% Options for packages loaded elsewhere
% Options for packages loaded elsewhere
\PassOptionsToPackage{unicode}{hyperref}
\PassOptionsToPackage{hyphens}{url}
\PassOptionsToPackage{dvipsnames,svgnames,x11names}{xcolor}
%
\documentclass[
  letterpaper,
  DIV=11,
  numbers=noendperiod]{scrreprt}
\usepackage{xcolor}
\usepackage{amsmath,amssymb}
\setcounter{secnumdepth}{5}
\usepackage{iftex}
\ifPDFTeX
  \usepackage[T1]{fontenc}
  \usepackage[utf8]{inputenc}
  \usepackage{textcomp} % provide euro and other symbols
\else % if luatex or xetex
  \usepackage{unicode-math} % this also loads fontspec
  \defaultfontfeatures{Scale=MatchLowercase}
  \defaultfontfeatures[\rmfamily]{Ligatures=TeX,Scale=1}
\fi
\usepackage{lmodern}
\ifPDFTeX\else
  % xetex/luatex font selection
\fi
% Use upquote if available, for straight quotes in verbatim environments
\IfFileExists{upquote.sty}{\usepackage{upquote}}{}
\IfFileExists{microtype.sty}{% use microtype if available
  \usepackage[]{microtype}
  \UseMicrotypeSet[protrusion]{basicmath} % disable protrusion for tt fonts
}{}
\makeatletter
\@ifundefined{KOMAClassName}{% if non-KOMA class
  \IfFileExists{parskip.sty}{%
    \usepackage{parskip}
  }{% else
    \setlength{\parindent}{0pt}
    \setlength{\parskip}{6pt plus 2pt minus 1pt}}
}{% if KOMA class
  \KOMAoptions{parskip=half}}
\makeatother
% Make \paragraph and \subparagraph free-standing
\makeatletter
\ifx\paragraph\undefined\else
  \let\oldparagraph\paragraph
  \renewcommand{\paragraph}{
    \@ifstar
      \xxxParagraphStar
      \xxxParagraphNoStar
  }
  \newcommand{\xxxParagraphStar}[1]{\oldparagraph*{#1}\mbox{}}
  \newcommand{\xxxParagraphNoStar}[1]{\oldparagraph{#1}\mbox{}}
\fi
\ifx\subparagraph\undefined\else
  \let\oldsubparagraph\subparagraph
  \renewcommand{\subparagraph}{
    \@ifstar
      \xxxSubParagraphStar
      \xxxSubParagraphNoStar
  }
  \newcommand{\xxxSubParagraphStar}[1]{\oldsubparagraph*{#1}\mbox{}}
  \newcommand{\xxxSubParagraphNoStar}[1]{\oldsubparagraph{#1}\mbox{}}
\fi
\makeatother

\usepackage{color}
\usepackage{fancyvrb}
\newcommand{\VerbBar}{|}
\newcommand{\VERB}{\Verb[commandchars=\\\{\}]}
\DefineVerbatimEnvironment{Highlighting}{Verbatim}{commandchars=\\\{\}}
% Add ',fontsize=\small' for more characters per line
\usepackage{framed}
\definecolor{shadecolor}{RGB}{241,243,245}
\newenvironment{Shaded}{\begin{snugshade}}{\end{snugshade}}
\newcommand{\AlertTok}[1]{\textcolor[rgb]{0.68,0.00,0.00}{#1}}
\newcommand{\AnnotationTok}[1]{\textcolor[rgb]{0.37,0.37,0.37}{#1}}
\newcommand{\AttributeTok}[1]{\textcolor[rgb]{0.40,0.45,0.13}{#1}}
\newcommand{\BaseNTok}[1]{\textcolor[rgb]{0.68,0.00,0.00}{#1}}
\newcommand{\BuiltInTok}[1]{\textcolor[rgb]{0.00,0.23,0.31}{#1}}
\newcommand{\CharTok}[1]{\textcolor[rgb]{0.13,0.47,0.30}{#1}}
\newcommand{\CommentTok}[1]{\textcolor[rgb]{0.37,0.37,0.37}{#1}}
\newcommand{\CommentVarTok}[1]{\textcolor[rgb]{0.37,0.37,0.37}{\textit{#1}}}
\newcommand{\ConstantTok}[1]{\textcolor[rgb]{0.56,0.35,0.01}{#1}}
\newcommand{\ControlFlowTok}[1]{\textcolor[rgb]{0.00,0.23,0.31}{\textbf{#1}}}
\newcommand{\DataTypeTok}[1]{\textcolor[rgb]{0.68,0.00,0.00}{#1}}
\newcommand{\DecValTok}[1]{\textcolor[rgb]{0.68,0.00,0.00}{#1}}
\newcommand{\DocumentationTok}[1]{\textcolor[rgb]{0.37,0.37,0.37}{\textit{#1}}}
\newcommand{\ErrorTok}[1]{\textcolor[rgb]{0.68,0.00,0.00}{#1}}
\newcommand{\ExtensionTok}[1]{\textcolor[rgb]{0.00,0.23,0.31}{#1}}
\newcommand{\FloatTok}[1]{\textcolor[rgb]{0.68,0.00,0.00}{#1}}
\newcommand{\FunctionTok}[1]{\textcolor[rgb]{0.28,0.35,0.67}{#1}}
\newcommand{\ImportTok}[1]{\textcolor[rgb]{0.00,0.46,0.62}{#1}}
\newcommand{\InformationTok}[1]{\textcolor[rgb]{0.37,0.37,0.37}{#1}}
\newcommand{\KeywordTok}[1]{\textcolor[rgb]{0.00,0.23,0.31}{\textbf{#1}}}
\newcommand{\NormalTok}[1]{\textcolor[rgb]{0.00,0.23,0.31}{#1}}
\newcommand{\OperatorTok}[1]{\textcolor[rgb]{0.37,0.37,0.37}{#1}}
\newcommand{\OtherTok}[1]{\textcolor[rgb]{0.00,0.23,0.31}{#1}}
\newcommand{\PreprocessorTok}[1]{\textcolor[rgb]{0.68,0.00,0.00}{#1}}
\newcommand{\RegionMarkerTok}[1]{\textcolor[rgb]{0.00,0.23,0.31}{#1}}
\newcommand{\SpecialCharTok}[1]{\textcolor[rgb]{0.37,0.37,0.37}{#1}}
\newcommand{\SpecialStringTok}[1]{\textcolor[rgb]{0.13,0.47,0.30}{#1}}
\newcommand{\StringTok}[1]{\textcolor[rgb]{0.13,0.47,0.30}{#1}}
\newcommand{\VariableTok}[1]{\textcolor[rgb]{0.07,0.07,0.07}{#1}}
\newcommand{\VerbatimStringTok}[1]{\textcolor[rgb]{0.13,0.47,0.30}{#1}}
\newcommand{\WarningTok}[1]{\textcolor[rgb]{0.37,0.37,0.37}{\textit{#1}}}

\usepackage{longtable,booktabs,array}
\newcounter{none} % for unnumbered tables
\usepackage{calc} % for calculating minipage widths
% Correct order of tables after \paragraph or \subparagraph
\usepackage{etoolbox}
\makeatletter
\patchcmd\longtable{\par}{\if@noskipsec\mbox{}\fi\par}{}{}
\makeatother
% Allow footnotes in longtable head/foot
\IfFileExists{footnotehyper.sty}{\usepackage{footnotehyper}}{\usepackage{footnote}}
\makesavenoteenv{longtable}
\usepackage{graphicx}
\makeatletter
\newsavebox\pandoc@box
\newcommand*\pandocbounded[1]{% scales image to fit in text height/width
  \sbox\pandoc@box{#1}%
  \Gscale@div\@tempa{\textheight}{\dimexpr\ht\pandoc@box+\dp\pandoc@box\relax}%
  \Gscale@div\@tempb{\linewidth}{\wd\pandoc@box}%
  \ifdim\@tempb\p@<\@tempa\p@\let\@tempa\@tempb\fi% select the smaller of both
  \ifdim\@tempa\p@<\p@\scalebox{\@tempa}{\usebox\pandoc@box}%
  \else\usebox{\pandoc@box}%
  \fi%
}
% Set default figure placement to htbp
\def\fps@figure{htbp}
\makeatother





\setlength{\emergencystretch}{3em} % prevent overfull lines

\providecommand{\tightlist}{%
  \setlength{\itemsep}{0pt}\setlength{\parskip}{0pt}}



 


\KOMAoption{captions}{tableheading}
\makeatletter
\@ifpackageloaded{bookmark}{}{\usepackage{bookmark}}
\makeatother
\makeatletter
\@ifpackageloaded{caption}{}{\usepackage{caption}}
\AtBeginDocument{%
\ifdefined\contentsname
  \renewcommand*\contentsname{Table of contents}
\else
  \newcommand\contentsname{Table of contents}
\fi
\ifdefined\listfigurename
  \renewcommand*\listfigurename{List of Figures}
\else
  \newcommand\listfigurename{List of Figures}
\fi
\ifdefined\listtablename
  \renewcommand*\listtablename{List of Tables}
\else
  \newcommand\listtablename{List of Tables}
\fi
\ifdefined\figurename
  \renewcommand*\figurename{Figure}
\else
  \newcommand\figurename{Figure}
\fi
\ifdefined\tablename
  \renewcommand*\tablename{Table}
\else
  \newcommand\tablename{Table}
\fi
}
\@ifpackageloaded{float}{}{\usepackage{float}}
\floatstyle{ruled}
\@ifundefined{c@chapter}{\newfloat{codelisting}{h}{lop}}{\newfloat{codelisting}{h}{lop}[chapter]}
\floatname{codelisting}{Listing}
\newcommand*\listoflistings{\listof{codelisting}{List of Listings}}
\makeatother
\makeatletter
\makeatother
\makeatletter
\@ifpackageloaded{caption}{}{\usepackage{caption}}
\@ifpackageloaded{subcaption}{}{\usepackage{subcaption}}
\makeatother
\makeatletter
\definecolor{QuartoInternalColor5}{rgb}{0.38,0.38,0.38}
\definecolor{QuartoInternalColor12}{rgb}{0.60,0.20,0.40}
\definecolor{QuartoInternalColor7}{rgb}{0.87,0.71,0.17}
\definecolor{QuartoInternalColor2}{rgb}{0,0,0}
\definecolor{QuartoInternalColor3}{rgb}{0.00,0.64,0.31}
\definecolor{QuartoInternalColor4}{rgb}{0.00,0.40,0.00}
\definecolor{QuartoInternalColor1}{rgb}{0.91,0.36,0.35}
\definecolor{QuartoInternalColor14}{rgb}{0.00,0.00,0.40}
\definecolor{QuartoInternalColor10}{rgb}{0.20,0.40,0.40}
\definecolor{QuartoInternalColor6}{rgb}{0.15,0.56,0.56}
\definecolor{QuartoInternalColor11}{rgb}{0.60,0.00,0.00}
\definecolor{QuartoInternalColor8}{rgb}{0.00,0.45,0.15}
\definecolor{QuartoInternalColor13}{rgb}{0.60,0.00,1.00}
\definecolor{QuartoInternalColor9}{rgb}{0.80,0.20,0.20}
\makeatother
\usepackage{bookmark}
\IfFileExists{xurl.sty}{\usepackage{xurl}}{} % add URL line breaks if available
\urlstyle{same}
\hypersetup{
  pdftitle={Métodos Computacionais da Astronomia},
  pdfauthor={Bruno E. Morgado; Luan M. Moraes (Org.)},
  colorlinks=true,
  linkcolor={blue},
  filecolor={Maroon},
  citecolor={Blue},
  urlcolor={Blue},
  pdfcreator={LaTeX via pandoc}}


\title{Métodos Computacionais da Astronomia}
\author{Bruno E. Morgado \and Luan M. Moraes (Org.)}
\date{}
\begin{document}
\maketitle

\renewcommand*\contentsname{Table of contents}
{
\setcounter{tocdepth}{2}
\tableofcontents
}

\bookmarksetup{startatroot}

\chapter*{Introdução}\label{introduuxe7uxe3o}
\addcontentsline{toc}{chapter}{Introdução}

\markboth{Introdução}{Introdução}

\part{Python Básico}

\chapter{Primeiros Passos}\label{primeiros-passos}

\section{Tipos de variáveis}\label{tipos-de-variuxe1veis}

Cada variável possui um tipo específico, temos:

\begin{itemize}
\tightlist
\item
  Número inteiro (\texttt{int})
\item
  Número real (\texttt{float})
\item
  Texto (\texttt{str}, se fala \emph{string})
\item
  Valores lógicos (\texttt{bool})
\end{itemize}

Podemos descobrir o tipo de uma variável usando a função
\texttt{type()}. Por exemplo:

\begin{Shaded}
\begin{Highlighting}[]
\NormalTok{\#| autorun: true}
\NormalTok{type(1000000)}
\end{Highlighting}
\end{Shaded}

\begin{Shaded}
\begin{Highlighting}[]
\NormalTok{\#| autorun: true}
\NormalTok{type(2.0)}
\end{Highlighting}
\end{Shaded}

\begin{Shaded}
\begin{Highlighting}[]
\NormalTok{\#| autorun: true}
\NormalTok{type(\textquotesingle{}blablabla teste\textquotesingle{})}
\end{Highlighting}
\end{Shaded}

\begin{Shaded}
\begin{Highlighting}[]
\NormalTok{\#| autorun: true}
\NormalTok{type(True), type(False)}
\end{Highlighting}
\end{Shaded}

Mas como o Python é uma linguagem \emph{case-sensitive}, ou seja, que
diferencia maiúsculas de minúsculas, é preciso tomar cuidado com
variáveis booleanas:

\begin{Shaded}
\begin{Highlighting}[]
\NormalTok{\#| autorun: true}
\NormalTok{\#| error: true}
\NormalTok{\# Atenção ao case}
\NormalTok{type(true)}
\end{Highlighting}
\end{Shaded}

\section{Operações Básicas}\label{operauxe7uxf5es-buxe1sicas}

\begin{itemize}
\tightlist
\item
  \texttt{+} \(\rightarrow\) Soma
\item
  \texttt{-} \(\rightarrow\) Subtração
\item
  \texttt{*} \(\rightarrow\) Produto
\item
  \texttt{/} \(\rightarrow\) Divisão
\end{itemize}

\begin{Shaded}
\begin{Highlighting}[]
\NormalTok{\#| autorun: true}
\NormalTok{2+5}
\end{Highlighting}
\end{Shaded}

\begin{Shaded}
\begin{Highlighting}[]
\NormalTok{\#| autorun: true}
\NormalTok{2+5.0}
\end{Highlighting}
\end{Shaded}

\begin{Shaded}
\begin{Highlighting}[]
\NormalTok{\#| autorun: true}
\NormalTok{3/2}
\end{Highlighting}
\end{Shaded}

Em especial, podemos dar valores para variáveis usando o sinal
\texttt{=}. Ele indica que a variável à esquerda \emph{recebe} o valor à
direita.

\begin{Shaded}
\begin{Highlighting}[]
\NormalTok{\#| autorun: true}
\NormalTok{variavel = 2}
\end{Highlighting}
\end{Shaded}

\begin{Shaded}
\begin{Highlighting}[]
\NormalTok{\#| autorun: true}
\NormalTok{variavel}
\end{Highlighting}
\end{Shaded}

E podemos operar então com ela:

\begin{Shaded}
\begin{Highlighting}[]
\NormalTok{\#| autorun: true}
\NormalTok{variavel + 2 }
\end{Highlighting}
\end{Shaded}

Mas só dá pra usar variáveis que você já declarou, ou seja, as que já
receberam algum valor.

\begin{Shaded}
\begin{Highlighting}[]
\NormalTok{\#| autorun: true}
\NormalTok{\#| error: true}
\NormalTok{banana {-} 3 }
\end{Highlighting}
\end{Shaded}

Operações com número são bastante intuitivas, porém podemos ir além. A
soma de \emph{strings} também pode ser feita e é uma concatenação:

\begin{Shaded}
\begin{Highlighting}[]
\NormalTok{\#| autorun: true}
\NormalTok{\# Válido para strings}
\NormalTok{\textquotesingle{}observatorio\textquotesingle{} + \textquotesingle{} do \textquotesingle{} + \textquotesingle{}valongo\textquotesingle{}}
\end{Highlighting}
\end{Shaded}

Multiplicar também é possível, o que vai dar em uma \emph{string}
repetida várias vezes.

\begin{Shaded}
\begin{Highlighting}[]
\NormalTok{\#| autorun: true}
\NormalTok{\textquotesingle{}observatorio\textquotesingle{} + \textquotesingle{} do \textquotesingle{} * 5 + \textquotesingle{}valongo\textquotesingle{}}
\end{Highlighting}
\end{Shaded}

Em compensação, substrair ou dividir \emph{strings} não é possível:

\begin{Shaded}
\begin{Highlighting}[]
\NormalTok{\#| autorun: true}
\NormalTok{\#| error: true}
\NormalTok{\textquotesingle{}casa\textquotesingle{} {-} \textquotesingle{}a\textquotesingle{}}
\end{Highlighting}
\end{Shaded}

\begin{Shaded}
\begin{Highlighting}[]
\NormalTok{\#| autorun: true}
\NormalTok{\#| error: true}
\NormalTok{\textquotesingle{}casa\textquotesingle{} / 2}
\end{Highlighting}
\end{Shaded}

\begin{Shaded}
\begin{Highlighting}[]
\NormalTok{\#| autorun: true}
\NormalTok{\# Válido para Booleanos}
\NormalTok{bool(True + False)}
\end{Highlighting}
\end{Shaded}

\begin{Shaded}
\begin{Highlighting}[]
\NormalTok{\#| autorun: true}
\NormalTok{\# 0 = False}
\NormalTok{\# 1 = True}

\NormalTok{bool(1 * 0)}
\end{Highlighting}
\end{Shaded}

\begin{Shaded}
\begin{Highlighting}[]
\NormalTok{\#| autorun: true}
\NormalTok{bool(True + True)}
\end{Highlighting}
\end{Shaded}

Aqui, a soma e o produto são operações lógicas, onde:

\begin{itemize}
\tightlist
\item
  \texttt{+} significa \emph{ou}
\item
  \texttt{*} significa \emph{e}
\end{itemize}

Se você não sabe ou não lembra como os operadores lógicos funcionam,
basta conferir o Appendix~\ref{sec-appendix-a}.

\begin{Shaded}
\begin{Highlighting}[]
\NormalTok{\#| autorun: true}
\NormalTok{teste\_1 = True}
\NormalTok{teste\_2 = False}

\NormalTok{\# Teste\_1 E teste\_2 são verdadeiros?}
\NormalTok{bool(teste\_1 * teste\_2)}
\end{Highlighting}
\end{Shaded}

\section{Condições}\label{condiuxe7uxf5es}

\begin{itemize}
\tightlist
\item
  \texttt{if} - Se for verdade
\item
  \texttt{elif} - E se for verdade
\item
  \texttt{if\ not} - Se não for verdade
\item
  \texttt{else} - Caso contrário
\end{itemize}

\begin{Shaded}
\begin{Highlighting}[]
\NormalTok{\#| autorun: true}
\NormalTok{if True:}
\NormalTok{    print(\textquotesingle{}oi\textquotesingle{})}
\end{Highlighting}
\end{Shaded}

\begin{Shaded}
\begin{Highlighting}[]
\NormalTok{\#| autorun: true}
\NormalTok{if False:}
\NormalTok{    print(\textquotesingle{}oi\textquotesingle{})}
\end{Highlighting}
\end{Shaded}

\begin{Shaded}
\begin{Highlighting}[]
\NormalTok{\#| autorun: true}
\NormalTok{if True:}
\NormalTok{    print(\textquotesingle{}oi\textquotesingle{})}
\NormalTok{else:}
\NormalTok{    print(\textquotesingle{}não deu\textquotesingle{})}
\end{Highlighting}
\end{Shaded}

\begin{Shaded}
\begin{Highlighting}[]
\NormalTok{\#| autorun: true}
\NormalTok{if not True:}
\NormalTok{    print(\textquotesingle{}oi\textquotesingle{})}
\NormalTok{else:}
\NormalTok{    print(\textquotesingle{}não deu\textquotesingle{})}
\end{Highlighting}
\end{Shaded}

\begin{Shaded}
\begin{Highlighting}[]
\NormalTok{\#| autorun: true}
\NormalTok{if  True * False:}
\NormalTok{    print(\textquotesingle{}oi\textquotesingle{})}

\NormalTok{elif True:}
\NormalTok{    print(\textquotesingle{}de novo\textquotesingle{})}

\NormalTok{else:}
\NormalTok{    print(\textquotesingle{}não deu\textquotesingle{})}
\end{Highlighting}
\end{Shaded}

O que resulta em um booleano? Uma pergunta!

\begin{itemize}
\tightlist
\item
  \texttt{==} \(\rightarrow\) É igual a?
\item
  \texttt{!=} \(\rightarrow\) É diferente de?
\item
  \texttt{\textgreater{}} \(\rightarrow\) É maior que?
\item
  \texttt{\textless{}} \(\rightarrow\) É menor que?
\item
  \texttt{\textgreater{}=} \(\rightarrow\) É maior ou igual a?
\item
  \texttt{\textless{}=} \(\rightarrow\) É menor ou igual a?
\end{itemize}

\begin{Shaded}
\begin{Highlighting}[]
\NormalTok{\#| autorun: true}
\NormalTok{2 == 2}
\end{Highlighting}
\end{Shaded}

\begin{Shaded}
\begin{Highlighting}[]
\NormalTok{\#| autorun: true}
\NormalTok{2 \textgreater{} 2}
\end{Highlighting}
\end{Shaded}

\begin{Shaded}
\begin{Highlighting}[]
\NormalTok{\#| autorun: true}
\NormalTok{2 \textgreater{}= 3}
\end{Highlighting}
\end{Shaded}

\begin{Shaded}
\begin{Highlighting}[]
\NormalTok{\#| autorun: true}
\NormalTok{3 != 2}
\end{Highlighting}
\end{Shaded}

\begin{Shaded}
\begin{Highlighting}[]
\NormalTok{\#| autorun: true}
\NormalTok{valor = 15}

\NormalTok{if valor \textgreater{} 5:}
\NormalTok{    print(\textquotesingle{}Valor é maior que 5\textquotesingle{})}
\NormalTok{elif valor == 5:}
\NormalTok{    print(\textquotesingle{}Valor é igual a cinco\textquotesingle{})}
\NormalTok{else:}
\NormalTok{    print(\textquotesingle{}Valor é menor que cinco\textquotesingle{})}
\end{Highlighting}
\end{Shaded}

Alterem a variável \texttt{valor}.

\begin{Shaded}
\begin{Highlighting}[]
\NormalTok{\#| autorun: true}
\NormalTok{int(1.9999999)}
\end{Highlighting}
\end{Shaded}

\section{Conjuntos de variáveis}\label{conjuntos-de-variuxe1veis}

\begin{itemize}
\tightlist
\item
  Tupla

  \begin{itemize}
  \tightlist
  \item
    Slice
  \end{itemize}
\item
  Dicionário
\end{itemize}

\begin{Shaded}
\begin{Highlighting}[]
\NormalTok{\#| autorun: true}
\NormalTok{conjunto = [1,2,3,4,5]}
\end{Highlighting}
\end{Shaded}

\begin{Shaded}
\begin{Highlighting}[]
\NormalTok{\#| autorun: true}
\NormalTok{type(conjunto)}
\end{Highlighting}
\end{Shaded}

\begin{Shaded}
\begin{Highlighting}[]
\NormalTok{\#| autorun: true}
\NormalTok{\# Quantos elementos tem esse conjunto de variaveis?}
\NormalTok{len(conjunto)}
\end{Highlighting}
\end{Shaded}

\begin{Shaded}
\begin{Highlighting}[]
\NormalTok{\#| autorun: true}
\NormalTok{\# Como criar um conjunto ordenado?}
\NormalTok{conjunto\_2 = range(10)}

\NormalTok{print(conjunto\_2)}
\end{Highlighting}
\end{Shaded}

\begin{Shaded}
\begin{Highlighting}[]
\NormalTok{\#| autorun: true}
\NormalTok{len(conjunto\_2)}
\end{Highlighting}
\end{Shaded}

\begin{Shaded}
\begin{Highlighting}[]
\NormalTok{\#| autorun: true}
\NormalTok{import numpy as np}
\end{Highlighting}
\end{Shaded}

\begin{Shaded}
\begin{Highlighting}[]
\CommentTok{\# IMPORTANTE: As funções têm Docstring!! (Ou deveriam ter)}
\BuiltInTok{int}\NormalTok{?}
\end{Highlighting}
\end{Shaded}

\begin{verbatim}
Init signature: int(self, /, *args, **kwargs)
Docstring:     
int([x]) -> integer
int(x, base=10) -> integer

Convert a number or string to an integer, or return 0 if no arguments
are given.  If x is a number, return x.__int__().  For floating-point
numbers, this truncates towards zero.

If x is not a number or if base is given, then x must be a string,
bytes, or bytearray instance representing an integer literal in the
given base.  The literal can be preceded by '+' or '-' and be surrounded
by whitespace.  The base defaults to 10.  Valid bases are 0 and 2-36.
Base 0 means to interpret the base from the string as an integer literal.
>>> int('0b100', base=0)
4
Type:           type
Subclasses:     bool, IntEnum, IntFlag, _NamedIntConstant, Handle, _CacheHashWrapper
\end{verbatim}

\begin{Shaded}
\begin{Highlighting}[]
\NormalTok{\#| autorun: true}
\NormalTok{\# como eu seleciono um elemento desse conjunto?}
\NormalTok{\# IMPORTANTE, Python sempre começa no zero!}

\NormalTok{conjunto\_3 = [\textquotesingle{}casa\textquotesingle{}, \textquotesingle{}valongo\textquotesingle{}, \textquotesingle{}trabalho\textquotesingle{}]}

\NormalTok{print(conjunto\_3[0])}
\NormalTok{print(conjunto\_3[1])}
\NormalTok{print(conjunto\_3[2])}
\end{Highlighting}
\end{Shaded}

\begin{Shaded}
\begin{Highlighting}[]
\NormalTok{\#| autorun: true}
\NormalTok{conjunto}
\end{Highlighting}
\end{Shaded}

\begin{Shaded}
\begin{Highlighting}[]
\NormalTok{\#| autorun: true}
\NormalTok{conjunto\_3.append(\textquotesingle{}extra\textquotesingle{})}
\end{Highlighting}
\end{Shaded}

\begin{Shaded}
\begin{Highlighting}[]
\NormalTok{\#| autorun: true}
\NormalTok{conjunto\_3.sort()}
\end{Highlighting}
\end{Shaded}

\begin{Shaded}
\begin{Highlighting}[]
\NormalTok{\#| autorun: true}
\NormalTok{conjunto\_3}
\end{Highlighting}
\end{Shaded}

\begin{Shaded}
\begin{Highlighting}[]
\NormalTok{\#| autorun: true}
\NormalTok{conjunto\_3.remove(\textquotesingle{}trabalho\textquotesingle{})}
\end{Highlighting}
\end{Shaded}

\begin{Shaded}
\begin{Highlighting}[]
\NormalTok{\#| autorun: true}
\NormalTok{conjunto\_3}
\end{Highlighting}
\end{Shaded}

\begin{Shaded}
\begin{Highlighting}[]
\NormalTok{\#| autorun: true}
\NormalTok{retorno = conjunto\_3.pop(2)}
\end{Highlighting}
\end{Shaded}

\begin{Shaded}
\begin{Highlighting}[]
\NormalTok{\#| autorun: true}
\NormalTok{retorno}
\end{Highlighting}
\end{Shaded}

\begin{Shaded}
\begin{Highlighting}[]
\NormalTok{\#| autorun: true}
\NormalTok{print(conjunto\_3)}
\end{Highlighting}
\end{Shaded}

Para pegar vários elementos pegamos um \texttt{slice}

\begin{Shaded}
\begin{Highlighting}[]
\NormalTok{\#| autorun: true}
\NormalTok{a = [0,1,2,3,4,5,6,7,8,9]}
\end{Highlighting}
\end{Shaded}

\begin{Shaded}
\begin{Highlighting}[]
\NormalTok{\#| autorun: true}
\NormalTok{a}
\end{Highlighting}
\end{Shaded}

Slice funciona assim:

{[}\texttt{inicio} : \texttt{final} : \texttt{passo}{]}

\begin{Shaded}
\begin{Highlighting}[]
\NormalTok{\#| autorun: true}
\NormalTok{\# Todos}
\NormalTok{print(a[:])}
\end{Highlighting}
\end{Shaded}

\begin{Shaded}
\begin{Highlighting}[]
\NormalTok{\#| autorun: true}
\NormalTok{\# até o quinto elemento}
\NormalTok{print(a[0:5])}
\NormalTok{\#| autorun: true}
\end{Highlighting}
\end{Shaded}

\begin{Shaded}
\begin{Highlighting}[]
\NormalTok{\#| autorun: true}
\NormalTok{\# a partir do quinto elemento}
\NormalTok{print(a[5:])}
\end{Highlighting}
\end{Shaded}

\begin{Shaded}
\begin{Highlighting}[]
\NormalTok{\#| autorun: true}
\NormalTok{a[::3]}
\end{Highlighting}
\end{Shaded}

\begin{Shaded}
\begin{Highlighting}[]
\NormalTok{\#| autorun: true}
\NormalTok{\# entre o 5 e o 7 elementos}
\NormalTok{print(a[5:7])}
\end{Highlighting}
\end{Shaded}

Uma \texttt{tupla} é apenas um conjunto de informações, porém a
informação também pode ser organizada na forma de um
\texttt{dicionário}.

\begin{Shaded}
\begin{Highlighting}[]
\NormalTok{\#| autorun: true}
\NormalTok{dic = \{\textquotesingle{}chave\textquotesingle{} : \textquotesingle{}valor\textquotesingle{}\}}
\end{Highlighting}
\end{Shaded}

\begin{Shaded}
\begin{Highlighting}[]
\NormalTok{\#| autorun: true}
\NormalTok{print(dic)}
\end{Highlighting}
\end{Shaded}

Em especial, em um \texttt{dicionário} podemos organizar informações de
diferentes tipos.

\begin{Shaded}
\begin{Highlighting}[]
\NormalTok{\#| autorun: true}
\NormalTok{prof = \{\textquotesingle{}Nome\textquotesingle{} : \textquotesingle{}Bruno\textquotesingle{},}
\NormalTok{        \textquotesingle{}Aniversario\textquotesingle{} : \textquotesingle{}03/05\textquotesingle{},}
\NormalTok{        \textquotesingle{}Tem gatos?\textquotesingle{} : True,}
\NormalTok{        \textquotesingle{}Gatos\textquotesingle{} : [1, 2, 3, 4, 5]\}}
\end{Highlighting}
\end{Shaded}

\begin{Shaded}
\begin{Highlighting}[]
\NormalTok{\#| autorun: true}
\NormalTok{print(prof)}
\end{Highlighting}
\end{Shaded}

\begin{Shaded}
\begin{Highlighting}[]
\NormalTok{\#| autorun: true}
\NormalTok{prof[\textquotesingle{}Nome\textquotesingle{}]}
\end{Highlighting}
\end{Shaded}

\section{Iterações}\label{iterauxe7uxf5es}

\begin{itemize}
\tightlist
\item
  \texttt{for} \ldots{} \texttt{in} \ldots{}
\item
  \texttt{while}
\item
  \texttt{enumerate}
\end{itemize}

\begin{Shaded}
\begin{Highlighting}[]
\NormalTok{\#| autorun: true}
\NormalTok{conjunto\_3}
\end{Highlighting}
\end{Shaded}

\begin{Shaded}
\begin{Highlighting}[]
\NormalTok{\#| autorun: true}
\NormalTok{n = 0}
\NormalTok{for a\_1 in conjunto\_3:}
\NormalTok{    print(a\_1)}
\NormalTok{    n = n + 1}
\end{Highlighting}
\end{Shaded}

\begin{Shaded}
\begin{Highlighting}[]
\NormalTok{\#| autorun: true}
\NormalTok{n}
\end{Highlighting}
\end{Shaded}

\begin{Shaded}
\begin{Highlighting}[]
\NormalTok{\#| autorun: true}
\NormalTok{for a\_2 in conjunto\_2:}
\NormalTok{    print(a\_2 + 50)}
\end{Highlighting}
\end{Shaded}

\begin{Shaded}
\begin{Highlighting}[]
\NormalTok{\#| autorun: true}
\NormalTok{banana = [10, 27]}
\end{Highlighting}
\end{Shaded}

\begin{Shaded}
\begin{Highlighting}[]
\NormalTok{\#| autorun: true}
\NormalTok{for banana\_1 in [10, 27, 32, 64, 200]:}
\NormalTok{    print(banana\_1)}
\end{Highlighting}
\end{Shaded}

\begin{Shaded}
\begin{Highlighting}[]
\NormalTok{\#| autorun: true}
\NormalTok{banana\_1}
\end{Highlighting}
\end{Shaded}

\begin{Shaded}
\begin{Highlighting}[]
\NormalTok{\#| autorun: true}
\NormalTok{conjunto\_4 = []}
\NormalTok{soma = 0}
\NormalTok{for a\_3 in conjunto\_2:}
\NormalTok{    conjunto\_4.append(a\_3 + 50)}
\NormalTok{    soma = soma + a\_3}
\end{Highlighting}
\end{Shaded}

\begin{Shaded}
\begin{Highlighting}[]
\NormalTok{\#| autorun: true}
\NormalTok{conjunto\_2}
\end{Highlighting}
\end{Shaded}

\begin{Shaded}
\begin{Highlighting}[]
\NormalTok{\#| autorun: true}
\NormalTok{conjunto\_4}
\end{Highlighting}
\end{Shaded}

\begin{Shaded}
\begin{Highlighting}[]
\NormalTok{\#| autorun: true}
\NormalTok{soma}
\end{Highlighting}
\end{Shaded}

\texttt{for} utiliza os valores de um determinado conjunto, por outro
lado, o \texttt{while} irá funcionar até uma condição ser satisfeita

\begin{Shaded}
\begin{Highlighting}[]
\NormalTok{\#| autorun: true}
\NormalTok{a\_4 = 0}
\NormalTok{while a\_4 \textless{} 10:}
\NormalTok{    a\_4 = a\_4 + 1}
\NormalTok{    print(a\_4)}
\end{Highlighting}
\end{Shaded}

\begin{Shaded}
\begin{Highlighting}[]
\NormalTok{\#| autorun: true}
\NormalTok{a\_5 = 0}
\NormalTok{while a\_5 \textless{} 10:}
\NormalTok{    print(a\_5)}
\NormalTok{    a\_5 = a\_5 + 1}
\NormalTok{print(a\_5)}
\end{Highlighting}
\end{Shaded}

\textbf{ATENÇÃO:} Aqui existe um grande perigo de entrar em um loop
infinito!!

\begin{Shaded}
\begin{Highlighting}[]
\NormalTok{\#| autorun: false}
\NormalTok{\#| timelimit: 1}
\NormalTok{\#| runbutton: true}
\NormalTok{a\_6 = 0}
\NormalTok{while a\_6 \textless{} 50:}
\NormalTok{    a\_6 = a\_6 {-} 1}
\NormalTok{print(a\_6)}
\end{Highlighting}
\end{Shaded}

\textbf{Extra:} Controle de Ansiedade

\begin{Shaded}
\begin{Highlighting}[]
\NormalTok{\#| autorun: false}
\NormalTok{\#| runbutton: true}
\NormalTok{\#| timelimit: 1}
\NormalTok{from tqdm import tqdm}
\NormalTok{for a\_7 in tqdm(range(100000000), \textquotesingle{}To testando\textquotesingle{}):}
\NormalTok{    a\_7 + 50}
\end{Highlighting}
\end{Shaded}

Uma outra forma de fazer loop é com o \texttt{enumerate}:

\begin{Shaded}
\begin{Highlighting}[]
\NormalTok{\#| autorun: true}
\NormalTok{for i, a\_8 in enumerate(range(100, 0, {-}1)):}
\NormalTok{    print(i, a\_8)}
\end{Highlighting}
\end{Shaded}

\section{Ler e Salvar arquivos}\label{ler-e-salvar-arquivos}

\begin{Shaded}
\begin{Highlighting}[]
\NormalTok{pwd}
\end{Highlighting}
\end{Shaded}

\begin{verbatim}
'D:\\OneDrive\\Desktop\\Livros Didaticos\\livro metcomp\\começo'
\end{verbatim}

\begin{Shaded}
\begin{Highlighting}[]
\NormalTok{\#| autorun: true}
\NormalTok{\#| edit: false}
\NormalTok{valor\_1 = []}
\NormalTok{f = open(\textquotesingle{}input.dat\textquotesingle{}, \textquotesingle{}r\textquotesingle{})}
\NormalTok{lines = f.readlines()  \# r = read}
\end{Highlighting}
\end{Shaded}

\begin{Shaded}
\begin{Highlighting}[]
\NormalTok{\#| autorun: true}
\NormalTok{lines}
\end{Highlighting}
\end{Shaded}

\begin{Shaded}
\begin{Highlighting}[]
\BuiltInTok{open}\NormalTok{?}
\end{Highlighting}
\end{Shaded}

\begin{verbatim}
Signature:
open(
    file,
    mode='r',
    buffering=-1,
    encoding=None,
    errors=None,
    newline=None,
    closefd=True,
    opener=None,
)
Docstring:
Open file and return a stream.  Raise OSError upon failure.

file is either a text or byte string giving the name (and the path
if the file isn't in the current working directory) of the file to
be opened or an integer file descriptor of the file to be
wrapped. (If a file descriptor is given, it is closed when the
returned I/O object is closed, unless closefd is set to False.)

mode is an optional string that specifies the mode in which the file
is opened. It defaults to 'r' which means open for reading in text
mode.  Other common values are 'w' for writing (truncating the file if
it already exists), 'x' for creating and writing to a new file, and
'a' for appending (which on some Unix systems, means that all writes
append to the end of the file regardless of the current seek position).
In text mode, if encoding is not specified the encoding used is platform
dependent: locale.getencoding() is called to get the current locale encoding.
(For reading and writing raw bytes use binary mode and leave encoding
unspecified.) The available modes are:

========= ===============================================================
Character Meaning
--------- ---------------------------------------------------------------
'r'       open for reading (default)
'w'       open for writing, truncating the file first
'x'       create a new file and open it for writing
'a'       open for writing, appending to the end of the file if it exists
'b'       binary mode
't'       text mode (default)
'+'       open a disk file for updating (reading and writing)
========= ===============================================================

The default mode is 'rt' (open for reading text). For binary random
access, the mode 'w+b' opens and truncates the file to 0 bytes, while
'r+b' opens the file without truncation. The 'x' mode implies 'w' and
raises an `FileExistsError` if the file already exists.

Python distinguishes between files opened in binary and text modes,
even when the underlying operating system doesn't. Files opened in
binary mode (appending 'b' to the mode argument) return contents as
bytes objects without any decoding. In text mode (the default, or when
't' is appended to the mode argument), the contents of the file are
returned as strings, the bytes having been first decoded using a
platform-dependent encoding or using the specified encoding if given.

buffering is an optional integer used to set the buffering policy.
Pass 0 to switch buffering off (only allowed in binary mode), 1 to select
line buffering (only usable in text mode), and an integer > 1 to indicate
the size of a fixed-size chunk buffer.  When no buffering argument is
given, the default buffering policy works as follows:

* Binary files are buffered in fixed-size chunks; the size of the buffer
  is chosen using a heuristic trying to determine the underlying device's
  "block size" and falling back on `io.DEFAULT_BUFFER_SIZE`.
  On many systems, the buffer will typically be 4096 or 8192 bytes long.

* "Interactive" text files (files for which isatty() returns True)
  use line buffering.  Other text files use the policy described above
  for binary files.

encoding is the name of the encoding used to decode or encode the
file. This should only be used in text mode. The default encoding is
platform dependent, but any encoding supported by Python can be
passed.  See the codecs module for the list of supported encodings.

errors is an optional string that specifies how encoding errors are to
be handled---this argument should not be used in binary mode. Pass
'strict' to raise a ValueError exception if there is an encoding error
(the default of None has the same effect), or pass 'ignore' to ignore
errors. (Note that ignoring encoding errors can lead to data loss.)
See the documentation for codecs.register or run 'help(codecs.Codec)'
for a list of the permitted encoding error strings.

newline controls how universal newlines works (it only applies to text
mode). It can be None, '', '\n', '\r', and '\r\n'.  It works as
follows:

* On input, if newline is None, universal newlines mode is
  enabled. Lines in the input can end in '\n', '\r', or '\r\n', and
  these are translated into '\n' before being returned to the
  caller. If it is '', universal newline mode is enabled, but line
  endings are returned to the caller untranslated. If it has any of
  the other legal values, input lines are only terminated by the given
  string, and the line ending is returned to the caller untranslated.

* On output, if newline is None, any '\n' characters written are
  translated to the system default line separator, os.linesep. If
  newline is '' or '\n', no translation takes place. If newline is any
  of the other legal values, any '\n' characters written are translated
  to the given string.

If closefd is False, the underlying file descriptor will be kept open
when the file is closed. This does not work when a file name is given
and must be True in that case.

A custom opener can be used by passing a callable as *opener*. The
underlying file descriptor for the file object is then obtained by
calling *opener* with (*file*, *flags*). *opener* must return an open
file descriptor (passing os.open as *opener* results in functionality
similar to passing None).

open() returns a file object whose type depends on the mode, and
through which the standard file operations such as reading and writing
are performed. When open() is used to open a file in a text mode ('w',
'r', 'wt', 'rt', etc.), it returns a TextIOWrapper. When used to open
a file in a binary mode, the returned class varies: in read binary
mode, it returns a BufferedReader; in write binary and append binary
modes, it returns a BufferedWriter, and in read/write mode, it returns
a BufferedRandom.

It is also possible to use a string or bytearray as a file for both
reading and writing. For strings StringIO can be used like a file
opened in a text mode, and for bytes a BytesIO can be used like a file
opened in a binary mode.
Type:      function
\end{verbatim}

\begin{Shaded}
\begin{Highlighting}[]
\NormalTok{\#| autorun: true}
\NormalTok{valor\_2 = []}
\NormalTok{f\_1 = open(\textquotesingle{}input.dat\textquotesingle{}, \textquotesingle{}r\textquotesingle{})}
\NormalTok{lines\_1 = f\_1.readlines()  \# r = read}
\NormalTok{for line in lines\_1:}
\NormalTok{    valor\_2.append(float(line))}
\NormalTok{f\_1.close()}
\end{Highlighting}
\end{Shaded}

\begin{Shaded}
\begin{Highlighting}[]
\NormalTok{\#| autorun: true}
\NormalTok{valor\_2}
\end{Highlighting}
\end{Shaded}

\begin{Shaded}
\begin{Highlighting}[]
\NormalTok{\#| autorun: true}
\NormalTok{valor\_2.sort()}
\end{Highlighting}
\end{Shaded}

\begin{Shaded}
\begin{Highlighting}[]
\NormalTok{\#| autorun: true}
\NormalTok{valor\_2}
\end{Highlighting}
\end{Shaded}

Exercício \(\rightarrow\) filtre e pegue apenas os valores positivos! e
crie uma nova tupla.

\begin{Shaded}
\begin{Highlighting}[]
\NormalTok{novos\_valores = []}

\NormalTok{\#...}
\NormalTok{\#...}
\NormalTok{\#...}
\NormalTok{\#...}
\NormalTok{\#...}
\end{Highlighting}
\end{Shaded}

Agora como salvar esses novos valores?

\begin{Shaded}
\begin{Highlighting}[]
\NormalTok{\#| autorun: true}
\NormalTok{f\_2 = open(\textquotesingle{}output.txt\textquotesingle{}, \textquotesingle{}w\textquotesingle{})  \# w = write}
\NormalTok{for v in valor\_2:}
\NormalTok{    f\_2.write(str(v) + \textquotesingle{}\textbackslash{}n\textquotesingle{})}
\NormalTok{f\_2.close()}
\end{Highlighting}
\end{Shaded}

\section{\texorpdfstring{Ferramentas úteis para
\emph{strings}}{Ferramentas úteis para strings}}\label{ferramentas-uxfateis-para-strings}

\begin{Shaded}
\begin{Highlighting}[]
\NormalTok{\#| autorun: true}
\NormalTok{\# Como fazer o print do valor junto com uma string?}
\NormalTok{print(\textquotesingle{}O valor é igual a \{:+8.2f\}\textquotesingle{}.format(valor\_2[20]))}
\end{Highlighting}
\end{Shaded}

\begin{Shaded}
\begin{Highlighting}[]
\NormalTok{\#| autorun: true}
\NormalTok{\# Pular a linha}
\NormalTok{print(\textquotesingle{}casa \textbackslash{}n do bruno\textquotesingle{})}
\end{Highlighting}
\end{Shaded}

\begin{Shaded}
\begin{Highlighting}[]
\NormalTok{\#| autorun: true}
\NormalTok{\# Adicionar tabulação}
\NormalTok{print(\textquotesingle{}teste 1 \textbackslash{}t = 2\textquotesingle{})}
\NormalTok{print(\textquotesingle{}teste 20 \textbackslash{}t = 2\textquotesingle{})}
\NormalTok{print(\textquotesingle{}teste 330 \textbackslash{}t = 2\textquotesingle{})}
\NormalTok{print(\textquotesingle{}teste 1000 \textbackslash{}t = 2\textquotesingle{})}
\end{Highlighting}
\end{Shaded}

\chapter{Python Orientado a Objetos}\label{python-orientado-a-objetos}

\section{Funções}\label{funuxe7uxf5es}

Quando precisamos repetir muitas vezes um pedaço de código, faz mais
sentido transformalo em uma função: - def -\textgreater{} Definir uma
função - return -\textgreater{} O que a função retorna

\begin{Shaded}
\begin{Highlighting}[]
\KeywordTok{def}\NormalTok{ soma(a, b):}
    \ControlFlowTok{return}\NormalTok{ a }\OperatorTok{+}\NormalTok{ b}
\end{Highlighting}
\end{Shaded}

\begin{Shaded}
\begin{Highlighting}[]
\NormalTok{soma(}\DecValTok{2}\NormalTok{,}\DecValTok{3}\NormalTok{)}
\end{Highlighting}
\end{Shaded}

\begin{verbatim}
5
\end{verbatim}

\begin{Shaded}
\begin{Highlighting}[]
\CommentTok{\# as variaveis dentro da função ficam dentro da função, a menos que você retorne elas}

\KeywordTok{def}\NormalTok{ soma(a, b):}
\NormalTok{    c }\OperatorTok{=} \DecValTok{10}
    \ControlFlowTok{return}\NormalTok{ a }\OperatorTok{+}\NormalTok{ b }\OperatorTok{+}\NormalTok{ c}
\end{Highlighting}
\end{Shaded}

\begin{Shaded}
\begin{Highlighting}[]
\NormalTok{soma(}\DecValTok{2}\NormalTok{,}\DecValTok{3}\NormalTok{)}
\end{Highlighting}
\end{Shaded}

\begin{verbatim}
15
\end{verbatim}

\begin{Shaded}
\begin{Highlighting}[]
\BuiltInTok{print}\NormalTok{(c)}
\end{Highlighting}
\end{Shaded}

\begin{Highlighting}
\textcolor{black}{NameError: name 'c' is not defined}
\textcolor{black}{}\textcolor{QuartoInternalColor1}{---------------------------------------------------------------------------}\textcolor{QuartoInternalColor2}{}
\textcolor{QuartoInternalColor2}{}\textcolor{QuartoInternalColor1}{NameError}\textcolor{QuartoInternalColor2}{                                 Traceback (most recent call last)}
\textcolor{QuartoInternalColor2}{Cell }\textcolor{QuartoInternalColor3}{In[5], line 1}\textcolor{QuartoInternalColor2}{}
\textcolor{QuartoInternalColor2}{}\textcolor{QuartoInternalColor3}{----> 1}\textcolor{QuartoInternalColor2}{ }\textcolor{QuartoInternalColor4}{print}\textcolor{QuartoInternalColor2}{(}\textcolor{QuartoInternalColor2}{c}\textcolor{QuartoInternalColor2}{)}
\textcolor{QuartoInternalColor2}{}\textcolor{QuartoInternalColor1}{NameError}\textcolor{QuartoInternalColor2}{: name 'c' is not defined}
\end{Highlighting}

O que acontece dentro da função fica dentro da função\ldots{}

\begin{Shaded}
\begin{Highlighting}[]
\CommentTok{\# Podemos dar valores predefinidos as variaveis}
\KeywordTok{def}\NormalTok{ soma(a, b}\OperatorTok{=}\DecValTok{5}\NormalTok{):}
    \ControlFlowTok{return}\NormalTok{ a }\OperatorTok{+}\NormalTok{ b}
\end{Highlighting}
\end{Shaded}

\begin{Shaded}
\begin{Highlighting}[]
\NormalTok{soma(}\DecValTok{5}\NormalTok{)}
\end{Highlighting}
\end{Shaded}

\begin{verbatim}
10
\end{verbatim}

\begin{Shaded}
\begin{Highlighting}[]
\NormalTok{soma(a}\OperatorTok{=}\DecValTok{5}\NormalTok{)}
\end{Highlighting}
\end{Shaded}

\begin{verbatim}
10
\end{verbatim}

\begin{Shaded}
\begin{Highlighting}[]
\NormalTok{soma(}\DecValTok{5}\NormalTok{, }\DecValTok{2}\NormalTok{)}
\end{Highlighting}
\end{Shaded}

\begin{verbatim}
7
\end{verbatim}

\begin{Shaded}
\begin{Highlighting}[]
\CommentTok{\# não necessariamente é necessário retornar algo}

\KeywordTok{def}\NormalTok{ soma\_1(a, b):}
    \BuiltInTok{print}\NormalTok{(a }\OperatorTok{+}\NormalTok{ b)}
    \ControlFlowTok{return}

\KeywordTok{def}\NormalTok{ soma\_2(a, b):}
    \ControlFlowTok{return}\NormalTok{ a }\OperatorTok{+}\NormalTok{ b}
\end{Highlighting}
\end{Shaded}

\begin{Shaded}
\begin{Highlighting}[]
\CommentTok{\# Qual a diferença?}

\NormalTok{soma\_1(}\DecValTok{2}\NormalTok{, }\DecValTok{5}\NormalTok{)}
\NormalTok{soma\_2(}\DecValTok{2}\NormalTok{, }\DecValTok{5}\NormalTok{)}
\end{Highlighting}
\end{Shaded}

\begin{verbatim}
7
\end{verbatim}

\begin{verbatim}
7
\end{verbatim}

\begin{Shaded}
\begin{Highlighting}[]
\NormalTok{valor }\OperatorTok{=}\NormalTok{ soma\_1(}\DecValTok{2}\NormalTok{, }\DecValTok{5}\NormalTok{)}
\end{Highlighting}
\end{Shaded}

\begin{verbatim}
7
\end{verbatim}

\begin{Shaded}
\begin{Highlighting}[]
\BuiltInTok{print}\NormalTok{(valor)}
\end{Highlighting}
\end{Shaded}

\begin{verbatim}
None
\end{verbatim}

\begin{Shaded}
\begin{Highlighting}[]
\NormalTok{valor }\OperatorTok{=}\NormalTok{ soma\_2(}\DecValTok{2}\NormalTok{, }\DecValTok{5}\NormalTok{)}
\end{Highlighting}
\end{Shaded}

\begin{Shaded}
\begin{Highlighting}[]
\BuiltInTok{print}\NormalTok{(valor)}
\end{Highlighting}
\end{Shaded}

\begin{verbatim}
7
\end{verbatim}

Além deles, também temos:

\begin{itemize}
\tightlist
\item
  args -\textgreater{} Argumentos das funções
\item
  kwargs -\textgreater{} Argumentos com Keyword
\end{itemize}

\begin{Shaded}
\begin{Highlighting}[]
\CommentTok{\#Vamos supor que queremos fazer a soma, mas o usuário pode escolher quantos valores ele quer colocar.}

\KeywordTok{def}\NormalTok{ soma(a,b):}
    \ControlFlowTok{return}\NormalTok{ a}\OperatorTok{+}\NormalTok{b}
\end{Highlighting}
\end{Shaded}

\begin{Shaded}
\begin{Highlighting}[]
\NormalTok{soma(}\DecValTok{2}\NormalTok{,}\DecValTok{3}\NormalTok{,}\DecValTok{5}\NormalTok{)}
\end{Highlighting}
\end{Shaded}

\begin{Highlighting}
\textcolor{black}{TypeError: soma() takes 2 positional arguments but 3 were given}
\textcolor{black}{}\textcolor{QuartoInternalColor1}{---------------------------------------------------------------------------}\textcolor{QuartoInternalColor2}{}
\textcolor{QuartoInternalColor2}{}\textcolor{QuartoInternalColor1}{TypeError}\textcolor{QuartoInternalColor2}{                                 Traceback (most recent call last)}
\textcolor{QuartoInternalColor2}{Cell }\textcolor{QuartoInternalColor3}{In[17], line 1}\textcolor{QuartoInternalColor2}{}
\textcolor{QuartoInternalColor2}{}\textcolor{QuartoInternalColor3}{----> 1}\textcolor{QuartoInternalColor2}{ }\textcolor{QuartoInternalColor2}{soma}\textcolor{QuartoInternalColor2}{}\textcolor{QuartoInternalColor2}{(}\textcolor{QuartoInternalColor2}{}\textcolor{QuartoInternalColor5}{2}\textcolor{QuartoInternalColor2}{}\textcolor{QuartoInternalColor2}{,}\textcolor{QuartoInternalColor2}{}\textcolor{QuartoInternalColor5}{3}\textcolor{QuartoInternalColor2}{}\textcolor{QuartoInternalColor2}{,}\textcolor{QuartoInternalColor2}{}\textcolor{QuartoInternalColor5}{5}\textcolor{QuartoInternalColor2}{}\textcolor{QuartoInternalColor2}{)}\textcolor{QuartoInternalColor2}{}
\textcolor{QuartoInternalColor2}{}\textcolor{QuartoInternalColor1}{TypeError}\textcolor{QuartoInternalColor2}{: soma() takes 2 positional arguments but 3 were given}
\end{Highlighting}

\begin{Shaded}
\begin{Highlighting}[]
\CommentTok{\# Solução, usar os argumentos}

\KeywordTok{def}\NormalTok{ soma(}\OperatorTok{*}\NormalTok{args):}
\NormalTok{    soma }\OperatorTok{=} \DecValTok{0}
    \ControlFlowTok{for}\NormalTok{ valor }\KeywordTok{in}\NormalTok{ args:}
\NormalTok{        soma }\OperatorTok{+=}\NormalTok{ valor}
    \ControlFlowTok{return}\NormalTok{ soma}
\end{Highlighting}
\end{Shaded}

\begin{Shaded}
\begin{Highlighting}[]
\NormalTok{soma(}\DecValTok{2}\NormalTok{,}\DecValTok{3}\NormalTok{)}
\end{Highlighting}
\end{Shaded}

\begin{verbatim}
5
\end{verbatim}

\begin{Shaded}
\begin{Highlighting}[]
\NormalTok{soma(}\DecValTok{2}\NormalTok{,}\DecValTok{3}\NormalTok{,}\DecValTok{5}\NormalTok{,}\DecValTok{6}\NormalTok{,}\DecValTok{7}\NormalTok{)}
\end{Highlighting}
\end{Shaded}

\begin{verbatim}
23
\end{verbatim}

Em resumo, você pode ter quantos argumentos você quiser, mas sem ter q
se preocupar em definir cada um deles.

Podemos inclusive colocar algumas informações extras

\begin{Shaded}
\begin{Highlighting}[]
\KeywordTok{def}\NormalTok{ soma(}\OperatorTok{*}\NormalTok{args, verbose}\OperatorTok{=}\VariableTok{False}\NormalTok{):}
\NormalTok{    soma }\OperatorTok{=} \DecValTok{0}
    \ControlFlowTok{for}\NormalTok{ valor }\KeywordTok{in}\NormalTok{ args:}
        \ControlFlowTok{if}\NormalTok{ verbose:}
            \BuiltInTok{print}\NormalTok{(}\StringTok{\textquotesingle{}somando }\SpecialCharTok{\{:+.1f\}}\StringTok{\textquotesingle{}}\NormalTok{.}\BuiltInTok{format}\NormalTok{(valor))}
\NormalTok{        soma }\OperatorTok{+=}\NormalTok{ valor}
    \ControlFlowTok{return}\NormalTok{ soma}
\end{Highlighting}
\end{Shaded}

\begin{Shaded}
\begin{Highlighting}[]
\NormalTok{valor }\OperatorTok{=}\NormalTok{ soma(}\DecValTok{2}\NormalTok{,}\DecValTok{3}\NormalTok{,}\DecValTok{5}\NormalTok{,}\DecValTok{6}\NormalTok{,}\DecValTok{7}\NormalTok{)}
\BuiltInTok{print}\NormalTok{(valor)}
\end{Highlighting}
\end{Shaded}

\begin{verbatim}
23
\end{verbatim}

\begin{Shaded}
\begin{Highlighting}[]
\NormalTok{valor }\OperatorTok{=}\NormalTok{ soma(}\DecValTok{2}\NormalTok{,}\DecValTok{3}\NormalTok{,}\DecValTok{5}\NormalTok{,}\DecValTok{6}\NormalTok{,}\DecValTok{7}\NormalTok{, verbose}\OperatorTok{=}\VariableTok{True}\NormalTok{)}
\BuiltInTok{print}\NormalTok{(valor)}
\end{Highlighting}
\end{Shaded}

\begin{verbatim}
somando +2.0
somando +3.0
somando +5.0
somando +6.0
somando +7.0
23
\end{verbatim}

Mas Bruno, eu preciso definir cada uma das variaveis nomeadas no inicio?

\begin{Shaded}
\begin{Highlighting}[]
\CommentTok{\# Aqui você pode usar os kwargs:}

\KeywordTok{def}\NormalTok{ soma(}\OperatorTok{*}\NormalTok{args, }\OperatorTok{**}\NormalTok{kwargs):}
\NormalTok{    soma }\OperatorTok{=} \DecValTok{0}
    \BuiltInTok{print}\NormalTok{(kwargs)}
    \ControlFlowTok{for}\NormalTok{ valor }\KeywordTok{in}\NormalTok{ args:}
\NormalTok{        soma }\OperatorTok{+=}\NormalTok{ valor}
    \ControlFlowTok{return}\NormalTok{ soma}
\end{Highlighting}
\end{Shaded}

\begin{Shaded}
\begin{Highlighting}[]
\NormalTok{soma(}\DecValTok{2}\NormalTok{,}\DecValTok{3}\NormalTok{,}\DecValTok{5}\NormalTok{,}\DecValTok{4}\NormalTok{, teste}\OperatorTok{=}\StringTok{\textquotesingle{}sim\textquotesingle{}}\NormalTok{, valor}\OperatorTok{=}\StringTok{\textquotesingle{}quero\textquotesingle{}}\NormalTok{, verbose}\OperatorTok{=}\VariableTok{True}\NormalTok{)}
\end{Highlighting}
\end{Shaded}

\begin{verbatim}
{'teste': 'sim', 'valor': 'quero', 'verbose': True}
\end{verbatim}

\begin{verbatim}
14
\end{verbatim}

O kwargs é um dicinário! Então como trabalhar com ele?

\begin{Shaded}
\begin{Highlighting}[]
\KeywordTok{def}\NormalTok{ soma(}\OperatorTok{*}\NormalTok{args, }\OperatorTok{**}\NormalTok{kwargs):}
\NormalTok{    soma }\OperatorTok{=} \DecValTok{0}
    \ControlFlowTok{if} \StringTok{\textquotesingle{}verbose\textquotesingle{}} \KeywordTok{in}\NormalTok{ kwargs:}
\NormalTok{        verbose }\OperatorTok{=}\NormalTok{ kwargs[}\StringTok{\textquotesingle{}verbose\textquotesingle{}}\NormalTok{]}
    \ControlFlowTok{for}\NormalTok{ valor }\KeywordTok{in}\NormalTok{ args:}
        \ControlFlowTok{if}\NormalTok{ verbose:}
            \BuiltInTok{print}\NormalTok{(}\StringTok{\textquotesingle{}somando }\SpecialCharTok{\{:+.1f\}}\StringTok{\textquotesingle{}}\NormalTok{.}\BuiltInTok{format}\NormalTok{(valor))}
\NormalTok{        soma }\OperatorTok{+=}\NormalTok{ valor}
    \ControlFlowTok{return}\NormalTok{ soma}
\end{Highlighting}
\end{Shaded}

\begin{Shaded}
\begin{Highlighting}[]
\NormalTok{soma(}\DecValTok{2}\NormalTok{,}\DecValTok{3}\NormalTok{,}\DecValTok{4}\NormalTok{,}\DecValTok{5}\NormalTok{, verbose}\OperatorTok{=}\VariableTok{True}\NormalTok{)}
\end{Highlighting}
\end{Shaded}

\begin{verbatim}
somando +2.0
somando +3.0
somando +4.0
somando +5.0
\end{verbatim}

\begin{verbatim}
14
\end{verbatim}

\begin{Shaded}
\begin{Highlighting}[]
\NormalTok{soma(}\DecValTok{2}\NormalTok{,}\DecValTok{3}\NormalTok{,}\DecValTok{4}\NormalTok{,}\DecValTok{5}\NormalTok{, verbose}\OperatorTok{=}\VariableTok{True}\NormalTok{, extra}\OperatorTok{=}\StringTok{\textquotesingle{}teste\textquotesingle{}}\NormalTok{)}
\end{Highlighting}
\end{Shaded}

\begin{verbatim}
somando +2.0
somando +3.0
somando +4.0
somando +5.0
\end{verbatim}

\begin{verbatim}
14
\end{verbatim}

O que aconteceu com esse extra? Nada!

\textbf{ATENÇÃO:} Na hora de fazer uma função é muito importante definir
sua Docstring

\begin{Shaded}
\begin{Highlighting}[]
\KeywordTok{def}\NormalTok{ soma(valor\_inicial, }\OperatorTok{*}\NormalTok{args, }\OperatorTok{**}\NormalTok{kwargs):}
    \CommentTok{"""Soma os valores adicionados.}

\CommentTok{    Parâmetros}
\CommentTok{    {-}{-}{-}{-}{-}{-}{-}{-}{-}{-}}
\CommentTok{    valor\_inicial : \textasciigrave{}int\textasciigrave{}, \textasciigrave{}float\textasciigrave{},}
\CommentTok{        Número inicial que será somado.}

\CommentTok{    args : \textasciigrave{}int\textasciigrave{}, \textasciigrave{}float\textasciigrave{},}
\CommentTok{        Números que serão somados a variavel valor.}

\CommentTok{    Kwargs}
\CommentTok{    {-}{-}{-}{-}{-}{-}{-}{-}{-}{-}}
\CommentTok{    verbose : \textasciigrave{}bool\textasciigrave{}}
\CommentTok{        Se \textasciigrave{}True\textasciigrave{}, irá mostrar informações durante as operações.}

\CommentTok{    Returns}
\CommentTok{    {-}{-}{-}{-}{-}{-}{-}}
\CommentTok{    soma : \textasciigrave{}int\textasciigrave{}, \textasciigrave{}float\textasciigrave{}}
\CommentTok{        Soma final dos parâmetros de entrada}
\CommentTok{    """}
\NormalTok{    soma }\OperatorTok{=}\NormalTok{ valor\_inicial}
    \ControlFlowTok{if} \StringTok{\textquotesingle{}verbose\textquotesingle{}} \KeywordTok{in}\NormalTok{ kwargs:}
\NormalTok{        verbose }\OperatorTok{=}\NormalTok{ kwargs[}\StringTok{\textquotesingle{}verbose\textquotesingle{}}\NormalTok{]}
    \ControlFlowTok{for}\NormalTok{ valor }\KeywordTok{in}\NormalTok{ args:}
        \ControlFlowTok{if}\NormalTok{ verbose:}
            \BuiltInTok{print}\NormalTok{(}\StringTok{\textquotesingle{}somando }\SpecialCharTok{\{:+.1f\}}\StringTok{\textquotesingle{}}\NormalTok{.}\BuiltInTok{format}\NormalTok{(valor))}
\NormalTok{        soma }\OperatorTok{+=}\NormalTok{ valor}
    \ControlFlowTok{return}\NormalTok{ soma}
\end{Highlighting}
\end{Shaded}

\begin{Shaded}
\begin{Highlighting}[]
\CommentTok{\# Agora eu consigo ler a docstring}
\NormalTok{soma?}
\end{Highlighting}
\end{Shaded}

\begin{Shaded}
\begin{Highlighting}[]
\CommentTok{\# também consigo ver o código interno}
\NormalTok{soma??}
\end{Highlighting}
\end{Shaded}

\section{Classes}\label{classes}

Quando apenas algumas funções não são o suficiente, criamos uma Classe
de objetos:

\begin{Shaded}
\begin{Highlighting}[]
\KeywordTok{class}\NormalTok{ Gatos:}
    \CommentTok{""" }
\CommentTok{    Dockstring blablabla}
\CommentTok{    """}

    \KeywordTok{def} \FunctionTok{\_\_init\_\_}\NormalTok{(}\VariableTok{self}\NormalTok{, }\OperatorTok{*}\NormalTok{, numero}\OperatorTok{=}\DecValTok{0}\NormalTok{):}
        \VariableTok{self}\NormalTok{.contagem }\OperatorTok{=}\NormalTok{ numero}
        \VariableTok{self}\NormalTok{.dono }\OperatorTok{=} \StringTok{\textquotesingle{}Você\textquotesingle{}}

    \AttributeTok{@property}
    \KeywordTok{def}\NormalTok{ quantos(}\VariableTok{self}\NormalTok{):}
        \ControlFlowTok{return} \VariableTok{self}\NormalTok{.contagem}

    \KeywordTok{def}\NormalTok{ adicionar(}\VariableTok{self}\NormalTok{, numero):}
        \CommentTok{"""}
\CommentTok{        adicionar número de gatos}
\CommentTok{        """}
        \BuiltInTok{print}\NormalTok{(}\StringTok{\textquotesingle{}Você tem mais gatos!\textquotesingle{}}\NormalTok{)}
        \VariableTok{self}\NormalTok{.contagem }\OperatorTok{+=}\NormalTok{ numero}
        \ControlFlowTok{return}\NormalTok{ numero}

    \KeywordTok{def}\NormalTok{ mudar\_dono(}\VariableTok{self}\NormalTok{, name):}
        \CommentTok{"""}
\CommentTok{        adiciona o nome do dono}
\CommentTok{        """}
        \VariableTok{self}\NormalTok{.dono }\OperatorTok{=}\NormalTok{ name}
        \ControlFlowTok{return}
\end{Highlighting}
\end{Shaded}

\begin{Shaded}
\begin{Highlighting}[]
\NormalTok{gt }\OperatorTok{=}\NormalTok{ Gatos()}
\end{Highlighting}
\end{Shaded}

\begin{Shaded}
\begin{Highlighting}[]
\BuiltInTok{type}\NormalTok{(gt)}
\end{Highlighting}
\end{Shaded}

\begin{verbatim}
__main__.Gatos
\end{verbatim}

\begin{Shaded}
\begin{Highlighting}[]
\NormalTok{gt.dono}
\end{Highlighting}
\end{Shaded}

\begin{verbatim}
'Você'
\end{verbatim}

\begin{Shaded}
\begin{Highlighting}[]
\NormalTok{gt.contagem}
\end{Highlighting}
\end{Shaded}

\begin{verbatim}
0
\end{verbatim}

\begin{Shaded}
\begin{Highlighting}[]
\NormalTok{gt.adicionar(}\DecValTok{2}\NormalTok{)}
\end{Highlighting}
\end{Shaded}

\begin{verbatim}
Você tem mais gatos!
\end{verbatim}

\begin{verbatim}
2
\end{verbatim}

\begin{Shaded}
\begin{Highlighting}[]
\NormalTok{gt.contagem}
\end{Highlighting}
\end{Shaded}

\begin{verbatim}
2
\end{verbatim}

\begin{Shaded}
\begin{Highlighting}[]
\NormalTok{gt.mudar\_dono(}\StringTok{\textquotesingle{}Bruno\textquotesingle{}}\NormalTok{)}
\end{Highlighting}
\end{Shaded}

\begin{Shaded}
\begin{Highlighting}[]
\NormalTok{gt.dono}
\end{Highlighting}
\end{Shaded}

\begin{verbatim}
'Bruno'
\end{verbatim}

\begin{Shaded}
\begin{Highlighting}[]
\KeywordTok{class}\NormalTok{ Gatos:}
    \CommentTok{""" }
\CommentTok{    Dockstring blablabla}
\CommentTok{    """}

    \KeywordTok{def} \FunctionTok{\_\_init\_\_}\NormalTok{(}\VariableTok{self}\NormalTok{, }\OperatorTok{*}\NormalTok{, numero}\OperatorTok{=}\DecValTok{0}\NormalTok{):}
        \VariableTok{self}\NormalTok{.contagem }\OperatorTok{=}\NormalTok{ numero}
        \VariableTok{self}\NormalTok{.dono }\OperatorTok{=} \StringTok{\textquotesingle{}Você\textquotesingle{}}

    \AttributeTok{@property}
    \KeywordTok{def}\NormalTok{ quantos(}\VariableTok{self}\NormalTok{):}
        \ControlFlowTok{return} \VariableTok{self}\NormalTok{.contagem}

    \KeywordTok{def}\NormalTok{ adicionar(}\VariableTok{self}\NormalTok{, numero):}
        \CommentTok{"""}
\CommentTok{        adicionar número de gatos}
\CommentTok{        """}
        \BuiltInTok{print}\NormalTok{(}\StringTok{\textquotesingle{}Você tem mais gatos!\textquotesingle{}}\NormalTok{)}
        \VariableTok{self}\NormalTok{.contagem }\OperatorTok{+=}\NormalTok{ numero}
        \ControlFlowTok{return}\NormalTok{ numero}

    \KeywordTok{def}\NormalTok{ mudar\_dono(}\VariableTok{self}\NormalTok{, name):}
        \CommentTok{"""}
\CommentTok{        adiciona o nome do dono}
\CommentTok{        """}
        \VariableTok{self}\NormalTok{.dono }\OperatorTok{=}\NormalTok{ name}
        \ControlFlowTok{return}

    \KeywordTok{def} \FunctionTok{\_\_repr\_\_}\NormalTok{(}\VariableTok{self}\NormalTok{):}
        \CommentTok{"""}
\CommentTok{        Representação na forma de String}
\CommentTok{        """}
        \ControlFlowTok{return} \StringTok{\textquotesingle{}}\SpecialCharTok{\{\}}\StringTok{ tem }\SpecialCharTok{\{\}}\StringTok{ gatos\textquotesingle{}}\NormalTok{.}\BuiltInTok{format}\NormalTok{(}\VariableTok{self}\NormalTok{.dono, }\VariableTok{self}\NormalTok{.contagem)}

    \KeywordTok{def} \FunctionTok{\_\_str\_\_}\NormalTok{(}\VariableTok{self}\NormalTok{):}
        \CommentTok{"""}
\CommentTok{        Representação na forma de string quando o print(obj) é usado}
\CommentTok{        """}
        \ControlFlowTok{return} \StringTok{\textquotesingle{}você está fazendo o print, e }\SpecialCharTok{\{\}}\StringTok{ tem }\SpecialCharTok{\{\}}\StringTok{ gatos\textquotesingle{}}\NormalTok{.}\BuiltInTok{format}\NormalTok{(}\VariableTok{self}\NormalTok{.dono, }\VariableTok{self}\NormalTok{.contagem)}
\end{Highlighting}
\end{Shaded}

\begin{Shaded}
\begin{Highlighting}[]
\NormalTok{gt }\OperatorTok{=}\NormalTok{ Gatos()}
\end{Highlighting}
\end{Shaded}

\begin{Shaded}
\begin{Highlighting}[]
\NormalTok{gt.dono}
\end{Highlighting}
\end{Shaded}

\begin{verbatim}
'Você'
\end{verbatim}

\begin{Shaded}
\begin{Highlighting}[]
\BuiltInTok{print}\NormalTok{(gt)}
\end{Highlighting}
\end{Shaded}

\begin{verbatim}
você está fazendo o print, e Você tem 0 gatos
\end{verbatim}

\begin{Shaded}
\begin{Highlighting}[]
\NormalTok{gt}
\end{Highlighting}
\end{Shaded}

\begin{verbatim}
Você tem 0 gatos
\end{verbatim}

\section{Pacotes}\label{pacotes}

Existe uma grande gama de pacotes python, como usar?

\begin{Shaded}
\begin{Highlighting}[]
\CommentTok{\# import numpy as np \# {-}\textgreater{} Importando o pacote inteiro e atribuindo um alias a ele}
\end{Highlighting}
\end{Shaded}

\begin{Shaded}
\begin{Highlighting}[]
\ImportTok{import}\NormalTok{ numpy }\CommentTok{\# {-}\textgreater{} Importando um pacote inteiro}
\end{Highlighting}
\end{Shaded}

\begin{Shaded}
\begin{Highlighting}[]
\NormalTok{numpy.array?}
\end{Highlighting}
\end{Shaded}

\begin{Shaded}
\begin{Highlighting}[]
\ImportTok{import}\NormalTok{ numpy }\ImportTok{as}\NormalTok{ np }\CommentTok{\# {-}\textgreater{} Importando o pacote inteiro e atribuindo um alias a ele}

\NormalTok{np.array?}
\end{Highlighting}
\end{Shaded}

\begin{Shaded}
\begin{Highlighting}[]
\NormalTok{numpy.array }\OperatorTok{==}\NormalTok{ np.array}
\end{Highlighting}
\end{Shaded}

\begin{verbatim}
True
\end{verbatim}

Não precisamos importar o pacote inteiro.

\begin{Shaded}
\begin{Highlighting}[]
\ImportTok{from}\NormalTok{ numpy }\ImportTok{import}\NormalTok{ random }\ImportTok{as}\NormalTok{ rd}
\end{Highlighting}
\end{Shaded}

\begin{Shaded}
\begin{Highlighting}[]
\NormalTok{rd.randint(low}\OperatorTok{=}\DecValTok{0}\NormalTok{, high}\OperatorTok{=}\DecValTok{10}\NormalTok{) }\CommentTok{\# {-}\textgreater{} Te devolve um número randomico entre o \textasciigrave{}low\textasciigrave{} e o \textasciigrave{}high\textasciigrave{} (Maiores detalhes na próxima aula)}
\end{Highlighting}
\end{Shaded}

\begin{verbatim}
1
\end{verbatim}

\begin{Shaded}
\begin{Highlighting}[]
\CommentTok{\# Podemos também pegar uma classe dentro de um módulo dentro de um pacote!}
\ImportTok{from}\NormalTok{ astropy.time }\ImportTok{import}\NormalTok{ Time}
\end{Highlighting}
\end{Shaded}

\begin{Shaded}
\begin{Highlighting}[]
\NormalTok{Time?}
\end{Highlighting}
\end{Shaded}

\textbf{Vamos criar o nosso próprio pacote!}

\begin{Shaded}
\begin{Highlighting}[]
\CommentTok{\#\#\#\# Lembrete, instalar o nov pacote, indo até o diretório onde encontramos o setup.py, e usar:}
\CommentTok{\#\#\#\# \textgreater{} pip install .}

\CommentTok{\#from bruno\_projeto import * \#Importando todas as funções/classes do pacote}
\end{Highlighting}
\end{Shaded}

\part{Pacotes Essenciais}

\chapter{OS e Numpy}\label{os-e-numpy}

\section{OS: Miscellaneous operating system
interface}\label{os-miscellaneous-operating-system-interface}

Aqui começamos com o pacote que nos permite navegar o sistema
operacional - Navegar os diretórios

\begin{Shaded}
\begin{Highlighting}[]
\NormalTok{import os}
\end{Highlighting}
\end{Shaded}

\begin{Shaded}
\begin{Highlighting}[]
\NormalTok{\# os.system(\textquotesingle{}ls\textquotesingle{}) \# Linux}
\NormalTok{os.system(\textquotesingle{}dir\textquotesingle{}) \# Windows }

\NormalTok{\# O que aconteceu? Deem uma olhada no terminal que está rodando o Jupyter!}
\end{Highlighting}
\end{Shaded}

\begin{Shaded}
\begin{Highlighting}[]
\NormalTok{\# os.system(\textquotesingle{}ls\textquotesingle{}) \# Linux}
\NormalTok{os.system(\textquotesingle{}dir\textquotesingle{}) \# Windows }

\NormalTok{\# O que aconteceu? Deem uma olhada no terminal que está rodando o Jupyter!}
\end{Highlighting}
\end{Shaded}

\begin{Shaded}
\begin{Highlighting}[]
\NormalTok{\# Informações uteis do seu computador }
\NormalTok{os.cpu\_count()}
\end{Highlighting}
\end{Shaded}

\begin{Shaded}
\begin{Highlighting}[]
\NormalTok{\# Lista os diretorios}
\NormalTok{os.listdir(\textquotesingle{}..\textquotesingle{})}
\end{Highlighting}
\end{Shaded}

\begin{Shaded}
\begin{Highlighting}[]
\NormalTok{\# buscar o diretório atual}
\NormalTok{os.getcwd()}
\end{Highlighting}
\end{Shaded}

\begin{Shaded}
\begin{Highlighting}[]
\NormalTok{\# Criando um novo diretório}
\NormalTok{os.mkdir(\textquotesingle{}Teste\_1\textquotesingle{})}
\end{Highlighting}
\end{Shaded}

\begin{Shaded}
\begin{Highlighting}[]
\NormalTok{os.listdir(\textquotesingle{}.\textquotesingle{})}
\end{Highlighting}
\end{Shaded}

\begin{Shaded}
\begin{Highlighting}[]
\NormalTok{\# Mudando de diretório}
\NormalTok{os.chdir(\textquotesingle{}Teste\_1/\textquotesingle{})}
\end{Highlighting}
\end{Shaded}

\begin{Shaded}
\begin{Highlighting}[]
\NormalTok{os.getcwd()}
\end{Highlighting}
\end{Shaded}

\begin{Shaded}
\begin{Highlighting}[]
\NormalTok{os.chdir(\textquotesingle{}..\textquotesingle{})}
\end{Highlighting}
\end{Shaded}

\begin{Shaded}
\begin{Highlighting}[]
\NormalTok{os.getcwd()}
\end{Highlighting}
\end{Shaded}

\begin{Shaded}
\begin{Highlighting}[]
\NormalTok{\# renomeando um diretório}
\NormalTok{os.rename(\textquotesingle{}Teste\_1/\textquotesingle{}, \textquotesingle{}Teste\_2/\textquotesingle{})}
\end{Highlighting}
\end{Shaded}

\begin{Shaded}
\begin{Highlighting}[]
\NormalTok{os.listdir(\textquotesingle{}.\textquotesingle{})}
\end{Highlighting}
\end{Shaded}

\begin{Shaded}
\begin{Highlighting}[]
\NormalTok{\# Removendo um diretório}
\NormalTok{os.rmdir(\textquotesingle{}Teste\_2/\textquotesingle{})}
\end{Highlighting}
\end{Shaded}

\begin{Shaded}
\begin{Highlighting}[]
\NormalTok{os.listdir(\textquotesingle{}.\textquotesingle{})}
\end{Highlighting}
\end{Shaded}

\begin{Shaded}
\begin{Highlighting}[]
\NormalTok{\# Vamos criar um arquivo de texto vazio}
\NormalTok{file = open(\textquotesingle{}texto.txt\textquotesingle{},\textquotesingle{}w\textquotesingle{})}

\NormalTok{file.close()}
\end{Highlighting}
\end{Shaded}

\begin{Shaded}
\begin{Highlighting}[]
\NormalTok{os.listdir(\textquotesingle{}.\textquotesingle{})}
\end{Highlighting}
\end{Shaded}

\begin{Shaded}
\begin{Highlighting}[]
\NormalTok{\# Mudando o nome de um arquivo}
\NormalTok{os.rename(\textquotesingle{}texto.txt\textquotesingle{},\textquotesingle{}texto\_novo.txt\textquotesingle{})}
\end{Highlighting}
\end{Shaded}

\begin{Shaded}
\begin{Highlighting}[]
\NormalTok{os.listdir()}
\end{Highlighting}
\end{Shaded}

\begin{Shaded}
\begin{Highlighting}[]
\NormalTok{\# Deletando um arquivo}
\NormalTok{os.remove(\textquotesingle{}texto\_novo.txt\textquotesingle{})}
\end{Highlighting}
\end{Shaded}

\begin{Shaded}
\begin{Highlighting}[]
\NormalTok{os.listdir()}
\end{Highlighting}
\end{Shaded}

\begin{Shaded}
\begin{Highlighting}[]
\NormalTok{os.listdir()}
\end{Highlighting}
\end{Shaded}

\textbf{Extra:} Se quiser receber só alguns arquivos específicos

\begin{Shaded}
\begin{Highlighting}[]
\NormalTok{import glob}
\end{Highlighting}
\end{Shaded}

\begin{Shaded}
\begin{Highlighting}[]
\NormalTok{glob.glob(\textquotesingle{}*.ipynb\textquotesingle{}) \# Igual faria no terminal}
\end{Highlighting}
\end{Shaded}

\section{Numpy: Array}\label{numpy-array}

Um pacote que visa automatizar e permite fazer operações matemáticas de
forma eficiente.

\begin{Shaded}
\begin{Highlighting}[]
\NormalTok{import numpy as np}
\end{Highlighting}
\end{Shaded}

\begin{Shaded}
\begin{Highlighting}[]
\NormalTok{\# Uma das grandes habilidades do Numpy é o np.array()}

\NormalTok{lista\_1 = [2,3,5,9,7]}
\NormalTok{lista\_2 = [6,7,2,1,3]}
\end{Highlighting}
\end{Shaded}

\begin{Shaded}
\begin{Highlighting}[]
\NormalTok{\# Há a concatenação!}
\NormalTok{lista\_1 + lista\_2}
\end{Highlighting}
\end{Shaded}

\begin{Shaded}
\begin{Highlighting}[]
\NormalTok{\# Mas eu querof azer operações! Neste caso:}

\NormalTok{lista\_1 = np.array([2,3,5,9,7])}
\NormalTok{lista\_2 = np.array([6,7,2,1,3])}
\end{Highlighting}
\end{Shaded}

\begin{Shaded}
\begin{Highlighting}[]
\NormalTok{type(lista\_1)}
\end{Highlighting}
\end{Shaded}

\begin{Shaded}
\begin{Highlighting}[]
\NormalTok{lista\_1 + lista\_2}
\end{Highlighting}
\end{Shaded}

Facilmente podemos fazer operações entre dois arrays de mesmo tamanho,
ou então 1 array e um valor

\begin{Shaded}
\begin{Highlighting}[]
\NormalTok{\# Vai somar cada elemento pelo menos valor}

\NormalTok{lista\_1 + 5}
\end{Highlighting}
\end{Shaded}

pyodide é orientado ao objeto! Existem diversas funções já inseridas no
objeto \texttt{numpy.ndarray}

\begin{Shaded}
\begin{Highlighting}[]
\NormalTok{\# média}
\NormalTok{lista\_1.mean()}
\end{Highlighting}
\end{Shaded}

\begin{Shaded}
\begin{Highlighting}[]
\NormalTok{\# soma}
\NormalTok{lista\_1.sum()}
\end{Highlighting}
\end{Shaded}

\begin{Shaded}
\begin{Highlighting}[]
\NormalTok{\# variancia}
\NormalTok{lista\_1.var()}
\end{Highlighting}
\end{Shaded}

\begin{Shaded}
\begin{Highlighting}[]
\NormalTok{\# desvio padrão}
\NormalTok{lista\_1.std()}
\end{Highlighting}
\end{Shaded}

\begin{Shaded}
\begin{Highlighting}[]
\NormalTok{lista\_1.max()}
\end{Highlighting}
\end{Shaded}

\begin{Shaded}
\begin{Highlighting}[]
\NormalTok{lista\_1.argmax()}
\end{Highlighting}
\end{Shaded}

\begin{Shaded}
\begin{Highlighting}[]
\NormalTok{lista\_1[3]}
\end{Highlighting}
\end{Shaded}

\begin{Shaded}
\begin{Highlighting}[]
\NormalTok{lista\_1.argmin()}
\end{Highlighting}
\end{Shaded}

\begin{Shaded}
\begin{Highlighting}[]
\NormalTok{lista\_1[0]}
\end{Highlighting}
\end{Shaded}

\begin{Shaded}
\begin{Highlighting}[]
\NormalTok{lista\_1.min()}
\end{Highlighting}
\end{Shaded}

\subsection{Agora, como criar um
array?}\label{agora-como-criar-um-array}

\begin{Shaded}
\begin{Highlighting}[]
\NormalTok{\# similar ao range, mas criando arrays}
\NormalTok{np.arange(10)}
\end{Highlighting}
\end{Shaded}

\begin{Shaded}
\begin{Highlighting}[]
\NormalTok{\#inicio, final, passo}

\NormalTok{np.arange(10, 100, 5) \# de 10, até 100, de 5 em 5}
\end{Highlighting}
\end{Shaded}

\begin{Shaded}
\begin{Highlighting}[]
\NormalTok{\#uma outra opção é usar o linspace}

\NormalTok{np.linspace(10,100, 7) \# aqui você define o primeiro, o ultimo e quantos números você quer}
\end{Highlighting}
\end{Shaded}

\begin{Shaded}
\begin{Highlighting}[]
\NormalTok{\# quero criar um conjunto de 0s}
\NormalTok{np.zeros(10)}
\end{Highlighting}
\end{Shaded}

\begin{Shaded}
\begin{Highlighting}[]
\NormalTok{\# quero criar um conjunto de 1s}
\NormalTok{np.ones(10)}
\end{Highlighting}
\end{Shaded}

\begin{Shaded}
\begin{Highlighting}[]
\NormalTok{\# quero repetir algo N vezes}

\NormalTok{np.repeat(100, 7)}
\end{Highlighting}
\end{Shaded}

\begin{Shaded}
\begin{Highlighting}[]
\NormalTok{\# também válido para outras coisas}

\NormalTok{np.repeat(\textquotesingle{}teste\textquotesingle{}, 7)}
\end{Highlighting}
\end{Shaded}

\begin{Shaded}
\begin{Highlighting}[]
\NormalTok{np.repeat(True, 7)}
\end{Highlighting}
\end{Shaded}

\begin{Shaded}
\begin{Highlighting}[]
\NormalTok{\# quero um número inteiro randomico}
\NormalTok{np.random.randint(low=0, high=10, size=5)}
\end{Highlighting}
\end{Shaded}

\begin{Shaded}
\begin{Highlighting}[]
\NormalTok{\# podemos receber números randômicos que seguem uma distribuição }
\NormalTok{np.random.normal(loc=0, scale=1, size=100)}
\end{Highlighting}
\end{Shaded}

\begin{Shaded}
\begin{Highlighting}[]
\NormalTok{\# teste outras distribuições!}
\NormalTok{\# ...}
\end{Highlighting}
\end{Shaded}

\section{Numpy: Matrizes}\label{numpy-matrizes}

Um dos principais tipos de dados na Astronomia -\textgreater{} Imagens
astronômicas Logo, não só de listas vive o astronômo.

\begin{Shaded}
\begin{Highlighting}[]
\NormalTok{print(lista\_1)}
\NormalTok{lista\_1.size}
\end{Highlighting}
\end{Shaded}

\begin{Shaded}
\begin{Highlighting}[]
\NormalTok{lista\_1.shape}
\end{Highlighting}
\end{Shaded}

\begin{Shaded}
\begin{Highlighting}[]
\NormalTok{\# mas se queremos combinar arrays}

\NormalTok{lista\_3 = np.array([lista\_1, lista\_2])}
\end{Highlighting}
\end{Shaded}

\begin{Shaded}
\begin{Highlighting}[]
\NormalTok{print(lista\_3)}

\NormalTok{lista\_3.shape}
\end{Highlighting}
\end{Shaded}

\begin{Shaded}
\begin{Highlighting}[]
\NormalTok{\# podemos criar essas matrizes igual criamos os arrays}

\NormalTok{np.ones([5,7])}
\end{Highlighting}
\end{Shaded}

\begin{Shaded}
\begin{Highlighting}[]
\NormalTok{np.zeros([7,5])}
\end{Highlighting}
\end{Shaded}

\begin{Shaded}
\begin{Highlighting}[]
\NormalTok{\# e podemos combinar dois arrays para criar uma matriz,algo como uma análise combinatório}

\NormalTok{lista\_1\_mx, lista\_2\_mx = np.meshgrid(lista\_1, lista\_2)}
\end{Highlighting}
\end{Shaded}

\begin{Shaded}
\begin{Highlighting}[]
\NormalTok{print(lista\_1\_mx)}
\end{Highlighting}
\end{Shaded}

\begin{Shaded}
\begin{Highlighting}[]
\NormalTok{print(lista\_2\_mx)}
\end{Highlighting}
\end{Shaded}

\begin{Shaded}
\begin{Highlighting}[]
\NormalTok{\# Como sempre, podemos operar diretamente}
\NormalTok{matrix = lista\_1\_mx + lista\_2\_mx}
\end{Highlighting}
\end{Shaded}

\begin{Shaded}
\begin{Highlighting}[]
\NormalTok{print(matrix)}
\end{Highlighting}
\end{Shaded}

\begin{Shaded}
\begin{Highlighting}[]
\NormalTok{\# podemos pegar os elementos da diagonal}
\NormalTok{matrix.diagonal()}
\end{Highlighting}
\end{Shaded}

\textbf{Agora, o que acontece se fazemos uso de algumas funções padrões
do Numpy}

\begin{Shaded}
\begin{Highlighting}[]
\NormalTok{matrix}
\end{Highlighting}
\end{Shaded}

\begin{Shaded}
\begin{Highlighting}[]
\NormalTok{\# Esse é o valor médio considerando TODOS os pontos}
\NormalTok{matrix.mean()}
\end{Highlighting}
\end{Shaded}

\begin{Shaded}
\begin{Highlighting}[]
\NormalTok{\#| collapsed: true}
\NormalTok{\#| jupyter: \{outputs\_hidden: true\}}
\NormalTok{\#| tags: []}
\NormalTok{np.mean?}
\end{Highlighting}
\end{Shaded}

\begin{Shaded}
\begin{Highlighting}[]
\NormalTok{matrix = np.meshgrid(np.arange(7),np.arange(10,15))}
\end{Highlighting}
\end{Shaded}

\begin{Shaded}
\begin{Highlighting}[]
\NormalTok{matrix = matrix[0] + 5*matrix[1]  }
\end{Highlighting}
\end{Shaded}

\begin{Shaded}
\begin{Highlighting}[]
\NormalTok{m3 = np.array([matrix,5*matrix])}
\end{Highlighting}
\end{Shaded}

\begin{Shaded}
\begin{Highlighting}[]
\NormalTok{m3.shape}
\end{Highlighting}
\end{Shaded}

\begin{Shaded}
\begin{Highlighting}[]
\NormalTok{\# porém podemos controlar em cima de que eixo queremos que a conta seja feita}
\NormalTok{m3.mean(axis=2)}
\end{Highlighting}
\end{Shaded}

\begin{Shaded}
\begin{Highlighting}[]
\NormalTok{matrix.mean(axis=1)}
\end{Highlighting}
\end{Shaded}

\begin{Shaded}
\begin{Highlighting}[]
\NormalTok{print(matrix.sum(axis=0))}
\NormalTok{print(matrix.sum(axis=1))}
\end{Highlighting}
\end{Shaded}

\begin{Shaded}
\begin{Highlighting}[]
\NormalTok{\# além disso, podemos alterar o shape da nossa matrix}
\NormalTok{print(matrix.shape, matrix.size)}
\end{Highlighting}
\end{Shaded}

\begin{Shaded}
\begin{Highlighting}[]
\NormalTok{matrix.reshape([35,1])}
\end{Highlighting}
\end{Shaded}

\begin{Shaded}
\begin{Highlighting}[]
\NormalTok{matrix.reshape([1,35])}
\end{Highlighting}
\end{Shaded}

\begin{Shaded}
\begin{Highlighting}[]
\NormalTok{\# Que é a mesma coisa que:}
\NormalTok{matrix.ravel()}
\end{Highlighting}
\end{Shaded}

\begin{Shaded}
\begin{Highlighting}[]
\NormalTok{\# também funciona de forma invertida}

\NormalTok{array = np.arange(0,20)}
\end{Highlighting}
\end{Shaded}

\begin{Shaded}
\begin{Highlighting}[]
\NormalTok{print(array.shape, array.size)}
\end{Highlighting}
\end{Shaded}

\begin{Shaded}
\begin{Highlighting}[]
\NormalTok{\#| error: true}
\NormalTok{array.reshape([2,10])}
\end{Highlighting}
\end{Shaded}

\begin{Shaded}
\begin{Highlighting}[]
\NormalTok{array.reshape([5,4])}
\end{Highlighting}
\end{Shaded}

\begin{Shaded}
\begin{Highlighting}[]
\NormalTok{\#| error: true}
\NormalTok{array.reshape([10,5])}
\end{Highlighting}
\end{Shaded}

\textbf{Importante: Em} \texttt{array.reshape({[}N,\ M{]})} \textbf{N vs
M tem que ser igual ao} \texttt{size}

\section{Numpy: Funções Especiais}\label{numpy-funuxe7uxf5es-especiais}

Além disso, o Numpy tem um conjunto gigantesco de funções que podemos
utilizar

\begin{Shaded}
\begin{Highlighting}[]
\NormalTok{\# Exemplo, queremos tirar a raiz quadrada}
\NormalTok{np.sqrt(4)}
\end{Highlighting}
\end{Shaded}

\begin{Shaded}
\begin{Highlighting}[]
\NormalTok{\# Exemplo, queremos trabalhar com trigonometria}
\NormalTok{angulo = 30}

\NormalTok{print(np.sin(angulo))}
\NormalTok{print(np.cos(angulo))}
\NormalTok{print(np.tan(angulo))}
\end{Highlighting}
\end{Shaded}

\textbf{Importante: o padrão é ângulos serem em radianos!}

\begin{Shaded}
\begin{Highlighting}[]
\NormalTok{np.sin?}
\end{Highlighting}
\end{Shaded}

\begin{Shaded}
\begin{Highlighting}[]
\NormalTok{\# Como saber o ângulo em radianos? Sempre podemos fazer a conta}

\NormalTok{print(np.deg2rad(30))}
\NormalTok{print(np.rad2deg(3.1415))}
\end{Highlighting}
\end{Shaded}

\begin{Shaded}
\begin{Highlighting}[]
\NormalTok{\# Exemplo, queremos trabalhar com trigonometria}
\NormalTok{angulo = np.deg2rad(30)}

\NormalTok{print(np.sin(angulo))}
\NormalTok{print(np.cos(angulo))}
\NormalTok{print(np.tan(angulo))}
\end{Highlighting}
\end{Shaded}

\begin{Shaded}
\begin{Highlighting}[]
\NormalTok{\# Inclusive, temos o valor de pi}
\NormalTok{print(np.pi)}
\NormalTok{print(np.sin(np.pi))}
\NormalTok{print(np.cos(np.pi))}
\end{Highlighting}
\end{Shaded}

\begin{Shaded}
\begin{Highlighting}[]
\NormalTok{\# Também temos os arcos}

\NormalTok{print ( np.arcsin (0.5) )}
\NormalTok{print ( np.arccos (0.5) )}
\end{Highlighting}
\end{Shaded}

\begin{Shaded}
\begin{Highlighting}[]
\NormalTok{np.round(np.pi, 5)}
\end{Highlighting}
\end{Shaded}

\begin{Shaded}
\begin{Highlighting}[]
\NormalTok{print(np.rad2deg(np.arcsin(0.5)))}
\NormalTok{print(np.rad2deg(np.arccos(0.5)))}
\end{Highlighting}
\end{Shaded}

\begin{Shaded}
\begin{Highlighting}[]
\NormalTok{\# Temos também o "e" (número de Euler)}
\NormalTok{np.e}
\end{Highlighting}
\end{Shaded}

\begin{Shaded}
\begin{Highlighting}[]
\NormalTok{np.log(100)}
\end{Highlighting}
\end{Shaded}

\textbf{Novamente, o padrão é logaritimo neperiano}

\begin{Shaded}
\begin{Highlighting}[]
\NormalTok{np.log(np.e**5)}
\end{Highlighting}
\end{Shaded}

\begin{Shaded}
\begin{Highlighting}[]
\NormalTok{\# Mas se eu quero o log de 10?}
\NormalTok{np.log10(10**7)}
\end{Highlighting}
\end{Shaded}

\begin{Shaded}
\begin{Highlighting}[]
\NormalTok{np.log2(2048)}
\end{Highlighting}
\end{Shaded}

\textbf{em vez de usar} \texttt{np.e**valor} \textbf{também podemos
utilizar} \texttt{np.exp()}

\begin{Shaded}
\begin{Highlighting}[]
\NormalTok{print(np.exp(2))}
\NormalTok{print(np.e**2)}
\end{Highlighting}
\end{Shaded}

\textbf{Exercício, escreva em pyodide a seguinte equação:}

\begin{itemize}
\tightlist
\item
  \(mu = 5\)
\item
  \(sigma = 1\)
\item
  \(x = 2\)
\end{itemize}

\section{\texorpdfstring{\$ y = \frac{1}{\sigma.\sqrt{2.\pi}} .
exp\left[ \frac{-(x - \mu)^2}{2\sigma^2}\right]\$}{\$ y =  . exp\$}}\label{y-.-exp}

\begin{Shaded}
\begin{Highlighting}[]
\NormalTok{sigma = 1.}
\NormalTok{mu = 5}
\NormalTok{x = 2.}

\NormalTok{y  = (1 / (sigma*np.sqrt(2*np.pi))) * np.exp({-}((x {-} mu)**2)/(2*(sigma**2)))}

\NormalTok{print(y)}
\end{Highlighting}
\end{Shaded}

\textbf{Vamos fazer um primeiro ajuste!}

Com o numpy podemos fazer de forma rápida um ajuste polinomial usando
\texttt{numpy.polyfit()}

\begin{Shaded}
\begin{Highlighting}[]
\NormalTok{x = np.arange(0, 100)}
\NormalTok{y = 5*x + 2.0}
\end{Highlighting}
\end{Shaded}

\begin{Shaded}
\begin{Highlighting}[]
\NormalTok{\#| collapsed: true}
\NormalTok{\#| jupyter: \{outputs\_hidden: true\}}
\NormalTok{\#| tags: []}
\NormalTok{np.polyfit?}
\end{Highlighting}
\end{Shaded}

\begin{Shaded}
\begin{Highlighting}[]
\NormalTok{y}
\end{Highlighting}
\end{Shaded}

\begin{Shaded}
\begin{Highlighting}[]
\NormalTok{poli = np.polyfit(x, y, deg=1)}
\end{Highlighting}
\end{Shaded}

\begin{Shaded}
\begin{Highlighting}[]
\NormalTok{poli}
\end{Highlighting}
\end{Shaded}

\begin{Shaded}
\begin{Highlighting}[]
\NormalTok{\#| collapsed: true}
\NormalTok{\#| jupyter: \{outputs\_hidden: true\}}
\NormalTok{\#| tags: []}
\NormalTok{np.polyfit?}
\end{Highlighting}
\end{Shaded}

Isso sigifica:

\(y = a[0].x^1 + a[1].x^0\)

\begin{Shaded}
\begin{Highlighting}[]
\NormalTok{\#| collapsed: true}
\NormalTok{\#| jupyter: \{outputs\_hidden: true\}}
\NormalTok{\#| tags: []}
\NormalTok{np.polyfit??}
\end{Highlighting}
\end{Shaded}

\begin{Shaded}
\begin{Highlighting}[]
\NormalTok{np.polyfit(x, y, deg=2)}
\end{Highlighting}
\end{Shaded}

Isso sigifica:

\(y = a[0].x^2 + a[1].x^1 + a[2].x^0\)

\begin{Shaded}
\begin{Highlighting}[]
\NormalTok{np.polyfit(x, y, deg=3)}
\end{Highlighting}
\end{Shaded}

Isso sigifica:

\(y = a[0].x^3 + a[1].x^2 + a[2].x^1 + a[3].x^0\)

\begin{Shaded}
\begin{Highlighting}[]
\NormalTok{\#Podemos também calcular a raiz de uma equação polinomial}

\NormalTok{polinomio = np.array([5., 2.])}
\end{Highlighting}
\end{Shaded}

\begin{Shaded}
\begin{Highlighting}[]
\NormalTok{np.roots?}
\end{Highlighting}
\end{Shaded}

\begin{Shaded}
\begin{Highlighting}[]
\NormalTok{np.roots(poli)}
\end{Highlighting}
\end{Shaded}

\begin{Shaded}
\begin{Highlighting}[]
\NormalTok{\# Válido também para casos mais complexos}

\NormalTok{polinomio = np.array([ 7.14701903e{-}19, {-}1.17901910e{-}16,  5.00000000e+00,  2.00000000e+00])}

\NormalTok{raizes = np.roots(polinomio)}

\NormalTok{print(raizes)}
\end{Highlighting}
\end{Shaded}

\begin{Shaded}
\begin{Highlighting}[]
\NormalTok{\# número complexos!}
\NormalTok{c = raizes[2]}
\NormalTok{print(c)}
\end{Highlighting}
\end{Shaded}

\begin{Shaded}
\begin{Highlighting}[]
\NormalTok{type(c)}
\end{Highlighting}
\end{Shaded}

\begin{Shaded}
\begin{Highlighting}[]
\NormalTok{float(c)}
\end{Highlighting}
\end{Shaded}

\begin{Shaded}
\begin{Highlighting}[]
\NormalTok{print(c.imag)}
\NormalTok{print(c.real)}
\end{Highlighting}
\end{Shaded}

\textbf{Para finalizar, uma coisa importante, como salvar e ler arquivos
usando Numpy}

\begin{Shaded}
\begin{Highlighting}[]
\NormalTok{\# Quero colocar meu x e y em um arquivo!}

\NormalTok{np.savetxt(\textquotesingle{}nome\_arquivo.txt\textquotesingle{}, x)}
\end{Highlighting}
\end{Shaded}

\begin{Shaded}
\begin{Highlighting}[]
\NormalTok{\# Pera, quero colocar o x E o Y}

\NormalTok{np.savetxt(\textquotesingle{}nome\_arquivo.txt\textquotesingle{}, np.array([x, y]))}
\end{Highlighting}
\end{Shaded}

\begin{Shaded}
\begin{Highlighting}[]
\NormalTok{\# Ficou no formato de duas linhas, com mil colunas! Quero invertido.}

\NormalTok{np.savetxt(\textquotesingle{}nome\_arquivo.txt\textquotesingle{}, np.array([x, y]).T)}
\end{Highlighting}
\end{Shaded}

\begin{Shaded}
\begin{Highlighting}[]
\NormalTok{\# Agora quero controlar o formato desses dados}

\NormalTok{np.savetxt?}
\end{Highlighting}
\end{Shaded}

\begin{Shaded}
\begin{Highlighting}[]
\NormalTok{np.savetxt(\textquotesingle{}nome\_arquivo.txt\textquotesingle{}, }
\NormalTok{           np.array([x, y]).T,}
\NormalTok{           fmt=\textquotesingle{}\%.6f\textquotesingle{})}
\end{Highlighting}
\end{Shaded}

\begin{Shaded}
\begin{Highlighting}[]
\NormalTok{\# Ok, sei salvar, como eu leio?}
\NormalTok{a = np.loadtxt(\textquotesingle{}nome\_arquivo.txt\textquotesingle{})}
\end{Highlighting}
\end{Shaded}

\begin{Shaded}
\begin{Highlighting}[]
\NormalTok{a}
\end{Highlighting}
\end{Shaded}

\begin{Shaded}
\begin{Highlighting}[]
\NormalTok{\# Pera, ficou no formato de 2 linhas por muitas colunas, quero invertido}

\NormalTok{a.T}
\end{Highlighting}
\end{Shaded}

\begin{Shaded}
\begin{Highlighting}[]
\NormalTok{a = np.loadtxt(\textquotesingle{}nome\_arquivo.txt\textquotesingle{}, unpack=True)}
\end{Highlighting}
\end{Shaded}

\begin{Shaded}
\begin{Highlighting}[]
\NormalTok{a[1]}
\end{Highlighting}
\end{Shaded}

\begin{Shaded}
\begin{Highlighting}[]
\NormalTok{\# Além, podemos designar variaveis diferentes para cada coluna}

\NormalTok{a, b = np.loadtxt(\textquotesingle{}nome\_arquivo.txt\textquotesingle{}, unpack=True)}
\end{Highlighting}
\end{Shaded}

\begin{Shaded}
\begin{Highlighting}[]
\NormalTok{\#| collapsed: true}
\NormalTok{\#| jupyter: \{outputs\_hidden: true\}}
\NormalTok{\#| tags: []}
\NormalTok{a}
\end{Highlighting}
\end{Shaded}

\begin{Shaded}
\begin{Highlighting}[]
\NormalTok{\#| collapsed: true}
\NormalTok{\#| jupyter: \{outputs\_hidden: true\}}
\NormalTok{\#| tags: []}
\NormalTok{b}
\end{Highlighting}
\end{Shaded}

\begin{Shaded}
\begin{Highlighting}[]
\NormalTok{\#| collapsed: true}
\NormalTok{\#| jupyter: \{outputs\_hidden: true\}}
\NormalTok{\#| tags: []}
\NormalTok{\# Além disso temos uma grande gama de opções}
\NormalTok{np.loadtxt?}
\end{Highlighting}
\end{Shaded}

\begin{Shaded}
\begin{Highlighting}[]
\NormalTok{\# Exemplo, quero usar apenas um das colunas e não todas}

\NormalTok{b = np.loadtxt(\textquotesingle{}nome\_arquivo.txt\textquotesingle{}, unpack=True, usecols=[1])}
\end{Highlighting}
\end{Shaded}

\begin{Shaded}
\begin{Highlighting}[]
\NormalTok{b}
\end{Highlighting}
\end{Shaded}

\begin{Shaded}
\begin{Highlighting}[]
\NormalTok{\# quero pular as 5 primeiras linhas}
\NormalTok{b = np.loadtxt(\textquotesingle{}nome\_arquivo.txt\textquotesingle{}, unpack=True, usecols=[1], skiprows=5)}

\NormalTok{\# Util quando o nosso arquivo de entrada possui um cabeçalho (header)!}
\end{Highlighting}
\end{Shaded}

\begin{Shaded}
\begin{Highlighting}[]
\NormalTok{b}
\end{Highlighting}
\end{Shaded}

\chapter{Isso é tudo}\label{isso-uxe9-tudo}

By BEM

\chapter{Matplotlib}\label{matplotlib}

\section{Gráficos}\label{gruxe1ficos}

\begin{Shaded}
\begin{Highlighting}[]
\NormalTok{\#| autorun: true}
\NormalTok{\# Iniciamos importando o Matplotlib}
\NormalTok{import matplotlib.pyplot as pl}
\NormalTok{import numpy as np}
\end{Highlighting}
\end{Shaded}

\begin{Shaded}
\begin{Highlighting}[]
\NormalTok{\#| autorun: true}
\NormalTok{\# Vamos criar dois arrays para servirem de dados ficticios}
\NormalTok{x = np.arange(0, 100, 10)}
\NormalTok{y = 5*x + 10}
\end{Highlighting}
\end{Shaded}

\begin{Shaded}
\begin{Highlighting}[]
\NormalTok{\#| autorun: true}
\NormalTok{print(x)}
\NormalTok{print(y)}
\end{Highlighting}
\end{Shaded}

\begin{Shaded}
\begin{Highlighting}[]
\NormalTok{\#| autorun: true}
\NormalTok{\# Podemos fazer um plot de maneira imediada}
\NormalTok{pl.plot(x, y)}
\NormalTok{pl.show()}
\end{Highlighting}
\end{Shaded}

\begin{Shaded}
\begin{Highlighting}[]
\CommentTok{\# Existem diferentes argumentos e Kwargs que podemos utilizar para modificar nossos plots}
\NormalTok{pl.plot?}
\end{Highlighting}
\end{Shaded}

\begin{verbatim}
Signature:
pl.plot(
    *args: 'float | ArrayLike | str',
    scalex: 'bool' = True,
    scaley: 'bool' = True,
    data: 'DataParamType' = None,
    **kwargs,
) -> 'list[Line2D]'
Docstring:
Plot y versus x as lines and/or markers.

Call signatures::

    plot([x], y, [fmt], *, data=None, **kwargs)
    plot([x], y, [fmt], [x2], y2, [fmt2], ..., **kwargs)

The coordinates of the points or line nodes are given by *x*, *y*.

The optional parameter *fmt* is a convenient way for defining basic
formatting like color, marker and linestyle. It's a shortcut string
notation described in the *Notes* section below.

>>> plot(x, y)        # plot x and y using default line style and color
>>> plot(x, y, 'bo')  # plot x and y using blue circle markers
>>> plot(y)           # plot y using x as index array 0..N-1
>>> plot(y, 'r+')     # ditto, but with red plusses

You can use `.Line2D` properties as keyword arguments for more
control on the appearance. Line properties and *fmt* can be mixed.
The following two calls yield identical results:

>>> plot(x, y, 'go--', linewidth=2, markersize=12)
>>> plot(x, y, color='green', marker='o', linestyle='dashed',
...      linewidth=2, markersize=12)

When conflicting with *fmt*, keyword arguments take precedence.


**Plotting labelled data**

There's a convenient way for plotting objects with labelled data (i.e.
data that can be accessed by index ``obj['y']``). Instead of giving
the data in *x* and *y*, you can provide the object in the *data*
parameter and just give the labels for *x* and *y*::

>>> plot('xlabel', 'ylabel', data=obj)

All indexable objects are supported. This could e.g. be a `dict`, a
`pandas.DataFrame` or a structured numpy array.


**Plotting multiple sets of data**

There are various ways to plot multiple sets of data.

- The most straight forward way is just to call `plot` multiple times.
  Example:

  >>> plot(x1, y1, 'bo')
  >>> plot(x2, y2, 'go')

- If *x* and/or *y* are 2D arrays, a separate data set will be drawn
  for every column. If both *x* and *y* are 2D, they must have the
  same shape. If only one of them is 2D with shape (N, m) the other
  must have length N and will be used for every data set m.

  Example:

  >>> x = [1, 2, 3]
  >>> y = np.array([[1, 2], [3, 4], [5, 6]])
  >>> plot(x, y)

  is equivalent to:

  >>> for col in range(y.shape[1]):
  ...     plot(x, y[:, col])

- The third way is to specify multiple sets of *[x]*, *y*, *[fmt]*
  groups::

  >>> plot(x1, y1, 'g^', x2, y2, 'g-')

  In this case, any additional keyword argument applies to all
  datasets. Also, this syntax cannot be combined with the *data*
  parameter.

By default, each line is assigned a different style specified by a
'style cycle'. The *fmt* and line property parameters are only
necessary if you want explicit deviations from these defaults.
Alternatively, you can also change the style cycle using
:rc:`axes.prop_cycle`.


Parameters
----------
x, y : array-like or float
    The horizontal / vertical coordinates of the data points.
    *x* values are optional and default to ``range(len(y))``.

    Commonly, these parameters are 1D arrays.

    They can also be scalars, or two-dimensional (in that case, the
    columns represent separate data sets).

    These arguments cannot be passed as keywords.

fmt : str, optional
    A format string, e.g. 'ro' for red circles. See the *Notes*
    section for a full description of the format strings.

    Format strings are just an abbreviation for quickly setting
    basic line properties. All of these and more can also be
    controlled by keyword arguments.

    This argument cannot be passed as keyword.

data : indexable object, optional
    An object with labelled data. If given, provide the label names to
    plot in *x* and *y*.

    .. note::
        Technically there's a slight ambiguity in calls where the
        second label is a valid *fmt*. ``plot('n', 'o', data=obj)``
        could be ``plt(x, y)`` or ``plt(y, fmt)``. In such cases,
        the former interpretation is chosen, but a warning is issued.
        You may suppress the warning by adding an empty format string
        ``plot('n', 'o', '', data=obj)``.

Returns
-------
list of `.Line2D`
    A list of lines representing the plotted data.

Other Parameters
----------------
scalex, scaley : bool, default: True
    These parameters determine if the view limits are adapted to the
    data limits. The values are passed on to
    `~.axes.Axes.autoscale_view`.

**kwargs : `~matplotlib.lines.Line2D` properties, optional
    *kwargs* are used to specify properties like a line label (for
    auto legends), linewidth, antialiasing, marker face color.
    Example::

    >>> plot([1, 2, 3], [1, 2, 3], 'go-', label='line 1', linewidth=2)
    >>> plot([1, 2, 3], [1, 4, 9], 'rs', label='line 2')

    If you specify multiple lines with one plot call, the kwargs apply
    to all those lines. In case the label object is iterable, each
    element is used as labels for each set of data.

    Here is a list of available `.Line2D` properties:

    Properties:
    agg_filter: a filter function, which takes a (m, n, 3) float array and a dpi value, and returns a (m, n, 3) array and two offsets from the bottom left corner of the image
    alpha: float or None
    animated: bool
    antialiased or aa: bool
    clip_box: `~matplotlib.transforms.BboxBase` or None
    clip_on: bool
    clip_path: Patch or (Path, Transform) or None
    color or c: :mpltype:`color`
    dash_capstyle: `.CapStyle` or {'butt', 'projecting', 'round'}
    dash_joinstyle: `.JoinStyle` or {'miter', 'round', 'bevel'}
    dashes: sequence of floats (on/off ink in points) or (None, None)
    data: (2, N) array or two 1D arrays
    drawstyle or ds: {'default', 'steps', 'steps-pre', 'steps-mid', 'steps-post'}, default: 'default'
    figure: `~matplotlib.figure.Figure` or `~matplotlib.figure.SubFigure`
    fillstyle: {'full', 'left', 'right', 'bottom', 'top', 'none'}
    gapcolor: :mpltype:`color` or None
    gid: str
    in_layout: bool
    label: object
    linestyle or ls: {'-', '--', '-.', ':', '', ...} or (offset, on-off-seq)
    linewidth or lw: float
    marker: marker style string, `~.path.Path` or `~.markers.MarkerStyle`
    markeredgecolor or mec: :mpltype:`color`
    markeredgewidth or mew: float
    markerfacecolor or mfc: :mpltype:`color`
    markerfacecoloralt or mfcalt: :mpltype:`color`
    markersize or ms: float
    markevery: None or int or (int, int) or slice or list[int] or float or (float, float) or list[bool]
    mouseover: bool
    path_effects: list of `.AbstractPathEffect`
    picker: float or callable[[Artist, Event], tuple[bool, dict]]
    pickradius: float
    rasterized: bool
    sketch_params: (scale: float, length: float, randomness: float)
    snap: bool or None
    solid_capstyle: `.CapStyle` or {'butt', 'projecting', 'round'}
    solid_joinstyle: `.JoinStyle` or {'miter', 'round', 'bevel'}
    transform: unknown
    url: str
    visible: bool
    xdata: 1D array
    ydata: 1D array
    zorder: float

See Also
--------
scatter : XY scatter plot with markers of varying size and/or color (
    sometimes also called bubble chart).

Notes
-----

.. note::

    This is the :ref:`pyplot wrapper <pyplot_interface>` for `.axes.Axes.plot`.

**Format Strings**

A format string consists of a part for color, marker and line::

    fmt = '[marker][line][color]'

Each of them is optional. If not provided, the value from the style
cycle is used. Exception: If ``line`` is given, but no ``marker``,
the data will be a line without markers.

Other combinations such as ``[color][marker][line]`` are also
supported, but note that their parsing may be ambiguous.

**Markers**

=============   ===============================
character       description
=============   ===============================
``'.'``         point marker
``','``         pixel marker
``'o'``         circle marker
``'v'``         triangle_down marker
``'^'``         triangle_up marker
``'<'``         triangle_left marker
``'>'``         triangle_right marker
``'1'``         tri_down marker
``'2'``         tri_up marker
``'3'``         tri_left marker
``'4'``         tri_right marker
``'8'``         octagon marker
``'s'``         square marker
``'p'``         pentagon marker
``'P'``         plus (filled) marker
``'*'``         star marker
``'h'``         hexagon1 marker
``'H'``         hexagon2 marker
``'+'``         plus marker
``'x'``         x marker
``'X'``         x (filled) marker
``'D'``         diamond marker
``'d'``         thin_diamond marker
``'|'``         vline marker
``'_'``         hline marker
=============   ===============================

**Line Styles**

=============    ===============================
character        description
=============    ===============================
``'-'``          solid line style
``'--'``         dashed line style
``'-.'``         dash-dot line style
``':'``          dotted line style
=============    ===============================

Example format strings::

    'b'    # blue markers with default shape
    'or'   # red circles
    '-g'   # green solid line
    '--'   # dashed line with default color
    '^k:'  # black triangle_up markers connected by a dotted line

**Colors**

The supported color abbreviations are the single letter codes

=============    ===============================
character        color
=============    ===============================
``'b'``          blue
``'g'``          green
``'r'``          red
``'c'``          cyan
``'m'``          magenta
``'y'``          yellow
``'k'``          black
``'w'``          white
=============    ===============================

and the ``'CN'`` colors that index into the default property cycle.

If the color is the only part of the format string, you can
additionally use any  `matplotlib.colors` spec, e.g. full names
(``'green'``) or hex strings (``'#008000'``).
File:      c:\users\luanx\appdata\roaming\python\python312\site-packages\matplotlib\pyplot.py
Type:      function
\end{verbatim}

\begin{Shaded}
\begin{Highlighting}[]
\NormalTok{\#| autorun: true}
\NormalTok{\# Por exemplo}

\NormalTok{pl.figure(figsize=[10, 3]) \# Podemos definir as dimensões das figuras}

\NormalTok{pl.plot(x, y, \textquotesingle{}\^{}{-}\textquotesingle{}, linewidth=3, color=\textquotesingle{}blue\textquotesingle{}) \# as cores de cada plot, largura de linha, simbolos utilizados, etc}
\NormalTok{pl.plot(x, 2*x + 5, \textquotesingle{}\textgreater{}:\textquotesingle{}, lw=0.5, color=\textquotesingle{}indigo\textquotesingle{})}
\NormalTok{pl.plot(x, {-}2*x + 300, \textquotesingle{}\textless{}{-}{-}\textquotesingle{}, color=\textquotesingle{}silver\textquotesingle{})}
\NormalTok{pl.text(50, 300, \textquotesingle{}Teste\textquotesingle{}, fontsize=20, color=\textquotesingle{}blue\textquotesingle{}) \# Podemos plotar texto}
\NormalTok{pl.grid(False)}

\NormalTok{pl.xlim(30, 60) \# Definir os limites dos plots (em x)}
\NormalTok{pl.ylim(100, 400) \# Definir os limites dos plots (em y)}

\NormalTok{pl.xlabel(\textquotesingle{}Valores X\textquotesingle{}, fontsize=15) \# Definir os labels (em x)}
\NormalTok{pl.ylabel(r\textquotesingle{}Valores $\textbackslash{}Delta$ [km]\textquotesingle{}, fontsize=15) \# Definir os labels (em y)}

\NormalTok{pl.xticks(np.array([35, 45, 55]),fontsize=15, rotation={-}60) \# Definir os ticks, (em x) suas posições, valores, e parâmetros}
\NormalTok{pl.yticks(np.array([150, 250, 350]),fontsize=15, rotation=0) \# Definir os ticks, (em x) suas posições, valores, e parâmetros}

\NormalTok{pl.show()}
\end{Highlighting}
\end{Shaded}

\begin{Shaded}
\begin{Highlighting}[]
\NormalTok{\#| autorun: true}
\NormalTok{\# Vamos criar um outro conjunto de dados, desta vez randomicos}
\NormalTok{a = np.random.randint(0, 300, 10)}
\end{Highlighting}
\end{Shaded}

\begin{Shaded}
\begin{Highlighting}[]
\NormalTok{\#| autorun: true}
\NormalTok{pl.bar(x, a, width=8) \# Criando gráficos de barra}
\NormalTok{pl.plot(x, a, color=\textquotesingle{}red\textquotesingle{}, lw=1, marker=\textquotesingle{}*\textquotesingle{})}
\NormalTok{\#pl.xlim(15, 85)}
\NormalTok{\#pl.ylim(200, 300)}
\NormalTok{pl.show()}
\end{Highlighting}
\end{Shaded}

\begin{Shaded}
\begin{Highlighting}[]
\NormalTok{\#| autorun: true}
\NormalTok{pl.barh(x, a, height=8) \# Barras horizontais}
\NormalTok{pl.plot(a, x, color=\textquotesingle{}red\textquotesingle{}, lw=1, marker=\textquotesingle{}*\textquotesingle{})}
\NormalTok{pl.show()}
\end{Highlighting}
\end{Shaded}

\begin{Shaded}
\begin{Highlighting}[]
\NormalTok{\#| autorun: true}
\NormalTok{pl.step(x, a) \# Plots Steps}
\NormalTok{pl.show()}
\end{Highlighting}
\end{Shaded}

\begin{Shaded}
\begin{Highlighting}[]
\NormalTok{pl.stem?}
\end{Highlighting}
\end{Shaded}

\begin{verbatim}
Signature:
pl.stem(
    *args: 'ArrayLike | str',
    linefmt: 'str | None' = None,
    markerfmt: 'str | None' = None,
    basefmt: 'str | None' = None,
    bottom: 'float' = 0,
    label: 'str | None' = None,
    orientation: "Literal['vertical', 'horizontal']" = 'vertical',
    data: 'DataParamType' = None,
) -> 'StemContainer'
Docstring:
Create a stem plot.

A stem plot draws lines perpendicular to a baseline at each location
*locs* from the baseline to *heads*, and places a marker there. For
vertical stem plots (the default), the *locs* are *x* positions, and
the *heads* are *y* values. For horizontal stem plots, the *locs* are
*y* positions, and the *heads* are *x* values.

Call signature::

  stem([locs,] heads, linefmt=None, markerfmt=None, basefmt=None)

The *locs*-positions are optional. *linefmt* may be provided as
positional, but all other formats must be provided as keyword
arguments.

Parameters
----------
locs : array-like, default: (0, 1, ..., len(heads) - 1)
    For vertical stem plots, the x-positions of the stems.
    For horizontal stem plots, the y-positions of the stems.

heads : array-like
    For vertical stem plots, the y-values of the stem heads.
    For horizontal stem plots, the x-values of the stem heads.

linefmt : str, optional
    A string defining the color and/or linestyle of the vertical lines:

    =========  =============
    Character  Line Style
    =========  =============
    ``'-'``    solid line
    ``'--'``   dashed line
    ``'-.'``   dash-dot line
    ``':'``    dotted line
    =========  =============

    Default: 'C0-', i.e. solid line with the first color of the color
    cycle.

    Note: Markers specified through this parameter (e.g. 'x') will be
    silently ignored. Instead, markers should be specified using
    *markerfmt*.

markerfmt : str, optional
    A string defining the color and/or shape of the markers at the stem
    heads. If the marker is not given, use the marker 'o', i.e. filled
    circles. If the color is not given, use the color from *linefmt*.

basefmt : str, default: 'C3-' ('C2-' in classic mode)
    A format string defining the properties of the baseline.

orientation : {'vertical', 'horizontal'}, default: 'vertical'
    The orientation of the stems.

bottom : float, default: 0
    The y/x-position of the baseline (depending on *orientation*).

label : str, optional
    The label to use for the stems in legends.

data : indexable object, optional
    If given, all parameters also accept a string ``s``, which is
    interpreted as ``data[s]`` if ``s`` is a key in ``data``.

Returns
-------
`.StemContainer`
    The container may be treated like a tuple
    (*markerline*, *stemlines*, *baseline*)

Notes
-----

.. note::

    This is the :ref:`pyplot wrapper <pyplot_interface>` for `.axes.Axes.stem`.

.. seealso::
    The MATLAB function
    `stem <https://www.mathworks.com/help/matlab/ref/stem.html>`_
    which inspired this method.
File:      c:\users\luanx\appdata\roaming\python\python312\site-packages\matplotlib\pyplot.py
Type:      function
\end{verbatim}

\begin{Shaded}
\begin{Highlighting}[]
\NormalTok{\#| autorun: true}
\NormalTok{pl.stem(x, a, markerfmt =\textquotesingle{}yo\textquotesingle{}, basefmt=\textquotesingle{}g.\textquotesingle{}, linefmt=\textquotesingle{}blue\textquotesingle{}) \# Plots stem}
\NormalTok{pl.stem(x, y, markerfmt =\textquotesingle{}ro\textquotesingle{}, basefmt=\textquotesingle{}g.\textquotesingle{}, linefmt=\textquotesingle{}blue\textquotesingle{})}
\NormalTok{pl.show()}
\end{Highlighting}
\end{Shaded}

\begin{Shaded}
\begin{Highlighting}[]
\NormalTok{\#| autorun: true}
\NormalTok{\# Criando novos conjuntos de dados}
\NormalTok{y2 = np.random.normal(0, 1, 100)}
\NormalTok{y3 = np.random.normal(0, 1, 100)}
\end{Highlighting}
\end{Shaded}

\begin{Shaded}
\begin{Highlighting}[]
\NormalTok{\#| autorun: true}
\NormalTok{pl.figure(figsize=[15,3])}
\NormalTok{pl.plot(y2, \textquotesingle{}+\textquotesingle{}, color=\textquotesingle{}gray\textquotesingle{}, alpha=1.0) \# Outro parâmetro importante é o alpha, que controla trsnaparencia, 0 = transparente, 1.0 = opaco}
\NormalTok{pl.show()}
\end{Highlighting}
\end{Shaded}

\begin{Shaded}
\begin{Highlighting}[]
\NormalTok{\#| autorun: true}
\NormalTok{\# Outro exemplo de transparência}

\NormalTok{pl.hist(y2, np.arange({-}4, +4.5, 0.5), color=\textquotesingle{}blue\textquotesingle{})}
\NormalTok{pl.hist(y3, np.arange({-}4, +4.5, 0.5), color=\textquotesingle{}red\textquotesingle{}, alpha=0.5) }
\NormalTok{pl.axvline(0, color=\textquotesingle{}k\textquotesingle{})}
\NormalTok{pl.axvline(+1, color=\textquotesingle{}gray\textquotesingle{})}
\NormalTok{pl.axvline({-}1, color=\textquotesingle{}gray\textquotesingle{})}
\NormalTok{pl.show()}
\end{Highlighting}
\end{Shaded}

\begin{Shaded}
\begin{Highlighting}[]
\NormalTok{\#| autorun: true}
\NormalTok{pl.plot(y2)}
\NormalTok{pl.show()}
\end{Highlighting}
\end{Shaded}

\begin{Shaded}
\begin{Highlighting}[]
\NormalTok{\#| autorun: true}
\NormalTok{\# Criando gráficos com barra de erros}
\NormalTok{pl.errorbar(x, y, yerr=a, xerr=2, linestyle=\textquotesingle{}\textquotesingle{}, marker=\textquotesingle{}.\textquotesingle{}, ecolor=\textquotesingle{}red\textquotesingle{}, color=\textquotesingle{}k\textquotesingle{})}
\NormalTok{pl.show()}
\end{Highlighting}
\end{Shaded}

\begin{Shaded}
\begin{Highlighting}[]
\NormalTok{\#| autorun: true}
\NormalTok{y2 = 2*y}
\end{Highlighting}
\end{Shaded}

\begin{Shaded}
\begin{Highlighting}[]
\NormalTok{\#| autorun: true}
\NormalTok{\# Pintando uma região, utilizando o fill\_between}
\NormalTok{pl.plot(x, y, color=\textquotesingle{}k\textquotesingle{})}
\NormalTok{pl.fill\_between(x, a, 2*a, alpha=0.5)}
\NormalTok{pl.show()}
\end{Highlighting}
\end{Shaded}

\begin{Shaded}
\begin{Highlighting}[]
\NormalTok{\#| autorun: true}
\NormalTok{\# Podemos utilizar o label para definir as legendas de cada elemento}
\NormalTok{pl.plot(x, y, color=\textquotesingle{}k\textquotesingle{}, label=\textquotesingle{}Area\textquotesingle{})}
\NormalTok{pl.fill\_between(x, a, 2*a, alpha=0.5, label=\textquotesingle{}Reta\textquotesingle{})}

\NormalTok{pl.legend(fontsize=15, loc=4) \# podemos definir o tamanho e localização da legenda}
\NormalTok{pl.show()}
\end{Highlighting}
\end{Shaded}

\begin{Shaded}
\begin{Highlighting}[]
\NormalTok{\#| autorun: true}
\NormalTok{\# Podemos também criar linhas retas (horizontais e na vertical)}

\NormalTok{pl.plot(x, y, color=\textquotesingle{}k\textquotesingle{}, label=\textquotesingle{}Area\textquotesingle{})}
\NormalTok{pl.fill\_between(x, a, 2*a, alpha=0.5, label=\textquotesingle{}Reta\textquotesingle{})}

\NormalTok{pl.axhline(400, color=\textquotesingle{}red\textquotesingle{}, label=r\textquotesingle{}$\textbackslash{}chi\^{}2$\textquotesingle{}) \# horizontal}
\NormalTok{pl.axvline(60, color=\textquotesingle{}blue\textquotesingle{}, linestyle=\textquotesingle{}{-}{-}\textquotesingle{}, label=r\textquotesingle{}$Bla\^{}2$\textquotesingle{}) \# vertical }

\NormalTok{pl.legend(fontsize=15, loc=4)}
\NormalTok{pl.show()}
\end{Highlighting}
\end{Shaded}

\begin{Shaded}
\begin{Highlighting}[]
\NormalTok{\#| autorun: true}
\NormalTok{\# Aqui podemos, ilustrar diferentes tamanhos, e diferentes cores para cada ponto  }
\NormalTok{pl.scatter(x, y, s=a, c=2*a, cmap=\textquotesingle{}cividis\textquotesingle{})}
\NormalTok{pl.show()}
\end{Highlighting}
\end{Shaded}

\begin{Shaded}
\begin{Highlighting}[]
\NormalTok{\#| autorun: true}
\NormalTok{\# Podemos também criar gráficos de Pizza}
\NormalTok{x2 = np.array([15, 10 , 30])}

\NormalTok{pl.pie(x2)}
\NormalTok{pl.show()}
\end{Highlighting}
\end{Shaded}

\begin{Shaded}
\begin{Highlighting}[]
\NormalTok{\#| autorun: true}
\NormalTok{\# Plotando matrizes (imagens!)}

\NormalTok{x, y = np.meshgrid(np.arange(1000), np.arange(1000))}

\NormalTok{imagem = 5*x + 10*y}
\end{Highlighting}
\end{Shaded}

\begin{Shaded}
\begin{Highlighting}[]
\NormalTok{\#| autorun: true}
\NormalTok{imagem}
\end{Highlighting}
\end{Shaded}

\begin{Shaded}
\begin{Highlighting}[]
\NormalTok{\#| autorun: true}
\NormalTok{\# Ao plotar imagens podemos escolher um mapa de cor e os pontos máximos e mínimos da coloração}

\NormalTok{pl.imshow(imagem, cmap=\textquotesingle{}terrain\textquotesingle{}, vmin=5000, vmax=10000)}
\NormalTok{pl.show()}
\end{Highlighting}
\end{Shaded}

\begin{Shaded}
\begin{Highlighting}[]
\NormalTok{\#| autorun: true}
\NormalTok{\# Podemos adicionar a grade de valores também se desejarmos}
\NormalTok{pl.imshow(imagem, cmap=\textquotesingle{}terrain\textquotesingle{}, vmin=5000, vmax=10000)}
\NormalTok{pl.grid(True)}
\end{Highlighting}
\end{Shaded}

\begin{Shaded}
\begin{Highlighting}[]
\NormalTok{\#| autorun: true}
\NormalTok{\# controlando o cmap}
\NormalTok{pl.imshow(imagem, cmap=\textquotesingle{}gray\textquotesingle{})}
\NormalTok{pl.show()}
\end{Highlighting}
\end{Shaded}

\textbf{Agora podemos combinar os diferentes plots utilizando uma grade
de plots}

\begin{Shaded}
\begin{Highlighting}[]
\NormalTok{\#| autorun: true}
\NormalTok{\# Importando o gridspec}

\NormalTok{import matplotlib.gridspec as gridspec}
\end{Highlighting}
\end{Shaded}

\begin{Shaded}
\begin{Highlighting}[]
\NormalTok{\#| autorun: true}
\NormalTok{fig = pl.figure(figsize=[10,5])}

\NormalTok{gs = gridspec.GridSpec(2, 2) \# Criando a nossa grade}


\NormalTok{\#\#\#\#\#\#\#\#\#\#\#\#\#\#\#\#\#\#\#\#\#\#\#\#\#\#\#\#\#\#\#\#\#\#\#\#\#\#\#\#\#\#\#\#\#\#\#\#\#\#\#\#\#\#\#\#\#\#\#\#\#\#\#\#\#\#\#\#\#\#\#\#\#\#\#\#\#\#}
\NormalTok{\# Criando o primeiro plot na posição [0, 0]}
\NormalTok{ax = fig.add\_subplot(gs[0, 0])}
\NormalTok{pl.imshow(imagem, cmap=\textquotesingle{}terrain\textquotesingle{}, vmin=5000, vmax=10000)}
\NormalTok{pl.xticks([])}

\NormalTok{\#\#\#\#\#\#\#\#\#\#\#\#\#\#\#\#\#\#\#\#\#\#\#\#\#\#\#\#\#\#\#\#\#\#\#\#\#\#\#\#\#\#\#\#\#\#\#\#\#\#\#\#\#\#\#\#\#\#\#\#\#\#\#\#\#\#\#\#\#\#\#\#\#\#\#\#\#\#}
\NormalTok{\# Criando o segundo plot na posição [0, 1]}
\NormalTok{ax = fig.add\_subplot(gs[0, 1])}
\NormalTok{pl.imshow(imagem, cmap=\textquotesingle{}terrain\textquotesingle{}, vmin=5000, vmax=10000)}
\NormalTok{pl.xticks([])}

\NormalTok{\#\#\#\#\#\#\#\#\#\#\#\#\#\#\#\#\#\#\#\#\#\#\#\#\#\#\#\#\#\#\#\#\#\#\#\#\#\#\#\#\#\#\#\#\#\#\#\#\#\#\#\#\#\#\#\#\#\#\#\#\#\#\#\#\#\#\#\#\#\#\#\#\#\#\#\#\#\#}
\NormalTok{\# Criando o terceiro plot na posição [1, 1]}
\NormalTok{ax = fig.add\_subplot(gs[1, 1])}
\NormalTok{pl.yticks([])}
\NormalTok{pl.scatter(np.arange(10), 5*np.arange(10) + 10, s=a, c=2*a, cmap=\textquotesingle{}cividis\textquotesingle{})}
\NormalTok{pl.show()}
\end{Highlighting}
\end{Shaded}

\begin{Shaded}
\begin{Highlighting}[]
\NormalTok{\#| autorun: true}
\NormalTok{np.random.normal()}
\end{Highlighting}
\end{Shaded}

\begin{Shaded}
\begin{Highlighting}[]
\NormalTok{\#| autorun: true}
\NormalTok{\#\#\# Exercício, Faça um gráfico contendo 3 plots, }
\NormalTok{\#\#\# o primeiro sendo um gráfico X vs Y, }
\NormalTok{\#\#\# o segundo, acima do primeiro, contendo o histograma de X, }
\NormalTok{\#\#\# e o terceiro, ao lado direito do primeiro, contendo o histograma de Y.}
\end{Highlighting}
\end{Shaded}

\begin{Shaded}
\begin{Highlighting}[]
\NormalTok{\#| autorun: true}
\NormalTok{x = np.arange(0, 100, 10)}
\NormalTok{y = 5*x + 10}
\end{Highlighting}
\end{Shaded}

\begin{Shaded}
\begin{Highlighting}[]
\NormalTok{\#| autorun: true}
\NormalTok{pl.plot(x, y)}
\NormalTok{pl.show()}
\end{Highlighting}
\end{Shaded}

\begin{Shaded}
\begin{Highlighting}[]
\NormalTok{\#| autorun: true}
\NormalTok{pl.bar(x, y)}
\NormalTok{pl.show()}
\end{Highlighting}
\end{Shaded}

\begin{Shaded}
\begin{Highlighting}[]
\NormalTok{\#| autorun: true}
\NormalTok{pl.barh(x, y)}
\NormalTok{pl.show()}
\end{Highlighting}
\end{Shaded}

\begin{Shaded}
\begin{Highlighting}[]
\NormalTok{\#| autorun: true}
\NormalTok{pl.plot(x)}
\NormalTok{pl.show()}
\end{Highlighting}
\end{Shaded}

\begin{Shaded}
\begin{Highlighting}[]
\NormalTok{\#| autorun: true}
\NormalTok{y2 = np.random.normal(0, 1, 10000)}
\end{Highlighting}
\end{Shaded}

\begin{Shaded}
\begin{Highlighting}[]
\NormalTok{\#| autorun: true}
\NormalTok{pl.hist(y2)}
\NormalTok{pl.show()}
\end{Highlighting}
\end{Shaded}

\begin{Shaded}
\begin{Highlighting}[]
\NormalTok{\#| autorun: true}
\NormalTok{pl.plot(y2)}
\NormalTok{pl.show()}
\end{Highlighting}
\end{Shaded}

\begin{Shaded}
\begin{Highlighting}[]
\NormalTok{\#| autorun: true}
\NormalTok{\# pl.scatter()}
\NormalTok{\# pl.errorbar()}
\NormalTok{\# pl.pie}
\end{Highlighting}
\end{Shaded}

\begin{Shaded}
\begin{Highlighting}[]
\NormalTok{\#| autorun: true}
\NormalTok{x2 = np.array([15, 10 , 30])}
\end{Highlighting}
\end{Shaded}

\begin{Shaded}
\begin{Highlighting}[]
\NormalTok{\#| autorun: true}
\NormalTok{pl.pie(x2)}
\NormalTok{pl.show()}
\end{Highlighting}
\end{Shaded}

\begin{Shaded}
\begin{Highlighting}[]
\NormalTok{\#| autorun: true}
\NormalTok{\# pl.text()}
\end{Highlighting}
\end{Shaded}

\section{Editando os plots}\label{editando-os-plots}

\begin{Shaded}
\begin{Highlighting}[]
\NormalTok{\#| autorun: true}
\NormalTok{\# figsize()}
\NormalTok{\# pl.xlim()}
\NormalTok{\# pl.ylim()}
\NormalTok{\# pl.xticks()}
\NormalTok{\# pl.yticks()}
\NormalTok{\# pl.axis()}
\NormalTok{\# pl.title()}
\NormalTok{\# pl.legend()}
\end{Highlighting}
\end{Shaded}

\chapter{Pandas}\label{pandas}

\begin{Shaded}
\begin{Highlighting}[]
\NormalTok{\#| autorun: true}
\NormalTok{import pandas as pd}
\NormalTok{import matplotlib.pyplot as pl}
\NormalTok{import numpy as np}
\end{Highlighting}
\end{Shaded}

\section{Iniciando o uso do Pandas}\label{iniciando-o-uso-do-pandas}

Vamos começar devagar, para que serve o Pandas?

\begin{Shaded}
\begin{Highlighting}[]
\NormalTok{\#| autorun: true}
\NormalTok{dic = \{\textquotesingle{}anos\textquotesingle{}:np.arange(2019, 2024), }
\NormalTok{       \textquotesingle{}alunos\textquotesingle{}:np.array([5,7,8,9,3]),}
\NormalTok{       \textquotesingle{}aprovados\textquotesingle{}:np.array([2,5,8,7,2])\}}
\end{Highlighting}
\end{Shaded}

\begin{Shaded}
\begin{Highlighting}[]
\NormalTok{\#| autorun: true}
\NormalTok{df = pd.DataFrame(dic)}
\end{Highlighting}
\end{Shaded}

\begin{Shaded}
\begin{Highlighting}[]
\NormalTok{\#| autorun: true}
\NormalTok{df}
\end{Highlighting}
\end{Shaded}

\begin{Shaded}
\begin{Highlighting}[]
\NormalTok{\#| autorun: true}
\NormalTok{df.iloc[2]}
\end{Highlighting}
\end{Shaded}

\begin{Shaded}
\begin{Highlighting}[]
\NormalTok{\#| autorun: true}
\NormalTok{type(df)}
\end{Highlighting}
\end{Shaded}

\begin{Shaded}
\begin{Highlighting}[]
\NormalTok{\#| autorun: true}
\NormalTok{df[\textquotesingle{}anos\textquotesingle{}]}
\end{Highlighting}
\end{Shaded}

\begin{Shaded}
\begin{Highlighting}[]
\NormalTok{\#| autorun: true}
\NormalTok{type(df[\textquotesingle{}alunos\textquotesingle{}])}
\end{Highlighting}
\end{Shaded}

\begin{Shaded}
\begin{Highlighting}[]
\NormalTok{\#| autorun: true}
\NormalTok{df.alunos}
\end{Highlighting}
\end{Shaded}

\begin{Shaded}
\begin{Highlighting}[]
\NormalTok{\#| autorun: true}
\NormalTok{df[\textquotesingle{}reprovados\textquotesingle{}] = df[\textquotesingle{}alunos\textquotesingle{}] {-} df[\textquotesingle{}aprovados\textquotesingle{}]}
\end{Highlighting}
\end{Shaded}

\begin{Shaded}
\begin{Highlighting}[]
\NormalTok{\#| autorun: true}
\NormalTok{df}
\end{Highlighting}
\end{Shaded}

\begin{Shaded}
\begin{Highlighting}[]
\NormalTok{\#| autorun: true}
\NormalTok{df[\textquotesingle{}blabla\textquotesingle{}] = df[\textquotesingle{}anos\textquotesingle{}]*10}
\end{Highlighting}
\end{Shaded}

\begin{Shaded}
\begin{Highlighting}[]
\NormalTok{\#| autorun: true}
\NormalTok{df.blabla}
\end{Highlighting}
\end{Shaded}

\section{Lendo um arquivo de entrada}\label{lendo-um-arquivo-de-entrada}

\begin{Shaded}
\begin{Highlighting}[]
\NormalTok{pd.read\_csv?}
\end{Highlighting}
\end{Shaded}

\begin{verbatim}
Signature:
pd.read_csv(
    filepath_or_buffer: 'FilePath | ReadCsvBuffer[bytes] | ReadCsvBuffer[str]',
    *,
    sep: 'str | None | lib.NoDefault' = <no_default>,
    delimiter: 'str | None | lib.NoDefault' = None,
    header: "int | Sequence[int] | None | Literal['infer']" = 'infer',
    names: 'Sequence[Hashable] | None | lib.NoDefault' = <no_default>,
    index_col: 'IndexLabel | Literal[False] | None' = None,
    usecols: 'UsecolsArgType' = None,
    dtype: 'DtypeArg | None' = None,
    engine: 'CSVEngine | None' = None,
    converters: 'Mapping[Hashable, Callable] | None' = None,
    true_values: 'list | None' = None,
    false_values: 'list | None' = None,
    skipinitialspace: 'bool' = False,
    skiprows: 'list[int] | int | Callable[[Hashable], bool] | None' = None,
    skipfooter: 'int' = 0,
    nrows: 'int | None' = None,
    na_values: 'Hashable | Iterable[Hashable] | Mapping[Hashable, Iterable[Hashable]] | None' = None,
    keep_default_na: 'bool' = True,
    na_filter: 'bool' = True,
    verbose: 'bool | lib.NoDefault' = <no_default>,
    skip_blank_lines: 'bool' = True,
    parse_dates: 'bool | Sequence[Hashable] | None' = None,
    infer_datetime_format: 'bool | lib.NoDefault' = <no_default>,
    keep_date_col: 'bool | lib.NoDefault' = <no_default>,
    date_parser: 'Callable | lib.NoDefault' = <no_default>,
    date_format: 'str | dict[Hashable, str] | None' = None,
    dayfirst: 'bool' = False,
    cache_dates: 'bool' = True,
    iterator: 'bool' = False,
    chunksize: 'int | None' = None,
    compression: 'CompressionOptions' = 'infer',
    thousands: 'str | None' = None,
    decimal: 'str' = '.',
    lineterminator: 'str | None' = None,
    quotechar: 'str' = '"',
    quoting: 'int' = 0,
    doublequote: 'bool' = True,
    escapechar: 'str | None' = None,
    comment: 'str | None' = None,
    encoding: 'str | None' = None,
    encoding_errors: 'str | None' = 'strict',
    dialect: 'str | csv.Dialect | None' = None,
    on_bad_lines: 'str' = 'error',
    delim_whitespace: 'bool | lib.NoDefault' = <no_default>,
    low_memory: 'bool' = True,
    memory_map: 'bool' = False,
    float_precision: "Literal['high', 'legacy'] | None" = None,
    storage_options: 'StorageOptions | None' = None,
    dtype_backend: 'DtypeBackend | lib.NoDefault' = <no_default>,
) -> 'DataFrame | TextFileReader'
Docstring:
Read a comma-separated values (csv) file into DataFrame.

Also supports optionally iterating or breaking of the file
into chunks.

Additional help can be found in the online docs for
`IO Tools <https://pandas.pydata.org/pandas-docs/stable/user_guide/io.html>`_.

Parameters
----------
filepath_or_buffer : str, path object or file-like object
    Any valid string path is acceptable. The string could be a URL. Valid
    URL schemes include http, ftp, s3, gs, and file. For file URLs, a host is
    expected. A local file could be: file://localhost/path/to/table.csv.

    If you want to pass in a path object, pandas accepts any ``os.PathLike``.

    By file-like object, we refer to objects with a ``read()`` method, such as
    a file handle (e.g. via builtin ``open`` function) or ``StringIO``.
sep : str, default ','
    Character or regex pattern to treat as the delimiter. If ``sep=None``, the
    C engine cannot automatically detect
    the separator, but the Python parsing engine can, meaning the latter will
    be used and automatically detect the separator from only the first valid
    row of the file by Python's builtin sniffer tool, ``csv.Sniffer``.
    In addition, separators longer than 1 character and different from
    ``'\s+'`` will be interpreted as regular expressions and will also force
    the use of the Python parsing engine. Note that regex delimiters are prone
    to ignoring quoted data. Regex example: ``'\r\t'``.
delimiter : str, optional
    Alias for ``sep``.
header : int, Sequence of int, 'infer' or None, default 'infer'
    Row number(s) containing column labels and marking the start of the
    data (zero-indexed). Default behavior is to infer the column names: if no ``names``
    are passed the behavior is identical to ``header=0`` and column
    names are inferred from the first line of the file, if column
    names are passed explicitly to ``names`` then the behavior is identical to
    ``header=None``. Explicitly pass ``header=0`` to be able to
    replace existing names. The header can be a list of integers that
    specify row locations for a :class:`~pandas.MultiIndex` on the columns
    e.g. ``[0, 1, 3]``. Intervening rows that are not specified will be
    skipped (e.g. 2 in this example is skipped). Note that this
    parameter ignores commented lines and empty lines if
    ``skip_blank_lines=True``, so ``header=0`` denotes the first line of
    data rather than the first line of the file.
names : Sequence of Hashable, optional
    Sequence of column labels to apply. If the file contains a header row,
    then you should explicitly pass ``header=0`` to override the column names.
    Duplicates in this list are not allowed.
index_col : Hashable, Sequence of Hashable or False, optional
  Column(s) to use as row label(s), denoted either by column labels or column
  indices.  If a sequence of labels or indices is given, :class:`~pandas.MultiIndex`
  will be formed for the row labels.

  Note: ``index_col=False`` can be used to force pandas to *not* use the first
  column as the index, e.g., when you have a malformed file with delimiters at
  the end of each line.
usecols : Sequence of Hashable or Callable, optional
    Subset of columns to select, denoted either by column labels or column indices.
    If list-like, all elements must either
    be positional (i.e. integer indices into the document columns) or strings
    that correspond to column names provided either by the user in ``names`` or
    inferred from the document header row(s). If ``names`` are given, the document
    header row(s) are not taken into account. For example, a valid list-like
    ``usecols`` parameter would be ``[0, 1, 2]`` or ``['foo', 'bar', 'baz']``.
    Element order is ignored, so ``usecols=[0, 1]`` is the same as ``[1, 0]``.
    To instantiate a :class:`~pandas.DataFrame` from ``data`` with element order
    preserved use ``pd.read_csv(data, usecols=['foo', 'bar'])[['foo', 'bar']]``
    for columns in ``['foo', 'bar']`` order or
    ``pd.read_csv(data, usecols=['foo', 'bar'])[['bar', 'foo']]``
    for ``['bar', 'foo']`` order.

    If callable, the callable function will be evaluated against the column
    names, returning names where the callable function evaluates to ``True``. An
    example of a valid callable argument would be ``lambda x: x.upper() in
    ['AAA', 'BBB', 'DDD']``. Using this parameter results in much faster
    parsing time and lower memory usage.
dtype : dtype or dict of {Hashable : dtype}, optional
    Data type(s) to apply to either the whole dataset or individual columns.
    E.g., ``{'a': np.float64, 'b': np.int32, 'c': 'Int64'}``
    Use ``str`` or ``object`` together with suitable ``na_values`` settings
    to preserve and not interpret ``dtype``.
    If ``converters`` are specified, they will be applied INSTEAD
    of ``dtype`` conversion.

    .. versionadded:: 1.5.0

        Support for ``defaultdict`` was added. Specify a ``defaultdict`` as input where
        the default determines the ``dtype`` of the columns which are not explicitly
        listed.
engine : {'c', 'python', 'pyarrow'}, optional
    Parser engine to use. The C and pyarrow engines are faster, while the python engine
    is currently more feature-complete. Multithreading is currently only supported by
    the pyarrow engine.

    .. versionadded:: 1.4.0

        The 'pyarrow' engine was added as an *experimental* engine, and some features
        are unsupported, or may not work correctly, with this engine.
converters : dict of {Hashable : Callable}, optional
    Functions for converting values in specified columns. Keys can either
    be column labels or column indices.
true_values : list, optional
    Values to consider as ``True`` in addition to case-insensitive variants of 'True'.
false_values : list, optional
    Values to consider as ``False`` in addition to case-insensitive variants of 'False'.
skipinitialspace : bool, default False
    Skip spaces after delimiter.
skiprows : int, list of int or Callable, optional
    Line numbers to skip (0-indexed) or number of lines to skip (``int``)
    at the start of the file.

    If callable, the callable function will be evaluated against the row
    indices, returning ``True`` if the row should be skipped and ``False`` otherwise.
    An example of a valid callable argument would be ``lambda x: x in [0, 2]``.
skipfooter : int, default 0
    Number of lines at bottom of file to skip (Unsupported with ``engine='c'``).
nrows : int, optional
    Number of rows of file to read. Useful for reading pieces of large files.
na_values : Hashable, Iterable of Hashable or dict of {Hashable : Iterable}, optional
    Additional strings to recognize as ``NA``/``NaN``. If ``dict`` passed, specific
    per-column ``NA`` values.  By default the following values are interpreted as
    ``NaN``: " ", "#N/A", "#N/A N/A", "#NA", "-1.#IND", "-1.#QNAN", "-NaN", "-nan",
    "1.#IND", "1.#QNAN", "<NA>", "N/A", "NA", "NULL", "NaN", "None",
    "n/a", "nan", "null ".

keep_default_na : bool, default True
    Whether or not to include the default ``NaN`` values when parsing the data.
    Depending on whether ``na_values`` is passed in, the behavior is as follows:

    * If ``keep_default_na`` is ``True``, and ``na_values`` are specified, ``na_values``
      is appended to the default ``NaN`` values used for parsing.
    * If ``keep_default_na`` is ``True``, and ``na_values`` are not specified, only
      the default ``NaN`` values are used for parsing.
    * If ``keep_default_na`` is ``False``, and ``na_values`` are specified, only
      the ``NaN`` values specified ``na_values`` are used for parsing.
    * If ``keep_default_na`` is ``False``, and ``na_values`` are not specified, no
      strings will be parsed as ``NaN``.

    Note that if ``na_filter`` is passed in as ``False``, the ``keep_default_na`` and
    ``na_values`` parameters will be ignored.
na_filter : bool, default True
    Detect missing value markers (empty strings and the value of ``na_values``). In
    data without any ``NA`` values, passing ``na_filter=False`` can improve the
    performance of reading a large file.
verbose : bool, default False
    Indicate number of ``NA`` values placed in non-numeric columns.

    .. deprecated:: 2.2.0
skip_blank_lines : bool, default True
    If ``True``, skip over blank lines rather than interpreting as ``NaN`` values.
parse_dates : bool, list of Hashable, list of lists or dict of {Hashable : list}, default False
    The behavior is as follows:

    * ``bool``. If ``True`` -> try parsing the index. Note: Automatically set to
      ``True`` if ``date_format`` or ``date_parser`` arguments have been passed.
    * ``list`` of ``int`` or names. e.g. If ``[1, 2, 3]`` -> try parsing columns 1, 2, 3
      each as a separate date column.
    * ``list`` of ``list``. e.g.  If ``[[1, 3]]`` -> combine columns 1 and 3 and parse
      as a single date column. Values are joined with a space before parsing.
    * ``dict``, e.g. ``{'foo' : [1, 3]}`` -> parse columns 1, 3 as date and call
      result 'foo'. Values are joined with a space before parsing.

    If a column or index cannot be represented as an array of ``datetime``,
    say because of an unparsable value or a mixture of timezones, the column
    or index will be returned unaltered as an ``object`` data type. For
    non-standard ``datetime`` parsing, use :func:`~pandas.to_datetime` after
    :func:`~pandas.read_csv`.

    Note: A fast-path exists for iso8601-formatted dates.
infer_datetime_format : bool, default False
    If ``True`` and ``parse_dates`` is enabled, pandas will attempt to infer the
    format of the ``datetime`` strings in the columns, and if it can be inferred,
    switch to a faster method of parsing them. In some cases this can increase
    the parsing speed by 5-10x.

    .. deprecated:: 2.0.0
        A strict version of this argument is now the default, passing it has no effect.

keep_date_col : bool, default False
    If ``True`` and ``parse_dates`` specifies combining multiple columns then
    keep the original columns.
date_parser : Callable, optional
    Function to use for converting a sequence of string columns to an array of
    ``datetime`` instances. The default uses ``dateutil.parser.parser`` to do the
    conversion. pandas will try to call ``date_parser`` in three different ways,
    advancing to the next if an exception occurs: 1) Pass one or more arrays
    (as defined by ``parse_dates``) as arguments; 2) concatenate (row-wise) the
    string values from the columns defined by ``parse_dates`` into a single array
    and pass that; and 3) call ``date_parser`` once for each row using one or
    more strings (corresponding to the columns defined by ``parse_dates``) as
    arguments.

    .. deprecated:: 2.0.0
       Use ``date_format`` instead, or read in as ``object`` and then apply
       :func:`~pandas.to_datetime` as-needed.
date_format : str or dict of column -> format, optional
    Format to use for parsing dates when used in conjunction with ``parse_dates``.
    The strftime to parse time, e.g. :const:`"%d/%m/%Y"`. See
    `strftime documentation
    <https://docs.python.org/3/library/datetime.html
    #strftime-and-strptime-behavior>`_ for more information on choices, though
    note that :const:`"%f"` will parse all the way up to nanoseconds.
    You can also pass:

    - "ISO8601", to parse any `ISO8601 <https://en.wikipedia.org/wiki/ISO_8601>`_
        time string (not necessarily in exactly the same format);
    - "mixed", to infer the format for each element individually. This is risky,
        and you should probably use it along with `dayfirst`.

    .. versionadded:: 2.0.0
dayfirst : bool, default False
    DD/MM format dates, international and European format.
cache_dates : bool, default True
    If ``True``, use a cache of unique, converted dates to apply the ``datetime``
    conversion. May produce significant speed-up when parsing duplicate
    date strings, especially ones with timezone offsets.

iterator : bool, default False
    Return ``TextFileReader`` object for iteration or getting chunks with
    ``get_chunk()``.
chunksize : int, optional
    Number of lines to read from the file per chunk. Passing a value will cause the
    function to return a ``TextFileReader`` object for iteration.
    See the `IO Tools docs
    <https://pandas.pydata.org/pandas-docs/stable/io.html#io-chunking>`_
    for more information on ``iterator`` and ``chunksize``.

compression : str or dict, default 'infer'
    For on-the-fly decompression of on-disk data. If 'infer' and 'filepath_or_buffer' is
    path-like, then detect compression from the following extensions: '.gz',
    '.bz2', '.zip', '.xz', '.zst', '.tar', '.tar.gz', '.tar.xz' or '.tar.bz2'
    (otherwise no compression).
    If using 'zip' or 'tar', the ZIP file must contain only one data file to be read in.
    Set to ``None`` for no decompression.
    Can also be a dict with key ``'method'`` set
    to one of {``'zip'``, ``'gzip'``, ``'bz2'``, ``'zstd'``, ``'xz'``, ``'tar'``} and
    other key-value pairs are forwarded to
    ``zipfile.ZipFile``, ``gzip.GzipFile``,
    ``bz2.BZ2File``, ``zstandard.ZstdDecompressor``, ``lzma.LZMAFile`` or
    ``tarfile.TarFile``, respectively.
    As an example, the following could be passed for Zstandard decompression using a
    custom compression dictionary:
    ``compression={'method': 'zstd', 'dict_data': my_compression_dict}``.

    .. versionadded:: 1.5.0
        Added support for `.tar` files.

    .. versionchanged:: 1.4.0 Zstandard support.

thousands : str (length 1), optional
    Character acting as the thousands separator in numerical values.
decimal : str (length 1), default '.'
    Character to recognize as decimal point (e.g., use ',' for European data).
lineterminator : str (length 1), optional
    Character used to denote a line break. Only valid with C parser.
quotechar : str (length 1), optional
    Character used to denote the start and end of a quoted item. Quoted
    items can include the ``delimiter`` and it will be ignored.
quoting : {0 or csv.QUOTE_MINIMAL, 1 or csv.QUOTE_ALL, 2 or csv.QUOTE_NONNUMERIC, 3 or csv.QUOTE_NONE}, default csv.QUOTE_MINIMAL
    Control field quoting behavior per ``csv.QUOTE_*`` constants. Default is
    ``csv.QUOTE_MINIMAL`` (i.e., 0) which implies that only fields containing special
    characters are quoted (e.g., characters defined in ``quotechar``, ``delimiter``,
    or ``lineterminator``.
doublequote : bool, default True
   When ``quotechar`` is specified and ``quoting`` is not ``QUOTE_NONE``, indicate
   whether or not to interpret two consecutive ``quotechar`` elements INSIDE a
   field as a single ``quotechar`` element.
escapechar : str (length 1), optional
    Character used to escape other characters.
comment : str (length 1), optional
    Character indicating that the remainder of line should not be parsed.
    If found at the beginning
    of a line, the line will be ignored altogether. This parameter must be a
    single character. Like empty lines (as long as ``skip_blank_lines=True``),
    fully commented lines are ignored by the parameter ``header`` but not by
    ``skiprows``. For example, if ``comment='#'``, parsing
    ``#empty\na,b,c\n1,2,3`` with ``header=0`` will result in ``'a,b,c'`` being
    treated as the header.
encoding : str, optional, default 'utf-8'
    Encoding to use for UTF when reading/writing (ex. ``'utf-8'``). `List of Python
    standard encodings
    <https://docs.python.org/3/library/codecs.html#standard-encodings>`_ .

encoding_errors : str, optional, default 'strict'
    How encoding errors are treated. `List of possible values
    <https://docs.python.org/3/library/codecs.html#error-handlers>`_ .

    .. versionadded:: 1.3.0

dialect : str or csv.Dialect, optional
    If provided, this parameter will override values (default or not) for the
    following parameters: ``delimiter``, ``doublequote``, ``escapechar``,
    ``skipinitialspace``, ``quotechar``, and ``quoting``. If it is necessary to
    override values, a ``ParserWarning`` will be issued. See ``csv.Dialect``
    documentation for more details.
on_bad_lines : {'error', 'warn', 'skip'} or Callable, default 'error'
    Specifies what to do upon encountering a bad line (a line with too many fields).
    Allowed values are :

    - ``'error'``, raise an Exception when a bad line is encountered.
    - ``'warn'``, raise a warning when a bad line is encountered and skip that line.
    - ``'skip'``, skip bad lines without raising or warning when they are encountered.

    .. versionadded:: 1.3.0

    .. versionadded:: 1.4.0

        - Callable, function with signature
          ``(bad_line: list[str]) -> list[str] | None`` that will process a single
          bad line. ``bad_line`` is a list of strings split by the ``sep``.
          If the function returns ``None``, the bad line will be ignored.
          If the function returns a new ``list`` of strings with more elements than
          expected, a ``ParserWarning`` will be emitted while dropping extra elements.
          Only supported when ``engine='python'``

    .. versionchanged:: 2.2.0

        - Callable, function with signature
          as described in `pyarrow documentation
          <https://arrow.apache.org/docs/python/generated/pyarrow.csv.ParseOptions.html
          #pyarrow.csv.ParseOptions.invalid_row_handler>`_ when ``engine='pyarrow'``

delim_whitespace : bool, default False
    Specifies whether or not whitespace (e.g. ``' '`` or ``'\t'``) will be
    used as the ``sep`` delimiter. Equivalent to setting ``sep='\s+'``. If this option
    is set to ``True``, nothing should be passed in for the ``delimiter``
    parameter.

    .. deprecated:: 2.2.0
        Use ``sep="\s+"`` instead.
low_memory : bool, default True
    Internally process the file in chunks, resulting in lower memory use
    while parsing, but possibly mixed type inference.  To ensure no mixed
    types either set ``False``, or specify the type with the ``dtype`` parameter.
    Note that the entire file is read into a single :class:`~pandas.DataFrame`
    regardless, use the ``chunksize`` or ``iterator`` parameter to return the data in
    chunks. (Only valid with C parser).
memory_map : bool, default False
    If a filepath is provided for ``filepath_or_buffer``, map the file object
    directly onto memory and access the data directly from there. Using this
    option can improve performance because there is no longer any I/O overhead.
float_precision : {'high', 'legacy', 'round_trip'}, optional
    Specifies which converter the C engine should use for floating-point
    values. The options are ``None`` or ``'high'`` for the ordinary converter,
    ``'legacy'`` for the original lower precision pandas converter, and
    ``'round_trip'`` for the round-trip converter.

storage_options : dict, optional
    Extra options that make sense for a particular storage connection, e.g.
    host, port, username, password, etc. For HTTP(S) URLs the key-value pairs
    are forwarded to ``urllib.request.Request`` as header options. For other
    URLs (e.g. starting with "s3://", and "gcs://") the key-value pairs are
    forwarded to ``fsspec.open``. Please see ``fsspec`` and ``urllib`` for more
    details, and for more examples on storage options refer `here
    <https://pandas.pydata.org/docs/user_guide/io.html?
    highlight=storage_options#reading-writing-remote-files>`_.

dtype_backend : {'numpy_nullable', 'pyarrow'}, default 'numpy_nullable'
    Back-end data type applied to the resultant :class:`DataFrame`
    (still experimental). Behaviour is as follows:

    * ``"numpy_nullable"``: returns nullable-dtype-backed :class:`DataFrame`
      (default).
    * ``"pyarrow"``: returns pyarrow-backed nullable :class:`ArrowDtype`
      DataFrame.

    .. versionadded:: 2.0

Returns
-------
DataFrame or TextFileReader
    A comma-separated values (csv) file is returned as two-dimensional
    data structure with labeled axes.

See Also
--------
DataFrame.to_csv : Write DataFrame to a comma-separated values (csv) file.
read_table : Read general delimited file into DataFrame.
read_fwf : Read a table of fixed-width formatted lines into DataFrame.

Examples
--------
>>> pd.read_csv('data.csv')  # doctest: +SKIP
File:      c:\python312\lib\site-packages\pandas\io\parsers\readers.py
Type:      function
\end{verbatim}

\section{Vamos começar lendo um
catálogo}\label{vamos-comeuxe7ar-lendo-um-catuxe1logo}

\begin{Shaded}
\begin{Highlighting}[]
\NormalTok{\#| autorun: true}
\NormalTok{df = pd.read\_csv(\textquotesingle{}nomad.csv\textquotesingle{}, sep=\textquotesingle{};\textquotesingle{}, comment=\textquotesingle{}\#\textquotesingle{}, }
\NormalTok{                 na\_values=\textquotesingle{}      \textquotesingle{})}
\end{Highlighting}
\end{Shaded}

\begin{Shaded}
\begin{Highlighting}[]
\NormalTok{\#| autorun: true}
\NormalTok{type(df)}
\end{Highlighting}
\end{Shaded}

\begin{Shaded}
\begin{Highlighting}[]
\NormalTok{\#| autorun: true}
\NormalTok{\#| tags: []}
\NormalTok{df}
\end{Highlighting}
\end{Shaded}

\begin{Shaded}
\begin{Highlighting}[]
\NormalTok{\#| autorun: true}
\NormalTok{df.dtypes}
\end{Highlighting}
\end{Shaded}

\begin{Shaded}
\begin{Highlighting}[]
\NormalTok{\#| autorun: true}
\NormalTok{df[\textquotesingle{}RAJ2000\textquotesingle{}]}
\end{Highlighting}
\end{Shaded}

\begin{Shaded}
\begin{Highlighting}[]
\NormalTok{\#| autorun: true}
\NormalTok{pl.plot(df.RAJ2000, df.DEJ2000, \textquotesingle{}r\^{}\textquotesingle{})}
\end{Highlighting}
\end{Shaded}

\begin{Shaded}
\begin{Highlighting}[]
\NormalTok{\#| autorun: true}
\NormalTok{df[\textquotesingle{}Bmag\textquotesingle{}]}
\end{Highlighting}
\end{Shaded}

\begin{Shaded}
\begin{Highlighting}[]
\NormalTok{\#| autorun: true}
\NormalTok{df}
\end{Highlighting}
\end{Shaded}

\begin{Shaded}
\begin{Highlighting}[]
\NormalTok{\#| autorun: true}
\NormalTok{\# Podemos aplicar filtros}
\NormalTok{df[df[\textquotesingle{}Bmag\textquotesingle{}] \textless{} 17]}
\end{Highlighting}
\end{Shaded}

\begin{Shaded}
\begin{Highlighting}[]
\NormalTok{\#| autorun: true}
\NormalTok{df[\textquotesingle{}Bmag\textquotesingle{}][\textasciitilde{}np.isnan(df[\textquotesingle{}Bmag\textquotesingle{}])]}
\end{Highlighting}
\end{Shaded}

\begin{Shaded}
\begin{Highlighting}[]
\NormalTok{\#| autorun: true}
\NormalTok{pl.plot(df.Bmag {-} df.Vmag, df.Vmag {-} df.Rmag, \textquotesingle{}k.\textquotesingle{})}
\NormalTok{pl.axis(\textquotesingle{}square\textquotesingle{})}
\NormalTok{pl.ylabel(\textquotesingle{}V {-} R\textquotesingle{},fontsize=12)}
\NormalTok{pl.xlabel(\textquotesingle{}B {-} V\textquotesingle{},fontsize=12)}
\NormalTok{pl.xlim({-}12, +12)}
\NormalTok{pl.ylim({-}12, +12)}

\NormalTok{pl.show()}
\end{Highlighting}
\end{Shaded}

\begin{Shaded}
\begin{Highlighting}[]
\NormalTok{\#| autorun: true}
\NormalTok{\#Podemos criar novas colunas de maneira imediata}

\NormalTok{print(df.shape)}
\NormalTok{df[\textquotesingle{}B{-}V\textquotesingle{}] = df.Bmag {-} df.Vmag}

\NormalTok{print(df.shape)}
\end{Highlighting}
\end{Shaded}

\begin{Shaded}
\begin{Highlighting}[]
\NormalTok{\#| autorun: true}
\NormalTok{df}
\end{Highlighting}
\end{Shaded}

\begin{Shaded}
\begin{Highlighting}[]
\NormalTok{\#| autorun: true}
\NormalTok{\# Podemos filtrar todas as colunas também de maneira simples}
\NormalTok{df2 = df[df.Vmag \textless{} 15]}
\end{Highlighting}
\end{Shaded}

\begin{Shaded}
\begin{Highlighting}[]
\NormalTok{\#| autorun: true}
\NormalTok{df2}
\end{Highlighting}
\end{Shaded}

\section{Vamos salvar esse novo
``catálogo''}\label{vamos-salvar-esse-novo-catuxe1logo}

\begin{Shaded}
\begin{Highlighting}[]
\NormalTok{\#| autorun: true}
\NormalTok{df2.to\_latex(\textquotesingle{}output.tex\textquotesingle{})}
\end{Highlighting}
\end{Shaded}

\begin{Shaded}
\begin{Highlighting}[]
\CommentTok{\# Podemos salvar o nosso arquivo de saída em diferentes formatos}

\CommentTok{\#df.to\_ \#tab}
\NormalTok{df.to\_csv?}
\end{Highlighting}
\end{Shaded}

\begin{verbatim}
Object `df.to_csv` not found.
\end{verbatim}

\begin{Shaded}
\begin{Highlighting}[]
\NormalTok{\#| autorun: true}
\NormalTok{df2.to\_csv(\textquotesingle{}output.csv\textquotesingle{})}
\end{Highlighting}
\end{Shaded}

\section{Agora vamos às análises}\label{agora-vamos-uxe0s-anuxe1lises}

\begin{itemize}
\tightlist
\item
  Análise geral e Filtros e Sort
\item
  Relações e Correlações entre os elementos
\item
  Agrupamentos
\end{itemize}

\begin{Shaded}
\begin{Highlighting}[]
\NormalTok{\#| autorun: true}
\NormalTok{\# Informações gerais sobre o seu catálogo / data frame}

\NormalTok{df5 = df.describe()}
\end{Highlighting}
\end{Shaded}

\begin{Shaded}
\begin{Highlighting}[]
\NormalTok{\#| autorun: true}

\NormalTok{np.mean(df[\textquotesingle{}pmRA\textquotesingle{}])}
\end{Highlighting}
\end{Shaded}

\begin{Shaded}
\begin{Highlighting}[]
\NormalTok{\#| autorun: true}
\NormalTok{\# Filtrando os valores}
\NormalTok{df[df.Vmag \textless{} 12]}
\end{Highlighting}
\end{Shaded}

\begin{Shaded}
\begin{Highlighting}[]
\NormalTok{\#| autorun: true}
\NormalTok{df.sort\_values(by=\textquotesingle{}Vmag\textquotesingle{})}
\end{Highlighting}
\end{Shaded}

\begin{Shaded}
\begin{Highlighting}[]
\NormalTok{\#| autorun: true}
\NormalTok{df.sort\_values(by=\textquotesingle{}Vmag\textquotesingle{},ascending=False)}
\end{Highlighting}
\end{Shaded}

\begin{Shaded}
\begin{Highlighting}[]
\NormalTok{\#| autorun: true}
\NormalTok{\# Qual a variancia e covariancia entre os parâmetros ?}

\NormalTok{df.cov(numeric\_only=True)}
\end{Highlighting}
\end{Shaded}

\begin{Shaded}
\begin{Highlighting}[]
\NormalTok{\#| autorun: true}
\NormalTok{\# Quais as correlações entre os parâmetros ?}

\NormalTok{df.corr(numeric\_only=True)}
\end{Highlighting}
\end{Shaded}

\begin{Shaded}
\begin{Highlighting}[]
\NormalTok{\#| autorun: true}
\NormalTok{df.hist(color="k", alpha=1.0, bins=10, figsize=[10,10])}
\NormalTok{pl.show()}
\end{Highlighting}
\end{Shaded}

\begin{Shaded}
\begin{Highlighting}[]
\NormalTok{\#| autorun: true}
\NormalTok{fig, ax = pl.subplots()}
\NormalTok{df.plot.scatter(x="Bmag", y="Kmag", ax=ax)}
\NormalTok{pl.show()}
\end{Highlighting}
\end{Shaded}

\begin{Shaded}
\begin{Highlighting}[]
\NormalTok{\#| autorun: true}
\NormalTok{fig, ax = pl.subplots()}
\NormalTok{df.plot.scatter(x="Bmag", }
\NormalTok{                y="Kmag", }
\NormalTok{                c=\textquotesingle{}Vmag\textquotesingle{},}
\NormalTok{                s=50, }
\NormalTok{                cmap=\textquotesingle{}terrain\textquotesingle{},}
\NormalTok{                ax=ax)}
\NormalTok{pl.show()}
\end{Highlighting}
\end{Shaded}

\begin{Shaded}
\begin{Highlighting}[]
\NormalTok{\#| autorun: true}
\NormalTok{\# listando as keys das colunas}
\NormalTok{df.keys()}
\end{Highlighting}
\end{Shaded}

\begin{Shaded}
\begin{Highlighting}[]
\NormalTok{\#| autorun: true}
\NormalTok{df10 = df[df[\textquotesingle{}Bmag\textquotesingle{}] \textless{} 15]}
\end{Highlighting}
\end{Shaded}

\begin{Shaded}
\begin{Highlighting}[]
\NormalTok{\#| autorun: true}
\NormalTok{df10}
\end{Highlighting}
\end{Shaded}

\begin{Shaded}
\begin{Highlighting}[]
\NormalTok{df.drop?}
\end{Highlighting}
\end{Shaded}

\begin{verbatim}
Object `df.drop` not found.
\end{verbatim}

\begin{Shaded}
\begin{Highlighting}[]
\NormalTok{\#| autorun: true}
\NormalTok{\# Vamos criar um novo DataFrame com menos colunas}

\NormalTok{df2 = df.drop(columns=[\textquotesingle{}\_RAJ2000\textquotesingle{}, \textquotesingle{}\_DEJ2000\textquotesingle{}, }
\NormalTok{                       \textquotesingle{}NOMAD1\textquotesingle{}, \textquotesingle{}YM\textquotesingle{}, \textquotesingle{}Bmag\textquotesingle{}, }
\NormalTok{                       \textquotesingle{}r\_Bmag\textquotesingle{}, \textquotesingle{}Vmag\textquotesingle{}, \textquotesingle{}r\_Vmag\textquotesingle{},}
\NormalTok{                       \textquotesingle{}Rmag\textquotesingle{}, \textquotesingle{}r\_Rmag\textquotesingle{}, \textquotesingle{}Jmag\textquotesingle{}, }
\NormalTok{                       \textquotesingle{}Hmag\textquotesingle{}, \textquotesingle{}Kmag\textquotesingle{}, \textquotesingle{}R\textquotesingle{}, }
\NormalTok{                       \textquotesingle{}B{-}V\textquotesingle{}, \textquotesingle{}r\textquotesingle{}, \textquotesingle{}e\_pmRA\textquotesingle{}, \textquotesingle{}e\_pmDE\textquotesingle{}])}
\end{Highlighting}
\end{Shaded}

\begin{Shaded}
\begin{Highlighting}[]
\NormalTok{\#| autorun: true}
\NormalTok{df2 = df[[\textquotesingle{}RAJ2000\textquotesingle{}, \textquotesingle{}DEJ2000\textquotesingle{}, \textquotesingle{}pmRA\textquotesingle{}, \textquotesingle{}pmDE\textquotesingle{}]]}
\end{Highlighting}
\end{Shaded}

\begin{Shaded}
\begin{Highlighting}[]
\NormalTok{\#| autorun: true}
\NormalTok{df2}
\end{Highlighting}
\end{Shaded}

\begin{Shaded}
\begin{Highlighting}[]
\NormalTok{\#| autorun: true}
\NormalTok{from pandas.plotting import scatter\_matrix}
\end{Highlighting}
\end{Shaded}

\begin{Shaded}
\begin{Highlighting}[]
\NormalTok{\#| autorun: true}
\NormalTok{scatter\_matrix(df2,  diagonal=\textquotesingle{}hist\textquotesingle{},alpha=0.5, figsize=(5, 5))}
\NormalTok{pl.show()}
\end{Highlighting}
\end{Shaded}

\begin{Shaded}
\begin{Highlighting}[]
\NormalTok{\#| autorun: true}
\NormalTok{fig, ax = pl.subplots()}
\NormalTok{df.plot.scatter(x="RAJ2000", y="DEJ2000", c="Vmag", s=50, figsize=[5,5], ax=ax)}
\NormalTok{pl.show()}
\end{Highlighting}
\end{Shaded}

\begin{Shaded}
\begin{Highlighting}[]
\NormalTok{\#| autorun: true}
\NormalTok{df[\textquotesingle{}r\textquotesingle{}]}
\end{Highlighting}
\end{Shaded}

\section{Olhando na descrição do
catálogo!}\label{olhando-na-descriuxe7uxe3o-do-catuxe1logo}

Note (3) : Sources of the values are defined as follows: - B = USNO-B1.0
(Monet et al., Cat. I/284) - C = UCAC2 (Zacharias et al., Cat. I/289) -
M = 2MASS catalog of point sources (Cutri et al., Cat. II/246) - Y = YB6
Catalog (USNO, unpublished) - T = Tycho-2 Catalog (Hog et al., 2000,
Cat. I/259) - H = Hipparcos catalog (Table I/239/hip\_main) - o =
Palomar-I blue (O) plate (for Bmag) - e = Palomar-I red (E) plate (for
Rmag)

\begin{Shaded}
\begin{Highlighting}[]
\NormalTok{\#| autorun: true}
\NormalTok{df3 = df[[\textquotesingle{}r\textquotesingle{},\textquotesingle{}Bmag\textquotesingle{}, \textquotesingle{}Vmag\textquotesingle{}, \textquotesingle{}Rmag\textquotesingle{}, \textquotesingle{}Jmag\textquotesingle{}, \textquotesingle{}Hmag\textquotesingle{}, \textquotesingle{}Kmag\textquotesingle{}]]}
\end{Highlighting}
\end{Shaded}

\begin{Shaded}
\begin{Highlighting}[]
\NormalTok{\#| autorun: true}
\NormalTok{df3}
\end{Highlighting}
\end{Shaded}

\begin{Shaded}
\begin{Highlighting}[]
\NormalTok{\#| autorun: true}
\NormalTok{df3.groupby([\textquotesingle{}r\textquotesingle{}]).agg([\textquotesingle{}count\textquotesingle{},\textquotesingle{}min\textquotesingle{},\textquotesingle{}mean\textquotesingle{}]) \#\textquotesingle{}median\textquotesingle{}, \textquotesingle{}sum\textquotesingle{}}
\end{Highlighting}
\end{Shaded}

\chapter{Astropy}\label{astropy}

\section{1) Visão Geral}\label{visuxe3o-geral}

Detalhes sobre o pacote Astropy podem ser encontrados em
\href{https://www.astropy.org/}{www.astropy.org}

\subsubsection{Data structures and
transformations}\label{data-structures-and-transformations}

\begin{itemize}
\tightlist
\item
  Constants (\texttt{astropy.constants})
\item
  Units and Quantities (\texttt{astropy.units})
\item
  N-Dimensional Datasets (\texttt{astropy.nddata})
\item
  Data Tables (\texttt{astropy.table})
\item
  Time and Dates (\texttt{astropy.time})
\item
  Time Series (\texttt{astropy.timeseries})
\item
  Astronomical Coordinate Systems (\texttt{astropy.coordinates})
\item
  World Coordinate System (\texttt{astropy.wcs})
\item
  Models and Fitting (\texttt{astropy.modeling})
\item
  Uncertainties and Distributions (\texttt{astropy.uncertainty})
\end{itemize}

\subsubsection{Files, I/O, and
Communication}\label{files-io-and-communication}

\begin{itemize}
\tightlist
\item
  Unified File Read/Write Interface
\item
  FITS File Handling (\texttt{astropy.io.fits})
\item
  ASCII Tables (\texttt{astropy.io.ascii})
\end{itemize}

\subsubsection{Computations and
utilities}\label{computations-and-utilities}

\begin{itemize}
\tightlist
\item
  Cosmological Calculations (\texttt{astropy.cosmology})
\item
  Convolution and Filtering (\texttt{astropy.convolution})
\item
  IERS data access (\texttt{astropy.utils.iers})
\item
  Data Visualization (\texttt{astropy.visualization})
\item
  Astrostatistics Tools (\texttt{astropy.stats})
\item
  Configuration System (\texttt{astropy.config})
\end{itemize}

\textbf{Infelizmente não temos o tempo necessário para cobrir cada um
dos itens, nosso foco será em mostrar alguns pontos chaves e como
iniciar}

\section{2) Astropy Units}\label{astropy-units}

Um módulo muito util para controlar as unidades das variaveis!

\begin{Shaded}
\begin{Highlighting}[]
\ImportTok{import}\NormalTok{ astropy.units }\ImportTok{as}\NormalTok{ u}
\end{Highlighting}
\end{Shaded}

\begin{Shaded}
\begin{Highlighting}[]
\NormalTok{teste }\OperatorTok{=} \DecValTok{10}\OperatorTok{*}\NormalTok{u.km}
\end{Highlighting}
\end{Shaded}

\begin{Shaded}
\begin{Highlighting}[]
\BuiltInTok{print}\NormalTok{(teste)}
\end{Highlighting}
\end{Shaded}

\begin{verbatim}
10.0 km
\end{verbatim}

\begin{Shaded}
\begin{Highlighting}[]
\BuiltInTok{type}\NormalTok{(teste)}
\end{Highlighting}
\end{Shaded}

\begin{verbatim}
astropy.units.quantity.Quantity
\end{verbatim}

\begin{Shaded}
\begin{Highlighting}[]
\CommentTok{\# Podemos ver qual o valor da variavel}
\NormalTok{teste.value}
\end{Highlighting}
\end{Shaded}

\begin{verbatim}
np.float64(10.0)
\end{verbatim}

\begin{Shaded}
\begin{Highlighting}[]
\CommentTok{\# Podemos ver qual a unidade da variavel}
\NormalTok{teste.unit}
\end{Highlighting}
\end{Shaded}

$\mathrm{km}$

\begin{Shaded}
\begin{Highlighting}[]
\CommentTok{\# podemos modificar a unidade da variavel de forma rápida}
\NormalTok{teste.to(}\StringTok{\textquotesingle{}m\textquotesingle{}}\NormalTok{)}
\end{Highlighting}
\end{Shaded}

$10000 \; \mathrm{m}$

\begin{Shaded}
\begin{Highlighting}[]
\NormalTok{teste.to(}\StringTok{\textquotesingle{}au\textquotesingle{}}\NormalTok{)}
\end{Highlighting}
\end{Shaded}

$6.6845871 \times 10^{-8} \; \mathrm{AU}$

\begin{Shaded}
\begin{Highlighting}[]
\NormalTok{teste.to(}\StringTok{\textquotesingle{}parsec\textquotesingle{}}\NormalTok{)}
\end{Highlighting}
\end{Shaded}

$3.2407793 \times 10^{-13} \; \mathrm{pc}$

\begin{Shaded}
\begin{Highlighting}[]
\NormalTok{teste.to(}\StringTok{\textquotesingle{}nm\textquotesingle{}}\NormalTok{)}
\end{Highlighting}
\end{Shaded}

$1 \times 10^{13} \; \mathrm{nm}$

\begin{Shaded}
\begin{Highlighting}[]
\CommentTok{\# Ao criar uma nova variavel ela vai carregar a unidade}
\NormalTok{teste\_2 }\OperatorTok{=} \DecValTok{2}\OperatorTok{*}\NormalTok{teste}
\BuiltInTok{print}\NormalTok{(teste\_2)}
\end{Highlighting}
\end{Shaded}

\begin{verbatim}
20.0 km
\end{verbatim}

\begin{Shaded}
\begin{Highlighting}[]
\CommentTok{\# O código automaticamente irá lidar com unidades relacionadas}
\DecValTok{2000}\OperatorTok{*}\NormalTok{u.m }\OperatorTok{{-}} \DecValTok{5}\OperatorTok{*}\NormalTok{u.km}
\end{Highlighting}
\end{Shaded}

$-3000 \; \mathrm{m}$

\begin{Shaded}
\begin{Highlighting}[]
\CommentTok{\#Porém dará erro com unidades sem relação}
\DecValTok{20}\OperatorTok{*}\NormalTok{u.s }\OperatorTok{+} \DecValTok{5}\OperatorTok{*}\NormalTok{u.km}
\end{Highlighting}
\end{Shaded}

\begin{Highlighting}
\textcolor{black}{UnitConversionError: Can only apply 'add' function to quantities with compatible dimensions}
\textcolor{black}{}\textcolor{QuartoInternalColor1}{---------------------------------------------------------------------------}\textcolor{QuartoInternalColor2}{}
\textcolor{QuartoInternalColor2}{}\textcolor{QuartoInternalColor1}{UnitConversionError}\textcolor{QuartoInternalColor2}{                       Traceback (most recent call last)}
\textcolor{QuartoInternalColor2}{File }\textcolor{QuartoInternalColor3}{\textasciitilde{}\textbackslash{}AppData\textbackslash{}Roaming\textbackslash{}Python\textbackslash{}Python312\textbackslash{}site-packages\textbackslash{}astropy\textbackslash{}units\textbackslash{}quantity\_helper\textbackslash{}helpers.py:94}\textcolor{QuartoInternalColor2}{, in }\textcolor{QuartoInternalColor6}{get\_converters\_and\_unit}\textcolor{QuartoInternalColor7}{(f, unit1, unit2)}\textcolor{QuartoInternalColor2}{}
\textcolor{QuartoInternalColor2}{}\textcolor{QuartoInternalColor8}{     93}\textcolor{QuartoInternalColor2}{ }\textcolor{QuartoInternalColor4}{try}\textcolor{QuartoInternalColor2}{:}
\textcolor{QuartoInternalColor2}{}\textcolor{QuartoInternalColor3}{---> 94}\textcolor{QuartoInternalColor2}{     converters[changeable] }\textcolor{QuartoInternalColor5}{=}\textcolor{QuartoInternalColor2}{ }\textcolor{QuartoInternalColor2}{get\_converter}\textcolor{QuartoInternalColor2}{}\textcolor{QuartoInternalColor2}{(}\textcolor{QuartoInternalColor2}{}\textcolor{QuartoInternalColor2}{unit2}\textcolor{QuartoInternalColor2}{}\textcolor{QuartoInternalColor2}{,}\textcolor{QuartoInternalColor2}{}\textcolor{QuartoInternalColor2}{ }\textcolor{QuartoInternalColor2}{}\textcolor{QuartoInternalColor2}{unit1}\textcolor{QuartoInternalColor2}{}\textcolor{QuartoInternalColor2}{)}\textcolor{QuartoInternalColor2}{}
\textcolor{QuartoInternalColor2}{}\textcolor{QuartoInternalColor8}{     95}\textcolor{QuartoInternalColor2}{ }\textcolor{QuartoInternalColor4}{except}\textcolor{QuartoInternalColor2}{ UnitsError:}
\textcolor{QuartoInternalColor2}{File }\textcolor{QuartoInternalColor3}{\textasciitilde{}\textbackslash{}AppData\textbackslash{}Roaming\textbackslash{}Python\textbackslash{}Python312\textbackslash{}site-packages\textbackslash{}astropy\textbackslash{}units\textbackslash{}quantity\_helper\textbackslash{}helpers.py:43}\textcolor{QuartoInternalColor2}{, in }\textcolor{QuartoInternalColor6}{get\_converter}\textcolor{QuartoInternalColor7}{(from\_unit, to\_unit)}\textcolor{QuartoInternalColor2}{}
\textcolor{QuartoInternalColor2}{}\textcolor{QuartoInternalColor8}{     42}\textcolor{QuartoInternalColor2}{ }\textcolor{QuartoInternalColor4}{try}\textcolor{QuartoInternalColor2}{:}
\textcolor{QuartoInternalColor2}{}\textcolor{QuartoInternalColor3}{---> 43}\textcolor{QuartoInternalColor2}{     converter }\textcolor{QuartoInternalColor5}{=}\textcolor{QuartoInternalColor2}{ }\textcolor{QuartoInternalColor2}{from\_unit}\textcolor{QuartoInternalColor2}{}\textcolor{QuartoInternalColor5}{.}\textcolor{QuartoInternalColor2}{}\textcolor{QuartoInternalColor2}{get\_converter}\textcolor{QuartoInternalColor2}{}\textcolor{QuartoInternalColor2}{(}\textcolor{QuartoInternalColor2}{}\textcolor{QuartoInternalColor2}{to\_unit}\textcolor{QuartoInternalColor2}{}\textcolor{QuartoInternalColor2}{)}\textcolor{QuartoInternalColor2}{}
\textcolor{QuartoInternalColor2}{}\textcolor{QuartoInternalColor8}{     44}\textcolor{QuartoInternalColor2}{ }\textcolor{QuartoInternalColor4}{except}\textcolor{QuartoInternalColor2}{ }\textcolor{QuartoInternalColor9}{AttributeError}\textcolor{QuartoInternalColor2}{ }\textcolor{QuartoInternalColor4}{as}\textcolor{QuartoInternalColor2}{ exc:}
\textcolor{QuartoInternalColor2}{}\textcolor{QuartoInternalColor8}{     45}\textcolor{QuartoInternalColor2}{     }\textcolor{QuartoInternalColor10}{# Check for lack of unit only now, to avoid delay for cases where a unit}\textcolor{QuartoInternalColor2}{}
\textcolor{QuartoInternalColor2}{}\textcolor{QuartoInternalColor8}{     46}\textcolor{QuartoInternalColor2}{     }\textcolor{QuartoInternalColor10}{# was present. Note that cases where dimensionless is expected are}\textcolor{QuartoInternalColor2}{}
\textcolor{QuartoInternalColor2}{}\textcolor{QuartoInternalColor8}{     47}\textcolor{QuartoInternalColor2}{     }\textcolor{QuartoInternalColor10}{# already short-circuited; here, we cover just the case where, e.g., the}\textcolor{QuartoInternalColor2}{}
\textcolor{QuartoInternalColor2}{}\textcolor{QuartoInternalColor8}{     48}\textcolor{QuartoInternalColor2}{     }\textcolor{QuartoInternalColor10}{# user has done u.add\_enabled\_equivalencies(u.dimensionless\_angles()).}\textcolor{QuartoInternalColor2}{}
\textcolor{QuartoInternalColor2}{File }\textcolor{QuartoInternalColor3}{\textasciitilde{}\textbackslash{}AppData\textbackslash{}Roaming\textbackslash{}Python\textbackslash{}Python312\textbackslash{}site-packages\textbackslash{}astropy\textbackslash{}units\textbackslash{}core.py:589}\textcolor{QuartoInternalColor2}{, in }\textcolor{QuartoInternalColor6}{UnitBase.get\_converter}\textcolor{QuartoInternalColor7}{(self, other, equivalencies)}\textcolor{QuartoInternalColor2}{}
\textcolor{QuartoInternalColor2}{}\textcolor{QuartoInternalColor8}{    587}\textcolor{QuartoInternalColor2}{                 }\textcolor{QuartoInternalColor4}{return}\textcolor{QuartoInternalColor2}{ }\textcolor{QuartoInternalColor4}{lambda}\textcolor{QuartoInternalColor2}{ v: b(converter(v))}
\textcolor{QuartoInternalColor2}{}\textcolor{QuartoInternalColor3}{--> 589}\textcolor{QuartoInternalColor2}{ }\textcolor{QuartoInternalColor4}{raise}\textcolor{QuartoInternalColor2}{ exc}
\textcolor{QuartoInternalColor2}{File }\textcolor{QuartoInternalColor3}{\textasciitilde{}\textbackslash{}AppData\textbackslash{}Roaming\textbackslash{}Python\textbackslash{}Python312\textbackslash{}site-packages\textbackslash{}astropy\textbackslash{}units\textbackslash{}core.py:572}\textcolor{QuartoInternalColor2}{, in }\textcolor{QuartoInternalColor6}{UnitBase.get\_converter}\textcolor{QuartoInternalColor7}{(self, other, equivalencies)}\textcolor{QuartoInternalColor2}{}
\textcolor{QuartoInternalColor2}{}\textcolor{QuartoInternalColor8}{    571}\textcolor{QuartoInternalColor2}{ }\textcolor{QuartoInternalColor4}{try}\textcolor{QuartoInternalColor2}{:}
\textcolor{QuartoInternalColor2}{}\textcolor{QuartoInternalColor3}{--> 572}\textcolor{QuartoInternalColor2}{     }\textcolor{QuartoInternalColor4}{return}\textcolor{QuartoInternalColor2}{ }\textcolor{QuartoInternalColor4}{self}\textcolor{QuartoInternalColor2}{}\textcolor{QuartoInternalColor5}{.}\textcolor{QuartoInternalColor2}{}\textcolor{QuartoInternalColor2}{\_apply\_equivalencies}\textcolor{QuartoInternalColor2}{}\textcolor{QuartoInternalColor2}{(}\textcolor{QuartoInternalColor2}{}
\textcolor{QuartoInternalColor2}{}\textcolor{QuartoInternalColor8}{    573}\textcolor{QuartoInternalColor2}{ }\textcolor{QuartoInternalColor2}{        }\textcolor{QuartoInternalColor2}{}\textcolor{QuartoInternalColor4}{self}\textcolor{QuartoInternalColor2}{}\textcolor{QuartoInternalColor2}{,}\textcolor{QuartoInternalColor2}{}\textcolor{QuartoInternalColor2}{ }\textcolor{QuartoInternalColor2}{}\textcolor{QuartoInternalColor2}{other}\textcolor{QuartoInternalColor2}{}\textcolor{QuartoInternalColor2}{,}\textcolor{QuartoInternalColor2}{}\textcolor{QuartoInternalColor2}{ }\textcolor{QuartoInternalColor2}{}\textcolor{QuartoInternalColor4}{self}\textcolor{QuartoInternalColor2}{}\textcolor{QuartoInternalColor5}{.}\textcolor{QuartoInternalColor2}{}\textcolor{QuartoInternalColor2}{\_normalize\_equivalencies}\textcolor{QuartoInternalColor2}{}\textcolor{QuartoInternalColor2}{(}\textcolor{QuartoInternalColor2}{}\textcolor{QuartoInternalColor2}{equivalencies}\textcolor{QuartoInternalColor2}{}\textcolor{QuartoInternalColor2}{)}\textcolor{QuartoInternalColor2}{}
\textcolor{QuartoInternalColor2}{}\textcolor{QuartoInternalColor8}{    574}\textcolor{QuartoInternalColor2}{ }\textcolor{QuartoInternalColor2}{    }\textcolor{QuartoInternalColor2}{}\textcolor{QuartoInternalColor2}{)}\textcolor{QuartoInternalColor2}{}
\textcolor{QuartoInternalColor2}{}\textcolor{QuartoInternalColor8}{    575}\textcolor{QuartoInternalColor2}{ }\textcolor{QuartoInternalColor4}{except}\textcolor{QuartoInternalColor2}{ UnitsError }\textcolor{QuartoInternalColor4}{as}\textcolor{QuartoInternalColor2}{ exc:}
\textcolor{QuartoInternalColor2}{}\textcolor{QuartoInternalColor8}{    576}\textcolor{QuartoInternalColor2}{     }\textcolor{QuartoInternalColor10}{# Last hope: maybe other knows how to do it?}\textcolor{QuartoInternalColor2}{}
\textcolor{QuartoInternalColor2}{}\textcolor{QuartoInternalColor8}{    577}\textcolor{QuartoInternalColor2}{     }\textcolor{QuartoInternalColor10}{# We assume the equivalencies have the unit itself as first item.}\textcolor{QuartoInternalColor2}{}
\textcolor{QuartoInternalColor2}{}\textcolor{QuartoInternalColor8}{    578}\textcolor{QuartoInternalColor2}{     }\textcolor{QuartoInternalColor10}{# TODO: maybe better for other to have a `\_back\_converter` method?}\textcolor{QuartoInternalColor2}{}
\textcolor{QuartoInternalColor2}{File }\textcolor{QuartoInternalColor3}{\textasciitilde{}\textbackslash{}AppData\textbackslash{}Roaming\textbackslash{}Python\textbackslash{}Python312\textbackslash{}site-packages\textbackslash{}astropy\textbackslash{}units\textbackslash{}core.py:523}\textcolor{QuartoInternalColor2}{, in }\textcolor{QuartoInternalColor6}{UnitBase.\_apply\_equivalencies}\textcolor{QuartoInternalColor7}{(self, unit, other, equivalencies)}\textcolor{QuartoInternalColor2}{}
\textcolor{QuartoInternalColor2}{}\textcolor{QuartoInternalColor8}{    521}\textcolor{QuartoInternalColor2}{ other\_str }\textcolor{QuartoInternalColor5}{=}\textcolor{QuartoInternalColor2}{ get\_err\_str(other)}
\textcolor{QuartoInternalColor2}{}\textcolor{QuartoInternalColor3}{--> 523}\textcolor{QuartoInternalColor2}{ }\textcolor{QuartoInternalColor4}{raise}\textcolor{QuartoInternalColor2}{ UnitConversionError(}\textcolor{QuartoInternalColor11}{f}\textcolor{QuartoInternalColor2}{}\textcolor{QuartoInternalColor11}{"}\textcolor{QuartoInternalColor2}{}\textcolor{QuartoInternalColor12}{\{}\textcolor{QuartoInternalColor2}{unit\_str}\textcolor{QuartoInternalColor12}{\}}\textcolor{QuartoInternalColor2}{}\textcolor{QuartoInternalColor11}{ and }\textcolor{QuartoInternalColor2}{}\textcolor{QuartoInternalColor12}{\{}\textcolor{QuartoInternalColor2}{other\_str}\textcolor{QuartoInternalColor12}{\}}\textcolor{QuartoInternalColor2}{}\textcolor{QuartoInternalColor11}{ are not convertible}\textcolor{QuartoInternalColor2}{}\textcolor{QuartoInternalColor11}{"}\textcolor{QuartoInternalColor2}{)}
\textcolor{QuartoInternalColor2}{}\textcolor{QuartoInternalColor1}{UnitConversionError}\textcolor{QuartoInternalColor2}{: 'km' (length) and 's' (time) are not convertible}
\textcolor{QuartoInternalColor2}{During handling of the above exception, another exception occurred:}
\textcolor{QuartoInternalColor2}{}\textcolor{QuartoInternalColor1}{UnitConversionError}\textcolor{QuartoInternalColor2}{                       Traceback (most recent call last)}
\textcolor{QuartoInternalColor2}{Cell }\textcolor{QuartoInternalColor3}{In[13], line 2}\textcolor{QuartoInternalColor2}{}
\textcolor{QuartoInternalColor2}{}\textcolor{QuartoInternalColor8}{      1}\textcolor{QuartoInternalColor2}{ }\textcolor{QuartoInternalColor10}{#Porém dará erro com unidades sem relação}\textcolor{QuartoInternalColor2}{}
\textcolor{QuartoInternalColor2}{}\textcolor{QuartoInternalColor3}{----> 2}\textcolor{QuartoInternalColor2}{ }\textcolor{QuartoInternalColor5}{20}\textcolor{QuartoInternalColor2}{}\textcolor{QuartoInternalColor5}{*}\textcolor{QuartoInternalColor2}{}\textcolor{QuartoInternalColor2}{u}\textcolor{QuartoInternalColor2}{}\textcolor{QuartoInternalColor5}{.}\textcolor{QuartoInternalColor2}{}\textcolor{QuartoInternalColor2}{s}\textcolor{QuartoInternalColor2}{}\textcolor{QuartoInternalColor2}{ }\textcolor{QuartoInternalColor2}{}\textcolor{QuartoInternalColor5}{+}\textcolor{QuartoInternalColor2}{}\textcolor{QuartoInternalColor2}{ }\textcolor{QuartoInternalColor2}{}\textcolor{QuartoInternalColor5}{5}\textcolor{QuartoInternalColor2}{}\textcolor{QuartoInternalColor5}{*}\textcolor{QuartoInternalColor2}{}\textcolor{QuartoInternalColor2}{u}\textcolor{QuartoInternalColor2}{}\textcolor{QuartoInternalColor5}{.}\textcolor{QuartoInternalColor2}{}\textcolor{QuartoInternalColor2}{km}\textcolor{QuartoInternalColor2}{}
\textcolor{QuartoInternalColor2}{File }\textcolor{QuartoInternalColor3}{\textasciitilde{}\textbackslash{}AppData\textbackslash{}Roaming\textbackslash{}Python\textbackslash{}Python312\textbackslash{}site-packages\textbackslash{}astropy\textbackslash{}units\textbackslash{}quantity.py:673}\textcolor{QuartoInternalColor2}{, in }\textcolor{QuartoInternalColor6}{Quantity.\_\_array\_ufunc\_\_}\textcolor{QuartoInternalColor7}{(self, function, method, *inputs, **kwargs)}\textcolor{QuartoInternalColor2}{}
\textcolor{QuartoInternalColor2}{}\textcolor{QuartoInternalColor8}{    671}\textcolor{QuartoInternalColor2}{     }\textcolor{QuartoInternalColor4}{return}\textcolor{QuartoInternalColor2}{ }\textcolor{QuartoInternalColor4}{NotImplemented}\textcolor{QuartoInternalColor2}{}
\textcolor{QuartoInternalColor2}{}\textcolor{QuartoInternalColor8}{    672}\textcolor{QuartoInternalColor2}{ }\textcolor{QuartoInternalColor4}{else}\textcolor{QuartoInternalColor2}{:}
\textcolor{QuartoInternalColor2}{}\textcolor{QuartoInternalColor3}{--> 673}\textcolor{QuartoInternalColor2}{     }\textcolor{QuartoInternalColor4}{raise}\textcolor{QuartoInternalColor2}{ e}
\textcolor{QuartoInternalColor2}{File }\textcolor{QuartoInternalColor3}{\textasciitilde{}\textbackslash{}AppData\textbackslash{}Roaming\textbackslash{}Python\textbackslash{}Python312\textbackslash{}site-packages\textbackslash{}astropy\textbackslash{}units\textbackslash{}quantity.py:618}\textcolor{QuartoInternalColor2}{, in }\textcolor{QuartoInternalColor6}{Quantity.\_\_array\_ufunc\_\_}\textcolor{QuartoInternalColor7}{(self, function, method, *inputs, **kwargs)}\textcolor{QuartoInternalColor2}{}
\textcolor{QuartoInternalColor2}{}\textcolor{QuartoInternalColor8}{    613}\textcolor{QuartoInternalColor2}{ }\textcolor{QuartoInternalColor10}{# Determine required conversion functions -- to bring the unit of the}\textcolor{QuartoInternalColor2}{}
\textcolor{QuartoInternalColor2}{}\textcolor{QuartoInternalColor8}{    614}\textcolor{QuartoInternalColor2}{ }\textcolor{QuartoInternalColor10}{# input to that expected (e.g., radian for np.sin), or to get}\textcolor{QuartoInternalColor2}{}
\textcolor{QuartoInternalColor2}{}\textcolor{QuartoInternalColor8}{    615}\textcolor{QuartoInternalColor2}{ }\textcolor{QuartoInternalColor10}{# consistent units between two inputs (e.g., in np.add) --}\textcolor{QuartoInternalColor2}{}
\textcolor{QuartoInternalColor2}{}\textcolor{QuartoInternalColor8}{    616}\textcolor{QuartoInternalColor2}{ }\textcolor{QuartoInternalColor10}{# and the unit of the result (or tuple of units for nout > 1).}\textcolor{QuartoInternalColor2}{}
\textcolor{QuartoInternalColor2}{}\textcolor{QuartoInternalColor8}{    617}\textcolor{QuartoInternalColor2}{ }\textcolor{QuartoInternalColor4}{try}\textcolor{QuartoInternalColor2}{:}
\textcolor{QuartoInternalColor2}{}\textcolor{QuartoInternalColor3}{--> 618}\textcolor{QuartoInternalColor2}{     converters, unit }\textcolor{QuartoInternalColor5}{=}\textcolor{QuartoInternalColor2}{ }\textcolor{QuartoInternalColor2}{converters\_and\_unit}\textcolor{QuartoInternalColor2}{}\textcolor{QuartoInternalColor2}{(}\textcolor{QuartoInternalColor2}{}\textcolor{QuartoInternalColor2}{function}\textcolor{QuartoInternalColor2}{}\textcolor{QuartoInternalColor2}{,}\textcolor{QuartoInternalColor2}{}\textcolor{QuartoInternalColor2}{ }\textcolor{QuartoInternalColor2}{}\textcolor{QuartoInternalColor2}{method}\textcolor{QuartoInternalColor2}{}\textcolor{QuartoInternalColor2}{,}\textcolor{QuartoInternalColor2}{}\textcolor{QuartoInternalColor2}{ }\textcolor{QuartoInternalColor2}{}\textcolor{QuartoInternalColor5}{*}\textcolor{QuartoInternalColor2}{}\textcolor{QuartoInternalColor2}{inputs}\textcolor{QuartoInternalColor2}{}\textcolor{QuartoInternalColor2}{)}\textcolor{QuartoInternalColor2}{}
\textcolor{QuartoInternalColor2}{}\textcolor{QuartoInternalColor8}{    620}\textcolor{QuartoInternalColor2}{     out }\textcolor{QuartoInternalColor5}{=}\textcolor{QuartoInternalColor2}{ kwargs}\textcolor{QuartoInternalColor5}{.}\textcolor{QuartoInternalColor2}{get(}\textcolor{QuartoInternalColor11}{"}\textcolor{QuartoInternalColor2}{}\textcolor{QuartoInternalColor11}{out}\textcolor{QuartoInternalColor2}{}\textcolor{QuartoInternalColor11}{"}\textcolor{QuartoInternalColor2}{)}
\textcolor{QuartoInternalColor2}{}\textcolor{QuartoInternalColor8}{    621}\textcolor{QuartoInternalColor2}{     }\textcolor{QuartoInternalColor10}{# Avoid loop back by turning any Quantity output into array views.}\textcolor{QuartoInternalColor2}{}
\textcolor{QuartoInternalColor2}{File }\textcolor{QuartoInternalColor3}{\textasciitilde{}\textbackslash{}AppData\textbackslash{}Roaming\textbackslash{}Python\textbackslash{}Python312\textbackslash{}site-packages\textbackslash{}astropy\textbackslash{}units\textbackslash{}quantity\_helper\textbackslash{}converters.py:177}\textcolor{QuartoInternalColor2}{, in }\textcolor{QuartoInternalColor6}{converters\_and\_unit}\textcolor{QuartoInternalColor7}{(function, method, *args)}\textcolor{QuartoInternalColor2}{}
\textcolor{QuartoInternalColor2}{}\textcolor{QuartoInternalColor8}{    174}\textcolor{QuartoInternalColor2}{ units }\textcolor{QuartoInternalColor5}{=}\textcolor{QuartoInternalColor2}{ [}\textcolor{QuartoInternalColor4}{getattr}\textcolor{QuartoInternalColor2}{(arg, }\textcolor{QuartoInternalColor11}{"}\textcolor{QuartoInternalColor2}{}\textcolor{QuartoInternalColor11}{unit}\textcolor{QuartoInternalColor2}{}\textcolor{QuartoInternalColor11}{"}\textcolor{QuartoInternalColor2}{, }\textcolor{QuartoInternalColor4}{None}\textcolor{QuartoInternalColor2}{) }\textcolor{QuartoInternalColor4}{for}\textcolor{QuartoInternalColor2}{ arg }\textcolor{QuartoInternalColor13}{in}\textcolor{QuartoInternalColor2}{ args]}
\textcolor{QuartoInternalColor2}{}\textcolor{QuartoInternalColor8}{    176}\textcolor{QuartoInternalColor2}{ }\textcolor{QuartoInternalColor10}{# Determine possible conversion functions, and the result unit.}\textcolor{QuartoInternalColor2}{}
\textcolor{QuartoInternalColor2}{}\textcolor{QuartoInternalColor3}{--> 177}\textcolor{QuartoInternalColor2}{ converters, result\_unit }\textcolor{QuartoInternalColor5}{=}\textcolor{QuartoInternalColor2}{ }\textcolor{QuartoInternalColor2}{ufunc\_helper}\textcolor{QuartoInternalColor2}{}\textcolor{QuartoInternalColor2}{(}\textcolor{QuartoInternalColor2}{}\textcolor{QuartoInternalColor2}{function}\textcolor{QuartoInternalColor2}{}\textcolor{QuartoInternalColor2}{,}\textcolor{QuartoInternalColor2}{}\textcolor{QuartoInternalColor2}{ }\textcolor{QuartoInternalColor2}{}\textcolor{QuartoInternalColor5}{*}\textcolor{QuartoInternalColor2}{}\textcolor{QuartoInternalColor2}{units}\textcolor{QuartoInternalColor2}{}\textcolor{QuartoInternalColor2}{)}\textcolor{QuartoInternalColor2}{}
\textcolor{QuartoInternalColor2}{}\textcolor{QuartoInternalColor8}{    179}\textcolor{QuartoInternalColor2}{ }\textcolor{QuartoInternalColor4}{if}\textcolor{QuartoInternalColor2}{ }\textcolor{QuartoInternalColor4}{any}\textcolor{QuartoInternalColor2}{(converter }\textcolor{QuartoInternalColor13}{is}\textcolor{QuartoInternalColor2}{ }\textcolor{QuartoInternalColor4}{False}\textcolor{QuartoInternalColor2}{ }\textcolor{QuartoInternalColor4}{for}\textcolor{QuartoInternalColor2}{ converter }\textcolor{QuartoInternalColor13}{in}\textcolor{QuartoInternalColor2}{ converters):}
\textcolor{QuartoInternalColor2}{}\textcolor{QuartoInternalColor8}{    180}\textcolor{QuartoInternalColor2}{     }\textcolor{QuartoInternalColor10}{# for multi-argument ufuncs with a quantity and a non-quantity,}\textcolor{QuartoInternalColor2}{}
\textcolor{QuartoInternalColor2}{}\textcolor{QuartoInternalColor8}{    181}\textcolor{QuartoInternalColor2}{     }\textcolor{QuartoInternalColor10}{# the quantity normally needs to be dimensionless, *except*}\textcolor{QuartoInternalColor2}{}
\textcolor{QuartoInternalColor2}{}\textcolor{QuartoInternalColor3}{   (...)}\textcolor{QuartoInternalColor2}{}
\textcolor{QuartoInternalColor2}{}\textcolor{QuartoInternalColor8}{    184}\textcolor{QuartoInternalColor2}{     }\textcolor{QuartoInternalColor10}{# can just have the unit of the quantity}\textcolor{QuartoInternalColor2}{}
\textcolor{QuartoInternalColor2}{}\textcolor{QuartoInternalColor8}{    185}\textcolor{QuartoInternalColor2}{     }\textcolor{QuartoInternalColor10}{# (this allows, e.g., `q > 0.` independent of unit)}\textcolor{QuartoInternalColor2}{}
\textcolor{QuartoInternalColor2}{}\textcolor{QuartoInternalColor8}{    186}\textcolor{QuartoInternalColor2}{     }\textcolor{QuartoInternalColor4}{try}\textcolor{QuartoInternalColor2}{:}
\textcolor{QuartoInternalColor2}{}\textcolor{QuartoInternalColor8}{    187}\textcolor{QuartoInternalColor2}{         }\textcolor{QuartoInternalColor10}{# Don't fold this loop in the test above: this rare case}\textcolor{QuartoInternalColor2}{}
\textcolor{QuartoInternalColor2}{}\textcolor{QuartoInternalColor8}{    188}\textcolor{QuartoInternalColor2}{         }\textcolor{QuartoInternalColor10}{# should not make the common case slower.}\textcolor{QuartoInternalColor2}{}
\textcolor{QuartoInternalColor2}{File }\textcolor{QuartoInternalColor3}{\textasciitilde{}\textbackslash{}AppData\textbackslash{}Roaming\textbackslash{}Python\textbackslash{}Python312\textbackslash{}site-packages\textbackslash{}astropy\textbackslash{}units\textbackslash{}quantity\_helper\textbackslash{}helpers.py:96}\textcolor{QuartoInternalColor2}{, in }\textcolor{QuartoInternalColor6}{get\_converters\_and\_unit}\textcolor{QuartoInternalColor7}{(f, unit1, unit2)}\textcolor{QuartoInternalColor2}{}
\textcolor{QuartoInternalColor2}{}\textcolor{QuartoInternalColor8}{     94}\textcolor{QuartoInternalColor2}{     converters[changeable] }\textcolor{QuartoInternalColor5}{=}\textcolor{QuartoInternalColor2}{ get\_converter(unit2, unit1)}
\textcolor{QuartoInternalColor2}{}\textcolor{QuartoInternalColor8}{     95}\textcolor{QuartoInternalColor2}{ }\textcolor{QuartoInternalColor4}{except}\textcolor{QuartoInternalColor2}{ UnitsError:}
\textcolor{QuartoInternalColor2}{}\textcolor{QuartoInternalColor3}{---> 96}\textcolor{QuartoInternalColor2}{     }\textcolor{QuartoInternalColor4}{raise}\textcolor{QuartoInternalColor2}{ UnitConversionError(}
\textcolor{QuartoInternalColor2}{}\textcolor{QuartoInternalColor8}{     97}\textcolor{QuartoInternalColor2}{         }\textcolor{QuartoInternalColor11}{f}\textcolor{QuartoInternalColor2}{}\textcolor{QuartoInternalColor11}{"}\textcolor{QuartoInternalColor2}{}\textcolor{QuartoInternalColor11}{Can only apply }\textcolor{QuartoInternalColor2}{}\textcolor{QuartoInternalColor11}{'}\textcolor{QuartoInternalColor2}{}\textcolor{QuartoInternalColor12}{\{}\textcolor{QuartoInternalColor2}{f}\textcolor{QuartoInternalColor5}{.}\textcolor{QuartoInternalColor2}{}\textcolor{QuartoInternalColor14}{\_\_name\_\_}\textcolor{QuartoInternalColor2}{}\textcolor{QuartoInternalColor12}{\}}\textcolor{QuartoInternalColor2}{}\textcolor{QuartoInternalColor11}{'}\textcolor{QuartoInternalColor2}{}\textcolor{QuartoInternalColor11}{ function to quantities }\textcolor{QuartoInternalColor2}{}\textcolor{QuartoInternalColor11}{"}\textcolor{QuartoInternalColor2}{}
\textcolor{QuartoInternalColor2}{}\textcolor{QuartoInternalColor8}{     98}\textcolor{QuartoInternalColor2}{         }\textcolor{QuartoInternalColor11}{"}\textcolor{QuartoInternalColor2}{}\textcolor{QuartoInternalColor11}{with compatible dimensions}\textcolor{QuartoInternalColor2}{}\textcolor{QuartoInternalColor11}{"}\textcolor{QuartoInternalColor2}{}
\textcolor{QuartoInternalColor2}{}\textcolor{QuartoInternalColor8}{     99}\textcolor{QuartoInternalColor2}{     )}
\textcolor{QuartoInternalColor2}{}\textcolor{QuartoInternalColor8}{    101}\textcolor{QuartoInternalColor2}{ }\textcolor{QuartoInternalColor4}{return}\textcolor{QuartoInternalColor2}{ converters, unit1}
\textcolor{QuartoInternalColor2}{}\textcolor{QuartoInternalColor1}{UnitConversionError}\textcolor{QuartoInternalColor2}{: Can only apply 'add' function to quantities with compatible dimensions}
\end{Highlighting}

\begin{Shaded}
\begin{Highlighting}[]
\CommentTok{\# Operações que fazem sentido irão funcionar}
\NormalTok{teste }\OperatorTok{=}\NormalTok{ (}\DecValTok{20}\OperatorTok{*}\NormalTok{u.km) }\OperatorTok{/}\NormalTok{ (}\DecValTok{10}\OperatorTok{*}\NormalTok{u.hour)}
\BuiltInTok{print}\NormalTok{(teste)}
\end{Highlighting}
\end{Shaded}

\begin{verbatim}
2.0 km / h
\end{verbatim}

\begin{Shaded}
\begin{Highlighting}[]
\DecValTok{20}\OperatorTok{*}\NormalTok{u.m }\OperatorTok{*} \DecValTok{40}\OperatorTok{*}\NormalTok{u.m}
\end{Highlighting}
\end{Shaded}

$800 \; \mathrm{m^{2}}$

\begin{Shaded}
\begin{Highlighting}[]
\CommentTok{\# De maneira similar podemos mudar as unidades}
\NormalTok{teste.to(}\StringTok{\textquotesingle{}m/s\textquotesingle{}}\NormalTok{)}
\end{Highlighting}
\end{Shaded}

$0.55555556 \; \mathrm{\frac{m}{s}}$

\begin{Shaded}
\begin{Highlighting}[]
\CommentTok{\# Em especial, isso nos auxilia lidar com ângulos}
\ImportTok{import}\NormalTok{ numpy }\ImportTok{as}\NormalTok{ np}

\NormalTok{np.tan(}\DecValTok{30}\OperatorTok{*}\NormalTok{u.deg)}
\end{Highlighting}
\end{Shaded}

$0.57735027 \; \mathrm{}$

\begin{Shaded}
\begin{Highlighting}[]
\CommentTok{\# Não precisamos trabalhar com apenas 1 valor, podemos lidar com arrays de forma similar}
\NormalTok{teste }\OperatorTok{=}\NormalTok{ np.linspace(}\DecValTok{0}\NormalTok{, }\DecValTok{360}\NormalTok{, }\DecValTok{36}\NormalTok{)}\OperatorTok{*}\NormalTok{u.deg}
\end{Highlighting}
\end{Shaded}

\begin{Shaded}
\begin{Highlighting}[]
\NormalTok{teste.to(}\StringTok{\textquotesingle{}hourangle\textquotesingle{}}\NormalTok{) }\CommentTok{\#+ 10*u.arcsec }
\end{Highlighting}
\end{Shaded}

$[0,~0.68571429,~1.3714286,~2.0571429,~2.7428571,~3.4285714,~4.1142857,~4.8,~5.4857143,~6.1714286,~6.8571429,~7.5428571,~8.2285714,~8.9142857,~9.6,~10.285714,~10.971429,~11.657143,~12.342857,~13.028571,~13.714286,~14.4,~15.085714,~15.771429,~16.457143,~17.142857,~17.828571,~18.514286,~19.2,~19.885714,~20.571429,~21.257143,~21.942857,~22.628571,~23.314286,~24]\mathrm{{}^{h}}$

\begin{Shaded}
\begin{Highlighting}[]
\BuiltInTok{print}\NormalTok{(teste)}
\end{Highlighting}
\end{Shaded}

\begin{verbatim}
[  0.          10.28571429  20.57142857  30.85714286  41.14285714
  51.42857143  61.71428571  72.          82.28571429  92.57142857
 102.85714286 113.14285714 123.42857143 133.71428571 144.
 154.28571429 164.57142857 174.85714286 185.14285714 195.42857143
 205.71428571 216.         226.28571429 236.57142857 246.85714286
 257.14285714 267.42857143 277.71428571 288.         298.28571429
 308.57142857 318.85714286 329.14285714 339.42857143 349.71428571
 360.        ] deg
\end{verbatim}

\begin{Shaded}
\begin{Highlighting}[]
\BuiltInTok{type}\NormalTok{(teste)}
\end{Highlighting}
\end{Shaded}

\begin{verbatim}
astropy.units.quantity.Quantity
\end{verbatim}

\begin{Shaded}
\begin{Highlighting}[]
\NormalTok{teste.to(}\StringTok{\textquotesingle{}arcmin\textquotesingle{}}\NormalTok{)}
\end{Highlighting}
\end{Shaded}

$[0,~617.14286,~1234.2857,~1851.4286,~2468.5714,~3085.7143,~3702.8571,~4320,~4937.1429,~5554.2857,~6171.4286,~6788.5714,~7405.7143,~8022.8571,~8640,~9257.1429,~9874.2857,~10491.429,~11108.571,~11725.714,~12342.857,~12960,~13577.143,~14194.286,~14811.429,~15428.571,~16045.714,~16662.857,~17280,~17897.143,~18514.286,~19131.429,~19748.571,~20365.714,~20982.857,~21600]\mathrm{{}^{\prime}}$

\begin{Shaded}
\begin{Highlighting}[]
\DecValTok{10}\OperatorTok{*}\NormalTok{u.}\BuiltInTok{min} \OperatorTok{+} \DecValTok{30}\OperatorTok{*}\NormalTok{u.s }\OperatorTok{+} \DecValTok{2}\OperatorTok{*}\NormalTok{u.hour}
\end{Highlighting}
\end{Shaded}

$130.5 \; \mathrm{min}$

\begin{Shaded}
\begin{Highlighting}[]
\NormalTok{u.dex?}
\end{Highlighting}
\end{Shaded}

\section{3) Astropy Constants}\label{astropy-constants}

Um módulo muito util para controlar constantes necessárias!

\begin{Shaded}
\begin{Highlighting}[]
\ImportTok{import}\NormalTok{ astropy.constants }\ImportTok{as}\NormalTok{ const}
\end{Highlighting}
\end{Shaded}

\begin{Shaded}
\begin{Highlighting}[]
\CommentTok{\# aqui podemos encontrar uma coletânea de constantes interessantes para a Astronomia}
\NormalTok{const.R\_sun.uncertainty}
\end{Highlighting}
\end{Shaded}

\begin{verbatim}
0.0
\end{verbatim}

\begin{Shaded}
\begin{Highlighting}[]
\CommentTok{\# As Docstrings cotém os detalhes dessas constantes}
\NormalTok{const.pc?}
\end{Highlighting}
\end{Shaded}

\begin{Shaded}
\begin{Highlighting}[]
\NormalTok{vel }\OperatorTok{=}\NormalTok{ const.c}
\BuiltInTok{print}\NormalTok{(vel)}
\end{Highlighting}
\end{Shaded}

\begin{verbatim}
  Name   = Speed of light in vacuum
  Value  = 299792458.0
  Uncertainty  = 0.0
  Unit  = m / s
  Reference = CODATA 2018
\end{verbatim}

\begin{Shaded}
\begin{Highlighting}[]
\BuiltInTok{type}\NormalTok{(vel)}
\end{Highlighting}
\end{Shaded}

\begin{verbatim}
astropy.constants.codata2018.CODATA2018
\end{verbatim}

\begin{Shaded}
\begin{Highlighting}[]
\CommentTok{\# Assim como as variaveis com unidades, podemos converter facilmente entre variaveis}
\NormalTok{vel.to(}\StringTok{\textquotesingle{}Rjup/day\textquotesingle{}}\NormalTok{)}
\end{Highlighting}
\end{Shaded}

$362307.23 \; \mathrm{\frac{R_{\rm J}}{d}}$

\begin{Shaded}
\begin{Highlighting}[]
\NormalTok{const.R\_jup?}
\end{Highlighting}
\end{Shaded}

\section{4) Astropy Time}\label{astropy-time}

Lidando com o tempo!

\begin{Shaded}
\begin{Highlighting}[]
\ImportTok{from}\NormalTok{ astropy.time }\ImportTok{import}\NormalTok{ Time}
\end{Highlighting}
\end{Shaded}

\begin{Shaded}
\begin{Highlighting}[]
\NormalTok{t }\OperatorTok{=}\NormalTok{ Time.now()}
\end{Highlighting}
\end{Shaded}

\begin{Shaded}
\begin{Highlighting}[]
\BuiltInTok{print}\NormalTok{(t)}
\end{Highlighting}
\end{Shaded}

\begin{verbatim}
2026-08-12 22:56:31.153605
\end{verbatim}

\begin{Shaded}
\begin{Highlighting}[]
\NormalTok{t.scale}
\end{Highlighting}
\end{Shaded}

\begin{verbatim}
'utc'
\end{verbatim}

\begin{Shaded}
\begin{Highlighting}[]
\CommentTok{\# agora podemos acessar esse tempo de diferentes formas}
\NormalTok{t.iso}
\end{Highlighting}
\end{Shaded}

\begin{verbatim}
'2026-08-12 22:56:31.154'
\end{verbatim}

\begin{Shaded}
\begin{Highlighting}[]
\CommentTok{\#Data Juliana}
\NormalTok{t.jd}
\end{Highlighting}
\end{Shaded}

\begin{verbatim}
np.float64(2461265.4559161295)
\end{verbatim}

\begin{Shaded}
\begin{Highlighting}[]
\CommentTok{\# Data Juliana modificada}
\NormalTok{t.mjd}
\end{Highlighting}
\end{Shaded}

\begin{verbatim}
np.float64(61264.955916129686)
\end{verbatim}

\begin{Shaded}
\begin{Highlighting}[]
\CommentTok{\# Besselian Epoch}
\NormalTok{t.byear}
\end{Highlighting}
\end{Shaded}

\begin{verbatim}
np.float64(2026.6150038261553)
\end{verbatim}

\begin{Shaded}
\begin{Highlighting}[]
\CommentTok{\# Além disso, o tempo tem uma escala associada}
\NormalTok{t.scale}
\end{Highlighting}
\end{Shaded}

\begin{verbatim}
'utc'
\end{verbatim}

\textbf{Segue a lista de escalas de tempo}

\begin{itemize}
\tightlist
\item
  tai -\textgreater{} International Atomic Time
\item
  tcb -\textgreater{} Barycentric Coordinate Time
\item
  tcg -\textgreater{} Geocentric Coordinate Time
\item
  tdb -\textgreater{} Barycentric Dynamical Time
\item
  tt -\textgreater{} Terrestrial Time
\item
  ut1 -\textgreater{} Universal Time
\item
  utc -\textgreater{} Coordinate Universal Time
\item
  local -\textgreater{} Local Time Scale
\end{itemize}

\begin{Shaded}
\begin{Highlighting}[]
\CommentTok{\# e podemos converter entre escalas de tempo}
\NormalTok{t.tdb.iso}
\end{Highlighting}
\end{Shaded}

\begin{verbatim}
'2026-08-12 22:57:40.337'
\end{verbatim}

\begin{Shaded}
\begin{Highlighting}[]
\NormalTok{t.tt.iso}
\end{Highlighting}
\end{Shaded}

\begin{verbatim}
'2026-08-12 22:57:40.338'
\end{verbatim}

\begin{Shaded}
\begin{Highlighting}[]
\NormalTok{t.utc.iso}
\end{Highlighting}
\end{Shaded}

\begin{verbatim}
'2026-08-12 22:56:31.154'
\end{verbatim}

\textbf{Cada escala de tempo tem diferentes aplicações na astronomia}

\begin{Shaded}
\begin{Highlighting}[]
\CommentTok{\# Além disso, podemos combinar o Time com o units de maneira automática}

\NormalTok{t1 }\OperatorTok{=}\NormalTok{ Time(}\StringTok{\textquotesingle{}2022{-}02{-}05 01:20\textquotesingle{}}\NormalTok{) }\OperatorTok{+} \FloatTok{2.3}\OperatorTok{*}\NormalTok{u.hour}

\BuiltInTok{print}\NormalTok{(t1.iso)}
\end{Highlighting}
\end{Shaded}

\begin{verbatim}
2022-02-05 03:38:00.000
\end{verbatim}

\begin{Shaded}
\begin{Highlighting}[]
\BuiltInTok{type}\NormalTok{(t1)}
\end{Highlighting}
\end{Shaded}

\begin{verbatim}
astropy.time.core.Time
\end{verbatim}

\begin{Shaded}
\begin{Highlighting}[]
\CommentTok{\# Uma coisa que estamos muito interessados é na passagem do tempo}

\NormalTok{t2 }\OperatorTok{=}\NormalTok{ Time.now()}
\end{Highlighting}
\end{Shaded}

\begin{Shaded}
\begin{Highlighting}[]
\NormalTok{dt }\OperatorTok{=}\NormalTok{ t2 }\OperatorTok{{-}}\NormalTok{ t}
\end{Highlighting}
\end{Shaded}

\begin{Shaded}
\begin{Highlighting}[]
\NormalTok{t2}
\end{Highlighting}
\end{Shaded}

\begin{verbatim}
<Time object: scale='utc' format='datetime' value=2026-08-12 22:56:31.429185>
\end{verbatim}

\begin{Shaded}
\begin{Highlighting}[]
\BuiltInTok{type}\NormalTok{(dt)}
\end{Highlighting}
\end{Shaded}

\begin{verbatim}
astropy.time.core.TimeDelta
\end{verbatim}

\begin{Shaded}
\begin{Highlighting}[]
\CommentTok{\# diferença em dias}
\NormalTok{dt.jd}
\end{Highlighting}
\end{Shaded}

\begin{verbatim}
np.float64(3.1895833333406287e-06)
\end{verbatim}

\begin{Shaded}
\begin{Highlighting}[]
\CommentTok{\# diferença em segundos}
\NormalTok{dt.sec}
\end{Highlighting}
\end{Shaded}

\begin{verbatim}
np.float64(0.2755800000006303)
\end{verbatim}

\begin{Shaded}
\begin{Highlighting}[]
\NormalTok{dt.datetime}
\end{Highlighting}
\end{Shaded}

\begin{verbatim}
datetime.timedelta(microseconds=275580)
\end{verbatim}

\section{5) Astropy coordinates}\label{astropy-coordinates}

Lidando com o coordenadas!

\begin{Shaded}
\begin{Highlighting}[]
\CommentTok{\# Vamos começar com localização na Terra}
\ImportTok{from}\NormalTok{ astropy.coordinates }\ImportTok{import}\NormalTok{ EarthLocation}
\end{Highlighting}
\end{Shaded}

\begin{Shaded}
\begin{Highlighting}[]
\NormalTok{bruno }\OperatorTok{=}\NormalTok{ EarthLocation(lon}\OperatorTok{=}\StringTok{\textquotesingle{}{-}43 11 11\textquotesingle{}}\NormalTok{, }
\NormalTok{                      lat}\OperatorTok{=}\StringTok{\textquotesingle{}{-}22 53 54\textquotesingle{}}\NormalTok{, }
\NormalTok{                      height}\OperatorTok{=}\DecValTok{38}\OperatorTok{*}\NormalTok{u.m)}
\end{Highlighting}
\end{Shaded}

\begin{Shaded}
\begin{Highlighting}[]
\BuiltInTok{type}\NormalTok{(bruno)}
\end{Highlighting}
\end{Shaded}

\begin{verbatim}
astropy.coordinates.earth.EarthLocation
\end{verbatim}

\begin{Shaded}
\begin{Highlighting}[]
\BuiltInTok{print}\NormalTok{(bruno)}
\end{Highlighting}
\end{Shaded}

\begin{verbatim}
(4286222.4313010145, -4023115.0608965205, -2466366.344181889) m
\end{verbatim}

\begin{Shaded}
\begin{Highlighting}[]
\NormalTok{bruno.geodetic}
\end{Highlighting}
\end{Shaded}

\begin{verbatim}
GeodeticLocation(lon=<Longitude -43.18638889 deg>, lat=<Latitude -22.89833333 deg>, height=<Quantity 38. m>)
\end{verbatim}

\begin{Shaded}
\begin{Highlighting}[]
\CommentTok{\# Vamos começar com localização na Terra}
\ImportTok{from}\NormalTok{ astropy.coordinates }\ImportTok{import}\NormalTok{ SkyCoord}
\end{Highlighting}
\end{Shaded}

\begin{Shaded}
\begin{Highlighting}[]
\NormalTok{SkyCoord?}
\end{Highlighting}
\end{Shaded}

\begin{Shaded}
\begin{Highlighting}[]
\NormalTok{star }\OperatorTok{=}\NormalTok{ SkyCoord(ra}\OperatorTok{=}\StringTok{\textquotesingle{}01 12 43.2\textquotesingle{}}\NormalTok{, }
\NormalTok{                dec}\OperatorTok{=}\StringTok{\textquotesingle{}+12 23 44\textquotesingle{}}\NormalTok{, }
\NormalTok{                distance}\OperatorTok{=}\DecValTok{10}\OperatorTok{*}\NormalTok{u.parsec, }
\NormalTok{                unit}\OperatorTok{=}\NormalTok{[u.hourangle, u.deg])}
\end{Highlighting}
\end{Shaded}

\begin{Shaded}
\begin{Highlighting}[]
\BuiltInTok{type}\NormalTok{(star)}
\end{Highlighting}
\end{Shaded}

\begin{verbatim}
astropy.coordinates.sky_coordinate.SkyCoord
\end{verbatim}

\begin{Shaded}
\begin{Highlighting}[]
\CommentTok{\# pegar só parte dos valores (só RA)}
\NormalTok{star.ra}
\end{Highlighting}
\end{Shaded}

$18^\circ10{}^\prime48{}^{\prime\prime}$

\begin{Shaded}
\begin{Highlighting}[]
\CommentTok{\# só Declinação}
\NormalTok{star.dec}
\end{Highlighting}
\end{Shaded}

$12^\circ23{}^\prime44{}^{\prime\prime}$

\begin{Shaded}
\begin{Highlighting}[]
\NormalTok{star.distance.to(}\StringTok{\textquotesingle{}au\textquotesingle{}}\NormalTok{)}
\end{Highlighting}
\end{Shaded}

$2062648.1 \; \mathrm{AU}$

\begin{Shaded}
\begin{Highlighting}[]
\CommentTok{\# Alterando a representação para cartesiana}
\NormalTok{star.cartesian}
\end{Highlighting}
\end{Shaded}

\begin{verbatim}
<CartesianRepresentation (x, y, z) in pc
    (9.27933616, 3.04730167, 2.14659566)>
\end{verbatim}

\begin{Shaded}
\begin{Highlighting}[]
\CommentTok{\# apenas um das coordenadas}
\NormalTok{star.cartesian.x}
\end{Highlighting}
\end{Shaded}

$9.2793362 \; \mathrm{pc}$

\begin{Shaded}
\begin{Highlighting}[]
\CommentTok{\# Agora podemos converter entre os sistemas de referência}
\NormalTok{star.barycentrictrueecliptic}
\end{Highlighting}
\end{Shaded}

\begin{verbatim}
<SkyCoord (BarycentricTrueEcliptic: equinox=J2000.000): (lon, lat, distance) in (deg, deg, pc)
    (21.46663903, 4.34325771, 10.)>
\end{verbatim}

\begin{Shaded}
\begin{Highlighting}[]
\NormalTok{star.galactic}
\end{Highlighting}
\end{Shaded}

\begin{verbatim}
<SkyCoord (Galactic): (l, b, distance) in (deg, deg, pc)
    (131.05544829, -50.14020949, 10.)>
\end{verbatim}

\begin{Shaded}
\begin{Highlighting}[]
\CommentTok{\# Podemos analisar diferentes estrelas}
\NormalTok{star\_2 }\OperatorTok{=}\NormalTok{ SkyCoord(ra}\OperatorTok{=}\StringTok{\textquotesingle{}01 12 43.2\textquotesingle{}}\NormalTok{, }
\NormalTok{                  dec}\OperatorTok{=}\StringTok{\textquotesingle{}{-}12 23 44\textquotesingle{}}\NormalTok{, }
\NormalTok{                  distance}\OperatorTok{=}\DecValTok{10}\OperatorTok{*}\NormalTok{u.parsec, }
\NormalTok{                  unit}\OperatorTok{=}\NormalTok{[u.hourangle, u.deg])}
\end{Highlighting}
\end{Shaded}

\begin{Shaded}
\begin{Highlighting}[]
\CommentTok{\# Podemos computar a separação (no céu) das duas estrelas}
\NormalTok{star.separation(star\_2)}
\end{Highlighting}
\end{Shaded}

$24^\circ47{}^\prime28{}^{\prime\prime}$

\chapter{Isso é tudo}\label{isso-uxe9-tudo-1}

By BEM

\section{1) Abrindo arquivos ASCII}\label{abrindo-arquivos-ascii}

\begin{Shaded}
\begin{Highlighting}[]
\ImportTok{import}\NormalTok{ astropy.io.}\BuiltInTok{ascii} \ImportTok{as}\NormalTok{ asc}
\end{Highlighting}
\end{Shaded}

\begin{Shaded}
\begin{Highlighting}[]
\NormalTok{asc.read?}
\end{Highlighting}
\end{Shaded}

\begin{Shaded}
\begin{Highlighting}[]
\NormalTok{tb }\OperatorTok{=}\NormalTok{ asc.read(}\StringTok{\textquotesingle{}teste.txt\textquotesingle{}}\NormalTok{)}
\end{Highlighting}
\end{Shaded}

\begin{Shaded}
\begin{Highlighting}[]
\NormalTok{tb}
\end{Highlighting}
\end{Shaded}

\begin{Shaded}
\begin{Highlighting}[]
\BuiltInTok{type}\NormalTok{(tb)}
\end{Highlighting}
\end{Shaded}

\begin{verbatim}
astropy.table.table.Table
\end{verbatim}

\begin{Shaded}
\begin{Highlighting}[]
\CommentTok{\# Temos algo bem parecido com o pandas}
\NormalTok{tb[}\StringTok{\textquotesingle{}obsid\textquotesingle{}}\NormalTok{]}
\end{Highlighting}
\end{Shaded}

{\def\LTcaptype{none} % do not increment counter
\begin{longtable}[]{@{}l@{}}
\toprule\noalign{}
\endhead
\bottomrule\noalign{}
\endlastfoot
3102 \\
877 \\
\end{longtable}
}

\begin{Shaded}
\begin{Highlighting}[]
\NormalTok{tb.write?}
\end{Highlighting}
\end{Shaded}

\begin{Shaded}
\begin{Highlighting}[]
\CommentTok{\# Inclusive podemos transformar em um objeto do Pandas}
\NormalTok{df }\OperatorTok{=}\NormalTok{ tb.to\_pandas()}
\end{Highlighting}
\end{Shaded}

\begin{Shaded}
\begin{Highlighting}[]
\NormalTok{df}
\end{Highlighting}
\end{Shaded}

{\def\LTcaptype{none} % do not increment counter
\begin{longtable}[]{@{}llllll@{}}
\toprule\noalign{}
& obsid & redshift & X & Y & object \\
\midrule\noalign{}
\endhead
\bottomrule\noalign{}
\endlastfoot
0 & 3102 & 0.32 & 4167 & 4085 & Q1250+568-A \\
1 & 877 & 0.22 & 4378 & 3892 & Source\_82 \\
\end{longtable}
}

\section{2) Abrindo imagens fits}\label{abrindo-imagens-fits}

\begin{Shaded}
\begin{Highlighting}[]
\ImportTok{import}\NormalTok{ astropy.io.fits }\ImportTok{as}\NormalTok{ fits}
\end{Highlighting}
\end{Shaded}

\begin{Shaded}
\begin{Highlighting}[]
\NormalTok{lista }\OperatorTok{=}\NormalTok{ fits.}\BuiltInTok{open}\NormalTok{(}\StringTok{\textquotesingle{}exemplo\_1.fits\textquotesingle{}}\NormalTok{)}
\end{Highlighting}
\end{Shaded}

\begin{Shaded}
\begin{Highlighting}[]
\BuiltInTok{print}\NormalTok{(lista)}
\end{Highlighting}
\end{Shaded}

\begin{verbatim}
[<astropy.io.fits.hdu.image.PrimaryHDU object at 0x0000025293BE9760>]
\end{verbatim}

\begin{Shaded}
\begin{Highlighting}[]
\BuiltInTok{type}\NormalTok{(lista)}
\end{Highlighting}
\end{Shaded}

\begin{verbatim}
astropy.io.fits.hdu.hdulist.HDUList
\end{verbatim}

\begin{Shaded}
\begin{Highlighting}[]
\BuiltInTok{type}\NormalTok{(lista[}\DecValTok{0}\NormalTok{])}
\end{Highlighting}
\end{Shaded}

\begin{verbatim}
astropy.io.fits.hdu.image.PrimaryHDU
\end{verbatim}

\begin{Shaded}
\begin{Highlighting}[]
\NormalTok{imagem }\OperatorTok{=}\NormalTok{ lista[}\DecValTok{0}\NormalTok{]}
\end{Highlighting}
\end{Shaded}

\begin{Shaded}
\begin{Highlighting}[]
\NormalTok{imagem.header}
\end{Highlighting}
\end{Shaded}

\begin{verbatim}
SIMPLE  =                    T / file does conform to FITS standard             
BITPIX  =                   16 / number of bits per data pixel                  
NAXIS   =                    2 / number of data axes                            
NAXIS1  =                 2048 / length of data axis 1                          
NAXIS2  =                 2048 / length of data axis 2                          
EXTEND  =                    T / FITS dataset may contain extensions            
COMMENT = 'Imagens para foco'  / Comentario                                     
COMMENT       and Astrophysics', volume 376, page 359; bibcode: 2001A&A...376..3
BZERO   =                32768 / offset data range to that of unsigned short    
BSCALE  =                    1 / default scaling factor                         
HEAD    = '        '           / Head Model                                     
ACQMODE = 'Single  '           / Acquisition Mode                               
READMODE= 'Image   '           / Readout Mode                                   
IMGRECT = '1, 2048, 2048, 1'   / Image Format                                   
HBIN    = '1       '           / Horizontal Binning                             
VBIN    = '1       '           / Vertical Binning                               
SUBRECT = '1, 2048,2048,1'     / Subimage Format                                
DATATYPE= 'Counts  '           / Data type                                      
XTYPE   = 'Pixel number'       / Calibration type                               
XUNIT   =                    0 / Calibration units                              
RAYWAVE =                   0. / Rayleigh Wavelength                            
CALBWVNM=                    0 / Wave calibration                               
TRIGGER = 'Internal'           / Trigger Mode                                   
CALIB   = '0.,1.,0.,0.'        / Calibration                                    
DLLVER  = '2.102.30024.0'      / Software Version                               
EXPOSURE= '10,00000'           / Total Exposure Time                            
TEMP    = '-999    '           / Temperature                                    
READTIME=         3.333333E-07 / Pixel readout time                             
OPERATN =                    4 / Type of system                                 
GAIN    = '1.0     '           / Ganho                                          
EMREALGN=                    0 / EM Real Gain                                   
VCLKAMP =                    0 / Vertical Clock Amplitude                       
VSHIFT  = '76,95   '           / Vertical Shift Speed                           
PREAMP  = '4x      '           / Pre Amplifier Gain                             
SERNO   = '10127   '           / Serial Number                                  
UNSTTEMP= '-63,0   '           / Unstabilized Temperature                       
BLCLAMP =                    F / Baseline Clamp                                 
PRECAN  =                    0 / Prescans                                       
FLIPX   =                    0 / Horizontally Flipped                           
FLIPY   =                    0 / Vertically Flipped                             
CNTCVTMD=                    0 / Count Convert Mode                             
CNTCVT  =                    0 / Count Convert                                  
DTNWLGTH=                 500. / Detection Wavelength                           
SNTVTY  =                   0. / Sensitivity                                    
SPSNFLTR=                    0 / Spurious Noise Filter Mode                     
THRSHLD =                   0. / Threshold                                      
PCNTENLD=                    0 / Photon Counting Enabled                        
NSETHSLD=                    0 / Number of Photon Counting Thresholds           
PTNTHLD1=                   0. / Photon Counting Threshold 1                    
PTNTHLD2=                   0. / Photon Counting Threshold 2                    
PTNTHLD3=                   0. / Photon Counting Threshold 3                    
PTNTHLD4=                   0. / Photon Counting Threshold 4                    
AVGFTRMD=                    0 / Averaging Filter Mode                          
AVGFCTR =                    1 / Averaging factor                               
FRMCNT  =                    1 / Frame Count                                    
PORTMODE=                   -1 / Port Readout Mode                              
LSHEIGHT=                    0 / Exposure Window Height                         
LSSPEED =                   0. / Line Scan Speed                                
LSALTDIR=                    0 / Alternating Readout Direction                  
LSCTRL  =                    0 / Scan Speed Control                             
LSDIR   =                    0 / Readout Direction                              
FKSMODE =                    0 / Fast Kinetics Storage Mode                     
FKTMODE =                    0 / Fast Kinetics Time Scan Mode                   
USERTXT1= '        '           / User text                                      
USERTXT2= '        '           / User text                                      
USERTXT3= '        '           / User text                                      
USERTXT4= '        '           / User text                                      
DATE    = '2019-09-16T23:29:10' / File Creation Date (YYYY-MM-HHThh:mm:ss)      
FRAME   = '2019-09-16T23:29:10.291' / Start of Frame Exposure                   
ESHTMODE=                    0 / Electronic Shuttering Mode                     
HIERARCH PREAMPGAINTEXT = '        ' / Pre-Amplifier Gain                       
HIERARCH SPECTROGRAPHSERIAL = '        ' / Spectrograph Serial                  
HIERARCH SHAMROCKISACTIVE =  0 / Shamrock is active                             
HIERARCH SPECTROGRAPHNAME = '        ' / Spectrograph name                      
HIERARCH SPECTROGRAPHISACTIVE = 0 / Spectrograph is active                      
HIERARCH IRIGDATAAVAILABLE = 0 / IRIG Availability                              
IRIGDATA= '        '           / IRIG Timestamp                                 
HIERARCH PIVENABLE =         0 / PIV Mode Status                                
HIERARCH EXTENDEDDYNAMICRANGE = 0 / Extended Dynamic Range Mode Status          
HIERARCH MILLISECONDS =      0 / Millisecond Count Status                       
HIERARCH DIRECTIRISPRESENT = 0 / Direct Iris Present                            
HIERARCH DIRECTIRISPOS =     0 / Direct Iris Position                           
HIERARCH SIDEIRISPRESENT =   0 / Side Iris Present                              
HIERARCH SIDEIRISPOS =       0 / Side Iris Position                             
DETECTOR= '        '           / Head Model                                     
EXPTIME = '10,00000'           / Total Exposure Time                            
OUTPTAMP= 'High Sensitivity'   / Output Amplifier                               
CAMGAIN = '4x      '           / Pre Amplifier Gain                             
DATE-OBS= '2019-09-16T23:29:10.291' / Start of Frame Exposure                   
RA      = '23:08:39'           / Alfa                                           
DEC     = '-03:08:08'          / Delta                                          
EPOCH   = '2000.0  '           / Epoca                                          
AIRMASS = '1.488   '           / Massa de Ar                                    
JD      = '2458743.47858'      / Data Juliana                                   
ST      = '20:09:16'           / Tempo Sideral                                  
HA      = '-3:00:20'           / Angulo Horario                                 
IMAGE   = 'Hektor_00001'       / Nome do arquivo                                
TELESCOP= '0.60m(BC)'          / Telescopio                                     
PLATESCL= '0.34    '           / Escala de Placa ("/pixel)                      
INSTRUME= 'cam 2+IkonL'        / Instrumento                                    
OBJECT  = 'Hektor  '           / Objeto                                         
OBSERVER= 'Bruno Morgado'      / Observador                                     
RDNOISE = '12.1    '           / Ruido de Leitura                               
FILTER  = 'Clear   '           / Filtro                                         
FOCUSVAL= '2117    '           / Valor do Foco                                  
W-BAR   = '760,9   '           / Pressao Atmosferica (mb)                       
W-TEMP  = '15,6    '           / Temperatura Ambiente                           
W-HUM   = '45      '           / Umidade (%)                                    
\end{verbatim}

\section{3) Lidando com o Header}\label{lidando-com-o-header}

\begin{Shaded}
\begin{Highlighting}[]
\NormalTok{im\_header }\OperatorTok{=}\NormalTok{ imagem.header}
\end{Highlighting}
\end{Shaded}

\begin{Shaded}
\begin{Highlighting}[]
\BuiltInTok{type}\NormalTok{(im\_header)}
\end{Highlighting}
\end{Shaded}

\begin{verbatim}
astropy.io.fits.header.Header
\end{verbatim}

\begin{Shaded}
\begin{Highlighting}[]
\NormalTok{im\_header}
\end{Highlighting}
\end{Shaded}

\begin{verbatim}
SIMPLE  =                    T / file does conform to FITS standard             
BITPIX  =                   16 / number of bits per data pixel                  
NAXIS   =                    2 / number of data axes                            
NAXIS1  =                 2048 / length of data axis 1                          
NAXIS2  =                 2048 / length of data axis 2                          
EXTEND  =                    T / FITS dataset may contain extensions            
COMMENT = 'Imagens para foco'  / Comentario                                     
COMMENT       and Astrophysics', volume 376, page 359; bibcode: 2001A&A...376..3
BZERO   =                32768 / offset data range to that of unsigned short    
BSCALE  =                    1 / default scaling factor                         
HEAD    = '        '           / Head Model                                     
ACQMODE = 'Single  '           / Acquisition Mode                               
READMODE= 'Image   '           / Readout Mode                                   
IMGRECT = '1, 2048, 2048, 1'   / Image Format                                   
HBIN    = '1       '           / Horizontal Binning                             
VBIN    = '1       '           / Vertical Binning                               
SUBRECT = '1, 2048,2048,1'     / Subimage Format                                
DATATYPE= 'Counts  '           / Data type                                      
XTYPE   = 'Pixel number'       / Calibration type                               
XUNIT   =                    0 / Calibration units                              
RAYWAVE =                   0. / Rayleigh Wavelength                            
CALBWVNM=                    0 / Wave calibration                               
TRIGGER = 'Internal'           / Trigger Mode                                   
CALIB   = '0.,1.,0.,0.'        / Calibration                                    
DLLVER  = '2.102.30024.0'      / Software Version                               
EXPOSURE= '10,00000'           / Total Exposure Time                            
TEMP    = '-999    '           / Temperature                                    
READTIME=         3.333333E-07 / Pixel readout time                             
OPERATN =                    4 / Type of system                                 
GAIN    = '1.0     '           / Ganho                                          
EMREALGN=                    0 / EM Real Gain                                   
VCLKAMP =                    0 / Vertical Clock Amplitude                       
VSHIFT  = '76,95   '           / Vertical Shift Speed                           
PREAMP  = '4x      '           / Pre Amplifier Gain                             
SERNO   = '10127   '           / Serial Number                                  
UNSTTEMP= '-63,0   '           / Unstabilized Temperature                       
BLCLAMP =                    F / Baseline Clamp                                 
PRECAN  =                    0 / Prescans                                       
FLIPX   =                    0 / Horizontally Flipped                           
FLIPY   =                    0 / Vertically Flipped                             
CNTCVTMD=                    0 / Count Convert Mode                             
CNTCVT  =                    0 / Count Convert                                  
DTNWLGTH=                 500. / Detection Wavelength                           
SNTVTY  =                   0. / Sensitivity                                    
SPSNFLTR=                    0 / Spurious Noise Filter Mode                     
THRSHLD =                   0. / Threshold                                      
PCNTENLD=                    0 / Photon Counting Enabled                        
NSETHSLD=                    0 / Number of Photon Counting Thresholds           
PTNTHLD1=                   0. / Photon Counting Threshold 1                    
PTNTHLD2=                   0. / Photon Counting Threshold 2                    
PTNTHLD3=                   0. / Photon Counting Threshold 3                    
PTNTHLD4=                   0. / Photon Counting Threshold 4                    
AVGFTRMD=                    0 / Averaging Filter Mode                          
AVGFCTR =                    1 / Averaging factor                               
FRMCNT  =                    1 / Frame Count                                    
PORTMODE=                   -1 / Port Readout Mode                              
LSHEIGHT=                    0 / Exposure Window Height                         
LSSPEED =                   0. / Line Scan Speed                                
LSALTDIR=                    0 / Alternating Readout Direction                  
LSCTRL  =                    0 / Scan Speed Control                             
LSDIR   =                    0 / Readout Direction                              
FKSMODE =                    0 / Fast Kinetics Storage Mode                     
FKTMODE =                    0 / Fast Kinetics Time Scan Mode                   
USERTXT1= '        '           / User text                                      
USERTXT2= '        '           / User text                                      
USERTXT3= '        '           / User text                                      
USERTXT4= '        '           / User text                                      
DATE    = '2019-09-16T23:29:10' / File Creation Date (YYYY-MM-HHThh:mm:ss)      
FRAME   = '2019-09-16T23:29:10.291' / Start of Frame Exposure                   
ESHTMODE=                    0 / Electronic Shuttering Mode                     
HIERARCH PREAMPGAINTEXT = '        ' / Pre-Amplifier Gain                       
HIERARCH SPECTROGRAPHSERIAL = '        ' / Spectrograph Serial                  
HIERARCH SHAMROCKISACTIVE =  0 / Shamrock is active                             
HIERARCH SPECTROGRAPHNAME = '        ' / Spectrograph name                      
HIERARCH SPECTROGRAPHISACTIVE = 0 / Spectrograph is active                      
HIERARCH IRIGDATAAVAILABLE = 0 / IRIG Availability                              
IRIGDATA= '        '           / IRIG Timestamp                                 
HIERARCH PIVENABLE =         0 / PIV Mode Status                                
HIERARCH EXTENDEDDYNAMICRANGE = 0 / Extended Dynamic Range Mode Status          
HIERARCH MILLISECONDS =      0 / Millisecond Count Status                       
HIERARCH DIRECTIRISPRESENT = 0 / Direct Iris Present                            
HIERARCH DIRECTIRISPOS =     0 / Direct Iris Position                           
HIERARCH SIDEIRISPRESENT =   0 / Side Iris Present                              
HIERARCH SIDEIRISPOS =       0 / Side Iris Position                             
DETECTOR= '        '           / Head Model                                     
EXPTIME = '10,00000'           / Total Exposure Time                            
OUTPTAMP= 'High Sensitivity'   / Output Amplifier                               
CAMGAIN = '4x      '           / Pre Amplifier Gain                             
DATE-OBS= '2019-09-16T23:29:10.291' / Start of Frame Exposure                   
RA      = '23:08:39'           / Alfa                                           
DEC     = '-03:08:08'          / Delta                                          
EPOCH   = '2000.0  '           / Epoca                                          
AIRMASS = '1.488   '           / Massa de Ar                                    
JD      = '2458743.47858'      / Data Juliana                                   
ST      = '20:09:16'           / Tempo Sideral                                  
HA      = '-3:00:20'           / Angulo Horario                                 
IMAGE   = 'Hektor_00001'       / Nome do arquivo                                
TELESCOP= '0.60m(BC)'          / Telescopio                                     
PLATESCL= '0.34    '           / Escala de Placa ("/pixel)                      
INSTRUME= 'cam 2+IkonL'        / Instrumento                                    
OBJECT  = 'Hektor  '           / Objeto                                         
OBSERVER= 'Bruno Morgado'      / Observador                                     
RDNOISE = '12.1    '           / Ruido de Leitura                               
FILTER  = 'Clear   '           / Filtro                                         
FOCUSVAL= '2117    '           / Valor do Foco                                  
W-BAR   = '760,9   '           / Pressao Atmosferica (mb)                       
W-TEMP  = '15,6    '           / Temperatura Ambiente                           
W-HUM   = '45      '           / Umidade (%)                                    
\end{verbatim}

\begin{Shaded}
\begin{Highlighting}[]
\CommentTok{\# Podemos buscar algum elemento especifico do Header}

\NormalTok{im\_header.get(}\StringTok{\textquotesingle{}DATE{-}OBS\textquotesingle{}}\NormalTok{)}
\end{Highlighting}
\end{Shaded}

\begin{verbatim}
'2019-09-16T23:29:10.291'
\end{verbatim}

\begin{Shaded}
\begin{Highlighting}[]
\CommentTok{\# Podemos adicionar elemenos ao Header}
\NormalTok{im\_header.append(}\StringTok{\textquotesingle{}ANALIST\textquotesingle{}}\NormalTok{)}
\end{Highlighting}
\end{Shaded}

\begin{Shaded}
\begin{Highlighting}[]
\NormalTok{im\_header}
\end{Highlighting}
\end{Shaded}

\begin{verbatim}
SIMPLE  =                    T / file does conform to FITS standard             
BITPIX  =                   16 / number of bits per data pixel                  
NAXIS   =                    2 / number of data axes                            
NAXIS1  =                 2048 / length of data axis 1                          
NAXIS2  =                 2048 / length of data axis 2                          
EXTEND  =                    T / FITS dataset may contain extensions            
COMMENT = 'Imagens para foco'  / Comentario                                     
COMMENT       and Astrophysics', volume 376, page 359; bibcode: 2001A&A...376..3
BZERO   =                32768 / offset data range to that of unsigned short    
BSCALE  =                    1 / default scaling factor                         
HEAD    = '        '           / Head Model                                     
ACQMODE = 'Single  '           / Acquisition Mode                               
READMODE= 'Image   '           / Readout Mode                                   
IMGRECT = '1, 2048, 2048, 1'   / Image Format                                   
HBIN    = '1       '           / Horizontal Binning                             
VBIN    = '1       '           / Vertical Binning                               
SUBRECT = '1, 2048,2048,1'     / Subimage Format                                
DATATYPE= 'Counts  '           / Data type                                      
XTYPE   = 'Pixel number'       / Calibration type                               
XUNIT   =                    0 / Calibration units                              
RAYWAVE =                   0. / Rayleigh Wavelength                            
CALBWVNM=                    0 / Wave calibration                               
TRIGGER = 'Internal'           / Trigger Mode                                   
CALIB   = '0.,1.,0.,0.'        / Calibration                                    
DLLVER  = '2.102.30024.0'      / Software Version                               
EXPOSURE= '10,00000'           / Total Exposure Time                            
TEMP    = '-999    '           / Temperature                                    
READTIME=         3.333333E-07 / Pixel readout time                             
OPERATN =                    4 / Type of system                                 
GAIN    = '1.0     '           / Ganho                                          
EMREALGN=                    0 / EM Real Gain                                   
VCLKAMP =                    0 / Vertical Clock Amplitude                       
VSHIFT  = '76,95   '           / Vertical Shift Speed                           
PREAMP  = '4x      '           / Pre Amplifier Gain                             
SERNO   = '10127   '           / Serial Number                                  
UNSTTEMP= '-63,0   '           / Unstabilized Temperature                       
BLCLAMP =                    F / Baseline Clamp                                 
PRECAN  =                    0 / Prescans                                       
FLIPX   =                    0 / Horizontally Flipped                           
FLIPY   =                    0 / Vertically Flipped                             
CNTCVTMD=                    0 / Count Convert Mode                             
CNTCVT  =                    0 / Count Convert                                  
DTNWLGTH=                 500. / Detection Wavelength                           
SNTVTY  =                   0. / Sensitivity                                    
SPSNFLTR=                    0 / Spurious Noise Filter Mode                     
THRSHLD =                   0. / Threshold                                      
PCNTENLD=                    0 / Photon Counting Enabled                        
NSETHSLD=                    0 / Number of Photon Counting Thresholds           
PTNTHLD1=                   0. / Photon Counting Threshold 1                    
PTNTHLD2=                   0. / Photon Counting Threshold 2                    
PTNTHLD3=                   0. / Photon Counting Threshold 3                    
PTNTHLD4=                   0. / Photon Counting Threshold 4                    
AVGFTRMD=                    0 / Averaging Filter Mode                          
AVGFCTR =                    1 / Averaging factor                               
FRMCNT  =                    1 / Frame Count                                    
PORTMODE=                   -1 / Port Readout Mode                              
LSHEIGHT=                    0 / Exposure Window Height                         
LSSPEED =                   0. / Line Scan Speed                                
LSALTDIR=                    0 / Alternating Readout Direction                  
LSCTRL  =                    0 / Scan Speed Control                             
LSDIR   =                    0 / Readout Direction                              
FKSMODE =                    0 / Fast Kinetics Storage Mode                     
FKTMODE =                    0 / Fast Kinetics Time Scan Mode                   
USERTXT1= '        '           / User text                                      
USERTXT2= '        '           / User text                                      
USERTXT3= '        '           / User text                                      
USERTXT4= '        '           / User text                                      
DATE    = '2019-09-16T23:29:10' / File Creation Date (YYYY-MM-HHThh:mm:ss)      
FRAME   = '2019-09-16T23:29:10.291' / Start of Frame Exposure                   
ESHTMODE=                    0 / Electronic Shuttering Mode                     
HIERARCH PREAMPGAINTEXT = '        ' / Pre-Amplifier Gain                       
HIERARCH SPECTROGRAPHSERIAL = '        ' / Spectrograph Serial                  
HIERARCH SHAMROCKISACTIVE =  0 / Shamrock is active                             
HIERARCH SPECTROGRAPHNAME = '        ' / Spectrograph name                      
HIERARCH SPECTROGRAPHISACTIVE = 0 / Spectrograph is active                      
HIERARCH IRIGDATAAVAILABLE = 0 / IRIG Availability                              
IRIGDATA= '        '           / IRIG Timestamp                                 
HIERARCH PIVENABLE =         0 / PIV Mode Status                                
HIERARCH EXTENDEDDYNAMICRANGE = 0 / Extended Dynamic Range Mode Status          
HIERARCH MILLISECONDS =      0 / Millisecond Count Status                       
HIERARCH DIRECTIRISPRESENT = 0 / Direct Iris Present                            
HIERARCH DIRECTIRISPOS =     0 / Direct Iris Position                           
HIERARCH SIDEIRISPRESENT =   0 / Side Iris Present                              
HIERARCH SIDEIRISPOS =       0 / Side Iris Position                             
DETECTOR= '        '           / Head Model                                     
EXPTIME = '10,00000'           / Total Exposure Time                            
OUTPTAMP= 'High Sensitivity'   / Output Amplifier                               
CAMGAIN = '4x      '           / Pre Amplifier Gain                             
DATE-OBS= '2019-09-16T23:29:10.291' / Start of Frame Exposure                   
RA      = '23:08:39'           / Alfa                                           
DEC     = '-03:08:08'          / Delta                                          
EPOCH   = '2000.0  '           / Epoca                                          
AIRMASS = '1.488   '           / Massa de Ar                                    
JD      = '2458743.47858'      / Data Juliana                                   
ST      = '20:09:16'           / Tempo Sideral                                  
HA      = '-3:00:20'           / Angulo Horario                                 
IMAGE   = 'Hektor_00001'       / Nome do arquivo                                
TELESCOP= '0.60m(BC)'          / Telescopio                                     
PLATESCL= '0.34    '           / Escala de Placa ("/pixel)                      
INSTRUME= 'cam 2+IkonL'        / Instrumento                                    
OBJECT  = 'Hektor  '           / Objeto                                         
OBSERVER= 'Bruno Morgado'      / Observador                                     
RDNOISE = '12.1    '           / Ruido de Leitura                               
FILTER  = 'Clear   '           / Filtro                                         
FOCUSVAL= '2117    '           / Valor do Foco                                  
W-BAR   = '760,9   '           / Pressao Atmosferica (mb)                       
W-TEMP  = '15,6    '           / Temperatura Ambiente                           
W-HUM   = '45      '           / Umidade (%)                                    
ANALIST =                                                                       
\end{verbatim}

\begin{Shaded}
\begin{Highlighting}[]
\CommentTok{\# Podemos editar elementos do Header}
\NormalTok{im\_header.}\BuiltInTok{set}\NormalTok{(}\StringTok{\textquotesingle{}ANALIST\textquotesingle{}}\NormalTok{, }\StringTok{\textquotesingle{}\textquotesingle{}}\NormalTok{, comment}\OperatorTok{=}\StringTok{\textquotesingle{}Editado em 27{-}05{-}2022\textquotesingle{}}\NormalTok{)}
\end{Highlighting}
\end{Shaded}

\begin{Shaded}
\begin{Highlighting}[]
\NormalTok{im\_header}
\end{Highlighting}
\end{Shaded}

\begin{verbatim}
SIMPLE  =                    T / file does conform to FITS standard             
BITPIX  =                   16 / number of bits per data pixel                  
NAXIS   =                    2 / number of data axes                            
NAXIS1  =                 2048 / length of data axis 1                          
NAXIS2  =                 2048 / length of data axis 2                          
EXTEND  =                    T / FITS dataset may contain extensions            
COMMENT = 'Imagens para foco'  / Comentario                                     
COMMENT       and Astrophysics', volume 376, page 359; bibcode: 2001A&A...376..3
BZERO   =                32768 / offset data range to that of unsigned short    
BSCALE  =                    1 / default scaling factor                         
HEAD    = '        '           / Head Model                                     
ACQMODE = 'Single  '           / Acquisition Mode                               
READMODE= 'Image   '           / Readout Mode                                   
IMGRECT = '1, 2048, 2048, 1'   / Image Format                                   
HBIN    = '1       '           / Horizontal Binning                             
VBIN    = '1       '           / Vertical Binning                               
SUBRECT = '1, 2048,2048,1'     / Subimage Format                                
DATATYPE= 'Counts  '           / Data type                                      
XTYPE   = 'Pixel number'       / Calibration type                               
XUNIT   =                    0 / Calibration units                              
RAYWAVE =                   0. / Rayleigh Wavelength                            
CALBWVNM=                    0 / Wave calibration                               
TRIGGER = 'Internal'           / Trigger Mode                                   
CALIB   = '0.,1.,0.,0.'        / Calibration                                    
DLLVER  = '2.102.30024.0'      / Software Version                               
EXPOSURE= '10,00000'           / Total Exposure Time                            
TEMP    = '-999    '           / Temperature                                    
READTIME=         3.333333E-07 / Pixel readout time                             
OPERATN =                    4 / Type of system                                 
GAIN    = '1.0     '           / Ganho                                          
EMREALGN=                    0 / EM Real Gain                                   
VCLKAMP =                    0 / Vertical Clock Amplitude                       
VSHIFT  = '76,95   '           / Vertical Shift Speed                           
PREAMP  = '4x      '           / Pre Amplifier Gain                             
SERNO   = '10127   '           / Serial Number                                  
UNSTTEMP= '-63,0   '           / Unstabilized Temperature                       
BLCLAMP =                    F / Baseline Clamp                                 
PRECAN  =                    0 / Prescans                                       
FLIPX   =                    0 / Horizontally Flipped                           
FLIPY   =                    0 / Vertically Flipped                             
CNTCVTMD=                    0 / Count Convert Mode                             
CNTCVT  =                    0 / Count Convert                                  
DTNWLGTH=                 500. / Detection Wavelength                           
SNTVTY  =                   0. / Sensitivity                                    
SPSNFLTR=                    0 / Spurious Noise Filter Mode                     
THRSHLD =                   0. / Threshold                                      
PCNTENLD=                    0 / Photon Counting Enabled                        
NSETHSLD=                    0 / Number of Photon Counting Thresholds           
PTNTHLD1=                   0. / Photon Counting Threshold 1                    
PTNTHLD2=                   0. / Photon Counting Threshold 2                    
PTNTHLD3=                   0. / Photon Counting Threshold 3                    
PTNTHLD4=                   0. / Photon Counting Threshold 4                    
AVGFTRMD=                    0 / Averaging Filter Mode                          
AVGFCTR =                    1 / Averaging factor                               
FRMCNT  =                    1 / Frame Count                                    
PORTMODE=                   -1 / Port Readout Mode                              
LSHEIGHT=                    0 / Exposure Window Height                         
LSSPEED =                   0. / Line Scan Speed                                
LSALTDIR=                    0 / Alternating Readout Direction                  
LSCTRL  =                    0 / Scan Speed Control                             
LSDIR   =                    0 / Readout Direction                              
FKSMODE =                    0 / Fast Kinetics Storage Mode                     
FKTMODE =                    0 / Fast Kinetics Time Scan Mode                   
USERTXT1= '        '           / User text                                      
USERTXT2= '        '           / User text                                      
USERTXT3= '        '           / User text                                      
USERTXT4= '        '           / User text                                      
DATE    = '2019-09-16T23:29:10' / File Creation Date (YYYY-MM-HHThh:mm:ss)      
FRAME   = '2019-09-16T23:29:10.291' / Start of Frame Exposure                   
ESHTMODE=                    0 / Electronic Shuttering Mode                     
HIERARCH PREAMPGAINTEXT = '        ' / Pre-Amplifier Gain                       
HIERARCH SPECTROGRAPHSERIAL = '        ' / Spectrograph Serial                  
HIERARCH SHAMROCKISACTIVE =  0 / Shamrock is active                             
HIERARCH SPECTROGRAPHNAME = '        ' / Spectrograph name                      
HIERARCH SPECTROGRAPHISACTIVE = 0 / Spectrograph is active                      
HIERARCH IRIGDATAAVAILABLE = 0 / IRIG Availability                              
IRIGDATA= '        '           / IRIG Timestamp                                 
HIERARCH PIVENABLE =         0 / PIV Mode Status                                
HIERARCH EXTENDEDDYNAMICRANGE = 0 / Extended Dynamic Range Mode Status          
HIERARCH MILLISECONDS =      0 / Millisecond Count Status                       
HIERARCH DIRECTIRISPRESENT = 0 / Direct Iris Present                            
HIERARCH DIRECTIRISPOS =     0 / Direct Iris Position                           
HIERARCH SIDEIRISPRESENT =   0 / Side Iris Present                              
HIERARCH SIDEIRISPOS =       0 / Side Iris Position                             
DETECTOR= '        '           / Head Model                                     
EXPTIME = '10,00000'           / Total Exposure Time                            
OUTPTAMP= 'High Sensitivity'   / Output Amplifier                               
CAMGAIN = '4x      '           / Pre Amplifier Gain                             
DATE-OBS= '2019-09-16T23:29:10.291' / Start of Frame Exposure                   
RA      = '23:08:39'           / Alfa                                           
DEC     = '-03:08:08'          / Delta                                          
EPOCH   = '2000.0  '           / Epoca                                          
AIRMASS = '1.488   '           / Massa de Ar                                    
JD      = '2458743.47858'      / Data Juliana                                   
ST      = '20:09:16'           / Tempo Sideral                                  
HA      = '-3:00:20'           / Angulo Horario                                 
IMAGE   = 'Hektor_00001'       / Nome do arquivo                                
TELESCOP= '0.60m(BC)'          / Telescopio                                     
PLATESCL= '0.34    '           / Escala de Placa ("/pixel)                      
INSTRUME= 'cam 2+IkonL'        / Instrumento                                    
OBJECT  = 'Hektor  '           / Objeto                                         
OBSERVER= 'Bruno Morgado'      / Observador                                     
RDNOISE = '12.1    '           / Ruido de Leitura                               
FILTER  = 'Clear   '           / Filtro                                         
FOCUSVAL= '2117    '           / Valor do Foco                                  
W-BAR   = '760,9   '           / Pressao Atmosferica (mb)                       
W-TEMP  = '15,6    '           / Temperatura Ambiente                           
W-HUM   = '45      '           / Umidade (%)                                    
ANALIST = '' / Editado em 27-05-2022                                            
\end{verbatim}

\begin{Shaded}
\begin{Highlighting}[]
\NormalTok{im\_header.get(}\StringTok{\textquotesingle{}ANALIST\textquotesingle{}}\NormalTok{)}
\end{Highlighting}
\end{Shaded}

\begin{verbatim}
''
\end{verbatim}

\begin{Shaded}
\begin{Highlighting}[]
\CommentTok{\# Podemos eliminar elementos do Header}
\NormalTok{im\_header.remove(}\StringTok{\textquotesingle{}ANALIST\textquotesingle{}}\NormalTok{)}
\end{Highlighting}
\end{Shaded}

\begin{Shaded}
\begin{Highlighting}[]
\NormalTok{im\_header.get(}\StringTok{\textquotesingle{}ANALIST\textquotesingle{}}\NormalTok{)}
\end{Highlighting}
\end{Shaded}

\begin{Shaded}
\begin{Highlighting}[]
\NormalTok{im\_header}
\end{Highlighting}
\end{Shaded}

\begin{verbatim}
SIMPLE  =                    T / file does conform to FITS standard             
BITPIX  =                   16 / number of bits per data pixel                  
NAXIS   =                    2 / number of data axes                            
NAXIS1  =                 2048 / length of data axis 1                          
NAXIS2  =                 2048 / length of data axis 2                          
EXTEND  =                    T / FITS dataset may contain extensions            
COMMENT = 'Imagens para foco'  / Comentario                                     
COMMENT       and Astrophysics', volume 376, page 359; bibcode: 2001A&A...376..3
BZERO   =                32768 / offset data range to that of unsigned short    
BSCALE  =                    1 / default scaling factor                         
HEAD    = '        '           / Head Model                                     
ACQMODE = 'Single  '           / Acquisition Mode                               
READMODE= 'Image   '           / Readout Mode                                   
IMGRECT = '1, 2048, 2048, 1'   / Image Format                                   
HBIN    = '1       '           / Horizontal Binning                             
VBIN    = '1       '           / Vertical Binning                               
SUBRECT = '1, 2048,2048,1'     / Subimage Format                                
DATATYPE= 'Counts  '           / Data type                                      
XTYPE   = 'Pixel number'       / Calibration type                               
XUNIT   =                    0 / Calibration units                              
RAYWAVE =                   0. / Rayleigh Wavelength                            
CALBWVNM=                    0 / Wave calibration                               
TRIGGER = 'Internal'           / Trigger Mode                                   
CALIB   = '0.,1.,0.,0.'        / Calibration                                    
DLLVER  = '2.102.30024.0'      / Software Version                               
EXPOSURE= '10,00000'           / Total Exposure Time                            
TEMP    = '-999    '           / Temperature                                    
READTIME=         3.333333E-07 / Pixel readout time                             
OPERATN =                    4 / Type of system                                 
GAIN    = '1.0     '           / Ganho                                          
EMREALGN=                    0 / EM Real Gain                                   
VCLKAMP =                    0 / Vertical Clock Amplitude                       
VSHIFT  = '76,95   '           / Vertical Shift Speed                           
PREAMP  = '4x      '           / Pre Amplifier Gain                             
SERNO   = '10127   '           / Serial Number                                  
UNSTTEMP= '-63,0   '           / Unstabilized Temperature                       
BLCLAMP =                    F / Baseline Clamp                                 
PRECAN  =                    0 / Prescans                                       
FLIPX   =                    0 / Horizontally Flipped                           
FLIPY   =                    0 / Vertically Flipped                             
CNTCVTMD=                    0 / Count Convert Mode                             
CNTCVT  =                    0 / Count Convert                                  
DTNWLGTH=                 500. / Detection Wavelength                           
SNTVTY  =                   0. / Sensitivity                                    
SPSNFLTR=                    0 / Spurious Noise Filter Mode                     
THRSHLD =                   0. / Threshold                                      
PCNTENLD=                    0 / Photon Counting Enabled                        
NSETHSLD=                    0 / Number of Photon Counting Thresholds           
PTNTHLD1=                   0. / Photon Counting Threshold 1                    
PTNTHLD2=                   0. / Photon Counting Threshold 2                    
PTNTHLD3=                   0. / Photon Counting Threshold 3                    
PTNTHLD4=                   0. / Photon Counting Threshold 4                    
AVGFTRMD=                    0 / Averaging Filter Mode                          
AVGFCTR =                    1 / Averaging factor                               
FRMCNT  =                    1 / Frame Count                                    
PORTMODE=                   -1 / Port Readout Mode                              
LSHEIGHT=                    0 / Exposure Window Height                         
LSSPEED =                   0. / Line Scan Speed                                
LSALTDIR=                    0 / Alternating Readout Direction                  
LSCTRL  =                    0 / Scan Speed Control                             
LSDIR   =                    0 / Readout Direction                              
FKSMODE =                    0 / Fast Kinetics Storage Mode                     
FKTMODE =                    0 / Fast Kinetics Time Scan Mode                   
USERTXT1= '        '           / User text                                      
USERTXT2= '        '           / User text                                      
USERTXT3= '        '           / User text                                      
USERTXT4= '        '           / User text                                      
DATE    = '2019-09-16T23:29:10' / File Creation Date (YYYY-MM-HHThh:mm:ss)      
FRAME   = '2019-09-16T23:29:10.291' / Start of Frame Exposure                   
ESHTMODE=                    0 / Electronic Shuttering Mode                     
HIERARCH PREAMPGAINTEXT = '        ' / Pre-Amplifier Gain                       
HIERARCH SPECTROGRAPHSERIAL = '        ' / Spectrograph Serial                  
HIERARCH SHAMROCKISACTIVE =  0 / Shamrock is active                             
HIERARCH SPECTROGRAPHNAME = '        ' / Spectrograph name                      
HIERARCH SPECTROGRAPHISACTIVE = 0 / Spectrograph is active                      
HIERARCH IRIGDATAAVAILABLE = 0 / IRIG Availability                              
IRIGDATA= '        '           / IRIG Timestamp                                 
HIERARCH PIVENABLE =         0 / PIV Mode Status                                
HIERARCH EXTENDEDDYNAMICRANGE = 0 / Extended Dynamic Range Mode Status          
HIERARCH MILLISECONDS =      0 / Millisecond Count Status                       
HIERARCH DIRECTIRISPRESENT = 0 / Direct Iris Present                            
HIERARCH DIRECTIRISPOS =     0 / Direct Iris Position                           
HIERARCH SIDEIRISPRESENT =   0 / Side Iris Present                              
HIERARCH SIDEIRISPOS =       0 / Side Iris Position                             
DETECTOR= '        '           / Head Model                                     
EXPTIME = '10,00000'           / Total Exposure Time                            
OUTPTAMP= 'High Sensitivity'   / Output Amplifier                               
CAMGAIN = '4x      '           / Pre Amplifier Gain                             
DATE-OBS= '2019-09-16T23:29:10.291' / Start of Frame Exposure                   
RA      = '23:08:39'           / Alfa                                           
DEC     = '-03:08:08'          / Delta                                          
EPOCH   = '2000.0  '           / Epoca                                          
AIRMASS = '1.488   '           / Massa de Ar                                    
JD      = '2458743.47858'      / Data Juliana                                   
ST      = '20:09:16'           / Tempo Sideral                                  
HA      = '-3:00:20'           / Angulo Horario                                 
IMAGE   = 'Hektor_00001'       / Nome do arquivo                                
TELESCOP= '0.60m(BC)'          / Telescopio                                     
PLATESCL= '0.34    '           / Escala de Placa ("/pixel)                      
INSTRUME= 'cam 2+IkonL'        / Instrumento                                    
OBJECT  = 'Hektor  '           / Objeto                                         
OBSERVER= 'Bruno Morgado'      / Observador                                     
RDNOISE = '12.1    '           / Ruido de Leitura                               
FILTER  = 'Clear   '           / Filtro                                         
FOCUSVAL= '2117    '           / Valor do Foco                                  
W-BAR   = '760,9   '           / Pressao Atmosferica (mb)                       
W-TEMP  = '15,6    '           / Temperatura Ambiente                           
W-HUM   = '45      '           / Umidade (%)                                    
\end{verbatim}

\section{4) Lidando com os Dados}\label{lidando-com-os-dados}

\begin{Shaded}
\begin{Highlighting}[]
\NormalTok{dados }\OperatorTok{=}\NormalTok{ imagem.data.astype(}\BuiltInTok{float}\NormalTok{)}
\end{Highlighting}
\end{Shaded}

\begin{Shaded}
\begin{Highlighting}[]
\NormalTok{dados}
\end{Highlighting}
\end{Shaded}

\begin{verbatim}
array([[2299., 2253., 2317., ..., 2296., 2303., 2277.],
       [2270., 2298., 2285., ..., 2301., 2301., 2301.],
       [2290., 2259., 2312., ..., 2289., 2321., 2304.],
       ...,
       [2308., 2291., 2269., ..., 2342., 2325., 2325.],
       [2291., 2275., 2309., ..., 2345., 2299., 2333.],
       [2283., 2298., 2275., ..., 2327., 2299., 2355.]],
      shape=(2048, 2048))
\end{verbatim}

\begin{Shaded}
\begin{Highlighting}[]
\ImportTok{import}\NormalTok{ numpy }\ImportTok{as}\NormalTok{ np}
\ImportTok{import}\NormalTok{ matplotlib.pylab }\ImportTok{as}\NormalTok{ pl}
\end{Highlighting}
\end{Shaded}

\begin{Shaded}
\begin{Highlighting}[]
\CommentTok{\# Podemos plotar a imagem}
\NormalTok{pl.figure(figsize}\OperatorTok{=}\NormalTok{[}\DecValTok{5}\NormalTok{,}\DecValTok{5}\NormalTok{])}

\NormalTok{pl.imshow(dados, cmap}\OperatorTok{=}\StringTok{\textquotesingle{}gray\textquotesingle{}}\NormalTok{)}

\NormalTok{pl.show()}
\end{Highlighting}
\end{Shaded}

\pandocbounded{\includegraphics[keepaspectratio]{astropy/astropy_files/figure-pdf/cell-103-output-1.pdf}}

\begin{Shaded}
\begin{Highlighting}[]
\NormalTok{dados.shape}
\end{Highlighting}
\end{Shaded}

\begin{verbatim}
(2048, 2048)
\end{verbatim}

\begin{Shaded}
\begin{Highlighting}[]
\CommentTok{\# Podemos controlar as intensidades}
\NormalTok{pl.figure(figsize}\OperatorTok{=}\NormalTok{[}\DecValTok{5}\NormalTok{,}\DecValTok{5}\NormalTok{], facecolor}\OperatorTok{=}\StringTok{\textquotesingle{}w\textquotesingle{}}\NormalTok{)}

\NormalTok{pl.imshow(dados, cmap}\OperatorTok{=}\StringTok{\textquotesingle{}gray\textquotesingle{}}\NormalTok{, vmin}\OperatorTok{=}\NormalTok{dados.mean() }\OperatorTok{{-}} \DecValTok{50}\NormalTok{, vmax}\OperatorTok{=}\NormalTok{dados.mean() }\OperatorTok{+}\DecValTok{50}\NormalTok{)}
\CommentTok{\#pl.show()}
\NormalTok{pl.savefig(}\StringTok{\textquotesingle{}imagem.png\textquotesingle{}}\NormalTok{)}
\end{Highlighting}
\end{Shaded}

\pandocbounded{\includegraphics[keepaspectratio]{astropy/astropy_files/figure-pdf/cell-105-output-1.pdf}}

\begin{Shaded}
\begin{Highlighting}[]
\CommentTok{\# Podemos controlar o colormap}
\NormalTok{pl.figure(figsize}\OperatorTok{=}\NormalTok{[}\DecValTok{5}\NormalTok{,}\DecValTok{5}\NormalTok{])}

\CommentTok{\#pl.imshow(dados, cmap=\textquotesingle{}viridis\textquotesingle{}, vmin=2300, vmax=2400)}
\NormalTok{pl.imshow(dados, cmap}\OperatorTok{=}\StringTok{\textquotesingle{}viridis\_r\textquotesingle{}}\NormalTok{, vmin}\OperatorTok{=}\DecValTok{2300}\NormalTok{, vmax}\OperatorTok{=}\DecValTok{2400}\NormalTok{)}

\NormalTok{pl.show()}

\CommentTok{\#\#\# https://matplotlib.org/3.5.0/tutorials/colors/colormaps.html}
\end{Highlighting}
\end{Shaded}

\pandocbounded{\includegraphics[keepaspectratio]{astropy/astropy_files/figure-pdf/cell-106-output-1.pdf}}

\begin{Shaded}
\begin{Highlighting}[]
\CommentTok{\# Como parte do matplotlib podemos combinar os diferentes plotes}

\NormalTok{pl.figure(figsize}\OperatorTok{=}\NormalTok{[}\DecValTok{10}\NormalTok{,}\DecValTok{10}\NormalTok{])}

\NormalTok{pl.imshow(dados, cmap}\OperatorTok{=}\StringTok{\textquotesingle{}viridis\textquotesingle{}}\NormalTok{, vmin}\OperatorTok{=}\DecValTok{2300}\NormalTok{, vmax}\OperatorTok{=}\DecValTok{2400}\NormalTok{)}

\NormalTok{pl.axhline(}\DecValTok{1024}\NormalTok{, color}\OperatorTok{=}\StringTok{\textquotesingle{}k\textquotesingle{}}\NormalTok{, ls}\OperatorTok{=}\StringTok{\textquotesingle{}{-}{-}\textquotesingle{}}\NormalTok{)}
\NormalTok{pl.axvline(}\DecValTok{1024}\NormalTok{, color}\OperatorTok{=}\StringTok{\textquotesingle{}r\textquotesingle{}}\NormalTok{, ls}\OperatorTok{=}\StringTok{\textquotesingle{}{-}{-}\textquotesingle{}}\NormalTok{)}

\CommentTok{\#pl.imshow(im\_data, cmap=\textquotesingle{}viridis\_r\textquotesingle{}, vmin=2300, vmax=2400)}

\NormalTok{pl.show()}
\end{Highlighting}
\end{Shaded}

\pandocbounded{\includegraphics[keepaspectratio]{astropy/astropy_files/figure-pdf/cell-107-output-1.pdf}}

\chapter{Operações básicas}\label{operauxe7uxf5es-buxe1sicas-1}

\begin{Shaded}
\begin{Highlighting}[]
\NormalTok{dados}
\end{Highlighting}
\end{Shaded}

\begin{verbatim}
array([[2299., 2253., 2317., ..., 2296., 2303., 2277.],
       [2270., 2298., 2285., ..., 2301., 2301., 2301.],
       [2290., 2259., 2312., ..., 2289., 2321., 2304.],
       ...,
       [2308., 2291., 2269., ..., 2342., 2325., 2325.],
       [2291., 2275., 2309., ..., 2345., 2299., 2333.],
       [2283., 2298., 2275., ..., 2327., 2299., 2355.]],
      shape=(2048, 2048))
\end{verbatim}

\begin{Shaded}
\begin{Highlighting}[]
\NormalTok{dados }\OperatorTok{+} \DecValTok{2000}
\end{Highlighting}
\end{Shaded}

\begin{verbatim}
array([[4299., 4253., 4317., ..., 4296., 4303., 4277.],
       [4270., 4298., 4285., ..., 4301., 4301., 4301.],
       [4290., 4259., 4312., ..., 4289., 4321., 4304.],
       ...,
       [4308., 4291., 4269., ..., 4342., 4325., 4325.],
       [4291., 4275., 4309., ..., 4345., 4299., 4333.],
       [4283., 4298., 4275., ..., 4327., 4299., 4355.]],
      shape=(2048, 2048))
\end{verbatim}

\begin{Shaded}
\begin{Highlighting}[]
\NormalTok{dados }\OperatorTok{{-}} \DecValTok{2000}
\end{Highlighting}
\end{Shaded}

\begin{verbatim}
array([[299., 253., 317., ..., 296., 303., 277.],
       [270., 298., 285., ..., 301., 301., 301.],
       [290., 259., 312., ..., 289., 321., 304.],
       ...,
       [308., 291., 269., ..., 342., 325., 325.],
       [291., 275., 309., ..., 345., 299., 333.],
       [283., 298., 275., ..., 327., 299., 355.]], shape=(2048, 2048))
\end{verbatim}

\begin{Shaded}
\begin{Highlighting}[]
\NormalTok{dados }\OperatorTok{*} \DecValTok{2}
\end{Highlighting}
\end{Shaded}

\begin{verbatim}
array([[4598., 4506., 4634., ..., 4592., 4606., 4554.],
       [4540., 4596., 4570., ..., 4602., 4602., 4602.],
       [4580., 4518., 4624., ..., 4578., 4642., 4608.],
       ...,
       [4616., 4582., 4538., ..., 4684., 4650., 4650.],
       [4582., 4550., 4618., ..., 4690., 4598., 4666.],
       [4566., 4596., 4550., ..., 4654., 4598., 4710.]],
      shape=(2048, 2048))
\end{verbatim}

\begin{Shaded}
\begin{Highlighting}[]
\NormalTok{dados }\OperatorTok{/} \DecValTok{2}
\end{Highlighting}
\end{Shaded}

\begin{verbatim}
array([[1149.5, 1126.5, 1158.5, ..., 1148. , 1151.5, 1138.5],
       [1135. , 1149. , 1142.5, ..., 1150.5, 1150.5, 1150.5],
       [1145. , 1129.5, 1156. , ..., 1144.5, 1160.5, 1152. ],
       ...,
       [1154. , 1145.5, 1134.5, ..., 1171. , 1162.5, 1162.5],
       [1145.5, 1137.5, 1154.5, ..., 1172.5, 1149.5, 1166.5],
       [1141.5, 1149. , 1137.5, ..., 1163.5, 1149.5, 1177.5]],
      shape=(2048, 2048))
\end{verbatim}

\chapter{Lendo apenas um linha}\label{lendo-apenas-um-linha}

\begin{Shaded}
\begin{Highlighting}[]
\NormalTok{dados[}\DecValTok{1024}\NormalTok{]}
\end{Highlighting}
\end{Shaded}

\begin{verbatim}
array([2330., 2373., 2339., ..., 2336., 2349., 2305.], shape=(2048,))
\end{verbatim}

\begin{Shaded}
\begin{Highlighting}[]
\NormalTok{pl.figure(figsize}\OperatorTok{=}\NormalTok{[}\DecValTok{10}\NormalTok{,}\DecValTok{3}\NormalTok{])}
\NormalTok{pl.plot(dados[}\DecValTok{790}\NormalTok{], }\StringTok{\textquotesingle{}k.{-}\textquotesingle{}}\NormalTok{)}
\end{Highlighting}
\end{Shaded}

\pandocbounded{\includegraphics[keepaspectratio]{astropy/astropy_files/figure-pdf/cell-114-output-1.pdf}}

\chapter{Lendo apenas uma coluna}\label{lendo-apenas-uma-coluna}

\begin{Shaded}
\begin{Highlighting}[]
\NormalTok{dados[:,}\DecValTok{1230}\NormalTok{]}
\end{Highlighting}
\end{Shaded}

\begin{verbatim}
array([2358., 2331., 2327., ..., 2334., 2353., 2342.], shape=(2048,))
\end{verbatim}

\begin{Shaded}
\begin{Highlighting}[]
\NormalTok{pl.figure(figsize}\OperatorTok{=}\NormalTok{[}\DecValTok{15}\NormalTok{,}\DecValTok{3}\NormalTok{])}
\NormalTok{pl.plot(dados[:,}\DecValTok{1230}\NormalTok{], }\StringTok{\textquotesingle{}k.{-}\textquotesingle{}}\NormalTok{)}
\end{Highlighting}
\end{Shaded}

\pandocbounded{\includegraphics[keepaspectratio]{astropy/astropy_files/figure-pdf/cell-116-output-1.pdf}}

\subsubsection{Relembrando slices}\label{relembrando-slices}

\begin{Shaded}
\begin{Highlighting}[]
\NormalTok{a }\OperatorTok{=}\NormalTok{ np.arange(}\DecValTok{0}\NormalTok{,}\DecValTok{10}\NormalTok{,}\DecValTok{1}\NormalTok{)}
\end{Highlighting}
\end{Shaded}

\begin{Shaded}
\begin{Highlighting}[]
\NormalTok{a}
\end{Highlighting}
\end{Shaded}

\begin{verbatim}
array([0, 1, 2, 3, 4, 5, 6, 7, 8, 9])
\end{verbatim}

\begin{Shaded}
\begin{Highlighting}[]
\NormalTok{a[}\DecValTok{0}\NormalTok{:}\DecValTok{4}\NormalTok{]}
\end{Highlighting}
\end{Shaded}

\begin{verbatim}
array([0, 1, 2, 3])
\end{verbatim}

\subsubsection{Slices em 2D}\label{slices-em-2d}

\begin{Shaded}
\begin{Highlighting}[]
\NormalTok{b, c }\OperatorTok{=}\NormalTok{ np.meshgrid(np.arange(}\DecValTok{0}\NormalTok{,}\DecValTok{10}\NormalTok{,}\DecValTok{1}\NormalTok{), np.arange(}\DecValTok{100}\NormalTok{,}\DecValTok{110}\NormalTok{,}\DecValTok{1}\NormalTok{))}

\NormalTok{d }\OperatorTok{=} \DecValTok{2}\OperatorTok{*}\NormalTok{b }\OperatorTok{+}\NormalTok{ c}
\end{Highlighting}
\end{Shaded}

\begin{Shaded}
\begin{Highlighting}[]
\NormalTok{d}
\end{Highlighting}
\end{Shaded}

\begin{verbatim}
array([[100, 102, 104, 106, 108, 110, 112, 114, 116, 118],
       [101, 103, 105, 107, 109, 111, 113, 115, 117, 119],
       [102, 104, 106, 108, 110, 112, 114, 116, 118, 120],
       [103, 105, 107, 109, 111, 113, 115, 117, 119, 121],
       [104, 106, 108, 110, 112, 114, 116, 118, 120, 122],
       [105, 107, 109, 111, 113, 115, 117, 119, 121, 123],
       [106, 108, 110, 112, 114, 116, 118, 120, 122, 124],
       [107, 109, 111, 113, 115, 117, 119, 121, 123, 125],
       [108, 110, 112, 114, 116, 118, 120, 122, 124, 126],
       [109, 111, 113, 115, 117, 119, 121, 123, 125, 127]])
\end{verbatim}

\begin{Shaded}
\begin{Highlighting}[]
\NormalTok{d[:,}\DecValTok{0}\NormalTok{]}
\end{Highlighting}
\end{Shaded}

\begin{verbatim}
array([100, 101, 102, 103, 104, 105, 106, 107, 108, 109])
\end{verbatim}

\begin{Shaded}
\begin{Highlighting}[]
\NormalTok{d[}\DecValTok{0}\NormalTok{:}\DecValTok{2}\NormalTok{]}
\end{Highlighting}
\end{Shaded}

\begin{verbatim}
array([[100, 102, 104, 106, 108, 110, 112, 114, 116, 118],
       [101, 103, 105, 107, 109, 111, 113, 115, 117, 119]])
\end{verbatim}

\begin{Shaded}
\begin{Highlighting}[]
\NormalTok{d[:,}\DecValTok{0}\NormalTok{:}\DecValTok{2}\NormalTok{]}
\end{Highlighting}
\end{Shaded}

\begin{verbatim}
array([[100, 102],
       [101, 103],
       [102, 104],
       [103, 105],
       [104, 106],
       [105, 107],
       [106, 108],
       [107, 109],
       [108, 110],
       [109, 111]])
\end{verbatim}

\begin{Shaded}
\begin{Highlighting}[]
\NormalTok{d[}\DecValTok{0}\NormalTok{:}\DecValTok{2}\NormalTok{, }\DecValTok{0}\NormalTok{:}\DecValTok{2}\NormalTok{]}
\end{Highlighting}
\end{Shaded}

\begin{verbatim}
array([[100, 102],
       [101, 103]])
\end{verbatim}

\begin{Shaded}
\begin{Highlighting}[]
\NormalTok{d[d }\OperatorTok{\textgreater{}} \DecValTok{110}\NormalTok{]}
\end{Highlighting}
\end{Shaded}

\begin{verbatim}
array([112, 114, 116, 118, 111, 113, 115, 117, 119, 112, 114, 116, 118,
       120, 111, 113, 115, 117, 119, 121, 112, 114, 116, 118, 120, 122,
       111, 113, 115, 117, 119, 121, 123, 112, 114, 116, 118, 120, 122,
       124, 111, 113, 115, 117, 119, 121, 123, 125, 112, 114, 116, 118,
       120, 122, 124, 126, 111, 113, 115, 117, 119, 121, 123, 125, 127])
\end{verbatim}

\chapter{Aplicando a imagens}\label{aplicando-a-imagens}

\begin{Shaded}
\begin{Highlighting}[]
\CommentTok{\# Podemos controlar as intensidades}
\NormalTok{pl.figure(figsize}\OperatorTok{=}\NormalTok{[}\DecValTok{5}\NormalTok{,}\DecValTok{5}\NormalTok{])}

\NormalTok{pl.imshow(dados, cmap}\OperatorTok{=}\StringTok{\textquotesingle{}gray\textquotesingle{}}\NormalTok{, vmin}\OperatorTok{=}\DecValTok{2300}\NormalTok{, vmax}\OperatorTok{=}\DecValTok{2400}\NormalTok{)}

\NormalTok{pl.show()}
\end{Highlighting}
\end{Shaded}

\pandocbounded{\includegraphics[keepaspectratio]{astropy/astropy_files/figure-pdf/cell-127-output-1.pdf}}

\begin{Shaded}
\begin{Highlighting}[]
\CommentTok{\# Podemos controlar as intensidades}
\NormalTok{pl.figure(figsize}\OperatorTok{=}\NormalTok{[}\DecValTok{10}\NormalTok{,}\DecValTok{10}\NormalTok{])}

\NormalTok{pl.imshow(dados[}\DecValTok{790}\OperatorTok{{-}}\DecValTok{50}\NormalTok{:}\DecValTok{790}\OperatorTok{+}\DecValTok{50}\NormalTok{], cmap}\OperatorTok{=}\StringTok{\textquotesingle{}gray\textquotesingle{}}\NormalTok{, vmin}\OperatorTok{=}\DecValTok{2300}\NormalTok{, vmax}\OperatorTok{=}\DecValTok{2400}\NormalTok{)}

\NormalTok{pl.show()}
\end{Highlighting}
\end{Shaded}

\pandocbounded{\includegraphics[keepaspectratio]{astropy/astropy_files/figure-pdf/cell-128-output-1.pdf}}

\begin{Shaded}
\begin{Highlighting}[]
\CommentTok{\# Podemos controlar as intensidades}
\NormalTok{pl.figure(figsize}\OperatorTok{=}\NormalTok{[}\DecValTok{10}\NormalTok{,}\DecValTok{10}\NormalTok{])}

\NormalTok{pl.imshow(dados[:,}\DecValTok{1230}\OperatorTok{{-}}\DecValTok{50}\NormalTok{:}\DecValTok{1230}\OperatorTok{+}\DecValTok{50}\NormalTok{], cmap}\OperatorTok{=}\StringTok{\textquotesingle{}gray\textquotesingle{}}\NormalTok{, vmin}\OperatorTok{=}\DecValTok{2300}\NormalTok{, vmax}\OperatorTok{=}\DecValTok{2400}\NormalTok{)}

\NormalTok{pl.show()}
\end{Highlighting}
\end{Shaded}

\pandocbounded{\includegraphics[keepaspectratio]{astropy/astropy_files/figure-pdf/cell-129-output-1.pdf}}

\begin{Shaded}
\begin{Highlighting}[]
\CommentTok{\# Podemos controlar as intensidades}
\NormalTok{pl.figure(figsize}\OperatorTok{=}\NormalTok{[}\DecValTok{5}\NormalTok{,}\DecValTok{5}\NormalTok{])}

\NormalTok{pl.imshow(dados[}\DecValTok{790}\OperatorTok{{-}}\DecValTok{50}\NormalTok{:}\DecValTok{790}\OperatorTok{+}\DecValTok{50}\NormalTok{, }\DecValTok{1230}\OperatorTok{{-}}\DecValTok{50}\NormalTok{:}\DecValTok{1230}\OperatorTok{+}\DecValTok{50}\NormalTok{], cmap}\OperatorTok{=}\StringTok{\textquotesingle{}gray\textquotesingle{}}\NormalTok{, vmin}\OperatorTok{=}\DecValTok{2300}\NormalTok{, vmax}\OperatorTok{=}\DecValTok{2400}\NormalTok{)}

\NormalTok{pl.axvline(}\DecValTok{50}\NormalTok{)}
\NormalTok{pl.axhline(}\DecValTok{60}\NormalTok{, color}\OperatorTok{=}\StringTok{\textquotesingle{}y\textquotesingle{}}\NormalTok{)}
\CommentTok{\#pl.ylim(0, 100)}
\NormalTok{pl.show()}
\end{Highlighting}
\end{Shaded}

\pandocbounded{\includegraphics[keepaspectratio]{astropy/astropy_files/figure-pdf/cell-130-output-1.pdf}}

\begin{Shaded}
\begin{Highlighting}[]
\NormalTok{new\_dado }\OperatorTok{=}\NormalTok{ dados[}\DecValTok{790}\OperatorTok{{-}}\DecValTok{50}\NormalTok{:}\DecValTok{790}\OperatorTok{+}\DecValTok{50}\NormalTok{, }\DecValTok{1230}\OperatorTok{{-}}\DecValTok{50}\NormalTok{:}\DecValTok{1230}\OperatorTok{+}\DecValTok{50}\NormalTok{]}
\end{Highlighting}
\end{Shaded}

\begin{Shaded}
\begin{Highlighting}[]
\NormalTok{dados[}\DecValTok{790}\OperatorTok{{-}}\DecValTok{50}\NormalTok{:}\DecValTok{790}\OperatorTok{+}\DecValTok{50}\NormalTok{, }\DecValTok{1230}\OperatorTok{{-}}\DecValTok{50}\NormalTok{:}\DecValTok{1230}\OperatorTok{+}\DecValTok{50}\NormalTok{].shape}
\end{Highlighting}
\end{Shaded}

\begin{verbatim}
(100, 100)
\end{verbatim}

\begin{Shaded}
\begin{Highlighting}[]
\CommentTok{\# Lendo só a linha}
\NormalTok{pl.figure(figsize}\OperatorTok{=}\NormalTok{[}\DecValTok{15}\NormalTok{,}\DecValTok{3}\NormalTok{])}
\NormalTok{pl.plot(new\_dado[}\DecValTok{50}\NormalTok{], }\StringTok{\textquotesingle{}b+{-}\textquotesingle{}}\NormalTok{)}
\end{Highlighting}
\end{Shaded}

\pandocbounded{\includegraphics[keepaspectratio]{astropy/astropy_files/figure-pdf/cell-133-output-1.pdf}}

\begin{Shaded}
\begin{Highlighting}[]
\CommentTok{\# Lendo só a coluna}
\NormalTok{pl.figure(figsize}\OperatorTok{=}\NormalTok{[}\DecValTok{15}\NormalTok{,}\DecValTok{3}\NormalTok{])}
\NormalTok{pl.plot(new\_dado[:,}\DecValTok{50}\NormalTok{], }\StringTok{\textquotesingle{}k.{-}\textquotesingle{}}\NormalTok{)}
\end{Highlighting}
\end{Shaded}

\pandocbounded{\includegraphics[keepaspectratio]{astropy/astropy_files/figure-pdf/cell-134-output-1.pdf}}

\begin{Shaded}
\begin{Highlighting}[]
\NormalTok{fits.writeto(}\StringTok{\textquotesingle{}new\_fig.fits\textquotesingle{}}\NormalTok{, new\_dado, overwrite}\OperatorTok{=}\VariableTok{True}\NormalTok{)}
\end{Highlighting}
\end{Shaded}

\chapter{Isso é tudo}\label{isso-uxe9-tudo-2}

By BEM

\chapter{Astroquery}\label{astroquery}

\subsubsection{by BEM, Last update:
19/05/2023}\label{by-bem-last-update-19052023}

Sumário: 1. \hyperref[pt1]{Buscando um catalogo} 2.
\hyperref[pt2]{Buscando um objeto} 3. \hyperref[pt3]{Buscando uma
região} 4. \hyperref[pt4]{Lidando com os Dados}

\textbf{Mais detalhes na documentação oficial do Astroquery
(https://astroquery.readthedocs.io/)}

\section{1) Buscando um catalogo}\label{buscando-um-catalogo}

\begin{Shaded}
\begin{Highlighting}[]
\ImportTok{import}\NormalTok{ numpy }\ImportTok{as}\NormalTok{ np}
\ImportTok{import}\NormalTok{ pandas }\ImportTok{as}\NormalTok{ pd}
\ImportTok{import}\NormalTok{ matplotlib.pyplot }\ImportTok{as}\NormalTok{ pl}
\end{Highlighting}
\end{Shaded}

\begin{Shaded}
\begin{Highlighting}[]
\CommentTok{\# $ pip install astroquery}

\ImportTok{from}\NormalTok{ astroquery.vizier }\ImportTok{import}\NormalTok{ Vizier}
\end{Highlighting}
\end{Shaded}

\begin{Shaded}
\begin{Highlighting}[]
\NormalTok{Vizier.find\_catalogs?}
\end{Highlighting}
\end{Shaded}

\begin{Shaded}
\begin{Highlighting}[]
\NormalTok{catalog\_list }\OperatorTok{=}\NormalTok{ Vizier.find\_catalogs(}\StringTok{\textquotesingle{}NOMAD\textquotesingle{}}\NormalTok{)}

\BuiltInTok{print}\NormalTok{(catalog\_list)}
\end{Highlighting}
\end{Shaded}

\begin{verbatim}
OrderedDict({'I/297/out': </>})
\end{verbatim}

\begin{Shaded}
\begin{Highlighting}[]
\BuiltInTok{type}\NormalTok{(catalog\_list)}
\end{Highlighting}
\end{Shaded}

\begin{verbatim}
collections.OrderedDict
\end{verbatim}

\begin{Shaded}
\begin{Highlighting}[]
\ControlFlowTok{for}\NormalTok{ k,v }\KeywordTok{in}\NormalTok{ catalog\_list.items():}
    \BuiltInTok{print}\NormalTok{(k, v.description)}
\end{Highlighting}
\end{Shaded}

\begin{verbatim}
I/297/out NOMAD Catalog (Zacharias+ 2005)
    Example of the an output of NOMAD-1
\end{verbatim}

\section{2) Buscando um objeto}\label{buscando-um-objeto}

\begin{Shaded}
\begin{Highlighting}[]
\NormalTok{result }\OperatorTok{=}\NormalTok{ Vizier.query\_object(}\StringTok{\textquotesingle{}NGC322\textquotesingle{}}\NormalTok{, catalog}\OperatorTok{=}\StringTok{\textquotesingle{}Gaia\textquotesingle{}}\NormalTok{)}

\BuiltInTok{print}\NormalTok{(result)}
\end{Highlighting}
\end{Shaded}

\begin{verbatim}
TableList with 23 tables:
    '0:I/324/igsl3' with 18 column(s) and 33 row(s) 
    '1:I/337/gaia' with 14 column(s) and 11 row(s) 
    '2:I/345/gaia2' with 32 column(s) and 10 row(s) 
    '3:I/347/gaia2dis' with 9 column(s) and 9 row(s) 
    '4:I/350/gaiaedr3' with 39 column(s) and 16 row(s) 
    '5:I/350/comscanl' with 13 column(s) and 3 row(s) 
    '6:I/352/gedr3dis' with 10 column(s) and 9 row(s) 
    '7:I/355/gaiadr3' with 57 column(s) and 16 row(s) 
    '8:I/355/paramp' with 24 column(s) and 16 row(s) 
    '9:I/355/paramsup' with 24 column(s) and 5 row(s) 
    '10:I/355/epphot' with 48 column(s) and 50 row(s) 
    '11:I/355/xpsummary' with 22 column(s) and 3 row(s) 
    '12:I/355/xpsample' with 6 column(s) and 50 row(s) 
    '13:I/355/xpcont' with 26 column(s) and 3 row(s) 
    '14:I/355/spectra' with 5 column(s) and 1 row(s) 
    '15:I/356/galcand' with 10 column(s) and 3 row(s) 
    '16:I/358/varisum' with 24 column(s) and 2 row(s) 
    '17:I/358/vclassre' with 6 column(s) and 2 row(s) 
    '18:I/360/syntphot' with 5 column(s) and 3 row(s) 
    '19:VI/137/gum_mw' with 17 column(s) and 22 row(s) 
    '20:VI/137/gum_gal' with 11 column(s) and 5 row(s) 
    '21:VI/145/attitude' with 18 column(s) and 2 row(s) 
    '22:J/A+A/674/A25/vspursig' with 40 column(s) and 2 row(s) 
\end{verbatim}

\begin{Shaded}
\begin{Highlighting}[]
\CommentTok{\# Acessando o catálogo de interesse, usando o indice}

\NormalTok{interesting\_table }\OperatorTok{=}\NormalTok{ result[}\DecValTok{7}\NormalTok{]}
\BuiltInTok{print}\NormalTok{(interesting\_table)}
\end{Highlighting}
\end{Shaded}

\begin{verbatim}
    RA_ICRS         DE_ICRS     ...     RAJ2000         DEJ2000    
      deg             deg       ...       deg             deg      
--------------- --------------- ... --------------- ---------------
 14.29273425407 -43.74559894270 ...  14.29286571292 -43.74554184280
 14.31429408969 -43.73940306075 ...  14.31417173601 -43.73939338952
 14.31899705927 -43.73169635937 ...  14.31894560169 -43.73156367775
 14.29571274662 -43.73317129831 ...  14.29571274662 -43.73317129831
 14.31877141341 -43.71781424104 ...  14.31873752752 -43.71779346020
 14.26853881591 -43.74216032776 ...  14.26837371808 -43.74210813273
 14.27464683852 -43.74076203091 ...  14.27464683852 -43.74076203091
 14.28813176441 -43.72978841053 ...  14.28813176441 -43.72978841053
 14.25026544433 -43.73389069171 ...  14.25028067512 -43.73385371996
 14.27238203528 -43.72041368068 ...  14.27225260354 -43.72051765787
 14.30604617647 -43.70841759485 ...  14.30604617647 -43.70841759485
 14.29176960133 -43.72706020178 ...  14.29176960133 -43.72706020178
 14.31396524408 -43.70267505024 ...  14.31396524408 -43.70267505024
 14.25962162355 -43.71234968924 ...  14.25961078061 -43.71232167401
 14.25433773055 -43.72243602775 ...  14.25433773055 -43.72243602775
 14.27873131081 -43.70595833204 ...  14.27855093648 -43.70595286209
\end{verbatim}

\begin{Shaded}
\begin{Highlighting}[]
\CommentTok{\# Acessando o catálogo de interesse, usando o identificador}

\NormalTok{interesting\_table }\OperatorTok{=}\NormalTok{ result[}\StringTok{\textquotesingle{}I/358/vclassre\textquotesingle{}}\NormalTok{]}
\BuiltInTok{print}\NormalTok{(interesting\_table)}
\end{Highlighting}
\end{Shaded}

\begin{verbatim}
       Source             SolID        ...     RA_ICRS         DE_ICRS    
                                       ...       deg             deg      
------------------- ------------------ ... --------------- ---------------
4982662345448240384 375316653866487564 ...  14.26853881591 -43.74216032776
4982662654685883520 375316653866487564 ...  14.27873131081 -43.70595833204
\end{verbatim}

\chapter{Buscando Keywords especificas de um catalogo e aplicando
filtros}\label{buscando-keywords-especificas-de-um-catalogo-e-aplicando-filtros}

\begin{Shaded}
\begin{Highlighting}[]
\NormalTok{v }\OperatorTok{=}\NormalTok{ Vizier(catalog}\OperatorTok{=}\StringTok{\textquotesingle{}I/297/out\textquotesingle{}}\NormalTok{, keywords}\OperatorTok{=}\NormalTok{[}\StringTok{\textquotesingle{}optical\textquotesingle{}}\NormalTok{])}

\NormalTok{result }\OperatorTok{=}\NormalTok{ v.query\_object(}\StringTok{\textquotesingle{}HD 226868\textquotesingle{}}\NormalTok{)}
\BuiltInTok{print}\NormalTok{(result)}
\end{Highlighting}
\end{Shaded}

\begin{verbatim}
TableList with 1 tables:
    '0:I/297/out' with 19 column(s) and 50 row(s) 
\end{verbatim}

\begin{Shaded}
\begin{Highlighting}[]
\BuiltInTok{type}\NormalTok{(result)}
\end{Highlighting}
\end{Shaded}

\begin{verbatim}
astroquery.utils.commons.TableList
\end{verbatim}

\begin{Shaded}
\begin{Highlighting}[]
\BuiltInTok{print}\NormalTok{(result[}\StringTok{\textquotesingle{}I/297/out\textquotesingle{}}\NormalTok{])}
\end{Highlighting}
\end{Shaded}

\begin{verbatim}
   NOMAD1      RAJ2000     DEJ2000    YM  r  ... r_Rmag  Jmag   Hmag   Kmag   R 
                 deg         deg             ...         mag    mag    mag      
------------ ----------- ----------- --- --- ... ------ ------ ------ ------ ---
1251-0389526 299.5704167  35.1752833  YM   B ...      B 13.657 12.775 12.420    
1251-0389529 299.5713194  35.1726750  .M   B ...      B 15.955 15.490 15.276    
1251-0389557 299.5756408  35.1774939  .M   M ...        15.789 14.757 14.566    
1251-0389568 299.5775083  35.1791417  .M   B ...      B 14.426 13.493 13.153    
1251-0389573 299.5782069  35.1764789  .M   M ...        16.146 14.342 15.153    
1251-0389597 299.5820333  35.1694750  .M   B ...      e 16.242 15.766 15.381    
1251-0389841 299.6278306  35.1904806  .M   B ...      B 15.349 14.455 14.214    
1251-0389558 299.5760889  35.1821167  .M   B ...      B 15.910 15.362 15.246    
1251-0389588 299.5802878  35.1768889  YM   C ...      B 11.691 11.412 11.297    
         ...         ...         ... ... ... ...    ...    ...    ...    ... ...
1251-0389700 299.6013947  35.1981117  YM   C ...      B 12.582 12.360 12.166    
1251-0389706 299.6020739  35.1816900  .M   M ...        15.586 14.509 14.144    
1251-0389714 299.6028944  35.1723444  .M   B ...      e 16.623 15.750 15.294    
1251-0389715 299.6031111  35.1943389  .M   B ...      B 15.835 15.400 15.362    
1251-0389716 299.6033972  35.1937806  .M   B ...      B 16.524 14.736 14.669    
1251-0389719 299.6041528  35.1826667  .M   B ...      B 15.640 14.869 14.736    
1251-0389720 299.6041547  35.1844178  .M   M ...        16.855 15.705 15.900    
1251-0389721 299.6041944  35.1815667  .M   B ...      B 16.150 15.564 14.461    
1251-0389722 299.6043361  35.1914000  ..   B ...            --     --     --    
1251-0389728 299.6053128  35.1957969  .M   M ...        16.448 15.839 15.647    
Length = 50 rows
\end{verbatim}

\begin{Shaded}
\begin{Highlighting}[]
\NormalTok{df }\OperatorTok{=}\NormalTok{ result[}\StringTok{\textquotesingle{}I/297/out\textquotesingle{}}\NormalTok{]}
\end{Highlighting}
\end{Shaded}

\begin{Shaded}
\begin{Highlighting}[]
\NormalTok{df}
\end{Highlighting}
\end{Shaded}

\begin{Shaded}
\begin{Highlighting}[]
\BuiltInTok{type}\NormalTok{(df)}
\end{Highlighting}
\end{Shaded}

\begin{verbatim}
astropy.table.table.Table
\end{verbatim}

\begin{Shaded}
\begin{Highlighting}[]
\NormalTok{df.keys()}
\end{Highlighting}
\end{Shaded}

\begin{verbatim}
['NOMAD1',
 'RAJ2000',
 'DEJ2000',
 'YM',
 'r',
 'pmRA',
 'e_pmRA',
 'pmDE',
 'e_pmDE',
 'Bmag',
 'r_Bmag',
 'Vmag',
 'r_Vmag',
 'Rmag',
 'r_Rmag',
 'Jmag',
 'Hmag',
 'Kmag',
 'R']
\end{verbatim}

\begin{Shaded}
\begin{Highlighting}[]
\NormalTok{df[}\StringTok{\textquotesingle{}Jmag\textquotesingle{}}\NormalTok{]}
\end{Highlighting}
\end{Shaded}

{\def\LTcaptype{none} % do not increment counter
\begin{longtable}[]{@{}l@{}}
\toprule\noalign{}
\endhead
\bottomrule\noalign{}
\endlastfoot
13.657 \\
15.955 \\
15.789 \\
14.426 \\
16.146 \\
16.242 \\
15.349 \\
15.910 \\
11.691 \\
12.863 \\
12.322 \\
15.891 \\
... \\
13.632 \\
12.582 \\
15.586 \\
16.623 \\
15.835 \\
16.524 \\
15.640 \\
16.855 \\
16.150 \\
-\/- \\
16.448 \\
\end{longtable}
}

\begin{Shaded}
\begin{Highlighting}[]
\CommentTok{\# Podemos escolher quais colunas queremos pegar e aplicar filtros}
\NormalTok{v }\OperatorTok{=}\NormalTok{ Vizier(columns}\OperatorTok{=}\NormalTok{[}\StringTok{\textquotesingle{}RAJ2000\textquotesingle{}}\NormalTok{, }\StringTok{\textquotesingle{}DEJ2000\textquotesingle{}}\NormalTok{, }\StringTok{\textquotesingle{}Vmag\textquotesingle{}}\NormalTok{, }\StringTok{\textquotesingle{}Bmag\textquotesingle{}}\NormalTok{,}\StringTok{\textquotesingle{}B{-}V\textquotesingle{}}\NormalTok{], catalog}\OperatorTok{=}\StringTok{\textquotesingle{}NOMAD\textquotesingle{}}\NormalTok{,}
\NormalTok{           column\_filters}\OperatorTok{=}\NormalTok{\{}\StringTok{\textquotesingle{}Vmag\textquotesingle{}}\NormalTok{:}\StringTok{\textquotesingle{}\textgreater{}10\textquotesingle{}}\NormalTok{\}, keywords}\OperatorTok{=}\NormalTok{[}\StringTok{\textquotesingle{}optical\textquotesingle{}}\NormalTok{])}
\end{Highlighting}
\end{Shaded}

\begin{Shaded}
\begin{Highlighting}[]
\NormalTok{result }\OperatorTok{=}\NormalTok{ v.query\_object(}\StringTok{\textquotesingle{}HD 226868\textquotesingle{}}\NormalTok{)}
\BuiltInTok{print}\NormalTok{(result)}
\end{Highlighting}
\end{Shaded}

\begin{verbatim}
TableList with 1 tables:
    '0:I/297/out' with 4 column(s) and 50 row(s) 
\end{verbatim}

\begin{Shaded}
\begin{Highlighting}[]
\BuiltInTok{print}\NormalTok{(result[}\StringTok{\textquotesingle{}I/297/out\textquotesingle{}}\NormalTok{])}
\end{Highlighting}
\end{Shaded}

\begin{verbatim}
  RAJ2000     DEJ2000    Vmag   Bmag 
    deg         deg      mag    mag  
----------- ----------- ------ ------
299.5704167  35.1752833 17.810 18.810
299.5802878  35.1768889 13.260 13.160
299.5825550  35.1801478 14.360 13.890
299.5825769  35.1852303 14.480 14.740
299.5878556  35.1741250 16.760 17.290
299.5895250  35.1903333 15.050     --
299.5939833  35.1759750 15.930     --
299.5940833  35.1715028 17.830 19.240
299.5941931  35.1799983 14.360 14.440
        ...         ...    ...    ...
299.5599139  35.1851250 16.450 17.080
299.5723806  35.1881667 16.440 16.930
299.5724361  35.1943167 15.600 15.630
299.5799303  35.1965647 15.460 14.970
299.5805278  35.1895722 17.630 17.170
299.5614972  35.2016922 15.240 15.680
299.5669472  35.2084528 17.630 18.350
299.5678417  35.2073083 17.590 18.000
299.5687194  35.2063333 17.550 18.010
299.5803000  35.2318639 17.770 19.640
Length = 50 rows
\end{verbatim}

\begin{Shaded}
\begin{Highlighting}[]
\NormalTok{df }\OperatorTok{=}\NormalTok{ result[}\StringTok{\textquotesingle{}I/297/out\textquotesingle{}}\NormalTok{]}
\end{Highlighting}
\end{Shaded}

\begin{Shaded}
\begin{Highlighting}[]
\NormalTok{df}
\end{Highlighting}
\end{Shaded}

\begin{Shaded}
\begin{Highlighting}[]
\NormalTok{df[}\StringTok{\textquotesingle{}B{-}V\textquotesingle{}}\NormalTok{] }\OperatorTok{=}\NormalTok{ df[}\StringTok{\textquotesingle{}Bmag\textquotesingle{}}\NormalTok{] }\OperatorTok{{-}}\NormalTok{ df[}\StringTok{\textquotesingle{}Vmag\textquotesingle{}}\NormalTok{]}
\end{Highlighting}
\end{Shaded}

\begin{Shaded}
\begin{Highlighting}[]
\NormalTok{df}
\end{Highlighting}
\end{Shaded}

\section{3) Buscando uma região}\label{buscando-uma-regiuxe3o}

\begin{Shaded}
\begin{Highlighting}[]
\NormalTok{v }\OperatorTok{=}\NormalTok{ Vizier(}\CommentTok{\#columns=[\textquotesingle{}\_RAJ2000\textquotesingle{}, \textquotesingle{}\_DEJ2000\textquotesingle{},\textquotesingle{}B{-}V\textquotesingle{}, \textquotesingle{}Vmag\textquotesingle{}, \textquotesingle{}Bmag\textquotesingle{}], }
\NormalTok{           catalog}\OperatorTok{=}\StringTok{\textquotesingle{}I/297/out\textquotesingle{}}\NormalTok{,}
\NormalTok{           column\_filters}\OperatorTok{=}\NormalTok{\{}\StringTok{\textquotesingle{}Vmag\textquotesingle{}}\NormalTok{:}\StringTok{\textquotesingle{}\textgreater{}10\textquotesingle{}}\NormalTok{\}, keywords}\OperatorTok{=}\NormalTok{[}\StringTok{\textquotesingle{}optical\textquotesingle{}}\NormalTok{])}
\end{Highlighting}
\end{Shaded}

\begin{Shaded}
\begin{Highlighting}[]
\NormalTok{v.query\_region?}
\end{Highlighting}
\end{Shaded}

\begin{Shaded}
\begin{Highlighting}[]
\ImportTok{import}\NormalTok{ astropy.units }\ImportTok{as}\NormalTok{ u}

\NormalTok{result }\OperatorTok{=}\NormalTok{ v.query\_region(}\StringTok{\textquotesingle{}299.561497200  35.201692200\textquotesingle{}}\NormalTok{, radius}\OperatorTok{=}\DecValTok{10}\OperatorTok{*}\NormalTok{u.arcmin)}
\BuiltInTok{print}\NormalTok{(result[}\StringTok{\textquotesingle{}I/297/out\textquotesingle{}}\NormalTok{])}
\end{Highlighting}
\end{Shaded}

\begin{verbatim}
   NOMAD1      RAJ2000     DEJ2000    YM  r  ... r_Rmag  Jmag   Hmag   Kmag   R 
                 deg         deg             ...         mag    mag    mag      
------------ ----------- ----------- --- --- ... ------ ------ ------ ------ ---
1250-0386686 299.6282917  35.0523361  YM   B ...      B 15.348 14.804 14.682    
1250-0386734 299.6363889  35.0511250  YM   B ...      B 14.240 13.528 13.311    
1250-0386760 299.6421417  35.0488111  YM   B ...      B 13.220 12.243 11.930    
1250-0386781 299.6459194  35.0619472  YM   B ...      B 14.730 14.490 14.371    
1250-0386789 299.6471806  35.0545278  YM   B ...      B 15.657 15.273 15.131    
1250-0386806 299.6518056  35.0523056  YM   B ...      B 15.829 15.176 15.212    
1250-0386796 299.6495444  35.0648028  YM   B ...      B 14.925 14.297 14.268    
1250-0386811 299.6533683  35.0579700  YM   C ...      B  9.277  8.732  8.532    
1250-0386829 299.6574433  35.0793014  YM   C ...      B 12.472 12.145 12.069    
         ...         ...         ... ... ... ...    ...    ...    ...    ... ...
1251-0390385 299.7315861  35.1156556  YM   B ...      B 13.765 13.398 13.185    
1251-0390404 299.7368083  35.1170222  YM   B ...      B 15.372 14.804 14.413    
1251-0390412 299.7378833  35.1181750  YM   B ...      B 13.291 12.278 11.898    
1251-0390421 299.7407881  35.1228969  YM   C ...      B 11.192 10.598 10.457    
1251-0390365 299.7275917  35.1275944  YM   B ...      B 13.837 13.298 13.112    
1251-0390369 299.7287511  35.1315561  YM   C ...      B 13.052 12.730 12.646    
1251-0390372 299.7288389  35.1233528  YM   B ...      B 15.441 14.771 14.490    
1251-0390390 299.7335194  35.1249083  YM   B ...      B 15.634 15.307 14.807    
1251-0390418 299.7386864  35.1283603  YM   C ...      B 14.047 13.653 13.491    
1251-0390434 299.7434722  35.1326694  YM   B ...      B 12.201 10.917 10.461    
Length = 50 rows
\end{verbatim}

\begin{verbatim}
WARNING: InputWarning: Coordinate string is being interpreted as an ICRS coordinate provided in degrees. [astroquery.utils.commons]
\end{verbatim}

\begin{Shaded}
\begin{Highlighting}[]
\NormalTok{result[}\StringTok{\textquotesingle{}I/297/out\textquotesingle{}}\NormalTok{].keys()}
\end{Highlighting}
\end{Shaded}

\begin{verbatim}
['NOMAD1',
 'RAJ2000',
 'DEJ2000',
 'YM',
 'r',
 'pmRA',
 'e_pmRA',
 'pmDE',
 'e_pmDE',
 'Bmag',
 'r_Bmag',
 'Vmag',
 'r_Vmag',
 'Rmag',
 'r_Rmag',
 'Jmag',
 'Hmag',
 'Kmag',
 'R']
\end{verbatim}

\begin{Shaded}
\begin{Highlighting}[]
\CommentTok{\# Note que o resultado é sempre uma tabela com 50 elementos, concidencia?}
\NormalTok{v.ROW\_LIMIT}
\end{Highlighting}
\end{Shaded}

\begin{verbatim}
50
\end{verbatim}

\begin{Shaded}
\begin{Highlighting}[]
\NormalTok{v.ROW\_LIMIT }\OperatorTok{=} \DecValTok{100}
\end{Highlighting}
\end{Shaded}

\begin{Shaded}
\begin{Highlighting}[]
\NormalTok{result }\OperatorTok{=}\NormalTok{ v.query\_region(}\StringTok{\textquotesingle{}299.561497200  35.201692200\textquotesingle{}}\NormalTok{, radius}\OperatorTok{=}\DecValTok{10}\OperatorTok{*}\NormalTok{u.arcmin)}
\BuiltInTok{print}\NormalTok{(result[}\StringTok{\textquotesingle{}I/297/out\textquotesingle{}}\NormalTok{])}
\end{Highlighting}
\end{Shaded}

\begin{verbatim}
   NOMAD1      RAJ2000     DEJ2000    YM  r  ... r_Rmag  Jmag   Hmag   Kmag   R 
                 deg         deg             ...         mag    mag    mag      
------------ ----------- ----------- --- --- ... ------ ------ ------ ------ ---
1250-0386686 299.6282917  35.0523361  YM   B ...      B 15.348 14.804 14.682    
1250-0386734 299.6363889  35.0511250  YM   B ...      B 14.240 13.528 13.311    
1250-0386760 299.6421417  35.0488111  YM   B ...      B 13.220 12.243 11.930    
1250-0386781 299.6459194  35.0619472  YM   B ...      B 14.730 14.490 14.371    
1250-0386789 299.6471806  35.0545278  YM   B ...      B 15.657 15.273 15.131    
1250-0386806 299.6518056  35.0523056  YM   B ...      B 15.829 15.176 15.212    
1250-0386796 299.6495444  35.0648028  YM   B ...      B 14.925 14.297 14.268    
1250-0386811 299.6533683  35.0579700  YM   C ...      B  9.277  8.732  8.532    
1250-0386829 299.6574433  35.0793014  YM   C ...      B 12.472 12.145 12.069    
         ...         ...         ... ... ... ...    ...    ...    ...    ... ...
1250-0386023 299.5096967  35.0662294  YM   C ...      B 14.339 14.023 13.941    
1250-0386024 299.5097167  35.0552278  YM   B ...      B 13.882 12.930 12.646    
1250-0386025 299.5097750  35.0429639  YM   B ...      B 15.605 15.029 14.829    
1250-0386040 299.5116528  35.0468056  YM   B ...      B 12.964 12.054 11.737    
1250-0386046 299.5128819  35.0493853  YM   C ...      B 10.387  9.838  9.704    
1250-0386071 299.5173361  35.0537500  YM   B ...      B 15.670 15.194 14.909    
1250-0386102 299.5238139  35.0422028  YM   B ...      B 12.484 11.547 11.237    
1250-0385830 299.4715083  35.0545222  YM   B ...      B 14.614 14.138 13.929    
1250-0385862 299.4783042  35.0718208  YM   C ...      B 12.448 12.194 12.175    
1250-0385866 299.4787361  35.0582389  YM   B ...      B 12.701 11.384 10.922    
Length = 100 rows
\end{verbatim}

\begin{verbatim}
WARNING: InputWarning: Coordinate string is being interpreted as an ICRS coordinate provided in degrees. [astroquery.utils.commons]
\end{verbatim}

\begin{Shaded}
\begin{Highlighting}[]
\NormalTok{v.ROW\_LIMIT }\OperatorTok{=} \DecValTok{10000}
\end{Highlighting}
\end{Shaded}

\begin{Shaded}
\begin{Highlighting}[]
\NormalTok{result }\OperatorTok{=}\NormalTok{ v.query\_region(}\StringTok{\textquotesingle{}299.561497200  35.201692200\textquotesingle{}}\NormalTok{, radius}\OperatorTok{=}\DecValTok{10}\OperatorTok{*}\NormalTok{u.arcmin)}
\BuiltInTok{print}\NormalTok{(result[}\StringTok{\textquotesingle{}I/297/out\textquotesingle{}}\NormalTok{])}
\end{Highlighting}
\end{Shaded}

\begin{verbatim}
WARNING
\end{verbatim}

\begin{verbatim}
   NOMAD1      RAJ2000     DEJ2000    YM  r  ... r_Rmag  Jmag   Hmag   Kmag   R 
                 deg         deg             ...         mag    mag    mag      
------------ ----------- ----------- --- --- ... ------ ------ ------ ------ ---
1250-0386686 299.6282917  35.0523361  YM   B ...      B 15.348 14.804 14.682    
1250-0386734 299.6363889  35.0511250  YM   B ...      B 14.240 13.528 13.311    
1250-0386760 299.6421417  35.0488111  YM   B ...      B 13.220 12.243 11.930    
1250-0386781 299.6459194  35.0619472  YM   B ...      B 14.730 14.490 14.371    
1250-0386789 299.6471806  35.0545278  YM   B ...      B 15.657 15.273 15.131    
1250-0386806 299.6518056  35.0523056  YM   B ...      B 15.829 15.176 15.212    
1250-0386796 299.6495444  35.0648028  YM   B ...      B 14.925 14.297 14.268    
1250-0386811 299.6533683  35.0579700  YM   C ...      B  9.277  8.732  8.532    
1250-0386829 299.6574433  35.0793014  YM   C ...      B 12.472 12.145 12.069    
         ...         ...         ... ... ... ...    ...    ...    ...    ... ...
1253-0392982 299.5073378  35.3557922  YM   C ...      B 11.706 10.745 10.449    
1253-0392992 299.5105194  35.3496944  YM   B ...      B 14.842 14.061 13.890    
1253-0393000 299.5120686  35.3485886  YM   C ...      B 13.367 13.104 13.006    
1253-0393003 299.5125917  35.3547667  YM   Y ...        16.295 15.593 15.217    
1253-0393005 299.5127139  35.3459583  YM   B ...      B 15.118 14.490 14.400    
1253-0393008 299.5131444  35.3528750  YM   Y ...        16.425 15.644 15.500    
1253-0393013 299.5133500  35.3589139  YM   B ...      B 12.537 12.199 12.115    
1253-0393026 299.5153528  35.3551333  YM   Y ...        11.452 10.621 10.351    
1253-0393041 299.5174222  35.3540861  YM   Y ...        12.874 12.611 12.535    
1253-0393096 299.5276778  35.3641861  YM   B ...      B 15.106 14.721 14.532    
Length = 1549 rows
\end{verbatim}

\begin{verbatim}
: InputWarning: Coordinate string is being interpreted as an ICRS coordinate provided in degrees. [astroquery.utils.commons]
\end{verbatim}

\begin{Shaded}
\begin{Highlighting}[]
\NormalTok{df2 }\OperatorTok{=}\NormalTok{ result[}\StringTok{\textquotesingle{}I/297/out\textquotesingle{}}\NormalTok{]}
\end{Highlighting}
\end{Shaded}

\begin{Shaded}
\begin{Highlighting}[]
\NormalTok{df2.write(}\StringTok{\textquotesingle{}output\_astroq.csv\textquotesingle{}}\NormalTok{, }\BuiltInTok{format}\OperatorTok{=}\StringTok{\textquotesingle{}csv\textquotesingle{}}\NormalTok{, overwrite}\OperatorTok{=}\VariableTok{True}\NormalTok{)}
\end{Highlighting}
\end{Shaded}

\begin{Shaded}
\begin{Highlighting}[]
\NormalTok{df2.write?}
\end{Highlighting}
\end{Shaded}

\begin{Shaded}
\begin{Highlighting}[]
\NormalTok{df5 }\OperatorTok{=}\NormalTok{ pd.read\_csv(}\StringTok{\textquotesingle{}output\_aq.csv\textquotesingle{}}\NormalTok{)}
\end{Highlighting}
\end{Shaded}

\begin{Shaded}
\begin{Highlighting}[]
\NormalTok{df5}
\end{Highlighting}
\end{Shaded}

{\def\LTcaptype{none} % do not increment counter
\begin{longtable}[]{@{}lllll@{}}
\toprule\noalign{}
& \_RAJ2000 & \_DEJ2000 & Vmag & Bmag \\
\midrule\noalign{}
\endhead
\bottomrule\noalign{}
\endlastfoot
0 & 299.361972 & 35.229758 & 17.58 & 18.70 \\
1 & 299.362706 & 35.212617 & 17.65 & 19.15 \\
2 & 299.366253 & 35.201208 & 14.05 & 14.54 \\
3 & 299.367789 & 35.248108 & 17.08 & 17.28 \\
4 & 299.367906 & 35.227378 & 17.58 & 18.80 \\
... & ... & ... & ... & ... \\
1544 & 299.692692 & 35.082156 & 17.02 & 17.66 \\
1545 & 299.699177 & 35.079201 & 14.56 & 14.74 \\
1546 & 299.701711 & 35.098582 & 13.30 & 13.55 \\
1547 & 299.706179 & 35.093528 & 15.10 & 15.19 \\
1548 & 299.711769 & 35.099983 & 17.74 & 15.03 \\
\end{longtable}
}

\begin{Shaded}
\begin{Highlighting}[]
\CommentTok{\# Levando o catálogo para o Pandas}
\NormalTok{df3 }\OperatorTok{=}\NormalTok{ df2.to\_pandas()}
\end{Highlighting}
\end{Shaded}

\begin{Shaded}
\begin{Highlighting}[]
\BuiltInTok{type}\NormalTok{(df3)}
\end{Highlighting}
\end{Shaded}

\begin{verbatim}
pandas.core.frame.DataFrame
\end{verbatim}

\begin{Shaded}
\begin{Highlighting}[]
\ImportTok{from}\NormalTok{ astroquery }\ImportTok{import}\NormalTok{ nasa\_ads }\ImportTok{as}\NormalTok{ na}
\end{Highlighting}
\end{Shaded}

\begin{Shaded}
\begin{Highlighting}[]
\NormalTok{na.ADS.TOKEN }\OperatorTok{=} \StringTok{\textquotesingle{}fgo81KRQVaxNngpL3zOmGgO8JraFaLSbzf7HwXQA\textquotesingle{}}
\end{Highlighting}
\end{Shaded}

\begin{Shaded}
\begin{Highlighting}[]
\NormalTok{results }\OperatorTok{=}\NormalTok{ na.ADS.query\_simple(}\StringTok{\textquotesingle{}author:"Morgado, Bruno", aff:"Valongo"\textquotesingle{}}\NormalTok{)}
\end{Highlighting}
\end{Shaded}

\begin{Shaded}
\begin{Highlighting}[]
\NormalTok{results}
\end{Highlighting}
\end{Shaded}

\chapter{Isso é tudo}\label{isso-uxe9-tudo-3}

By BEM

\part{Ferramentas Computacionais}

\chapter{LaTeX}\label{latex}

\chapter{Git}\label{git}

\part{Ferramentas Matemáticas}

\chapter{Raízes de Funções}\label{rauxedzes-de-funuxe7uxf5es}

\subsubsection{by BEM, Last update:
10/07/2023}\label{by-bem-last-update-10072023}

\begin{Shaded}
\begin{Highlighting}[]
\ImportTok{import}\NormalTok{ numpy }\ImportTok{as}\NormalTok{ np}
\ImportTok{import}\NormalTok{ matplotlib.pyplot }\ImportTok{as}\NormalTok{ pl}
\end{Highlighting}
\end{Shaded}

\chapter{\texorpdfstring{\(f(x) = \sqrt{x} - 5e^{-x}\)}{f(x) = \textbackslash sqrt\{x\} - 5e\^{}\{-x\}}}\label{fx-sqrtx---5e-x}

\begin{Shaded}
\begin{Highlighting}[]
\KeywordTok{def}\NormalTok{ func(x):}
\NormalTok{    f }\OperatorTok{=}\NormalTok{ np.sqrt(x) }\OperatorTok{{-}} \DecValTok{5}\OperatorTok{*}\NormalTok{(np.e}\OperatorTok{**}\NormalTok{(}\OperatorTok{{-}}\NormalTok{x))}
    \ControlFlowTok{return}\NormalTok{ f}
\end{Highlighting}
\end{Shaded}

\begin{Shaded}
\begin{Highlighting}[]
\CommentTok{\# Avaliando o sinal de f(x), para x entre 0 e 10}

\NormalTok{x }\OperatorTok{=}\NormalTok{ np.arange(}\DecValTok{10}\NormalTok{)}

\ControlFlowTok{for}\NormalTok{ i }\KeywordTok{in} \BuiltInTok{range}\NormalTok{(}\BuiltInTok{len}\NormalTok{(x)):}
\NormalTok{    f }\OperatorTok{=}\NormalTok{ func(x[i])}
    \ControlFlowTok{if}\NormalTok{ f }\OperatorTok{\textless{}} \DecValTok{0}\NormalTok{:}
        \BuiltInTok{print}\NormalTok{(x[i], }\StringTok{\textquotesingle{}{-}\textquotesingle{}}\NormalTok{)}
    \ControlFlowTok{else}\NormalTok{:}
        \BuiltInTok{print}\NormalTok{(x[i], }\StringTok{\textquotesingle{}+\textquotesingle{}}\NormalTok{)}
\end{Highlighting}
\end{Shaded}

\begin{verbatim}
0 -
1 -
2 +
3 +
4 +
5 +
6 +
7 +
8 +
9 +
\end{verbatim}

\begin{Shaded}
\begin{Highlighting}[]
\CommentTok{\# Colocando este sinal em uma variavel}

\NormalTok{x }\OperatorTok{=}\NormalTok{ np.arange(}\DecValTok{0}\NormalTok{,}\DecValTok{10}\NormalTok{, }\DecValTok{1}\NormalTok{)}
\NormalTok{sinal }\OperatorTok{=}\NormalTok{ np.array([]) }\CommentTok{\#\#\#\#\#\#\#\#\#\#\#\#\#\#\#\#\#\#\#\#\#\#\#\#\#\#\#\#\#}

\ControlFlowTok{for}\NormalTok{ i }\KeywordTok{in} \BuiltInTok{range}\NormalTok{(}\BuiltInTok{len}\NormalTok{(x)):}
\NormalTok{    f }\OperatorTok{=}\NormalTok{ np.sqrt(x[i]) }\OperatorTok{{-}} \DecValTok{5}\OperatorTok{*}\NormalTok{(np.e}\OperatorTok{**}\NormalTok{(}\OperatorTok{{-}}\NormalTok{x[i]))}
    \ControlFlowTok{if}\NormalTok{ f }\OperatorTok{\textless{}} \DecValTok{0}\NormalTok{:}
\NormalTok{        sinal }\OperatorTok{=}\NormalTok{ np.append(sinal, }\OperatorTok{{-}}\DecValTok{1}\NormalTok{) }\CommentTok{\#\#\#\#\#\#\#\#\#\#\#\#\#\#\#\#\#\#\#\#\#\#\#\#\#\#\#\#\#}
        \BuiltInTok{print}\NormalTok{(x[i], }\StringTok{\textquotesingle{}{-}\textquotesingle{}}\NormalTok{)}
    \ControlFlowTok{else}\NormalTok{:}
\NormalTok{        sinal }\OperatorTok{=}\NormalTok{ np.append(sinal, }\OperatorTok{+}\DecValTok{1}\NormalTok{) }\CommentTok{\#\#\#\#\#\#\#\#\#\#\#\#\#\#\#\#\#\#\#\#\#\#\#\#\#\#\#\#\#}
        \BuiltInTok{print}\NormalTok{(x[i], }\StringTok{\textquotesingle{}+\textquotesingle{}}\NormalTok{)}
\end{Highlighting}
\end{Shaded}

\begin{verbatim}
0 -
1 -
2 +
3 +
4 +
5 +
6 +
7 +
8 +
9 +
\end{verbatim}

\begin{Shaded}
\begin{Highlighting}[]
\BuiltInTok{len}\NormalTok{(sinal)}
\end{Highlighting}
\end{Shaded}

\begin{verbatim}
10
\end{verbatim}

\begin{Shaded}
\begin{Highlighting}[]
\CommentTok{\# Onde há troca de sinal?}

\ControlFlowTok{for}\NormalTok{ i }\KeywordTok{in} \BuiltInTok{range}\NormalTok{(}\BuiltInTok{len}\NormalTok{(x)}\OperatorTok{{-}}\DecValTok{1}\NormalTok{):}
    \ControlFlowTok{if}\NormalTok{ sinal[i]}\OperatorTok{*}\NormalTok{sinal[i}\OperatorTok{+}\DecValTok{1}\NormalTok{] }\OperatorTok{\textless{}} \DecValTok{0}\NormalTok{:}
        \BuiltInTok{print}\NormalTok{(x[i],x[i}\OperatorTok{+}\DecValTok{1}\NormalTok{], sinal[i]}\OperatorTok{*}\NormalTok{sinal[i}\OperatorTok{+}\DecValTok{1}\NormalTok{])}
\NormalTok{        a }\OperatorTok{=}\NormalTok{ x[i]}
\NormalTok{        b }\OperatorTok{=}\NormalTok{ x[i}\OperatorTok{+}\DecValTok{1}\NormalTok{]}
\end{Highlighting}
\end{Shaded}

\begin{verbatim}
1 2 -1.0
\end{verbatim}

\begin{Shaded}
\begin{Highlighting}[]
\CommentTok{\# A raiz desejada está entre a e b!}
\BuiltInTok{print}\NormalTok{(a, b)}
\end{Highlighting}
\end{Shaded}

\begin{verbatim}
1 2
\end{verbatim}

\begin{Shaded}
\begin{Highlighting}[]
\CommentTok{\# vamos começar chutando o valor central}

\NormalTok{x0 }\OperatorTok{=}\NormalTok{ (a }\OperatorTok{+}\NormalTok{ b)}\OperatorTok{/}\DecValTok{2}
\end{Highlighting}
\end{Shaded}

\begin{Shaded}
\begin{Highlighting}[]
\KeywordTok{def}\NormalTok{ funcao(x):}
    \ControlFlowTok{return}\NormalTok{ np.sqrt(x) }\OperatorTok{{-}} \DecValTok{5}\OperatorTok{*}\NormalTok{(np.e}\OperatorTok{**}\NormalTok{(}\OperatorTok{{-}}\NormalTok{x))}
\end{Highlighting}
\end{Shaded}

\begin{Shaded}
\begin{Highlighting}[]
\CommentTok{\# Avaliando nosso chute e escolhendo novos chutes}

\BuiltInTok{print}\NormalTok{(funcao(}\DecValTok{1}\NormalTok{))}
\BuiltInTok{print}\NormalTok{(funcao(}\FloatTok{1.25}\NormalTok{))}
\BuiltInTok{print}\NormalTok{(funcao(}\FloatTok{1.5}\NormalTok{))}
\end{Highlighting}
\end{Shaded}

\begin{verbatim}
-0.8393972058572117
-0.3144899955510556
0.10909407064943966
\end{verbatim}

\begin{Shaded}
\begin{Highlighting}[]
\ImportTok{from}\NormalTok{ scipy.optimize }\ImportTok{import}\NormalTok{ fsolve}
\end{Highlighting}
\end{Shaded}

\begin{Shaded}
\begin{Highlighting}[]
\NormalTok{fsolve(func, }\DecValTok{0}\NormalTok{)}
\end{Highlighting}
\end{Shaded}

\begin{verbatim}
array([1.43044509])
\end{verbatim}

\begin{Shaded}
\begin{Highlighting}[]
\CommentTok{\# A troca de sinal ocorre entre a e x0, logo:}

\NormalTok{a }\OperatorTok{=}\NormalTok{ a}
\NormalTok{b }\OperatorTok{=}\NormalTok{ x0}

\NormalTok{x0 }\OperatorTok{=}\NormalTok{ (a }\OperatorTok{+}\NormalTok{ b)}\OperatorTok{/}\DecValTok{2}
\BuiltInTok{print}\NormalTok{(b}\OperatorTok{{-}}\NormalTok{a)}
\end{Highlighting}
\end{Shaded}

\begin{verbatim}
0.5
\end{verbatim}

\begin{Shaded}
\begin{Highlighting}[]
\CommentTok{\# Avaliando nosso chute e escolhendo novos chutes}

\BuiltInTok{print}\NormalTok{(funcao(a ))}
\BuiltInTok{print}\NormalTok{(funcao(x0))}
\BuiltInTok{print}\NormalTok{(funcao(b ))}
\end{Highlighting}
\end{Shaded}

\begin{verbatim}
-0.8393972058572117
-0.3144899955510556
0.10909407064943966
\end{verbatim}

\begin{Shaded}
\begin{Highlighting}[]
\CommentTok{\# Agora a troca de sinal ocorre entre x0 e b, logo:}

\NormalTok{a }\OperatorTok{=}\NormalTok{ x0}
\NormalTok{b }\OperatorTok{=}\NormalTok{ b}

\NormalTok{x0 }\OperatorTok{=}\NormalTok{ (a }\OperatorTok{+}\NormalTok{ b)}\OperatorTok{/}\DecValTok{2}
\BuiltInTok{print}\NormalTok{(b}\OperatorTok{{-}}\NormalTok{a)}
\end{Highlighting}
\end{Shaded}

\begin{verbatim}
0.25
\end{verbatim}

\begin{Shaded}
\begin{Highlighting}[]
\BuiltInTok{print}\NormalTok{(funcao(a ))}
\BuiltInTok{print}\NormalTok{(funcao(x0))}
\BuiltInTok{print}\NormalTok{(funcao(b ))}
\end{Highlighting}
\end{Shaded}

\begin{verbatim}
-0.3144899955510556
-0.09159403906787511
0.10909407064943966
\end{verbatim}

Podemos continuar fazendo até que a \textbar b - a\textbar{} seja menor
do que a precisão desejada

\begin{Shaded}
\begin{Highlighting}[]
\CommentTok{\# Ou de forma mais automática, podemos fazer}

\NormalTok{a }\OperatorTok{=} \DecValTok{1} \CommentTok{\# Chute inicial}
\NormalTok{b }\OperatorTok{=} \DecValTok{2} \CommentTok{\# Chute inicial}
\NormalTok{x0 }\OperatorTok{=}\NormalTok{ (a }\OperatorTok{+}\NormalTok{ b)}\OperatorTok{/}\DecValTok{2}
\NormalTok{err }\OperatorTok{=} \BuiltInTok{abs}\NormalTok{(b }\OperatorTok{{-}}\NormalTok{ a)}

\ControlFlowTok{while}\NormalTok{ err }\OperatorTok{\textgreater{}} \FloatTok{0.001}\NormalTok{:}
    \ControlFlowTok{if}\NormalTok{ (funcao(a) }\OperatorTok{\textless{}} \DecValTok{0}\NormalTok{) }\OperatorTok{*}\NormalTok{ (funcao(b) }\OperatorTok{\textgreater{}} \DecValTok{0}\NormalTok{) }\OperatorTok{*}\NormalTok{ (funcao(x0) }\OperatorTok{\textgreater{}} \DecValTok{0}\NormalTok{):}
\NormalTok{        a }\OperatorTok{=}\NormalTok{ a}
\NormalTok{        b }\OperatorTok{=}\NormalTok{ x0}
    \ControlFlowTok{if}\NormalTok{ (funcao(a) }\OperatorTok{\textless{}} \DecValTok{0}\NormalTok{) }\OperatorTok{*}\NormalTok{ (funcao(b) }\OperatorTok{\textgreater{}} \DecValTok{0}\NormalTok{) }\OperatorTok{*}\NormalTok{ (funcao(x0) }\OperatorTok{\textless{}} \DecValTok{0}\NormalTok{):}
\NormalTok{        a }\OperatorTok{=}\NormalTok{ x0}
\NormalTok{        b }\OperatorTok{=}\NormalTok{ b}
    \ControlFlowTok{if}\NormalTok{ (funcao(a) }\OperatorTok{\textgreater{}} \DecValTok{0}\NormalTok{) }\OperatorTok{*}\NormalTok{ (funcao(b) }\OperatorTok{\textless{}} \DecValTok{0}\NormalTok{) }\OperatorTok{*}\NormalTok{ (funcao(x0) }\OperatorTok{\textless{}} \DecValTok{0}\NormalTok{):}
\NormalTok{        a }\OperatorTok{=}\NormalTok{ a}
\NormalTok{        b }\OperatorTok{=}\NormalTok{ x0}
    \ControlFlowTok{if}\NormalTok{ (funcao(a) }\OperatorTok{\textgreater{}} \DecValTok{0}\NormalTok{) }\OperatorTok{*}\NormalTok{ (funcao(b) }\OperatorTok{\textless{}} \DecValTok{0}\NormalTok{) }\OperatorTok{*}\NormalTok{ (funcao(x0) }\OperatorTok{\textgreater{}} \DecValTok{0}\NormalTok{):}
\NormalTok{        a }\OperatorTok{=}\NormalTok{ x0}
\NormalTok{        b }\OperatorTok{=}\NormalTok{ b}

\NormalTok{    x0 }\OperatorTok{=}\NormalTok{ (a }\OperatorTok{+}\NormalTok{ b)}\OperatorTok{/}\DecValTok{2}
\NormalTok{    err }\OperatorTok{=} \BuiltInTok{abs}\NormalTok{(b }\OperatorTok{{-}}\NormalTok{ a)}
    
\BuiltInTok{print}\NormalTok{(x0, err)}
\end{Highlighting}
\end{Shaded}

\begin{verbatim}
1.43017578125 0.0009765625
\end{verbatim}

\subsubsection{Como fazer com o scipy?}\label{como-fazer-com-o-scipy}

\begin{Shaded}
\begin{Highlighting}[]
\ImportTok{from}\NormalTok{ scipy.optimize }\ImportTok{import}\NormalTok{ fsolve}
\end{Highlighting}
\end{Shaded}

\begin{Shaded}
\begin{Highlighting}[]
\NormalTok{fsolve?}
\end{Highlighting}
\end{Shaded}

\begin{Shaded}
\begin{Highlighting}[]
\NormalTok{fsolve(funcao, [}\DecValTok{1}\NormalTok{])}
\end{Highlighting}
\end{Shaded}

\begin{verbatim}
array([1.43044509])
\end{verbatim}

\begin{Shaded}
\begin{Highlighting}[]
\NormalTok{xxx }\OperatorTok{=}\NormalTok{ np.arange(}\DecValTok{0}\NormalTok{,}\DecValTok{5}\NormalTok{,}\FloatTok{0.01}\NormalTok{)}
\NormalTok{pl.plot(xxx, funcao(xxx))}
\NormalTok{pl.axhline(}\DecValTok{0}\NormalTok{, ls}\OperatorTok{=}\StringTok{\textquotesingle{}{-}{-}\textquotesingle{}}\NormalTok{, color}\OperatorTok{=}\StringTok{\textquotesingle{}r\textquotesingle{}}\NormalTok{)}
\end{Highlighting}
\end{Shaded}

\pandocbounded{\includegraphics[keepaspectratio]{raizes_files/figure-pdf/cell-22-output-1.pdf}}

\chapter{Isso é tudo}\label{isso-uxe9-tudo-4}

By BEM

\chapter{Ajuste de Funções}\label{ajuste-de-funuxe7uxf5es}

\subsubsection{by BEM, Last update:
10/07/2023}\label{by-bem-last-update-10072023-1}

\begin{Shaded}
\begin{Highlighting}[]
\ImportTok{import}\NormalTok{ numpy }\ImportTok{as}\NormalTok{ np}
\ImportTok{import}\NormalTok{ matplotlib.pyplot }\ImportTok{as}\NormalTok{ pl}
\end{Highlighting}
\end{Shaded}

\chapter{Caso particular da Reta}\label{caso-particular-da-reta}

\begin{Shaded}
\begin{Highlighting}[]
\CommentTok{\# Vamos Ler o arquivo de entrada}
\NormalTok{x, y }\OperatorTok{=}\NormalTok{ np.loadtxt(}\StringTok{\textquotesingle{}exercicio\_1.dat\textquotesingle{}}\NormalTok{, unpack}\OperatorTok{=}\VariableTok{True}\NormalTok{)}
\end{Highlighting}
\end{Shaded}

\begin{Shaded}
\begin{Highlighting}[]
\CommentTok{\# Visualizando nossos dados}

\NormalTok{pl.plot(x, y, }\StringTok{\textquotesingle{}k.{-}\textquotesingle{}}\NormalTok{)}
\NormalTok{pl.show()}
\end{Highlighting}
\end{Shaded}

\pandocbounded{\includegraphics[keepaspectratio]{ajuste/ajuste_files/figure-pdf/cell-4-output-1.pdf}}

Podemos computar os coeficientes utilizando:

\$ \frac{\partial S}{\partial a} = \sum\_\{i = 1\}\^{}\{n\}2.x\_i.(ax\_i
+ b - y\_i) = 0\$

\$ \frac{\partial S}{\partial b} = \sum\_\{i = 1\}\^{}\{n\}2.(ax\_i + b
- y\_i) = 0\$

\(\therefore\)

\(a = \frac{\sum_{i = 1}^{n} x_i (y_i - \bar{y})}{\sum_{i = 1}^{n} x_i (x_i - \bar{x})}\)

\(b = \bar{y} - a\bar{x}\)

onde

\(\bar{x} = \frac{1}{N}\sum_{i = 1}^{n} x_i\)

\(\bar{y} = \frac{1}{N}\sum_{i = 1}^{n} y_i\)

\begin{Shaded}
\begin{Highlighting}[]
\NormalTok{a }\OperatorTok{=}\NormalTok{ np.}\BuiltInTok{sum}\NormalTok{(x}\OperatorTok{*}\NormalTok{(y }\OperatorTok{{-}}\NormalTok{ y.mean())) }\OperatorTok{/}\NormalTok{ np.}\BuiltInTok{sum}\NormalTok{(x}\OperatorTok{*}\NormalTok{(x }\OperatorTok{{-}}\NormalTok{ x.mean()))}
\BuiltInTok{print}\NormalTok{(}\StringTok{\textquotesingle{}Coeficiente Angular: \textquotesingle{}}\NormalTok{, a)}
\end{Highlighting}
\end{Shaded}

\begin{verbatim}
Coeficiente Angular:  0.3093444729928444
\end{verbatim}

\begin{Shaded}
\begin{Highlighting}[]
\NormalTok{b }\OperatorTok{=}\NormalTok{ y.mean() }\OperatorTok{{-}}\NormalTok{ a}\OperatorTok{*}\NormalTok{x.mean()}
\BuiltInTok{print}\NormalTok{(}\StringTok{\textquotesingle{}Coeficiente linear: \textquotesingle{}}\NormalTok{, b)}
\end{Highlighting}
\end{Shaded}

\begin{verbatim}
Coeficiente linear:  17.56814695035778
\end{verbatim}

\begin{Shaded}
\begin{Highlighting}[]
\CommentTok{\# Visualizando nosso resultado}

\NormalTok{pl.plot(x, y, }\StringTok{\textquotesingle{}k.{-}\textquotesingle{}}\NormalTok{)}
\NormalTok{pl.plot(x, a}\OperatorTok{*}\NormalTok{x }\OperatorTok{+}\NormalTok{ b, }\StringTok{\textquotesingle{}r{-}\textquotesingle{}}\NormalTok{)}
\NormalTok{pl.show()}
\end{Highlighting}
\end{Shaded}

\pandocbounded{\includegraphics[keepaspectratio]{ajuste/ajuste_files/figure-pdf/cell-7-output-1.pdf}}

Esse foi um exemplo de um ajuste por mínimos quadrados para o caso de um
ajuste simples que pode ser resolvido algebricamente. Mas como fazer de
forma numérica?

\subsection{Vamos fazer um Monte
Carlo}\label{vamos-fazer-um-monte-carlo}

\begin{Shaded}
\begin{Highlighting}[]
\CommentTok{\# Primeiro, definir o nosso problema}

\KeywordTok{def}\NormalTok{ reta(x, a, b):}
    \ControlFlowTok{return}\NormalTok{ a}\OperatorTok{*}\NormalTok{x }\OperatorTok{+}\NormalTok{ b}
\end{Highlighting}
\end{Shaded}

\begin{Shaded}
\begin{Highlighting}[]
\CommentTok{\# Definimos alguns chutes iniciais }

\NormalTok{pl.plot(x, y, }\StringTok{\textquotesingle{}k.{-}\textquotesingle{}}\NormalTok{)}
\NormalTok{pl.plot(x, reta(x, }\FloatTok{0.5}\NormalTok{, }\DecValTok{10}\NormalTok{), }\StringTok{\textquotesingle{}r{-}\textquotesingle{}}\NormalTok{)}
\NormalTok{pl.show()}
\end{Highlighting}
\end{Shaded}

\pandocbounded{\includegraphics[keepaspectratio]{ajuste/ajuste_files/figure-pdf/cell-9-output-1.pdf}}

\begin{Shaded}
\begin{Highlighting}[]
\CommentTok{\# Vamos criar valores randômicos para testar}
\NormalTok{param\_a }\OperatorTok{=}\NormalTok{ np.random.uniform(}\DecValTok{0}\NormalTok{, }\DecValTok{1}\NormalTok{, }\DecValTok{100000}\NormalTok{)}
\NormalTok{param\_b }\OperatorTok{=}\NormalTok{ np.random.uniform(}\DecValTok{0}\NormalTok{, }\DecValTok{30}\NormalTok{, }\DecValTok{100000}\NormalTok{)}
\end{Highlighting}
\end{Shaded}

\begin{Shaded}
\begin{Highlighting}[]
\CommentTok{\# Como estamos preenchendo o espaço de parâmetros?}

\NormalTok{pl.plot(param\_a, param\_b, }\StringTok{\textquotesingle{}k.\textquotesingle{}}\NormalTok{)}
\NormalTok{pl.xlabel(}\StringTok{\textquotesingle{}a\textquotesingle{}}\NormalTok{, fontsize}\OperatorTok{=}\DecValTok{15}\NormalTok{)}
\NormalTok{pl.ylabel(}\StringTok{\textquotesingle{}b\textquotesingle{}}\NormalTok{, fontsize}\OperatorTok{=}\DecValTok{15}\NormalTok{)}
\NormalTok{pl.show()}

\CommentTok{\# Cada ponto preto é uma simulação que vamos fazer de forma iterativa}
\end{Highlighting}
\end{Shaded}

\pandocbounded{\includegraphics[keepaspectratio]{ajuste/ajuste_files/figure-pdf/cell-11-output-1.pdf}}

\begin{Shaded}
\begin{Highlighting}[]
\NormalTok{pl.plot(param\_a, param\_b, }\StringTok{\textquotesingle{}k.\textquotesingle{}}\NormalTok{)}
\NormalTok{pl.plot(param\_a[}\DecValTok{0}\NormalTok{], param\_b[}\DecValTok{0}\NormalTok{], }\StringTok{\textquotesingle{}r.\textquotesingle{}}\NormalTok{)}
\NormalTok{pl.xlabel(}\StringTok{\textquotesingle{}a\textquotesingle{}}\NormalTok{, fontsize}\OperatorTok{=}\DecValTok{15}\NormalTok{)}
\NormalTok{pl.ylabel(}\StringTok{\textquotesingle{}b\textquotesingle{}}\NormalTok{, fontsize}\OperatorTok{=}\DecValTok{15}\NormalTok{)}
\NormalTok{pl.show()}

\BuiltInTok{print}\NormalTok{(param\_a[}\DecValTok{0}\NormalTok{], param\_b[}\DecValTok{0}\NormalTok{])}
\end{Highlighting}
\end{Shaded}

\pandocbounded{\includegraphics[keepaspectratio]{ajuste/ajuste_files/figure-pdf/cell-12-output-1.pdf}}

\begin{verbatim}
0.6615403366447085 26.80812683250735
\end{verbatim}

Computando um modelo e calculando o valor de:

\(S = \sum_{i = 1}^{N} (y_i - f(x_i))^2\)

\begin{Shaded}
\begin{Highlighting}[]
\NormalTok{y\_model }\OperatorTok{=}\NormalTok{ reta(x, a}\OperatorTok{=}\NormalTok{param\_a[}\DecValTok{0}\NormalTok{], b}\OperatorTok{=}\NormalTok{param\_b[}\DecValTok{0}\NormalTok{])}

\NormalTok{S }\OperatorTok{=}\NormalTok{ np.}\BuiltInTok{sum}\NormalTok{((y }\OperatorTok{{-}}\NormalTok{ y\_model)}\OperatorTok{**}\DecValTok{2}\NormalTok{)}
\BuiltInTok{print}\NormalTok{(}\StringTok{\textquotesingle{}Soma dos quadrados dos Resíduos: \textquotesingle{}}\NormalTok{, S)}

\CommentTok{\# Quanto maior o S, pior é o ajuste, vamos calcular S para todos os valores simulados}
\end{Highlighting}
\end{Shaded}

\begin{verbatim}
Soma dos quadrados dos Resíduos:  25157.79080745567
\end{verbatim}

\begin{Shaded}
\begin{Highlighting}[]
\NormalTok{S }\OperatorTok{=}\NormalTok{ np.array([])}

\ControlFlowTok{for}\NormalTok{ i }\KeywordTok{in} \BuiltInTok{range}\NormalTok{(}\BuiltInTok{len}\NormalTok{(param\_a)):}
\NormalTok{    y\_model }\OperatorTok{=}\NormalTok{ reta(x, a}\OperatorTok{=}\NormalTok{param\_a[i], b}\OperatorTok{=}\NormalTok{param\_b[i])}
\NormalTok{    S }\OperatorTok{=}\NormalTok{ np.append(S, np.}\BuiltInTok{sum}\NormalTok{((y }\OperatorTok{{-}}\NormalTok{ y\_model)}\OperatorTok{**}\DecValTok{2}\NormalTok{))}
\end{Highlighting}
\end{Shaded}

\begin{Shaded}
\begin{Highlighting}[]
\BuiltInTok{print}\NormalTok{(}\StringTok{\textquotesingle{}Valor mínimo da soma dos quadrados dos resíduos: \textquotesingle{}}\NormalTok{, S.}\BuiltInTok{min}\NormalTok{())}
\BuiltInTok{print}\NormalTok{(}\StringTok{\textquotesingle{}Melhor valor de a: \textquotesingle{}}\NormalTok{, param\_a[S.argmin()])}
\BuiltInTok{print}\NormalTok{(}\StringTok{\textquotesingle{}Melhor valor de b: \textquotesingle{}}\NormalTok{, param\_b[S.argmin()])}
\end{Highlighting}
\end{Shaded}

\begin{verbatim}
Valor mínimo da soma dos quadrados dos resíduos:  215.6692339877744
Melhor valor de a:  0.3074223639137542
Melhor valor de b:  17.673651041569038
\end{verbatim}

\begin{Shaded}
\begin{Highlighting}[]
\CommentTok{\# Definimos alguns chutes iniciais }

\NormalTok{pl.plot(x, y, }\StringTok{\textquotesingle{}k.{-}\textquotesingle{}}\NormalTok{)}
\NormalTok{pl.plot(x, reta(x, param\_a[S.argmin()], param\_b[S.argmin()]), }\StringTok{\textquotesingle{}r{-}\textquotesingle{}}\NormalTok{)}
\NormalTok{pl.show()}
\end{Highlighting}
\end{Shaded}

\pandocbounded{\includegraphics[keepaspectratio]{ajuste/ajuste_files/figure-pdf/cell-16-output-1.pdf}}

\begin{Shaded}
\begin{Highlighting}[]
\NormalTok{pl.plot(param\_a, S, }\StringTok{\textquotesingle{}k.\textquotesingle{}}\NormalTok{)}
\NormalTok{pl.axvline(param\_a[S.argmin()], color}\OperatorTok{=}\StringTok{\textquotesingle{}red\textquotesingle{}}\NormalTok{, ls}\OperatorTok{=}\StringTok{\textquotesingle{}{-}{-}\textquotesingle{}}\NormalTok{)}
\NormalTok{pl.ylim(}\DecValTok{0}\NormalTok{, }\DecValTok{1000}\NormalTok{)}
\end{Highlighting}
\end{Shaded}

\pandocbounded{\includegraphics[keepaspectratio]{ajuste/ajuste_files/figure-pdf/cell-17-output-1.pdf}}

\begin{Shaded}
\begin{Highlighting}[]
\NormalTok{pl.plot(param\_b, S, }\StringTok{\textquotesingle{}k.\textquotesingle{}}\NormalTok{)}
\NormalTok{pl.axvline(param\_b[S.argmin()], color}\OperatorTok{=}\StringTok{\textquotesingle{}red\textquotesingle{}}\NormalTok{, ls}\OperatorTok{=}\StringTok{\textquotesingle{}{-}{-}\textquotesingle{}}\NormalTok{)}
\NormalTok{pl.ylim(}\DecValTok{0}\NormalTok{, }\DecValTok{1000}\NormalTok{)}
\end{Highlighting}
\end{Shaded}

\pandocbounded{\includegraphics[keepaspectratio]{ajuste/ajuste_files/figure-pdf/cell-18-output-1.pdf}}

\chapter{Exercício 2}\label{exercuxedcio-2}

Vamos ajustar os parâmetros da função

\subsubsection{\texorpdfstring{\(f(x) = e^{\frac{-x}{\beta_1}} + \beta_2\)}{f(x) = e\^{}\{\textbackslash frac\{-x\}\{\textbackslash beta\_1\}\} + \textbackslash beta\_2}}\label{fx-efrac-xbeta_1-beta_2}

\begin{Shaded}
\begin{Highlighting}[]
\CommentTok{\# Vamos Ler o arquivo de entrada}
\NormalTok{x, y }\OperatorTok{=}\NormalTok{ np.loadtxt(}\StringTok{\textquotesingle{}exercicio\_2.dat\textquotesingle{}}\NormalTok{, unpack}\OperatorTok{=}\VariableTok{True}\NormalTok{)}
\end{Highlighting}
\end{Shaded}

\begin{Shaded}
\begin{Highlighting}[]
\CommentTok{\# Vamos criar a função que queremos ajustar}
\KeywordTok{def}\NormalTok{ funcao(x, beta\_1, beta\_2):}
    \ControlFlowTok{return}\NormalTok{ np.exp(}\OperatorTok{{-}}\NormalTok{x}\OperatorTok{/}\NormalTok{beta\_1) }\OperatorTok{+}\NormalTok{ beta\_2}
\end{Highlighting}
\end{Shaded}

\begin{Shaded}
\begin{Highlighting}[]
\CommentTok{\# Visuaizando os dados}
\NormalTok{pl.plot(x, y, }\StringTok{\textquotesingle{}k.{-}\textquotesingle{}}\NormalTok{)}
\NormalTok{pl.show()}
\end{Highlighting}
\end{Shaded}

\pandocbounded{\includegraphics[keepaspectratio]{ajuste/ajuste_files/figure-pdf/cell-21-output-1.pdf}}

\begin{Shaded}
\begin{Highlighting}[]
\CommentTok{\# Visualizando os dados e um modelo qualquer}
\NormalTok{pl.plot(x, y, }\StringTok{\textquotesingle{}k.{-}\textquotesingle{}}\NormalTok{)}

\NormalTok{y\_model }\OperatorTok{=}\NormalTok{ funcao(x, beta\_1}\OperatorTok{=}\DecValTok{1}\NormalTok{, beta\_2}\OperatorTok{={-}}\DecValTok{4}\NormalTok{)}

\CommentTok{\# Computando o valor dos mínimos quadrados}
\NormalTok{S }\OperatorTok{=}\NormalTok{ np.}\BuiltInTok{sum}\NormalTok{((y }\OperatorTok{{-}}\NormalTok{ y\_model)}\OperatorTok{**}\DecValTok{2}\NormalTok{)}
\BuiltInTok{print}\NormalTok{(S)}
\NormalTok{pl.plot(x, y\_model, }\StringTok{\textquotesingle{}r{-}\textquotesingle{}}\NormalTok{)}
\NormalTok{pl.show()}
\end{Highlighting}
\end{Shaded}

\begin{verbatim}
1539.7626219543715
\end{verbatim}

\pandocbounded{\includegraphics[keepaspectratio]{ajuste/ajuste_files/figure-pdf/cell-22-output-2.pdf}}

\begin{Shaded}
\begin{Highlighting}[]
\NormalTok{beta\_1 }\OperatorTok{=}\NormalTok{ np.random.uniform(}\OperatorTok{{-}}\DecValTok{100}\NormalTok{, }\DecValTok{100}\NormalTok{, }\DecValTok{100000}\NormalTok{)}
\NormalTok{beta\_2 }\OperatorTok{=}\NormalTok{ np.random.uniform(}\OperatorTok{{-}}\DecValTok{100}\NormalTok{, }\DecValTok{100}\NormalTok{, }\DecValTok{100000}\NormalTok{)}
\end{Highlighting}
\end{Shaded}

\begin{Shaded}
\begin{Highlighting}[]
\NormalTok{pl.plot(beta\_1, beta\_2, }\StringTok{\textquotesingle{}k.\textquotesingle{}}\NormalTok{)}
\end{Highlighting}
\end{Shaded}

\pandocbounded{\includegraphics[keepaspectratio]{ajuste/ajuste_files/figure-pdf/cell-24-output-1.pdf}}

\begin{Shaded}
\begin{Highlighting}[]
\ImportTok{from}\NormalTok{ tqdm }\ImportTok{import}\NormalTok{ tqdm}

\NormalTok{S }\OperatorTok{=}\NormalTok{ np.array([])}

\ControlFlowTok{for}\NormalTok{ i }\KeywordTok{in}\NormalTok{ tqdm(}\BuiltInTok{range}\NormalTok{(}\BuiltInTok{len}\NormalTok{(beta\_1))): }
\NormalTok{    y\_model }\OperatorTok{=}\NormalTok{ funcao(x, beta\_1}\OperatorTok{=}\NormalTok{beta\_1[i], beta\_2}\OperatorTok{=}\NormalTok{beta\_2[i])}
\NormalTok{    S }\OperatorTok{=}\NormalTok{ np.append(S, np.}\BuiltInTok{sum}\NormalTok{((y }\OperatorTok{{-}}\NormalTok{ y\_model)}\OperatorTok{**}\DecValTok{2}\NormalTok{))}
\end{Highlighting}
\end{Shaded}

\begin{verbatim}

  0%|          | 0/100000 [00:00<?, ?it/s]C:\Users\luanx\AppData\Local\Temp\ipykernel_13228\825112861.py:3: RuntimeWarning:

overflow encountered in exp

C:\Users\luanx\AppData\Local\Temp\ipykernel_13228\3067340706.py:7: RuntimeWarning:

overflow encountered in square


  7%|▋         | 7002/100000 [00:00<00:01, 69836.58it/s]
 15%|█▍        | 14955/100000 [00:00<00:01, 75416.36it/s]
 23%|██▎       | 22854/100000 [00:00<00:01, 76856.54it/s]
 31%|███       | 30540/100000 [00:00<00:00, 76374.62it/s]
 38%|███▊      | 38178/100000 [00:00<00:00, 71864.08it/s]
 45%|████▌     | 45406/100000 [00:00<00:00, 66237.87it/s]
 52%|█████▏    | 52116/100000 [00:00<00:00, 63554.98it/s]
 59%|█████▊    | 58532/100000 [00:00<00:00, 59774.86it/s]
 65%|██████▍   | 64569/100000 [00:01<00:00, 56893.51it/s]
 70%|███████   | 70303/100000 [00:01<00:00, 52431.88it/s]
 76%|███████▌  | 75610/100000 [00:01<00:00, 47473.77it/s]
 80%|████████  | 80448/100000 [00:01<00:00, 42129.95it/s]
 85%|████████▍ | 84783/100000 [00:01<00:00, 36549.50it/s]
 89%|████████▊ | 88596/100000 [00:01<00:00, 33693.98it/s]
 92%|█████████▏| 92072/100000 [00:01<00:00, 31875.69it/s]
 95%|█████████▌| 95318/100000 [00:01<00:00, 29974.40it/s]
 98%|█████████▊| 98346/100000 [00:02<00:00, 28571.87it/s]
100%|██████████| 100000/100000 [00:02<00:00, 45847.05it/s]
\end{verbatim}

\begin{Shaded}
\begin{Highlighting}[]
\NormalTok{pl.plot(beta\_1, S, }\StringTok{\textquotesingle{}k.\textquotesingle{}}\NormalTok{)}
\NormalTok{pl.ylim(}\DecValTok{0}\NormalTok{, }\DecValTok{10}\NormalTok{)}
\CommentTok{\#pl.xlim(0, 50)}
\end{Highlighting}
\end{Shaded}

\pandocbounded{\includegraphics[keepaspectratio]{ajuste/ajuste_files/figure-pdf/cell-26-output-1.pdf}}

\begin{Shaded}
\begin{Highlighting}[]
\NormalTok{pl.plot(beta\_2, S, }\StringTok{\textquotesingle{}k.\textquotesingle{}}\NormalTok{)}
\NormalTok{pl.ylim(}\DecValTok{0}\NormalTok{, }\DecValTok{2}\NormalTok{)}
\NormalTok{pl.xlim(}\DecValTok{0}\NormalTok{, }\DecValTok{5}\NormalTok{)}
\end{Highlighting}
\end{Shaded}

\pandocbounded{\includegraphics[keepaspectratio]{ajuste/ajuste_files/figure-pdf/cell-27-output-1.pdf}}

\begin{Shaded}
\begin{Highlighting}[]
\BuiltInTok{print}\NormalTok{(beta\_1[S.argmin()])}
\BuiltInTok{print}\NormalTok{(beta\_2[S.argmin()])}
\end{Highlighting}
\end{Shaded}

\begin{verbatim}
17.841378154564282
3.031074431670149
\end{verbatim}

\begin{Shaded}
\begin{Highlighting}[]
\NormalTok{pl.plot(x, y, }\StringTok{\textquotesingle{}k.{-}\textquotesingle{}}\NormalTok{)}

\NormalTok{y\_model }\OperatorTok{=}\NormalTok{ funcao(x, beta\_1}\OperatorTok{=}\FloatTok{15.446324928570547}\NormalTok{, beta\_2}\OperatorTok{=}\FloatTok{2.987769552921577}\NormalTok{)}

\NormalTok{S }\OperatorTok{=}\NormalTok{ np.}\BuiltInTok{sum}\NormalTok{((y }\OperatorTok{{-}}\NormalTok{ y\_model)}\OperatorTok{**}\DecValTok{2}\NormalTok{)}
\BuiltInTok{print}\NormalTok{(S)}
\NormalTok{pl.plot(x, y\_model, }\StringTok{\textquotesingle{}r{-}\textquotesingle{}}\NormalTok{)}
\NormalTok{pl.show()}
\end{Highlighting}
\end{Shaded}

\begin{verbatim}
0.10738497702461611
\end{verbatim}

\pandocbounded{\includegraphics[keepaspectratio]{ajuste/ajuste_files/figure-pdf/cell-29-output-2.pdf}}

O ajuste parece OK, mas será que estamos de fato nos melhores valores?

\begin{Shaded}
\begin{Highlighting}[]
\CommentTok{\# Refazer, agora mais próximo da solução}

\NormalTok{beta\_1 }\OperatorTok{=}\NormalTok{ np.random.uniform(}\DecValTok{0}\NormalTok{, }\DecValTok{100}\NormalTok{, }\DecValTok{100000}\NormalTok{)}
\NormalTok{beta\_2 }\OperatorTok{=}\NormalTok{ np.random.uniform(}\OperatorTok{{-}}\DecValTok{10}\NormalTok{, }\DecValTok{10}\NormalTok{, }\DecValTok{100000}\NormalTok{)}
\end{Highlighting}
\end{Shaded}

\begin{Shaded}
\begin{Highlighting}[]
\ImportTok{from}\NormalTok{ tqdm }\ImportTok{import}\NormalTok{ tqdm}

\NormalTok{S }\OperatorTok{=}\NormalTok{ np.array([])}

\ControlFlowTok{for}\NormalTok{ i }\KeywordTok{in}\NormalTok{ tqdm(}\BuiltInTok{range}\NormalTok{(}\BuiltInTok{len}\NormalTok{(beta\_1))): }
\NormalTok{    y\_model }\OperatorTok{=}\NormalTok{ funcao(x, beta\_1}\OperatorTok{=}\NormalTok{beta\_1[i], beta\_2}\OperatorTok{=}\NormalTok{beta\_2[i])}
\NormalTok{    S }\OperatorTok{=}\NormalTok{ np.append(S, np.}\BuiltInTok{sum}\NormalTok{((y }\OperatorTok{{-}}\NormalTok{ y\_model)}\OperatorTok{**}\DecValTok{2}\NormalTok{))}
\end{Highlighting}
\end{Shaded}

\begin{verbatim}

  0%|          | 0/100000 [00:00<?, ?it/s]
 10%|█         | 10395/100000 [00:00<00:00, 103525.81it/s]
 21%|██        | 20748/100000 [00:00<00:00, 91741.49it/s] 
 30%|███       | 30008/100000 [00:00<00:00, 86828.49it/s]
 39%|███▊      | 38741/100000 [00:00<00:00, 76154.60it/s]
 47%|████▋     | 46505/100000 [00:00<00:00, 69013.54it/s]
 54%|█████▎    | 53541/100000 [00:00<00:00, 65327.08it/s]
 60%|██████    | 60150/100000 [00:00<00:00, 57040.07it/s]
 66%|██████▌   | 66006/100000 [00:01<00:00, 50367.95it/s]
 71%|███████   | 71205/100000 [00:01<00:00, 48582.99it/s]
 76%|███████▌  | 76151/100000 [00:01<00:00, 45412.75it/s]
 81%|████████  | 80746/100000 [00:01<00:00, 42248.35it/s]
 85%|████████▌ | 85006/100000 [00:01<00:00, 38808.70it/s]
 89%|████████▉ | 88919/100000 [00:01<00:00, 35927.91it/s]
 93%|█████████▎| 92534/100000 [00:01<00:00, 34516.99it/s]
 96%|█████████▌| 95988/100000 [00:01<00:00, 32955.85it/s]
 99%|█████████▉| 99276/100000 [00:02<00:00, 31762.93it/s]
100%|██████████| 100000/100000 [00:02<00:00, 48356.46it/s]
\end{verbatim}

\begin{Shaded}
\begin{Highlighting}[]
\NormalTok{pl.plot(beta\_1, S, }\StringTok{\textquotesingle{}k.\textquotesingle{}}\NormalTok{)}
\NormalTok{pl.ylim(}\DecValTok{0}\NormalTok{, }\DecValTok{1}\NormalTok{)}
\end{Highlighting}
\end{Shaded}

\pandocbounded{\includegraphics[keepaspectratio]{ajuste/ajuste_files/figure-pdf/cell-32-output-1.pdf}}

\begin{Shaded}
\begin{Highlighting}[]
\NormalTok{pl.plot(beta\_2, S, }\StringTok{\textquotesingle{}k.\textquotesingle{}}\NormalTok{)}
\NormalTok{pl.ylim(}\DecValTok{0}\NormalTok{, }\DecValTok{1}\NormalTok{)}
\end{Highlighting}
\end{Shaded}

\pandocbounded{\includegraphics[keepaspectratio]{ajuste/ajuste_files/figure-pdf/cell-33-output-1.pdf}}

\begin{Shaded}
\begin{Highlighting}[]
\BuiltInTok{print}\NormalTok{(beta\_1[S.argmin()])}
\BuiltInTok{print}\NormalTok{(beta\_2[S.argmin()])}
\end{Highlighting}
\end{Shaded}

\begin{verbatim}
18.406046127938
3.000448007550089
\end{verbatim}

\begin{Shaded}
\begin{Highlighting}[]
\NormalTok{pl.plot(x, y, }\StringTok{\textquotesingle{}k.{-}\textquotesingle{}}\NormalTok{)}

\NormalTok{y\_model }\OperatorTok{=}\NormalTok{ funcao(x, beta\_1}\OperatorTok{=}\FloatTok{18.77266155859759}\NormalTok{, beta\_2}\OperatorTok{=}\FloatTok{2.9917787435681813}\NormalTok{)}

\NormalTok{S }\OperatorTok{=}\NormalTok{ np.}\BuiltInTok{sum}\NormalTok{((y }\OperatorTok{{-}}\NormalTok{ y\_model)}\OperatorTok{**}\DecValTok{2}\NormalTok{)}
\BuiltInTok{print}\NormalTok{(S)}
\NormalTok{pl.plot(x, y\_model, }\StringTok{\textquotesingle{}r{-}\textquotesingle{}}\NormalTok{)}
\NormalTok{pl.show()}
\end{Highlighting}
\end{Shaded}

\begin{verbatim}
0.04337814021210659
\end{verbatim}

\pandocbounded{\includegraphics[keepaspectratio]{ajuste/ajuste_files/figure-pdf/cell-35-output-2.pdf}}

Quanto mais modelos fizermos mais próximos do resultado real estaremos!

\chapter{Como fazer com o Scipy}\label{como-fazer-com-o-scipy-1}

\begin{Shaded}
\begin{Highlighting}[]
\ImportTok{from}\NormalTok{ scipy.optimize }\ImportTok{import}\NormalTok{ curve\_fit}
\end{Highlighting}
\end{Shaded}

\begin{Shaded}
\begin{Highlighting}[]
\CommentTok{\# Funçao Curve\_fit fará o ajuste para você de forma eficiente}
\NormalTok{param, error }\OperatorTok{=}\NormalTok{ curve\_fit(funcao, x, y)}
\end{Highlighting}
\end{Shaded}

\begin{Shaded}
\begin{Highlighting}[]
\BuiltInTok{print}\NormalTok{(}\StringTok{\textquotesingle{}Valores: \textquotesingle{}}\NormalTok{, param)}
\BuiltInTok{print}\NormalTok{(}\StringTok{\textquotesingle{}Incertezas: \textquotesingle{}}\NormalTok{, np.sqrt(np.diagonal(error)))}
\end{Highlighting}
\end{Shaded}

\begin{verbatim}
Valores:  [18.4336622  3.0006474]
Incertezas:  [1.08368818 0.01240034]
\end{verbatim}

\chapter{Isso é tudo}\label{isso-uxe9-tudo-5}

By BEM

\chapter{Simulações de Montecarlo}\label{simulauxe7uxf5es-de-montecarlo}

\chapter{Cálculo Numérico}\label{cuxe1lculo-numuxe9rico}

\begin{Shaded}
\begin{Highlighting}[]
\ImportTok{import}\NormalTok{ numpy }\ImportTok{as}\NormalTok{ np}
\ImportTok{import}\NormalTok{ matplotlib.pyplot }\ImportTok{as}\NormalTok{ pl}
\end{Highlighting}
\end{Shaded}

\chapter{Calcular numéricamente a Derivada de
Funções}\label{calcular-numuxe9ricamente-a-derivada-de-funuxe7uxf5es}

\begin{Shaded}
\begin{Highlighting}[]
\CommentTok{\# Vamos definir a função}
\KeywordTok{def}\NormalTok{ funcao(x):}
    \ControlFlowTok{return}\NormalTok{ x}\OperatorTok{**}\DecValTok{2} \OperatorTok{+} \DecValTok{10} 
\end{Highlighting}
\end{Shaded}

\begin{Shaded}
\begin{Highlighting}[]
\CommentTok{\# Para visualização, como a curva se comporta}
\NormalTok{x }\OperatorTok{=}\NormalTok{ np.arange(}\OperatorTok{{-}}\DecValTok{10}\NormalTok{, }\OperatorTok{+}\DecValTok{10}\NormalTok{, }\FloatTok{0.0001}\NormalTok{)}
\NormalTok{pl.plot(x, funcao(x), }\StringTok{\textquotesingle{}k{-}\textquotesingle{}}\NormalTok{)}
\NormalTok{pl.show()}
\end{Highlighting}
\end{Shaded}

\pandocbounded{\includegraphics[keepaspectratio]{calculo/calculo_files/figure-pdf/cell-4-output-1.pdf}}

\begin{Shaded}
\begin{Highlighting}[]
\CommentTok{\# Definir um passo e determinar a derivada}

\NormalTok{h }\OperatorTok{=} \FloatTok{0.001}
\NormalTok{dx }\OperatorTok{=}\NormalTok{ (funcao(x }\OperatorTok{+}\NormalTok{ h) }\OperatorTok{{-}}\NormalTok{ funcao(x }\OperatorTok{{-}}\NormalTok{ h))}\OperatorTok{/}\NormalTok{(}\DecValTok{2}\OperatorTok{*}\NormalTok{h)}
\end{Highlighting}
\end{Shaded}

\begin{Shaded}
\begin{Highlighting}[]
\CommentTok{\# Vamos visualizar a derivada}
\NormalTok{pl.plot(x, dx, }\StringTok{\textquotesingle{}r{-}\textquotesingle{}}\NormalTok{)}
\end{Highlighting}
\end{Shaded}

\pandocbounded{\includegraphics[keepaspectratio]{calculo/calculo_files/figure-pdf/cell-6-output-1.pdf}}

\begin{Shaded}
\begin{Highlighting}[]
\CommentTok{\# Como o que obtivemos numéricamente se assemelha com o resultado algebrico?}

\NormalTok{pl.plot(x, }\DecValTok{2}\OperatorTok{*}\NormalTok{x,color}\OperatorTok{=}\StringTok{\textquotesingle{}g\textquotesingle{}}\NormalTok{, lw}\OperatorTok{=}\DecValTok{5}\NormalTok{, ls}\OperatorTok{=}\StringTok{\textquotesingle{}{-}{-}\textquotesingle{}}\NormalTok{)}
\NormalTok{pl.plot(x, dx, }\StringTok{\textquotesingle{}r{-}\textquotesingle{}}\NormalTok{)}

\CommentTok{\# Perfeitamente!}
\end{Highlighting}
\end{Shaded}

\pandocbounded{\includegraphics[keepaspectratio]{calculo/calculo_files/figure-pdf/cell-7-output-1.pdf}}

\subsection{Como fazer a derivada numérica de
dados?}\label{como-fazer-a-derivada-numuxe9rica-de-dados}

\begin{Shaded}
\begin{Highlighting}[]
\CommentTok{\# Lendo nossos dados}

\NormalTok{x, y }\OperatorTok{=}\NormalTok{ np.loadtxt(}\StringTok{\textquotesingle{}exercicio\_3.dat\textquotesingle{}}\NormalTok{, unpack}\OperatorTok{=}\VariableTok{True}\NormalTok{)}
\end{Highlighting}
\end{Shaded}

\begin{Shaded}
\begin{Highlighting}[]
\CommentTok{\# Visualizando}

\NormalTok{pl.plot(x, y, }\StringTok{\textquotesingle{}k{-}\textquotesingle{}}\NormalTok{)}
\NormalTok{pl.show()}
\end{Highlighting}
\end{Shaded}

\pandocbounded{\includegraphics[keepaspectratio]{calculo/calculo_files/figure-pdf/cell-9-output-1.pdf}}

\begin{Shaded}
\begin{Highlighting}[]
\CommentTok{\# Computando o menor passo possível, a diferença entre dois pontos}
\BuiltInTok{print}\NormalTok{(x[}\DecValTok{2}\NormalTok{] }\OperatorTok{{-}}\NormalTok{ x[}\DecValTok{0}\NormalTok{] )}
\BuiltInTok{print}\NormalTok{(x[}\DecValTok{3}\NormalTok{] }\OperatorTok{{-}}\NormalTok{ x[}\DecValTok{1}\NormalTok{] )}
\end{Highlighting}
\end{Shaded}

\begin{verbatim}
0.1
0.09999999999999999
\end{verbatim}

\begin{Shaded}
\begin{Highlighting}[]
\CommentTok{\# Determinando a derivada}

\NormalTok{derivada }\OperatorTok{=}\NormalTok{ np.array([])}

\ControlFlowTok{for}\NormalTok{ i }\KeywordTok{in} \BuiltInTok{range}\NormalTok{(}\DecValTok{1}\NormalTok{, }\BuiltInTok{len}\NormalTok{(x) }\OperatorTok{{-}} \DecValTok{1}\NormalTok{):}
\NormalTok{    dx }\OperatorTok{=}\NormalTok{ (y[i}\OperatorTok{+}\DecValTok{1}\NormalTok{] }\OperatorTok{{-}}\NormalTok{ y[i}\OperatorTok{{-}}\DecValTok{1}\NormalTok{])}\OperatorTok{/}\NormalTok{(x[i}\OperatorTok{+}\DecValTok{1}\NormalTok{] }\OperatorTok{{-}}\NormalTok{ x[i}\OperatorTok{{-}}\DecValTok{1}\NormalTok{])}
\NormalTok{    derivada }\OperatorTok{=}\NormalTok{ np.append(derivada, dx)}
\end{Highlighting}
\end{Shaded}

\begin{Shaded}
\begin{Highlighting}[]
\CommentTok{\# Visualizando a derivada}
\NormalTok{pl.plot(x[}\DecValTok{1}\NormalTok{:}\OperatorTok{{-}}\DecValTok{1}\NormalTok{], derivada, }\StringTok{\textquotesingle{}r{-}\textquotesingle{}}\NormalTok{)}
\NormalTok{pl.show()}
\end{Highlighting}
\end{Shaded}

\pandocbounded{\includegraphics[keepaspectratio]{calculo/calculo_files/figure-pdf/cell-12-output-1.pdf}}

\subsubsection{Como fazer com o scipy?}\label{como-fazer-com-o-scipy-2}

\begin{Shaded}
\begin{Highlighting}[]
\ImportTok{from}\NormalTok{ scipy.differentiate }\ImportTok{import}\NormalTok{ derivative}
\end{Highlighting}
\end{Shaded}

\begin{Shaded}
\begin{Highlighting}[]
\NormalTok{resultado }\OperatorTok{=}\NormalTok{ derivative(funcao, x)}
\NormalTok{dx }\OperatorTok{=}\NormalTok{ resultado.df}
\end{Highlighting}
\end{Shaded}

\begin{Shaded}
\begin{Highlighting}[]
\NormalTok{pl.plot(x, dx)}
\end{Highlighting}
\end{Shaded}

\pandocbounded{\includegraphics[keepaspectratio]{calculo/calculo_files/figure-pdf/cell-15-output-1.pdf}}

\chapter{Calcular numéricamente a Integral de
Funções}\label{calcular-numuxe9ricamente-a-integral-de-funuxe7uxf5es}

\begin{Shaded}
\begin{Highlighting}[]
\CommentTok{\# Vamos definir a função}
\KeywordTok{def}\NormalTok{ funcao(x):}
    \ControlFlowTok{return}\NormalTok{ x}\OperatorTok{**}\DecValTok{2} \OperatorTok{+} \DecValTok{10} 
\end{Highlighting}
\end{Shaded}

\begin{Shaded}
\begin{Highlighting}[]
\NormalTok{x }\OperatorTok{=}\NormalTok{ np.arange(}\OperatorTok{{-}}\DecValTok{10}\NormalTok{, }\OperatorTok{+}\DecValTok{10}\NormalTok{, }\FloatTok{0.0001}\NormalTok{)}

\NormalTok{h }\OperatorTok{=} \FloatTok{0.0001}
\NormalTok{integral }\OperatorTok{=} \DecValTok{0}

\ControlFlowTok{for}\NormalTok{ i }\KeywordTok{in} \BuiltInTok{range}\NormalTok{(}\BuiltInTok{len}\NormalTok{(x) }\OperatorTok{{-}} \DecValTok{1}\NormalTok{):}
\NormalTok{    integral }\OperatorTok{+=}\NormalTok{ (h}\OperatorTok{/}\DecValTok{2}\NormalTok{)}\OperatorTok{*}\NormalTok{(funcao(x[i]) }\OperatorTok{+}\NormalTok{ funcao(x[i}\OperatorTok{+}\DecValTok{1}\NormalTok{]))}
    
\BuiltInTok{print}\NormalTok{(integral)}
\end{Highlighting}
\end{Shaded}

\begin{verbatim}
866.655666796912
\end{verbatim}

Algebricamente o resultado seria:

\(\frac{x^3}{3} + 10x\), começando em x = -10, até x = 10

\begin{Shaded}
\begin{Highlighting}[]
\NormalTok{xx }\OperatorTok{=}\NormalTok{ [}\OperatorTok{{-}}\DecValTok{10}\NormalTok{, }\DecValTok{10}\NormalTok{]}
\NormalTok{((xx[}\DecValTok{1}\NormalTok{]}\OperatorTok{**}\DecValTok{3}\NormalTok{)}\OperatorTok{/}\DecValTok{3} \OperatorTok{+} \DecValTok{10}\OperatorTok{*}\NormalTok{xx[}\DecValTok{1}\NormalTok{]) }\OperatorTok{{-}}\NormalTok{ ((xx[}\DecValTok{0}\NormalTok{]}\OperatorTok{**}\DecValTok{3}\NormalTok{)}\OperatorTok{/}\DecValTok{3} \OperatorTok{+} \DecValTok{10}\OperatorTok{*}\NormalTok{xx[}\DecValTok{0}\NormalTok{])}
\end{Highlighting}
\end{Shaded}

\begin{verbatim}
866.6666666666666
\end{verbatim}

Quanto menor for o passo mais próximo do resultado correto chegaremos

\subsection{Como fazer a integral numérica de
dados?}\label{como-fazer-a-integral-numuxe9rica-de-dados}

\begin{Shaded}
\begin{Highlighting}[]
\CommentTok{\# Lendo o arquivo de entrada}
\NormalTok{x, y }\OperatorTok{=}\NormalTok{ np.loadtxt(}\StringTok{\textquotesingle{}exercicio\_4.dat\textquotesingle{}}\NormalTok{, unpack}\OperatorTok{=}\VariableTok{True}\NormalTok{)}
\end{Highlighting}
\end{Shaded}

\begin{Shaded}
\begin{Highlighting}[]
\CommentTok{\# Visualizando}
\NormalTok{pl.plot(x, y)}
\end{Highlighting}
\end{Shaded}

\pandocbounded{\includegraphics[keepaspectratio]{calculo/calculo_files/figure-pdf/cell-20-output-1.pdf}}

\begin{Shaded}
\begin{Highlighting}[]
\NormalTok{integral }\OperatorTok{=} \DecValTok{0}

\ControlFlowTok{for}\NormalTok{ i }\KeywordTok{in} \BuiltInTok{range}\NormalTok{(}\BuiltInTok{len}\NormalTok{(x) }\OperatorTok{{-}} \DecValTok{1}\NormalTok{):}
\NormalTok{    integral }\OperatorTok{+=}\NormalTok{ ((x[i}\OperatorTok{+}\DecValTok{1}\NormalTok{] }\OperatorTok{{-}}\NormalTok{ x[i])}\OperatorTok{/}\DecValTok{2}\NormalTok{)}\OperatorTok{*}\NormalTok{(y[i] }\OperatorTok{+}\NormalTok{ y[i}\OperatorTok{+}\DecValTok{1}\NormalTok{])}
    
\BuiltInTok{print}\NormalTok{(integral)}
\end{Highlighting}
\end{Shaded}

\begin{verbatim}
0.9999961999999989
\end{verbatim}

\subsubsection{Como fazer com o scipy?}\label{como-fazer-com-o-scipy-3}

\begin{Shaded}
\begin{Highlighting}[]
\ImportTok{from}\NormalTok{ scipy.integrate }\ImportTok{import}\NormalTok{ trapezoid}

\CommentTok{\#\#\# Ou}

\CommentTok{\# from scipy.integrate import trapz \# No caso de Python 2}
\end{Highlighting}
\end{Shaded}

\begin{Shaded}
\begin{Highlighting}[]
\NormalTok{pl.plot(x, y)}
\NormalTok{pl.fill\_between(x, y, }\DecValTok{0}\NormalTok{)}
\end{Highlighting}
\end{Shaded}

\pandocbounded{\includegraphics[keepaspectratio]{calculo/calculo_files/figure-pdf/cell-23-output-1.pdf}}

\begin{Shaded}
\begin{Highlighting}[]
\NormalTok{trapezoid(y, x)}
\end{Highlighting}
\end{Shaded}

\begin{verbatim}
np.float64(0.9999961999999998)
\end{verbatim}

\chapter{Isso é tudo}\label{isso-uxe9-tudo-6}

By BEM

\subsubsection{by BEM, Last update:
10/07/2023}\label{by-bem-last-update-10072023-2}

\begin{Shaded}
\begin{Highlighting}[]
\ImportTok{import}\NormalTok{ numpy }\ImportTok{as}\NormalTok{ np}
\ImportTok{import}\NormalTok{ matplotlib.pyplot }\ImportTok{as}\NormalTok{ pl}
\end{Highlighting}
\end{Shaded}

\chapter{Calcular numéricamente a Derivada de
Funções}\label{calcular-numuxe9ricamente-a-derivada-de-funuxe7uxf5es-1}

\begin{Shaded}
\begin{Highlighting}[]
\CommentTok{\# Vamos definir a função}
\KeywordTok{def}\NormalTok{ funcao(x):}
    \ControlFlowTok{return}\NormalTok{ x}\OperatorTok{**}\DecValTok{2} \OperatorTok{+} \DecValTok{10} 
\end{Highlighting}
\end{Shaded}

\begin{Shaded}
\begin{Highlighting}[]
\CommentTok{\# Para visualização, como a curva se comporta}
\NormalTok{x }\OperatorTok{=}\NormalTok{ np.arange(}\OperatorTok{{-}}\DecValTok{10}\NormalTok{, }\OperatorTok{+}\DecValTok{10}\NormalTok{, }\FloatTok{0.0001}\NormalTok{)}
\NormalTok{pl.plot(x, funcao(x), }\StringTok{\textquotesingle{}k{-}\textquotesingle{}}\NormalTok{)}
\NormalTok{pl.show()}
\end{Highlighting}
\end{Shaded}

\pandocbounded{\includegraphics[keepaspectratio]{calculo/calculo_files/figure-pdf/cell-27-output-1.pdf}}

\begin{Shaded}
\begin{Highlighting}[]
\CommentTok{\# Definir um passo e determinar a derivada}

\NormalTok{h }\OperatorTok{=} \FloatTok{0.001}
\NormalTok{dx }\OperatorTok{=}\NormalTok{ (funcao(x }\OperatorTok{+}\NormalTok{ h) }\OperatorTok{{-}}\NormalTok{ funcao(x }\OperatorTok{{-}}\NormalTok{ h))}\OperatorTok{/}\NormalTok{(}\DecValTok{2}\OperatorTok{*}\NormalTok{h)}
\end{Highlighting}
\end{Shaded}

\begin{Shaded}
\begin{Highlighting}[]
\CommentTok{\# Vamos visualizar a derivada}
\NormalTok{pl.plot(x, dx, }\StringTok{\textquotesingle{}r{-}\textquotesingle{}}\NormalTok{)}
\end{Highlighting}
\end{Shaded}

\pandocbounded{\includegraphics[keepaspectratio]{calculo/calculo_files/figure-pdf/cell-29-output-1.pdf}}

\begin{Shaded}
\begin{Highlighting}[]
\CommentTok{\# Como o que obtivemos numéricamente se assemelha com o resultado algebrico?}

\NormalTok{pl.plot(x, }\DecValTok{2}\OperatorTok{*}\NormalTok{x,color}\OperatorTok{=}\StringTok{\textquotesingle{}g\textquotesingle{}}\NormalTok{, lw}\OperatorTok{=}\DecValTok{5}\NormalTok{, ls}\OperatorTok{=}\StringTok{\textquotesingle{}{-}{-}\textquotesingle{}}\NormalTok{)}
\NormalTok{pl.plot(x, dx, }\StringTok{\textquotesingle{}r{-}\textquotesingle{}}\NormalTok{)}

\CommentTok{\# Perfeitamente!}
\end{Highlighting}
\end{Shaded}

\pandocbounded{\includegraphics[keepaspectratio]{calculo/calculo_files/figure-pdf/cell-30-output-1.pdf}}

\subsection{Como fazer a derivada numérica de
dados?}\label{como-fazer-a-derivada-numuxe9rica-de-dados-1}

\begin{Shaded}
\begin{Highlighting}[]
\CommentTok{\# Lendo nossos dados}

\NormalTok{x, y }\OperatorTok{=}\NormalTok{ np.loadtxt(}\StringTok{\textquotesingle{}exercicio\_3.dat\textquotesingle{}}\NormalTok{, unpack}\OperatorTok{=}\VariableTok{True}\NormalTok{)}
\end{Highlighting}
\end{Shaded}

\begin{Shaded}
\begin{Highlighting}[]
\CommentTok{\# Visualizando}

\NormalTok{pl.plot(x, y, }\StringTok{\textquotesingle{}k{-}\textquotesingle{}}\NormalTok{)}
\NormalTok{pl.show()}
\end{Highlighting}
\end{Shaded}

\pandocbounded{\includegraphics[keepaspectratio]{calculo/calculo_files/figure-pdf/cell-32-output-1.pdf}}

\begin{Shaded}
\begin{Highlighting}[]
\CommentTok{\# Computando o menor passo possível, a diferença entre dois pontos}
\BuiltInTok{print}\NormalTok{(x[}\DecValTok{2}\NormalTok{] }\OperatorTok{{-}}\NormalTok{ x[}\DecValTok{0}\NormalTok{] )}
\BuiltInTok{print}\NormalTok{(x[}\DecValTok{3}\NormalTok{] }\OperatorTok{{-}}\NormalTok{ x[}\DecValTok{1}\NormalTok{] )}
\end{Highlighting}
\end{Shaded}

\begin{verbatim}
0.1
0.09999999999999999
\end{verbatim}

\begin{Shaded}
\begin{Highlighting}[]
\CommentTok{\# Determinando a derivada}

\NormalTok{derivada }\OperatorTok{=}\NormalTok{ np.array([])}

\ControlFlowTok{for}\NormalTok{ i }\KeywordTok{in} \BuiltInTok{range}\NormalTok{(}\DecValTok{1}\NormalTok{, }\BuiltInTok{len}\NormalTok{(x) }\OperatorTok{{-}} \DecValTok{1}\NormalTok{):}
\NormalTok{    dx }\OperatorTok{=}\NormalTok{ (y[i}\OperatorTok{+}\DecValTok{1}\NormalTok{] }\OperatorTok{{-}}\NormalTok{ y[i}\OperatorTok{{-}}\DecValTok{1}\NormalTok{])}\OperatorTok{/}\NormalTok{(x[i}\OperatorTok{+}\DecValTok{1}\NormalTok{] }\OperatorTok{{-}}\NormalTok{ x[i}\OperatorTok{{-}}\DecValTok{1}\NormalTok{])}
\NormalTok{    derivada }\OperatorTok{=}\NormalTok{ np.append(derivada, dx)}
\end{Highlighting}
\end{Shaded}

\begin{Shaded}
\begin{Highlighting}[]
\CommentTok{\# Visualizando a derivada}
\NormalTok{pl.plot(x[}\DecValTok{1}\NormalTok{:}\OperatorTok{{-}}\DecValTok{1}\NormalTok{], derivada, }\StringTok{\textquotesingle{}r{-}\textquotesingle{}}\NormalTok{)}
\NormalTok{pl.show()}
\end{Highlighting}
\end{Shaded}

\pandocbounded{\includegraphics[keepaspectratio]{calculo/calculo_files/figure-pdf/cell-35-output-1.pdf}}

\subsubsection{Como fazer com o scipy?}\label{como-fazer-com-o-scipy-4}

\begin{Shaded}
\begin{Highlighting}[]
\ImportTok{from}\NormalTok{ scipy.differentiate }\ImportTok{import}\NormalTok{ derivative}
\end{Highlighting}
\end{Shaded}

\begin{Shaded}
\begin{Highlighting}[]
\NormalTok{resultado }\OperatorTok{=}\NormalTok{ derivative(funcao, x)}
\NormalTok{dx }\OperatorTok{=}\NormalTok{ resultado.df}
\end{Highlighting}
\end{Shaded}

\begin{Shaded}
\begin{Highlighting}[]
\NormalTok{pl.plot(x, dx)}
\end{Highlighting}
\end{Shaded}

\pandocbounded{\includegraphics[keepaspectratio]{calculo/calculo_files/figure-pdf/cell-38-output-1.pdf}}

\chapter{Calcular numéricamente a Integral de
Funções}\label{calcular-numuxe9ricamente-a-integral-de-funuxe7uxf5es-1}

\begin{Shaded}
\begin{Highlighting}[]
\CommentTok{\# Vamos definir a função}
\KeywordTok{def}\NormalTok{ funcao(x):}
    \ControlFlowTok{return}\NormalTok{ x}\OperatorTok{**}\DecValTok{2} \OperatorTok{+} \DecValTok{10} 
\end{Highlighting}
\end{Shaded}

\begin{Shaded}
\begin{Highlighting}[]
\NormalTok{x }\OperatorTok{=}\NormalTok{ np.arange(}\OperatorTok{{-}}\DecValTok{10}\NormalTok{, }\OperatorTok{+}\DecValTok{10}\NormalTok{, }\FloatTok{0.0001}\NormalTok{)}

\NormalTok{h }\OperatorTok{=} \FloatTok{0.0001}
\NormalTok{integral }\OperatorTok{=} \DecValTok{0}

\ControlFlowTok{for}\NormalTok{ i }\KeywordTok{in} \BuiltInTok{range}\NormalTok{(}\BuiltInTok{len}\NormalTok{(x) }\OperatorTok{{-}} \DecValTok{1}\NormalTok{):}
\NormalTok{    integral }\OperatorTok{+=}\NormalTok{ (h}\OperatorTok{/}\DecValTok{2}\NormalTok{)}\OperatorTok{*}\NormalTok{(funcao(x[i]) }\OperatorTok{+}\NormalTok{ funcao(x[i}\OperatorTok{+}\DecValTok{1}\NormalTok{]))}
    
\BuiltInTok{print}\NormalTok{(integral)}
\end{Highlighting}
\end{Shaded}

\begin{verbatim}
866.655666796912
\end{verbatim}

Algebricamente o resultado seria:

\(\frac{x^3}{3} + 10x\), começando em x = -10, até x = 10

\begin{Shaded}
\begin{Highlighting}[]
\NormalTok{xx }\OperatorTok{=}\NormalTok{ [}\OperatorTok{{-}}\DecValTok{10}\NormalTok{, }\DecValTok{10}\NormalTok{]}
\NormalTok{((xx[}\DecValTok{1}\NormalTok{]}\OperatorTok{**}\DecValTok{3}\NormalTok{)}\OperatorTok{/}\DecValTok{3} \OperatorTok{+} \DecValTok{10}\OperatorTok{*}\NormalTok{xx[}\DecValTok{1}\NormalTok{]) }\OperatorTok{{-}}\NormalTok{ ((xx[}\DecValTok{0}\NormalTok{]}\OperatorTok{**}\DecValTok{3}\NormalTok{)}\OperatorTok{/}\DecValTok{3} \OperatorTok{+} \DecValTok{10}\OperatorTok{*}\NormalTok{xx[}\DecValTok{0}\NormalTok{])}
\end{Highlighting}
\end{Shaded}

\begin{verbatim}
866.6666666666666
\end{verbatim}

Quanto menor for o passo mais próximo do resultado correto chegaremos

\subsection{Como fazer a integral numérica de
dados?}\label{como-fazer-a-integral-numuxe9rica-de-dados-1}

\begin{Shaded}
\begin{Highlighting}[]
\CommentTok{\# Lendo o arquivo de entrada}
\NormalTok{x, y }\OperatorTok{=}\NormalTok{ np.loadtxt(}\StringTok{\textquotesingle{}exercicio\_4.dat\textquotesingle{}}\NormalTok{, unpack}\OperatorTok{=}\VariableTok{True}\NormalTok{)}
\end{Highlighting}
\end{Shaded}

\begin{Shaded}
\begin{Highlighting}[]
\CommentTok{\# Visualizando}
\NormalTok{pl.plot(x, y)}
\end{Highlighting}
\end{Shaded}

\pandocbounded{\includegraphics[keepaspectratio]{calculo/calculo_files/figure-pdf/cell-43-output-1.pdf}}

\begin{Shaded}
\begin{Highlighting}[]
\NormalTok{integral }\OperatorTok{=} \DecValTok{0}

\ControlFlowTok{for}\NormalTok{ i }\KeywordTok{in} \BuiltInTok{range}\NormalTok{(}\BuiltInTok{len}\NormalTok{(x) }\OperatorTok{{-}} \DecValTok{1}\NormalTok{):}
\NormalTok{    integral }\OperatorTok{+=}\NormalTok{ ((x[i}\OperatorTok{+}\DecValTok{1}\NormalTok{] }\OperatorTok{{-}}\NormalTok{ x[i])}\OperatorTok{/}\DecValTok{2}\NormalTok{)}\OperatorTok{*}\NormalTok{(y[i] }\OperatorTok{+}\NormalTok{ y[i}\OperatorTok{+}\DecValTok{1}\NormalTok{])}
    
\BuiltInTok{print}\NormalTok{(integral)}
\end{Highlighting}
\end{Shaded}

\begin{verbatim}
0.9999961999999989
\end{verbatim}

\subsubsection{Como fazer com o scipy?}\label{como-fazer-com-o-scipy-5}

\begin{Shaded}
\begin{Highlighting}[]
\ImportTok{from}\NormalTok{ scipy.integrate }\ImportTok{import}\NormalTok{ trapezoid}

\CommentTok{\#\#\# Ou}

\CommentTok{\# from scipy.integrate import trapz \# No caso de Python 2}
\end{Highlighting}
\end{Shaded}

\begin{Shaded}
\begin{Highlighting}[]
\NormalTok{pl.plot(x, y)}
\NormalTok{pl.fill\_between(x, y, }\DecValTok{0}\NormalTok{)}
\end{Highlighting}
\end{Shaded}

\pandocbounded{\includegraphics[keepaspectratio]{calculo/calculo_files/figure-pdf/cell-46-output-1.pdf}}

\begin{Shaded}
\begin{Highlighting}[]
\NormalTok{trapezoid(y, x)}
\end{Highlighting}
\end{Shaded}

\begin{verbatim}
np.float64(0.9999961999999998)
\end{verbatim}

\chapter{Isso é tudo}\label{isso-uxe9-tudo-7}

By BEM

\chapter{Equações Diferenciais Ordinárias
(EDOs)}\label{equauxe7uxf5es-diferenciais-ordinuxe1rias-edos}

\subsubsection{by BEM, Last update:
07/07/2023}\label{by-bem-last-update-07072023}

\begin{Shaded}
\begin{Highlighting}[]
\ImportTok{import}\NormalTok{ numpy }\ImportTok{as}\NormalTok{ np}
\ImportTok{import}\NormalTok{ pylab }\ImportTok{as}\NormalTok{ pl}
\end{Highlighting}
\end{Shaded}

\chapter{Método de Euler}\label{muxe9todo-de-euler}

\section{\texorpdfstring{\(y' = y\)}{y\textquotesingle{} = y}}\label{y-y}

\section{\texorpdfstring{\(y(x_0) = y_0\)}{y(x\_0) = y\_0}}\label{yx_0-y_0}

\begin{itemize}
\tightlist
\item
  \(x_0\) = 1.0
\item
  \(y_0\) = 4.6
\end{itemize}

\begin{Shaded}
\begin{Highlighting}[]
\KeywordTok{def}\NormalTok{ f(x, c):}
    \ControlFlowTok{return}\NormalTok{ c}\OperatorTok{*}\NormalTok{(np.exp(x))}
\end{Highlighting}
\end{Shaded}

\begin{Shaded}
\begin{Highlighting}[]
\NormalTok{x0 }\OperatorTok{=} \FloatTok{1.0}
\NormalTok{y0 }\OperatorTok{=} \FloatTok{4.6}
\end{Highlighting}
\end{Shaded}

\begin{Shaded}
\begin{Highlighting}[]
\NormalTok{pl.plot(x0, y0, }\StringTok{\textquotesingle{}ko\textquotesingle{}}\NormalTok{)}
\NormalTok{pl.plot([x0, x0], [}\DecValTok{0}\NormalTok{, y0], }\StringTok{\textquotesingle{}k{-}{-}\textquotesingle{}}\NormalTok{)}
\NormalTok{pl.plot([}\DecValTok{0}\NormalTok{, x0], [y0, y0], }\StringTok{\textquotesingle{}k{-}{-}\textquotesingle{}}\NormalTok{)}
\NormalTok{pl.ylim(}\DecValTok{0}\NormalTok{, }\DecValTok{100}\NormalTok{)}
\NormalTok{pl.xlim(}\DecValTok{0}\NormalTok{, }\DecValTok{5}\NormalTok{)}
\NormalTok{pl.show()}
\end{Highlighting}
\end{Shaded}

\pandocbounded{\includegraphics[keepaspectratio]{edos_files/figure-pdf/cell-5-output-1.pdf}}

\chapter{Qual vai ser o próximo
ponto?}\label{qual-vai-ser-o-pruxf3ximo-ponto}

\begin{Shaded}
\begin{Highlighting}[]
\NormalTok{dy0 }\OperatorTok{=}\NormalTok{ y0}
\NormalTok{h }\OperatorTok{=} \FloatTok{0.5}

\NormalTok{y1 }\OperatorTok{=}\NormalTok{ y0 }\OperatorTok{+}\NormalTok{ h}\OperatorTok{*}\NormalTok{dy0}
\NormalTok{x1 }\OperatorTok{=}\NormalTok{ x0 }\OperatorTok{+}\NormalTok{ h }
\end{Highlighting}
\end{Shaded}

\begin{Shaded}
\begin{Highlighting}[]
\BuiltInTok{print}\NormalTok{(x1, y1)}
\end{Highlighting}
\end{Shaded}

\begin{verbatim}
1.5 6.8999999999999995
\end{verbatim}

\begin{Shaded}
\begin{Highlighting}[]
\NormalTok{pl.plot(x0, y0, }\StringTok{\textquotesingle{}ko\textquotesingle{}}\NormalTok{)}
\NormalTok{pl.plot([x0, x0], [}\DecValTok{0}\NormalTok{, y0], }\StringTok{\textquotesingle{}k{-}{-}\textquotesingle{}}\NormalTok{)}
\NormalTok{pl.plot([}\DecValTok{0}\NormalTok{, x0], [y0, y0], }\StringTok{\textquotesingle{}k{-}{-}\textquotesingle{}}\NormalTok{)}

\NormalTok{pl.plot(x1, y1, }\StringTok{\textquotesingle{}ro\textquotesingle{}}\NormalTok{)}
\NormalTok{pl.plot([x1, x1], [}\DecValTok{0}\NormalTok{, y1], }\StringTok{\textquotesingle{}r{-}{-}\textquotesingle{}}\NormalTok{)}
\NormalTok{pl.plot([}\DecValTok{0}\NormalTok{, x1], [y1, y1], }\StringTok{\textquotesingle{}r{-}{-}\textquotesingle{}}\NormalTok{)}


\NormalTok{pl.ylim(}\DecValTok{0}\NormalTok{, }\DecValTok{100}\NormalTok{)}
\NormalTok{pl.xlim(}\DecValTok{0}\NormalTok{, }\DecValTok{5}\NormalTok{)}
\NormalTok{pl.show()}
\end{Highlighting}
\end{Shaded}

\pandocbounded{\includegraphics[keepaspectratio]{edos_files/figure-pdf/cell-8-output-1.pdf}}

\chapter{Qual será o próximo
ponto?}\label{qual-seruxe1-o-pruxf3ximo-ponto}

\begin{Shaded}
\begin{Highlighting}[]
\NormalTok{dy1 }\OperatorTok{=}\NormalTok{ y1}
\NormalTok{h }\OperatorTok{=} \FloatTok{0.5}

\NormalTok{y2 }\OperatorTok{=}\NormalTok{ y1 }\OperatorTok{+}\NormalTok{ h}\OperatorTok{*}\NormalTok{dy1}
\NormalTok{x2 }\OperatorTok{=}\NormalTok{ x1 }\OperatorTok{+}\NormalTok{ h }
\end{Highlighting}
\end{Shaded}

\begin{Shaded}
\begin{Highlighting}[]
\NormalTok{pl.plot(x0, y0, }\StringTok{\textquotesingle{}ko\textquotesingle{}}\NormalTok{)}
\NormalTok{pl.plot([x0, x0], [}\DecValTok{0}\NormalTok{, y0], }\StringTok{\textquotesingle{}k{-}{-}\textquotesingle{}}\NormalTok{)}
\NormalTok{pl.plot([}\DecValTok{0}\NormalTok{, x0], [y0, y0], }\StringTok{\textquotesingle{}k{-}{-}\textquotesingle{}}\NormalTok{)}

\NormalTok{pl.plot(x1, y1, }\StringTok{\textquotesingle{}ro\textquotesingle{}}\NormalTok{)}
\NormalTok{pl.plot([x1, x1], [}\DecValTok{0}\NormalTok{, y1], }\StringTok{\textquotesingle{}r{-}{-}\textquotesingle{}}\NormalTok{)}
\NormalTok{pl.plot([}\DecValTok{0}\NormalTok{, x1], [y1, y1], }\StringTok{\textquotesingle{}r{-}{-}\textquotesingle{}}\NormalTok{)}

\NormalTok{pl.plot(x2, y2, }\StringTok{\textquotesingle{}bo\textquotesingle{}}\NormalTok{)}
\NormalTok{pl.plot([x2, x2], [}\DecValTok{0}\NormalTok{, y1], }\StringTok{\textquotesingle{}b{-}{-}\textquotesingle{}}\NormalTok{)}
\NormalTok{pl.plot([}\DecValTok{0}\NormalTok{, x2], [y2, y2], }\StringTok{\textquotesingle{}b{-}{-}\textquotesingle{}}\NormalTok{)}

\NormalTok{pl.ylim(}\DecValTok{0}\NormalTok{, }\DecValTok{100}\NormalTok{)}
\NormalTok{pl.xlim(}\DecValTok{0}\NormalTok{, }\DecValTok{5}\NormalTok{)}
\NormalTok{pl.show()}
\end{Highlighting}
\end{Shaded}

\pandocbounded{\includegraphics[keepaspectratio]{edos_files/figure-pdf/cell-10-output-1.pdf}}

\chapter{Agora como organizar isso de forma
eficiente?}\label{agora-como-organizar-isso-de-forma-eficiente}

\begin{Shaded}
\begin{Highlighting}[]
\KeywordTok{def}\NormalTok{ edo(x, y):}
    \ControlFlowTok{return}\NormalTok{ y}
\end{Highlighting}
\end{Shaded}

\begin{Shaded}
\begin{Highlighting}[]
\NormalTok{x }\OperatorTok{=}\NormalTok{ np.array([}\FloatTok{1.0}\NormalTok{])}
\NormalTok{y }\OperatorTok{=}\NormalTok{ np.array([}\FloatTok{4.6}\NormalTok{])}
\end{Highlighting}
\end{Shaded}

\begin{Shaded}
\begin{Highlighting}[]
\NormalTok{h }\OperatorTok{=} \FloatTok{0.5}

\ControlFlowTok{for}\NormalTok{ i }\KeywordTok{in} \BuiltInTok{range}\NormalTok{(}\DecValTok{1}\NormalTok{,}\DecValTok{10}\NormalTok{):}
    \BuiltInTok{print}\NormalTok{(i)}
\NormalTok{    dy }\OperatorTok{=}\NormalTok{ edo(x, y)}
\NormalTok{    y }\OperatorTok{=}\NormalTok{ np.append(y, y[i}\OperatorTok{{-}}\DecValTok{1}\NormalTok{] }\OperatorTok{+}\NormalTok{ h}\OperatorTok{*}\NormalTok{dy[i}\OperatorTok{{-}}\DecValTok{1}\NormalTok{])}
\NormalTok{    x }\OperatorTok{=}\NormalTok{ np.append(x, x[i}\OperatorTok{{-}}\DecValTok{1}\NormalTok{] }\OperatorTok{+}\NormalTok{ h)}
\end{Highlighting}
\end{Shaded}

\begin{verbatim}
1
2
3
4
5
6
7
8
9
\end{verbatim}

\begin{Shaded}
\begin{Highlighting}[]
\NormalTok{pl.plot(x, y, }\StringTok{\textquotesingle{}k.\textquotesingle{}}\NormalTok{) }\CommentTok{\# Solução numérica}
\end{Highlighting}
\end{Shaded}

\pandocbounded{\includegraphics[keepaspectratio]{edos_files/figure-pdf/cell-14-output-1.pdf}}

\chapter{Qual a solução
algebrica?}\label{qual-a-soluuxe7uxe3o-algebrica}

\section{\texorpdfstring{\(y(x) = y_0~.~e^{x - x0}\)}{y(x) = y\_0\textasciitilde.\textasciitilde e\^{}\{x - x0\}}}\label{yx-y_0.ex---x0}

\begin{Shaded}
\begin{Highlighting}[]
\KeywordTok{def}\NormalTok{ f(x):}
    \ControlFlowTok{return} \FloatTok{4.6}\OperatorTok{*}\NormalTok{np.exp(x }\OperatorTok{{-}} \FloatTok{1.0}\NormalTok{)}
\end{Highlighting}
\end{Shaded}

\begin{Shaded}
\begin{Highlighting}[]
\NormalTok{pl.plot(x, y, }\StringTok{\textquotesingle{}k.\textquotesingle{}}\NormalTok{) }\CommentTok{\# Solução numérica}
\NormalTok{pl.plot(x, f(x), }\StringTok{\textquotesingle{}r.{-}\textquotesingle{}}\NormalTok{) }\CommentTok{\# Solução algebrica}
\NormalTok{pl.ylim(}\DecValTok{0}\NormalTok{, }\DecValTok{100}\NormalTok{)}
\NormalTok{pl.xlim(}\DecValTok{0}\NormalTok{, }\DecValTok{5}\NormalTok{)}
\end{Highlighting}
\end{Shaded}

\pandocbounded{\includegraphics[keepaspectratio]{edos_files/figure-pdf/cell-16-output-1.pdf}}

\chapter{Não parece que fomos tão bem assim\ldots.
Pq?}\label{nuxe3o-parece-que-fomos-tuxe3o-bem-assim.-pq}

\begin{Shaded}
\begin{Highlighting}[]
\NormalTok{x }\OperatorTok{=}\NormalTok{ np.array([}\FloatTok{1.0}\NormalTok{])}
\NormalTok{y }\OperatorTok{=}\NormalTok{ np.array([}\FloatTok{4.6}\NormalTok{])}

\NormalTok{h }\OperatorTok{=} \FloatTok{0.05}

\ControlFlowTok{for}\NormalTok{ i }\KeywordTok{in} \BuiltInTok{range}\NormalTok{(}\DecValTok{1}\NormalTok{,}\DecValTok{100}\NormalTok{):}
\NormalTok{    dy }\OperatorTok{=}\NormalTok{ edo(x, y)}
\NormalTok{    y }\OperatorTok{=}\NormalTok{ np.append(y, y[i}\OperatorTok{{-}}\DecValTok{1}\NormalTok{] }\OperatorTok{+}\NormalTok{ h}\OperatorTok{*}\NormalTok{dy[i}\OperatorTok{{-}}\DecValTok{1}\NormalTok{])}
\NormalTok{    x }\OperatorTok{=}\NormalTok{ np.append(x, x[i}\OperatorTok{{-}}\DecValTok{1}\NormalTok{] }\OperatorTok{+}\NormalTok{ h)}
\end{Highlighting}
\end{Shaded}

\begin{Shaded}
\begin{Highlighting}[]
\NormalTok{pl.plot(x, y, }\StringTok{\textquotesingle{}k.\textquotesingle{}}\NormalTok{) }\CommentTok{\# Solução numérica}
\NormalTok{pl.plot(x, f(x), }\StringTok{\textquotesingle{}r.{-}\textquotesingle{}}\NormalTok{) }\CommentTok{\# Solução numérica}

\NormalTok{pl.ylim(}\DecValTok{0}\NormalTok{, }\DecValTok{100}\NormalTok{)}
\NormalTok{pl.xlim(}\DecValTok{0}\NormalTok{, }\DecValTok{5}\NormalTok{)}
\end{Highlighting}
\end{Shaded}

\pandocbounded{\includegraphics[keepaspectratio]{edos_files/figure-pdf/cell-18-output-1.pdf}}

\chapter{Diminuindo ainda mais o
passo\ldots{}}\label{diminuindo-ainda-mais-o-passo}

\begin{Shaded}
\begin{Highlighting}[]
\NormalTok{x }\OperatorTok{=}\NormalTok{ np.array([}\FloatTok{1.0}\NormalTok{])}
\NormalTok{y }\OperatorTok{=}\NormalTok{ np.array([}\FloatTok{4.6}\NormalTok{])}

\NormalTok{h }\OperatorTok{=} \FloatTok{0.005}

\ControlFlowTok{for}\NormalTok{ i }\KeywordTok{in} \BuiltInTok{range}\NormalTok{(}\DecValTok{1}\NormalTok{,}\DecValTok{1000}\NormalTok{):}
\NormalTok{    dy }\OperatorTok{=}\NormalTok{ edo(x, y)}
\NormalTok{    y }\OperatorTok{=}\NormalTok{ np.append(y, y[i}\OperatorTok{{-}}\DecValTok{1}\NormalTok{] }\OperatorTok{+}\NormalTok{ h}\OperatorTok{*}\NormalTok{dy[i}\OperatorTok{{-}}\DecValTok{1}\NormalTok{])}
\NormalTok{    x }\OperatorTok{=}\NormalTok{ np.append(x, x[i}\OperatorTok{{-}}\DecValTok{1}\NormalTok{] }\OperatorTok{+}\NormalTok{ h)}
\end{Highlighting}
\end{Shaded}

\begin{Shaded}
\begin{Highlighting}[]
\NormalTok{pl.plot(x, y, }\StringTok{\textquotesingle{}k.\textquotesingle{}}\NormalTok{) }\CommentTok{\# Solução numérica}
\NormalTok{pl.plot(x, f(x), }\StringTok{\textquotesingle{}r{-}\textquotesingle{}}\NormalTok{) }\CommentTok{\# Solução numérica}

\NormalTok{pl.ylim(}\DecValTok{0}\NormalTok{, }\DecValTok{100}\NormalTok{)}
\NormalTok{pl.xlim(}\DecValTok{0}\NormalTok{, }\DecValTok{5}\NormalTok{)}
\end{Highlighting}
\end{Shaded}

\pandocbounded{\includegraphics[keepaspectratio]{edos_files/figure-pdf/cell-20-output-1.pdf}}

\chapter{Quanto menor o passo, melhor o nosso
resultado!}\label{quanto-menor-o-passo-melhor-o-nosso-resultado}

\chapter{Método de Runge-Kutta}\label{muxe9todo-de-runge-kutta}

\begin{itemize}
\tightlist
\item
  Ordem 2
\end{itemize}

\begin{Shaded}
\begin{Highlighting}[]
\KeywordTok{def}\NormalTok{ edo(x, y):}
    \ControlFlowTok{return}\NormalTok{ y}
\end{Highlighting}
\end{Shaded}

\begin{Shaded}
\begin{Highlighting}[]
\NormalTok{x0 }\OperatorTok{=} \FloatTok{1.0}
\NormalTok{y0 }\OperatorTok{=} \FloatTok{4.6}

\NormalTok{h }\OperatorTok{=} \FloatTok{0.5}
\end{Highlighting}
\end{Shaded}

\begin{Shaded}
\begin{Highlighting}[]
\NormalTok{k1 }\OperatorTok{=}\NormalTok{ edo(x0, y0)}
\NormalTok{k2 }\OperatorTok{=}\NormalTok{ edo(x0 }\OperatorTok{+}\NormalTok{ h, y0 }\OperatorTok{+}\NormalTok{ h}\OperatorTok{*}\NormalTok{k1)}

\NormalTok{x1 }\OperatorTok{=}\NormalTok{ x0 }\OperatorTok{+}\NormalTok{ h}
\NormalTok{y1 }\OperatorTok{=}\NormalTok{ y0 }\OperatorTok{+}\NormalTok{ (h}\OperatorTok{/}\DecValTok{2}\NormalTok{)}\OperatorTok{*}\NormalTok{(k1 }\OperatorTok{+}\NormalTok{ k2)}
\end{Highlighting}
\end{Shaded}

\begin{Shaded}
\begin{Highlighting}[]
\NormalTok{pl.plot(x0, y0, }\StringTok{\textquotesingle{}ko\textquotesingle{}}\NormalTok{)}
\NormalTok{pl.plot([x0, x0], [}\DecValTok{0}\NormalTok{, y0], }\StringTok{\textquotesingle{}k{-}{-}\textquotesingle{}}\NormalTok{)}
\NormalTok{pl.plot([}\DecValTok{0}\NormalTok{, x0], [y0, y0], }\StringTok{\textquotesingle{}k{-}{-}\textquotesingle{}}\NormalTok{)}

\NormalTok{pl.plot(x1, y1, }\StringTok{\textquotesingle{}ro\textquotesingle{}}\NormalTok{)}
\NormalTok{pl.plot([x1, x1], [}\DecValTok{0}\NormalTok{, y1], }\StringTok{\textquotesingle{}r{-}{-}\textquotesingle{}}\NormalTok{)}
\NormalTok{pl.plot([}\DecValTok{0}\NormalTok{, x1], [y1, y1], }\StringTok{\textquotesingle{}r{-}{-}\textquotesingle{}}\NormalTok{)}


\NormalTok{pl.ylim(}\DecValTok{0}\NormalTok{, }\DecValTok{100}\NormalTok{)}
\NormalTok{pl.xlim(}\DecValTok{0}\NormalTok{, }\DecValTok{5}\NormalTok{)}
\NormalTok{pl.show()}
\end{Highlighting}
\end{Shaded}

\pandocbounded{\includegraphics[keepaspectratio]{edos_files/figure-pdf/cell-24-output-1.pdf}}

\begin{Shaded}
\begin{Highlighting}[]
\NormalTok{k1 }\OperatorTok{=}\NormalTok{ edo(x1, y1)}
\NormalTok{k2 }\OperatorTok{=}\NormalTok{ edo(x1 }\OperatorTok{+}\NormalTok{ h, y1 }\OperatorTok{+}\NormalTok{ h}\OperatorTok{*}\NormalTok{k1)}

\NormalTok{x2 }\OperatorTok{=}\NormalTok{ x1 }\OperatorTok{+}\NormalTok{ h}
\NormalTok{y2 }\OperatorTok{=}\NormalTok{ y1 }\OperatorTok{+}\NormalTok{ (h}\OperatorTok{/}\DecValTok{2}\NormalTok{)}\OperatorTok{*}\NormalTok{(k1 }\OperatorTok{+}\NormalTok{ k2)}
\end{Highlighting}
\end{Shaded}

\begin{Shaded}
\begin{Highlighting}[]
\NormalTok{pl.plot(x0, y0, }\StringTok{\textquotesingle{}ko\textquotesingle{}}\NormalTok{)}
\NormalTok{pl.plot([x0, x0], [}\DecValTok{0}\NormalTok{, y0], }\StringTok{\textquotesingle{}k{-}{-}\textquotesingle{}}\NormalTok{)}
\NormalTok{pl.plot([}\DecValTok{0}\NormalTok{, x0], [y0, y0], }\StringTok{\textquotesingle{}k{-}{-}\textquotesingle{}}\NormalTok{)}

\NormalTok{pl.plot(x1, y1, }\StringTok{\textquotesingle{}ro\textquotesingle{}}\NormalTok{)}
\NormalTok{pl.plot([x1, x1], [}\DecValTok{0}\NormalTok{, y1], }\StringTok{\textquotesingle{}r{-}{-}\textquotesingle{}}\NormalTok{)}
\NormalTok{pl.plot([}\DecValTok{0}\NormalTok{, x1], [y1, y1], }\StringTok{\textquotesingle{}r{-}{-}\textquotesingle{}}\NormalTok{)}

\NormalTok{pl.plot(x2, y2, }\StringTok{\textquotesingle{}bo\textquotesingle{}}\NormalTok{)}
\NormalTok{pl.plot([x2, x2], [}\DecValTok{0}\NormalTok{, y1], }\StringTok{\textquotesingle{}b{-}{-}\textquotesingle{}}\NormalTok{)}
\NormalTok{pl.plot([}\DecValTok{0}\NormalTok{, x2], [y2, y2], }\StringTok{\textquotesingle{}b{-}{-}\textquotesingle{}}\NormalTok{)}

\NormalTok{pl.ylim(}\DecValTok{0}\NormalTok{, }\DecValTok{100}\NormalTok{)}
\NormalTok{pl.xlim(}\DecValTok{0}\NormalTok{, }\DecValTok{5}\NormalTok{)}
\NormalTok{pl.show()}
\end{Highlighting}
\end{Shaded}

\pandocbounded{\includegraphics[keepaspectratio]{edos_files/figure-pdf/cell-26-output-1.pdf}}

\chapter{Agora como organizar isso de forma
eficiente?}\label{agora-como-organizar-isso-de-forma-eficiente-1}

\begin{Shaded}
\begin{Highlighting}[]
\NormalTok{xk }\OperatorTok{=}\NormalTok{ np.array([}\FloatTok{1.0}\NormalTok{])}
\NormalTok{yk }\OperatorTok{=}\NormalTok{ np.array([}\FloatTok{4.6}\NormalTok{])}
\end{Highlighting}
\end{Shaded}

\begin{Shaded}
\begin{Highlighting}[]
\NormalTok{h }\OperatorTok{=} \FloatTok{0.05}

\ControlFlowTok{for}\NormalTok{ i }\KeywordTok{in} \BuiltInTok{range}\NormalTok{(}\DecValTok{1}\NormalTok{,}\DecValTok{100}\NormalTok{):}
    \BuiltInTok{print}\NormalTok{(i)}
\NormalTok{    k1 }\OperatorTok{=}\NormalTok{ edo(xk[i}\OperatorTok{{-}}\DecValTok{1}\NormalTok{], yk[i}\OperatorTok{{-}}\DecValTok{1}\NormalTok{])}
\NormalTok{    k2 }\OperatorTok{=}\NormalTok{ edo(xk[i}\OperatorTok{{-}}\DecValTok{1}\NormalTok{] }\OperatorTok{+}\NormalTok{ h, yk[i}\OperatorTok{{-}}\DecValTok{1}\NormalTok{] }\OperatorTok{+}\NormalTok{ h}\OperatorTok{*}\NormalTok{k1)}

\NormalTok{    xk }\OperatorTok{=}\NormalTok{ np.append(xk, xk[i}\OperatorTok{{-}}\DecValTok{1}\NormalTok{] }\OperatorTok{+}\NormalTok{ h)}
\NormalTok{    yk }\OperatorTok{=}\NormalTok{ np.append(yk, yk[i}\OperatorTok{{-}}\DecValTok{1}\NormalTok{] }\OperatorTok{+}\NormalTok{ (h}\OperatorTok{/}\DecValTok{2}\NormalTok{)}\OperatorTok{*}\NormalTok{(k1 }\OperatorTok{+}\NormalTok{ k2))}
\end{Highlighting}
\end{Shaded}

\begin{verbatim}
1
2
3
4
5
6
7
8
9
10
11
12
13
14
15
16
17
18
19
20
21
22
23
24
25
26
27
28
29
30
31
32
33
34
35
36
37
38
39
40
41
42
43
44
45
46
47
48
49
50
51
52
53
54
55
56
57
58
59
60
61
62
63
64
65
66
67
68
69
70
71
72
73
74
75
76
77
78
79
80
81
82
83
84
85
86
87
88
89
90
91
92
93
94
95
96
97
98
99
\end{verbatim}

\begin{Shaded}
\begin{Highlighting}[]
\NormalTok{pl.plot(xk, yk, }\StringTok{\textquotesingle{}k.\textquotesingle{}}\NormalTok{) }\CommentTok{\# Solução numérica}
\end{Highlighting}
\end{Shaded}

\pandocbounded{\includegraphics[keepaspectratio]{edos_files/figure-pdf/cell-29-output-1.pdf}}

\begin{Shaded}
\begin{Highlighting}[]
\NormalTok{pl.plot(x, y, }\StringTok{\textquotesingle{}k{-}.\textquotesingle{}}\NormalTok{) }\CommentTok{\# Solução numérica}
\NormalTok{pl.plot(xk, yk, }\StringTok{\textquotesingle{}b{-}+\textquotesingle{}}\NormalTok{) }\CommentTok{\# Solução numérica}
\NormalTok{pl.plot(x, f(x), }\StringTok{\textquotesingle{}r.{-}\textquotesingle{}}\NormalTok{) }\CommentTok{\# Solução algebrica}
\NormalTok{pl.ylim(}\DecValTok{0}\NormalTok{, }\DecValTok{100}\NormalTok{)}
\NormalTok{pl.xlim(}\DecValTok{0}\NormalTok{, }\DecValTok{5}\NormalTok{)}
\end{Highlighting}
\end{Shaded}

\pandocbounded{\includegraphics[keepaspectratio]{edos_files/figure-pdf/cell-30-output-1.pdf}}

\chapter{Mesmo com passo grande, temos uma acuracia muito
melhor!}\label{mesmo-com-passo-grande-temos-uma-acuracia-muito-melhor}

\begin{Shaded}
\begin{Highlighting}[]
\KeywordTok{def}\NormalTok{ edo(x, y):}
    \ControlFlowTok{return}\NormalTok{ (x}\OperatorTok{**}\DecValTok{2}\NormalTok{) }\OperatorTok{+}\NormalTok{ (y}\OperatorTok{**}\DecValTok{2}\NormalTok{)}
\end{Highlighting}
\end{Shaded}

\begin{Shaded}
\begin{Highlighting}[]
\NormalTok{xk }\OperatorTok{=}\NormalTok{ np.array([}\DecValTok{0}\NormalTok{])}
\NormalTok{yk }\OperatorTok{=}\NormalTok{ np.array([}\DecValTok{1}\NormalTok{])}

\NormalTok{h }\OperatorTok{=} \FloatTok{0.05}

\ControlFlowTok{for}\NormalTok{ i }\KeywordTok{in} \BuiltInTok{range}\NormalTok{(}\DecValTok{1}\NormalTok{,}\DecValTok{21}\NormalTok{):}
    \BuiltInTok{print}\NormalTok{(i)}
\NormalTok{    k1 }\OperatorTok{=}\NormalTok{ edo(xk[i}\OperatorTok{{-}}\DecValTok{1}\NormalTok{], yk[i}\OperatorTok{{-}}\DecValTok{1}\NormalTok{])}
\NormalTok{    k2 }\OperatorTok{=}\NormalTok{ edo(xk[i}\OperatorTok{{-}}\DecValTok{1}\NormalTok{] }\OperatorTok{+}\NormalTok{ h}\OperatorTok{/}\DecValTok{2}\NormalTok{, yk[i}\OperatorTok{{-}}\DecValTok{1}\NormalTok{] }\OperatorTok{+}\NormalTok{ h}\OperatorTok{*}\NormalTok{k1}\OperatorTok{/}\DecValTok{2}\NormalTok{)}
\NormalTok{    k3 }\OperatorTok{=}\NormalTok{ edo(xk[i}\OperatorTok{{-}}\DecValTok{1}\NormalTok{] }\OperatorTok{+}\NormalTok{ h}\OperatorTok{/}\DecValTok{2}\NormalTok{, yk[i}\OperatorTok{{-}}\DecValTok{1}\NormalTok{] }\OperatorTok{+}\NormalTok{ h}\OperatorTok{*}\NormalTok{k2}\OperatorTok{/}\DecValTok{2}\NormalTok{)}
\NormalTok{    k4 }\OperatorTok{=}\NormalTok{ edo(xk[i}\OperatorTok{{-}}\DecValTok{1}\NormalTok{] }\OperatorTok{+}\NormalTok{ h, yk[i}\OperatorTok{{-}}\DecValTok{1}\NormalTok{] }\OperatorTok{+}\NormalTok{ h}\OperatorTok{*}\NormalTok{k3)}

\NormalTok{    xk }\OperatorTok{=}\NormalTok{ np.append(xk, xk[i}\OperatorTok{{-}}\DecValTok{1}\NormalTok{] }\OperatorTok{+}\NormalTok{ h)}
\NormalTok{    yk }\OperatorTok{=}\NormalTok{ np.append(yk, yk[i}\OperatorTok{{-}}\DecValTok{1}\NormalTok{] }\OperatorTok{+}\NormalTok{ (h}\OperatorTok{/}\DecValTok{6}\NormalTok{)}\OperatorTok{*}\NormalTok{(k1 }\OperatorTok{+}\NormalTok{ k2 }\OperatorTok{+}\NormalTok{ k3 }\OperatorTok{+}\NormalTok{ k4))}
\end{Highlighting}
\end{Shaded}

\begin{verbatim}
1
2
3
4
5
6
7
8
9
10
11
12
13
14
15
16
17
18
19
20
\end{verbatim}

\begin{Shaded}
\begin{Highlighting}[]
\NormalTok{pl.plot(xk, yk, }\StringTok{\textquotesingle{}k.{-}\textquotesingle{}}\NormalTok{)}
\NormalTok{pl.axvline(}\DecValTok{1}\NormalTok{)}
\end{Highlighting}
\end{Shaded}

\pandocbounded{\includegraphics[keepaspectratio]{edos_files/figure-pdf/cell-33-output-1.pdf}}

\bookmarksetup{startatroot}

\chapter*{Referências}\label{referuxeancias}
\addcontentsline{toc}{chapter}{Referências}

\markboth{Referências}{Referências}

\protect\phantomsection\label{refs}

\cleardoublepage
\phantomsection
\addcontentsline{toc}{part}{Apêndices}
\appendix

\chapter{Lógica Booleana}\label{sec-appendix-a}




\end{document}
